%========================================================
%  main.tex  — Triple Interface to (Long) Goldbach
%========================================================
\DeclareUnicodeCharacter{2032}{\ensuremath{'}}
\documentclass[11pt]{article}

%------------ Page & fonts
\usepackage[margin=1.1in]{geometry}
\usepackage[T1]{fontenc}
\usepackage[utf8]{inputenc}
\usepackage{microtype}
\usepackage{lmodern}
\usepackage[backend=biber]{biblatex}
\addbibresource{bibliography.bib}
\providecommand{\Z}{\mathbb{Z}}

%------------ Math & theorems
\usepackage{amsmath, amssymb, amsthm, amsfonts, mathtools, tabularx,enumitem}
\DeclarePairedDelimiter{\angles}{\langle}{\rangle}
\usepackage{thmtools}
\newcommand{\F}{\mathcal{F}}
\usepackage{bm}
\usepackage{commath}
\usepackage{enumitem}
\usepackage{physics}
\usepackage{graphicx}
\declaretheorem[title=Theorem F,numberwithin=section]{theoremF}
\allowdisplaybreaks
\DeclareMathOperator{\dist}{dist}

%------------ Custom abbreviations

\newcommand{\ee}[1]{\mathrm{e}^{#1}}
\newcommand{\ol}[1]{\overline{#1}}
\newcommand{\ppart}{_+}
\newcommand{\pospart}[1]{\left(#1\right)_{+}}
\renewcommand{\Z}{\mathbb{Z}}
\newcommand{\TFA}{\textup{TFA}}
\newcommand{\AO}{\textup{AO}}
\newcommand{\BG}{\textup{BG}}
\newcommand{\PSB}{\textup{PSB}}
\newcommand{\SSU}{\textup{SSU}}
\newcommand{\R}{\mathbb{R}}

%------------ Graphics (rare)
\usepackage{graphicx}
\usepackage{xcolor}

%------------ Hyperlinks & clever references
\usepackage[colorlinks,allcolors=blue!60!black]{hyperref}
\usepackage[nameinlink,noabbrev]{cleveref}


%------------ Bibliography (use your .bib file or comment out)
% \usepackage[backend=biber,style=alphabetic,maxbibnames=9]{biblatex}
% \addbibresource{ti_goldbach.bib}

%========================================================
%  Theorem environments
%========================================================
\numberwithin{equation}{section}
\declaretheorem[name=Theorem,numberwithin=section]{theorem}
\declaretheorem[name=Proposition,sibling=theorem]{proposition}
\declaretheorem[name=Lemma,sibling=theorem]{lemma}
\declaretheorem[name=Corollary,sibling=theorem]{corollary}
\declaretheorem[name=Claim,sibling=theorem]{claim}
\declaretheorem[name=Definition,sibling=theorem]{definition}
\declaretheorem[name=Remark,sibling=theorem,style=remark]{remark}
\declaretheorem[name=Example,sibling=theorem,style=remark]{example}

%========================================================
%  Custom commands & operators (match section files)
%========================================================
\newcommand{\llangle}{\langle\!\langle}
\newcommand{\rrangle}{\rangle\!\rangle}
\newcommand{\e}{\mathrm{e}}
\newcommand{\ii}{\mathrm{i}}
% \newcommand{\dd}{\mathrm{d}}
\newcommand{\1}{\mathbf{1}}
\newcommand{\eps}{\varepsilon}
\newcommand{\sgn}{\operatorname{sgn}}
\renewcommand{\Res}{\operatorname{Res}}
\newcommand{\Var}{\operatorname{Var}}
\newcommand{\supp}{\operatorname{supp}}

\newcommand{\TT}{\mathbb{T}}
\newcommand{\T}{\mathbb{T}}
\newcommand{\NN}{\mathbb{N}}
\newcommand{\RR}{\mathbb{R}}

\newcommand{\asympX}{\,\asymp\,}
\renewcommand{\qq}{\quad}
\newcommand{\qqand}{\quad\text{and}\quad}
\newcommand{\avg}{\mathbb{E}}

\DeclareMathOperator{\Sym}{Sym}
\DeclareMathOperator{\Varop}{Var}

% Short exponential
\newcommand{\ex}[1]{\mathrm{e}^{2\pi\ii\, #1}}

% Von Mangoldt and related
\newcommand{\Lam}{\Lambda}
\newcommand{\Lams}{\Lambda^\sharp}

% Bank/window notations
\newcommand{\wh}{\widehat}
\newcommand{\whw}{\widehat{w}_{\tau^\ast}}
\newcommand{\whK}{\widehat{K}}
\newcommand{\whKh}{\widehat{K}_{H}}

% Smoothing scale
\newcommand{\HH}{H}
\newcommand{\QQ}{Q}

\newcommand{\cA}{\mathcal{A}}
\newcommand{\PA}{P_{\mathcal{A}}} % Fourier projection to the bank

% ==== Appendix macros (safe to \input into your main doc preamble or before the appendices) ====

\usepackage{amssymb,amsmath,amsthm,mathtools}

\newcommand{\PreErrTypeI}{\mathcal{E}_{\mathrm{pre}}}
\newcommand{\N}{\mathbb{N}}
\newcommand{\SingSeriesBankPosSpec}[1]{\mathrm{SingSeriesBankPosSpec}(#1)}

\DeclareMathOperator{\LB}{LB}
\DeclareMathOperator{\Margin}{Margin}
\DeclareMathOperator{\TailBudget}{TailBudget}
\DeclareMathOperator{\MainSym}{MainSym}
\DeclareMathOperator{\MainTail}{MainTail}
\DeclareMathOperator{\MainPP}{MainPP}
\DeclareMathOperator{\ErrBudget}{ErrBudget}
\DeclareMathOperator{\ErrTypeI}{ErrTypeI}
\DeclareMathOperator{\ErrSSU}{ErrSSU}
\DeclareMathOperator{\ErrTail}{ErrTail}
\DeclareMathOperator{\wrapdist}{wrapdist}

\newcommand{\KH}{K_H}

\newcommand{\czero}{c_0}
\newcommand{\kappaFour}{\kappa_4}
\newcommand{\Aexp}{A}
\newcommand{\gexp}{\gamma}

\newcommand{\TTstar}{\textup{TT}^{\!*}}
\newcommand{\guard}{\mathrm{guard}}
\newcommand{\halfwidth}{w}

\theoremstyle{plain}
\newcommand{\TailCore}{\mathrm{TailCore}}
\newcommand{\TailOff}{\mathrm{TailOff}}

%========================================================
%  Title & metadata are in a separate file
%========================================================

\title{Goldbach and the Triple Interface Method:\\Banded TFA, Bilinear Geometry, and AO Dispersion}
\author{Edited by A. Quiller}
\date{\today}

%========================================================
%  Document
%========================================================
\newcommand{\bbT}{\mathbb{T}}
\begin{document}
\maketitle


%--------------------------------------------------------
%  00_frontmatter.tex — Title, authors, abstract, metadata
%--------------------------------------------------------

\begin{abstract}
We develop a research program to prove a bank--local closure for the binary Goldbach correlation using three ingredients that are unconditional at polylogarithmic scales: (i) a deterministic arc--occupancy (AO) dispersion bound for a Farey bank of moduli $q\le Q=(\log X)^{A\gamma}$, stating that no subfamily of $L$ arcs captures more than its proportional share $L/m$ of the bank energy up to $O(Q^{-1+\varepsilon})$; (ii) a two--axis short--shift uncertainty (SSU) estimate that yields a square--root saving $\big(H/X\big)^{1/2}$ for the minor--arc variance of bank–projected statistics, with $H=(\log X)^A$; and (iii) a Type--I leakage bound at the banked scale which, combined with AO dispersion and alias suppression, contributes $X/(H Q^2)$ to the global variance.

Let $H=(\log X)^A$ with $A\ge 10$ and $Q=H^\gamma$ with any fixed $0<\gamma<\tfrac12$. Uniform major--arc asymptotics and positivity of the singular series give a bank--projected mean $\gg (\log X)^{-2}$. The minor--arc variance is bounded by
\[
\mathrm{Var}\ \ll\ C_2\Big(\tfrac{H}{X}\Big)^{1/2}\ +\ \frac{C_3}{H\,Q^2}\,,
\]
for absolute $C_2,C_3\ge 1$. Hence, for each fixed $A$ and $\gamma\in(0,\tfrac12)$ there exists an explicit $X_0(A,\gamma)$ such that Paley--Zygmund yields positivity in every dyadic $[X,2X]$ for $X\ge X_0(A,\gamma)$. In our ledger the closure is governed solely by the two bank--local levers $(H/X)^{1/2}$ and $1/(H Q^2)$ together with the uniform main term. Center-uniformization completes the pointwise step.

\end{abstract}

% If you want the TOC on the title page or not; usually after \maketitle
\maketitle

% Optional: outline/roadmap paragraph
\noindent\textbf{This} is a research plan to solve the Goldbach conjecture.
Section~\ref{sec:intro} states the TI framework and the main theorem.
Sections~\ref{sec:TFA}--\ref{sec:AO-deterministic} develop TFA and AO.
Sections~\ref{sec:BG}--\ref{sec:ssu} establish BG (non-flat) and the slope--shift uncertainty.
Section~\ref{sec:typeI} treats type-I leakage. Section~\ref{sec:PSB} proves the PSB step.
We pass from dyadic mean positivity to pointwise positivity by \emph{two} routes.
\emph{Route A} is a quantitative Sampling+Bernstein (Logvinenko–Sereda) argument at bandwidth $1/H$.
\emph{Route B} is a bank–projection inequality: for the Fourier projector (or a smoothed variant) $P_{\cA}$ onto the bank,
\[
R_2(N)\ \ge\ (P_{\cA} M)(N) - \|P_{\cA}E\|_2 - \|P_{\cA^c}R_2\|_2,
\]
and the bank–summed singular series yields $(P_{\cA} M)(N)\gg 1/\log^2 X$ uniformly, while the two $L^2$ errors inherit the $ (\log X)^C(H/X)^{1/2}$ decay from SSU and Type–I.
Either route suffices; we present both to emphasize robustness.
Section~\ref{sec:center-uniform} carries out the arithmetic routing. Section~\ref{sec:final-assembly} assembles the provisional proof of Goldbach and records parameter choices. Appendices contain kernel details and auxiliary bounds and other treats.

% If you want the table of contents right after frontmatter:
% \setcounter{tocdepth}{2}
% \tableofcontents        % title, authors, abstract, MSC, keywords
\tableofcontents

%--------------------------------------------------------
%  Introductory materials
%--------------------------------------------------------
%========================================================
% sections/intro.tex
%========================================================
\section{Introduction and statement of results}\label{sec:intro}

\subsection{Parameters and overview}\label{ass:bank}
Let $X\to\infty$ be a large parameter. Fix $0<\varepsilon\le 1/10$ and $A\ge 10$, and set 
\[
H := (\log X)^A,\qquad Q := H^\gamma\quad\text{with}\quad 0<\gamma<\tfrac12.
\]
We work on dyadic blocks $[X,2X]$ and use two dyadic auxiliary scales $D,N$ satisfying
\[
X^\varepsilon \le D,N \le X^{1-\varepsilon},\qquad DN\asymp X.
\]
A fixed smooth bump $W\in C_c^\infty([1/2,2]^2)$ (with bounded derivatives up to order~$3$) is used to localize divisors in the bilinear pieces. Throughout, implied constants may depend on $(\varepsilon,A,\gamma,W)$ only.

\medskip
This paper develops and assembles a three–leg harmonic-analytic framework (the \emph{Triple Interface}, abbreviated \TFA, \AO, \BG) that suppresses short-shift correlations in arithmetic sums and feeds them into a $TT^\ast$ energy identity to produce \emph{Primes from Smooth Bilinear} (\PSB). The key quantitative scale is the \emph{short-shift window} of width $H$, which in our applications is polylogarithmic in $X$.

On the frequency side we deploy an \emph{arithmetic Fej\'er bank} (\TFA), a positive, smoothly-tapered union of $\asymp Q^2$ disjoint arcs of width $\asymp 1/H$ centred at rationals $a/q$ with $q\le Q$. Averaging over a deterministic family of small-modulus characters enforces \emph{dispersion} over the bank (\AO). On the physical side, a \emph{bilinear geometry} (\BG) analysis---based on a shearing of the multiplicative hyperbola and a slope--shift uncertainty (\SSU) bound on tubes---yields a square-root short-shift saving $(H/X)^{1/2}$ for Type-II structures \emph{without} assuming coefficient flatness (via a positive-operator, two-weight argument). A standard Type-I leakage estimate under the bank closes the $TT^\ast$ bound.

The output is a full Goldbach theorem: every interval of length $C(\log x)^A$ contains a \emph{Goldbach even}. We also give two ``uniformizers'' that turn band-limited mean positivity into a uniform lower bound at \emph{each} centre, providing a clean route to pointwise upgrades.

\subsection{Main theorem: full Goldbach (TI version)}
Write $R_{2,c}(n)$ for the smoothed number of representations of an even $n$ as $n=p_1+p_2$ with $p_i\equiv c_i \pmod W$ in a fixed admissible pair of residue classes modulo $W=(\log X)^B$, with the standard disposal of prime powers and a smooth dyadic weight (all fixed below in \S\ref{sec:AP}). The precise normalization does not affect the statement.

\begin{theorem}[Full Goldbach via TI]\label{thm:long-interval}
Fix $\varepsilon>0$ and $A\ge 10$. There exists $B=B(\varepsilon,A)$ and $C=C(\varepsilon,A)>0$ such that for all sufficiently large $x$ the interval
\[
[x,\,x+C(\log x)^A]
\]
contains an even integer $n$ with $R_{2,c}(n)>0$, hence representable as a sum of two primes in the prescribed residue classes modulo $W=(\log x)^B$. In particular, for the unrestricted Goldbach problem (taking the union over classes) every interval $[x,\,x+C(\log x)^A]$ contains a Goldbach even.
\end{theorem}

\begin{remark}
The proof is unconditional and uses only polylogarithmic moduli in the arithmetic localization (Siegel--Walfisz strength). The constants $C,B$ are effective. Our methods also yield density-one statements for the centres and admit natural sharpening via uniformizers in \S\ref{sec:CUNI}.
\end{remark}

\subsection{Method outline}
The argument proceeds in nine steps, each proved in later sections:

\begin{enumerate}[leftmargin=*,itemsep=0.25em]
  \item \textbf{\TFA} (\S\ref{sec:TFA}): build a positive Fej\'er-banked spectral window supported on $\asymp Q^2$ disjoint arcs of width $\asymp 1/H$, with total mass $\asymp Q^2/H$ and second-order notches that enforce alias suppression at small rationals.
  \item \textbf{\AO} (\S\ref{sec:AO-deterministic}): average over a deterministic family of additive characters mod $q\le Q$ to force dispersion on the bank: energy on any $L$ arcs is $\ll (\log X)^C(L/Q^2)\,$ of total.
  \item \textbf{\BG} (\S\ref{sec:BG}): a sheared hyperbola/tube decomposition plus the \SSU\ operator bound on each tube gives a $(H/X)^{1/2}$ short-shift saving for Type-II sums. A two-weight (Sawyer) criterion removes any coefficient flatness assumption.
  \item \textbf{Type-I leakage} (\S\ref{sec:typeI}): under the bank, Type-I energy is $\ll X^{1+o(1)}/(HQ^2)$, negligible after \AO.
  
  \item \textbf{\PSB} (\S\ref{sec:PSB}): Define the nonnegative extractor with the Fejér extractor $J_H\ge 0$, $\supp\widehat J_H\subset[-1/H,1/H]$, and $\widehat J_H(0)=1$.
Bank localization is imposed by the projector $P_{\mathfrak A}$ in frequency; neither $J_H$ nor the analysis kernel is multiplied by the bank.
By Lemma~\ref{lem:extractor}, $S^{\mathrm{ext}}_{n_0}\ge 0$ for all $n_0$, and $S^{\mathrm{ext}}_{n_0}>0$ implies $R_2(n)>0$ for some $|n-n_0|\le H$.

By Lemma~\ref{lem:kernel-stability}, the mean/variance bounds from \ref{sec:BG}-\ref{sec:typeI} also hold for $S^{\mathrm{ext}}_{n_0}$.
  
  \item \textbf{Uniformizers} (\S\ref{sec:CUNI}): sampling/LS + Bernstein and a banked projection convert band-limited mean positivity into uniform every-window means.
  \item \textbf{AP localization \& prime-power disposal} (\S\ref{sec:AP}): log-free Siegel--Walfisz at modulus $W=(\log X)^B$ and elimination of $p^k,\,k\ge2$ with $o(H/\log^2 X)$ loss.
  \item \textbf{Assembly} (\S\ref{sec:assembly}): combine the above to prove Theorem~\ref{thm:long-interval}.
\end{enumerate}

% ============================
% Methods: LLM use and authorship disclosure
% ============================
\subsection{Use of a large language model (LLM) and witness disclosure}\label{sec:llm-methods}

\paragraph{Warning and scope.}
Substantial portions of this manuscript (including preliminary drafts of statements, proofs, and expository text) were generated with the assistance of a large language model (“LLM”). This LLM, MathGPT (by Pulsr), operates with a Python backend that checks for mathematical accuracy.

While every displayed claim that remains in the final version is intended to be mathematically correct, LLMs are known to produce fluent but invalid arguments. Acute skepticism is therefore warranted.
Readers and referees should treat all arguments as provisional until independently verified. Only statements that survive explicit human checking (or external peer review) should be relied upon as certified mathematics.

\paragraph{Provenance and controls.}
Preliminary drafts were little more than wanderings related to the editor's naive intuitions. After a general research plan was created, workflow shifted to a dialectical ("good cop, bad cop") method, as different LLMs were asked to either cooperatively or adversarially critique the manuscript. In mature drafts, the LLM was used as a drafting and refactoring tool under a falsification-first workflow: proposed steps were stress-tested against explicit “falsifiers,” and several originally claimed routes were removed or replaced. 
Where possible, the paper points to a single audit inequality that collects all quantitative inputs; this is included to facilitate independent checking.
Conversation transcripts and intermediate artifacts are archived and can be provided on request to support provenance and reproduction.

\paragraph{Why there is a human \emph{editor} and not a human \emph{author}.}
A conventional author is expected to (i) originate the key arguments, (ii) comprehend and vouch for each step, and (iii) assume responsibility for mathematical correctness. 
In this project, the human contributor acted as philosopher-editor–curator: specifying goals, setting falsifiers, selecting among LLM-proposed routes, removing defective steps, and repeated LLM-based checking, but \emph{not} independently verifying every derivation or supplying the full mathematical substance personally. 
Out of intellectual honesty, the human contributor is therefore recorded as a \emph{witness/editor} to the process rather than an author asserting personal authorship or comprehensive verification. 
This choice is also motivated by accountability: an LLM cannot bear scholarly responsibility, and the witness declines to claim it on the LLM’s behalf.

\paragraph{Citation and responsibility.}
If this manuscript is cited, it should be clear in the citation text that an LLM was used as a drafting tool and that the human participant is listed as witness/editor. 
All responsibility for correctness lies with the mathematical community’s verification process; any errors discovered should be documented publicly so that the record is clear and future versions can be corrected. Any credit of the argument belongs to the mathematical community on the whole, since the LLM is trained on their labor.
              % introduction + statement of main theorems
%========================================================
%  Notation and standing conventions
%========================================================
% Notation section
\providecommand{\Z}{\mathbb{Z}}
\providecommand{\C}{\mathbb{C}}
\providecommand{\wh}{\widehat}

\section{Notation and standing conventions}\label{sec:notation}

\begin{description}[style=nextline]

\item[Epistemology and proof.]
Terms like "unconditional proof" and "proof" will be used throughout the paper draft to describe the state of deductive closure absent human verification. As such, strictly speaking, they are only provisionally proven.

\item[Transforms and convolution on $\mathbb{Z}_M$.]
For $f:\{0,\dots,M-1\}\to\mathbb{C}$,
\[
\wh f(k)=\sum_{n=0}^{M-1} f(n)\,e^{-2\pi i kn/M},\qquad
f(n)=\frac{1}{M}\sum_{k=0}^{M-1}\wh f(k)\,e^{2\pi i kn/M}.
\]
Convolution: $(f*g)(n)=\sum_m f(m)\,g(n-m)$. Parseval:
$\sum_n |f(n)|^2=\frac{1}{M}\sum_k |\wh f(k)|^2$.

\item[Asymptotic symbols.]
We write $A\ll B$ or $A=O(B)$ if $|A|\le C\,B$ for some constant $C>0$ depending only on the fixed parameters indicated (by subscripts when needed), e.g.\ $A\ll_{\varepsilon} B$.  
We write $A\asymp B$ if $A\ll B$ and $B\ll A$.  
We write $A=X^{o(1)}$ to mean: for every $\eta>0$, $|A|\le X^{\eta}$ for all sufficiently large $X$ (with the same fixed parameters).  
All implied constants may depend on $(\varepsilon,A,\gamma,C_W)$ unless otherwise stated.

\item[Global scale and dyadics.]
$X\to\infty$ is the main asymptotic parameter.  
Dyadic variables satisfy $X^{\varepsilon}\le D,N\le X^{1-\varepsilon}$ with $DN\asymp X$.  
We use the short--shift (smoothing) scale
\[
H=(\log X)^A,\qquad A\ge 10,
\]
and the bank parameter $Q=H^{\gamma}$ with fixed $0<\gamma<\tfrac12$.

\item[Arithmetic functions.]
$\mathbf{1}_E$ is the indicator of an event/set $E$.  
$\mu$ is the M\"obius function, $\varphi$ Euler’s totient.  
$\tau(n)$ is the divisor function.  
$\Lambda(n)$ is von Mangoldt; $\Lambda^{\sharp}(n)=\log p$ if $n=p$ is prime, otherwise $0$.  
We use $p$ for primes and $p^k$ for prime powers ($k\ge 2$).

\item[Exponential notation and Fourier transforms.]
\end{description}

$e(z):=e^{2\pi i z}$.  
On $\R$, for $f\in L^1(\R)$ we set
\[
\widehat f(\xi)=\int_{\R} f(t)\,e(-\xi t)\,dt,\qquad
f(t)=\int_{\R} \widehat f(\xi)\,e(\xi t)\,d\xi.
\]
On $\Z$, for $F:\Z\to\C$ we use the Fourier series on $\mathbb T=\R/\Z$:
\[
\widehat F(\xi)=\sum_{n\in\Z} F(n)\,e(-\xi n),\qquad
F(n)=\int_{\mathbb T}\widehat F(\xi)\,e(\xi n)\,d\xi.
\]
We write $\|\cdot\|_p$ for $L^p$ or $\ell^p$ norms as context dictates.

\paragraph{Bank, projectors, and multipliers.}
Let $\cA\subset\{0,\dots,M-1\}$ be the (alias-suppressed) bank: disjoint arcs of width $\asymp M/H$ around reduced $a/q$ with $q\le Q$.
\begin{itemize}
  \item \emph{Sharp bank projector.} $\PA$ is the orthogonal Fourier projection to $\cA$:
  $\wh{(\PA f)}(k)=\mathbf{1}_{\cA}(k)\,\wh f(k)$. In time, $\PA f=K_{\cA}*f$ with
  \[
  K_{\cA}(n)=\frac{1}{M}\sum_{k\in\cA} e^{2\pi i kn/M}\qquad\text{(generally signed)}.
  \]
  \item \emph{Smoothed bank operator.} Choose $0\le \psi\le 1$ with $\supp\psi\subset\cA$ and $\psi\equiv 1$ on the middle third of each arc. Define $T_\psi$ by $\wh{(T_\psi f)}=\psi\,\wh f$. Then $\|T_\psi\|_{\ell^2\to \ell^2}\le 1$ and $T_\psi f=K_\psi * f$ with
  $K_\psi(n)=\frac{1}{M}\sum_k \psi(k)\,e^{2\pi i kn/M}$ (signed).
\end{itemize}
We never assume time-domain kernels are nonnegative; all pointwise steps use the projector inequalities below.

\textbf{Bank and arcs (TFA$^{+}$, AO$^{+}$).}
$\mathfrak A=\bigcup_{j=1}^{J}\mathfrak a_j\subset\mathbb T$ is the union of disjoint “bank arcs” centered at rationals $a_j/q_j$ with $1\le q_j\le Q$, each of length $\asymp 1/H$.  
$J\asymp Q^2$.  
$\widehat w_{\tau^\ast}$ is the (nonnegative) banked spectral window supported on $\mathfrak A$ with total mass $\asymp Q^2/H$ (TFA$^{+}$).  
When we partition $\mathfrak A$, we write $\psi_{\mathfrak a}\in C_c^\infty(\mathfrak a)$ for a smooth partition of unity.

\paragraph{Local arc detectors (AO).}
For each arc $I_i\subset\cA$, fix a nonnegative bump $\phi_i$ with $\supp\phi_i\subset I_i$ and $\int\phi_i\asymp |I_i|$. Set
\[
\psi_i:=\phi_i^{1/2}\big/\|\phi_i^{1/2}\|_{L^2},\qquad \|\psi_i\|_{2}=1.
\]
Their Gram matrix $G_{ij}=\angles{\psi_i,\psi_j}$ is diagonally dominant with off-diagonals
$|G_{ij}|\ll Q^{-1+\varepsilon}$ (near-orthogonality). The $\psi_i$’s are \emph{frequency} bumps; they are not the global $\psi$ used in $T_\psi$.

\paragraph{Fejér window at bandwidth $1/H$.}
Let $\wh \Phi_H(\xi)$ be a nonnegative Fejér/triangle window supported on $|\xi|\lesssim 1/H$ with total mass $\asymp 1$.
For the dyadic shells $\mathcal D_j=\{\xi:2^{-j-1}/H<|\xi|\le 2^{-j}/H\}$ define
\[
W_j := \sum_{\xi\in \mathcal D_j} |\wh \Phi_H(\xi)|,\qquad
\sum_{j\ge 0}W_j\ll 1,\qquad \sum_{j\ge 0}2^{j/2}W_j\ll \log H.
\]
These are the only properties used in the SSU dyadic split.


\begin{description}
\item[Type–II (bilinear) sums.]
Coefficient sequences $(\alpha_d)$, $(\beta_n)$ are supported on $[D,2D]$, $[N,2N]$, respectively.  
The Type–II aggregate at level $k$ is
\[
A_k \ :=\ \sum_{dn=k}\alpha_d\,\beta_n\,
W\!\Big(\tfrac dD,\tfrac nN\Big).
\]
We write
\[
\textstyle \mathrm{Corr}_H:=\frac{1}{H}\sum_{0<|\Delta|<H}\sum_k |A_{k+\Delta}\,\overline{A_k}\,|
\]
for the short–shift correlation average.

\end{description}
\begin{description}
\item[Sheared coordinates and tubes (BG).]
For a reduced slope $a/q$ we use the unimodular linear map
\[
\Phi_{a/q}:\ (d,n)\mapsto (u,v)=(qd-an,\ ad+qn),
\]
so that $d'n-d n'=(u'v-uv')/(a^2+q^2)$.  
A (sheared) tube $T_\tau$ is a set of lattice points in $[D,2D]\times[N,2N]$ satisfying a slope constraint $|n/d-a/q|\ll H/X$ together with a short–shift constraint $|d'n-d n'|\ll H$.  
We use $S_\tau=\Phi_{a/q}(T_\tau)$ for its image.

\item[Short–shift operator and SSU.]
The short–shift kernel in space is
\[
\mathcal K_H\big((d,n),(d',n')\big):=K_H(d'n-d n')\,W\!\Big(\tfrac dD,\tfrac nN\Big)\,W\!\Big(\tfrac{d'}D,\tfrac{n'}N\Big),
\]
and $(\mathbf T F)(d',n'):=\sum_{d,n}\mathcal K_H((d,n),(d',n'))F(d,n)$.  
“SSU” (slope–shift uncertainty) denotes the key tube estimate
\[
\sum_{(d',n'),(d,n)\in T_\tau} F(d,n)\,K_H(d'n-d n')\,\overline{F(d',n')}
\ \ll\ (\log X)^C\Big(\frac{H}{X}\Big)^{\!1/2}\ \sum_{T_\tau}|F|^2,
\]
uniformly over tubes $T_\tau$.

\item[Circle method, major/minor arcs.]
For $R_{2,c}(n)$ (smoothed Goldbach count in a residue class $c\bmod W$) we write its Fourier transform $\widehat{R_{2,c}}(\xi)$ on $\mathbb T$.  
On each major arc $\mathfrak a_{a/q}\subset\mathfrak A$ we use the standard factorization
\[
\widehat{R_{2,c}}(\xi)\ =\ \mathfrak S_{a/q,c}\ \mathcal V(\xi-a/q)\ +\ \mathrm{Err}_{a/q}(\xi),
\]
where $\mathfrak S_{a/q,c}>0$ is the (localized) singular series and $\mathcal V\ge 0$ is the archimedean profile; $\mathrm{Err}_{a/q}$ is controlled by Type–I leakage, BG, and AO dispersion.  
“Minor arcs’’ means $\mathbb T\setminus \mathfrak A$.

\item[Smoothed local statistic.]
The band-limited local count is
\[
\mathcal S(y)\ :=\ \sum_{n\in\Z} R_{2,c}(n)\,K_H(n-y),\qquad
\mathcal T(y)\ :=\ \frac{\log^2 X}{H}\,\mathcal S(y).
\]
We frequently evaluate $\mathcal T$ on a grid $\mathcal G=\{X+j\Delta\}\cap[X,2X]$ with $\Delta=H/C_0$.

\item[AP localization.]
$W=(\log X)^B$ is the smooth modulus used for arithmetic progression localization (Siegel–Walfisz range).  
Residues $c\bmod W$ are assumed admissible.  
We write averages over classes as $\mathbb E_{a \bmod q}$ and over bank packets (AO) as $\mathbb E_{b\in\mathfrak B_Q}$.

\end{description}

\paragraph{Projection inequalities (used for pointwise steps).}
Writing $R_2=M+E$ one has, for every even $N$,
\[
R_2(N)\ \ge\ (\PA M)(N)\ -\ \|\PA E\|_2\ -\ \|P_{\cA^c}R_2\|_2,
\]
and, with the smoothed multiplier $T_\psi$,
\[
R_2(N)\ \ge\ (T_\psi M)(N)\ -\ \|T_\psi E\|_2\ -\ \|(I-T_\psi)R_2\|_2.
\]
Lower bounds for $(\PA M)$ or $(T_\psi M)$ come from the \emph{bank-summed} singular series; the two $L^2$ terms are controlled by the BG/SSU and Type–I sections.
           % notation & standing conventions

%--------------------------------------------------------
%  Analytic interfaces
%--------------------------------------------------------
%========================================================
% sections/tfa.tex — TFA+: Arithmetic Fejér Bank
%========================================================
\section{Time Frequency Analysis (TFA)}\label{sec:TFA}

The purpose of this section is to record the frequency-space
objects used throughout the paper: the \emph{bank} of arcs at scale $1/H$,
and the associated short--shift kernel $K_H$.  Later arguments (AO, BG, SSU)
only need to cite the properties listed here.

\subsection{Construction of the bank}

Fix parameters $A\ge 10$, $0<\gamma<\tfrac12$, $H=(\log X)^A$, and
$Q=H^\gamma$.  
Let
\[
\mathfrak A\ :=\ \bigcup_{1\le q\le Q}\ \bigcup_{\substack{1\le a\le q\\(a,q)=1}}
\ \mathfrak a_{a/q},
\]
where each $\mathfrak a_{a/q}$ is an interval of length $\asymp 1/H$
centered at $a/q$.  
We choose a smooth cutoff 
$\phi_{a/q}\in C_c^\infty(\mathfrak a_{a/q})$, $0\le \phi_{a/q}\le 1$, with
\[
\int_{\mathfrak a_{a/q}} \phi_{a/q}(\xi)\,d\xi\ \asymp\ \tfrac{1}{H}, 
\qquad 
\|\phi_{a/q}^{(m)}\|_\infty \ \ll\ H^m\ \ (m\le 3).
\]
Finally, define the \emph{bank weight}
\[
\widehat w_{\tau^\ast}(\xi)\ :=\ \sum_{a/q}\phi_{a/q}(\xi).
\]
Thus $\widehat w_{\tau^\ast}$ is supported on $\mathfrak A$, satisfies
$0\le \widehat w_{\tau^\ast}\le 1$, and has mass $\int \widehat w_{\tau^\ast}\asymp Q^2/H$.

\paragraph{Terminology: bank envelope vs detector width.}
We make a fixed distinction.

\emph{Bank envelope.} The family $\{\phi_{a/q}\}_{q\le Q,(a,q)=1}$ forms the coarse envelope: each $\phi_{a/q}$ has height $\asymp 1/H$ on its support. Define the total \emph{bank mass}
\[
\mathsf{Mass}_{\mathrm{bank}}
\ :=\ 
\sum_{q\le Q}\ \sum_{\substack{a\!\!\!\pmod q\\(a,q)=1}}\ \int_{\mathbb T}\phi_{a/q}(\alpha)\,d\alpha.
\]
Since each support has length $\asymp 1/H$ and $\sum_{q\le Q}\varphi(q)=\tfrac{3}{\pi^2}Q^2+O(Q\log Q)$, we have
\begin{equation}\label{eq:bank-mass}
\mathsf{Mass}_{\mathrm{bank}}\ \asymp\ \frac{Q^2}{H}.
\end{equation}

\emph{Detector windows.} The fine detectors $w_{a/q}$ localize at length scale
\begin{equation}\label{eq:detector-width}
\mathrm{width}(w_{a/q})\ \asymp\ \frac{1}{qH},
\end{equation}
centered at $a/q$. This is the scale at which we apply AO dispersion and large-sieve estimates.

\begin{lemma}[Bookkeeping of scales]\label{lem:bank-vs-detector}
(i) For each fixed $\alpha\in\mathbb T$,
\(
\sum_{q\le Q}\sum_{(a,q)=1}\phi_{a/q}(\alpha)\ \ll\ Q^2/H.
\)
(ii) The total detector measure satisfies
\(
\sum_{q\le Q}\sum_{(a,q)=1}\operatorname{length}(\operatorname{supp} w_{a/q})\asymp Q^2/H.
\)
\end{lemma}

\begin{remark}[Reader map]\label{rem:bank-detector-flag}
When we refer to “bank mass $\asymp Q^2/H$,” we mean \eqref{eq:bank-mass} (envelope counting/measure).
When we refer to “detector width $\asymp (qH)^{-1}$,” we mean \eqref{eq:detector-width} (analytic localization scale).
Downstream, we explicitly indicate which layer is being used whenever we switch between them.
\end{remark}


\subsection{The kernel \texorpdfstring{$K_H$}{KH}}

Define the time--side kernel $K_H$ by
\[
K_H(n) := \int_{\mathbb R} \widehat K_H(\xi) e(\xi n)\,d\xi,
\qquad \widehat K_H(\xi) := H\,\phi(H\xi),
\]
where $\phi$ is a fixed, even, nonnegative bump supported on $[-2/H,2/H]$
with $\phi(0)=1$.  

Then $K_H$ enjoys the following properties.

\begin{lemma}
[Signed analysis kernel, bank via projector]\label{ass:signed-pin}
Fix $A\ge 10$ and set $H=(\log X)^A$. 
Let $K_H:\mathbb Z\to\mathbb C$ be a \emph{signed} time--kernel with
\[
\supp\widehat K_H \subset [-1/H,\,1/H],\qquad 
\sum_{n\in\mathbb Z}K_H(n)=\widehat K_H(0)= H\,(1+O(H^{-1})),
\]
and kernel moments
\[
\sum_{n}|K_H(n)|\ \ll\ H,\qquad \sum_{n}|n|\,|K_H(n)|\ \ll\ H\log H.
\]
We never use $K_H\ge 0$. Bank localization is imposed \emph{only} by the frequency projector
\[
P_{\mathfrak A}f := \sum_{\substack{1\le q\le Q\\(a,q)=1}} (w_{a/q}\,\widehat f)^{\vee}\!,\qquad 
Q=H^\gamma,\ 0<\gamma<\tfrac12,
\]
with Fejér--type windows $w_{a/q}\ge0$ supported on arcs of width $\asymp (qH)^{-1}$.
\end{lemma}

\begin{lemma}[Model construction of $K_H$]\label{lem:model-K}
Let $\phi\in C_c^\infty(\mathbb R)$ be even, supported in $[-2,2]$, $\phi\equiv 1$ on $[-1,1]$, and set 
$\widehat K_H(\xi):= H\,\phi(H\xi)$. 
Then $K_H$ satisfies Assumption~\ref{ass:signed-pin}. 
Moreover $\widehat K_H\ge 0$ (so $K_H$ is positive--definite) but $K_H$ may change sign; no pointwise positivity is claimed or used.
\end{lemma}

% ==========================================================
% Section 3.2 — Bank normalization and alias suppression
% (positive bank for main term; balanced bank for TT* tests)
% ==========================================================
\subsection{Bank normalization and alias suppression}\label{subsec:bank-alias}

Let $Q=H^\gamma$ with $0<\gamma<\frac12$. For each reduced $a/q$ ($q\le Q$), let
$\Phi_H$ be a fixed nonnegative Fejér window of bandwidth $1/H$ on $\mathbb T$
and set $w_{a/q}(\alpha):=\Phi_H(\alpha-a/q)$. Define the \emph{positive bank weight}
\begin{equation}\label{eq:wbank-def}
\widehat w_{\mathrm{bank}}(\alpha)\ :=\ \sum_{q\le Q}\ \sum_{\substack{a\pmod q\\(a,q)=1}} w_{a/q}(\alpha),
\qquad\text{so that}\quad \supp \widehat w_{\mathrm{bank}}\subset \cA.
\end{equation}
We also introduce a \emph{balanced alias–suppressed test weight} by subtracting a small multiple
of a global Fejér window $u_H$ of the same bandwidth:
\begin{equation}\label{eq:vbank-def}
\widehat v_{\mathrm{bank}}(\alpha)\ :=\ \widehat w_{\mathrm{bank}}(\alpha)\ -\ \kappa\ u_H(\alpha),
\quad\text{where }\ \kappa=\frac{\int_{\mathbb T}\widehat w_{\mathrm{bank}}}{\int_{\mathbb T} u_H}.
\end{equation}
By construction $\int_{\mathbb T}\widehat v_{\mathrm{bank}}=0$ (mean–zero). We use
$\widehat w_{\mathrm{bank}}\ge 0$ for \emph{main–term lower bounds} and $\widehat v_{\mathrm{bank}}$ for
\emph{TT$^\ast$ tests}, where the mean–zero property removes the constant Fourier mode (“alias suppression $\delta=1$”).
A higher–order version with $\int \alpha\,\widehat v_{\mathrm{bank}}=0$ can be arranged as well, but we do not need it.

\paragraph{Mass capture (positive bank).}
\begin{lemma}[Bank mass and normalization]\label{lem:bank-mass}
Let $M_H:=\int_{\mathbb T}\Phi_H(\alpha)\,d\alpha\asymp H^{-1}$. Then
\[
\int_{\mathbb T}\widehat w_{\mathrm{bank}}(\alpha)\,d\alpha
=\Big(\sum_{q\le Q}\varphi(q)\Big)\,M_H
\asymp \frac{Q^2}{H}.
\]
Consequently, the normalized projector $P_{\cA}$ (sharp cutoff on $\supp \widehat w_{\mathrm{bank}}$) captures a fixed positive fraction of the major–arc mass.
\end{lemma}
\begin{proof}
$\int w_{a/q}=M_H$ for each $(a,q)=1$; there are $\varphi(q)$ such $a$; and
$\sum_{q\le Q}\varphi(q)\asymp Q^2$.
\end{proof}

\paragraph{Alias suppression (balanced bank).}
\begin{lemma}[Mean–zero alias suppression]\label{lem:alias-meanzero}
For any trigonometric polynomial $F(\alpha)=\sum_{m\in\mathbb Z} \widehat F(m)\,e(m\alpha)$,
\[
\int_{\mathbb T}\widehat v_{\mathrm{bank}}(\alpha)\,|F(\alpha)|^2\,d\alpha
=\sum_{m\ne 0} \widehat F(m)\,\overline{\widehat F(m)}\, \widehat v_{\mathrm{bank}}^{\ }(m),
\]
i.e. the $m=0$ (constant) mode drops out. In particular, when $F=S$ is a Dirichlet polynomial $S(\alpha)=\sum_{n\le X}\lambda(n)e(\alpha n)$,
the $n=m$ diagonal contributes $0$ to the bank–averaged energy with $\widehat v_{\mathrm{bank}}$.
\end{lemma}
\begin{proof}
Expand $|F|^2$ in Fourier modes and use $\int \widehat v_{\mathrm{bank}}=0$.
\end{proof}

\begin{proposition}[Hybrid large sieve with alias suppression]\label{prop:hls-alias}
Let $S(\alpha)=\sum_{n\le X}\lambda(n)\,e(\alpha n)$ with $\sum_{n\le X}|\lambda(n)|^2\ll X(\log X)^{C_0}$.
Define
\[
\mathsf{Leak}_I^{\mathrm{bal}}\ :=\ \frac{H}{Q^2}
\sum_{q\le Q}\sum_{\substack{a\pmod q\\(a,q)=1}}
\int \widehat v_{\mathrm{bank}}(\alpha)\,|S(\alpha)|^2\,d\alpha.
\]
Then
\[
\mathsf{Leak}_I^{\mathrm{bal}}\ \ll\ (\log X)^{C}\,\frac{X}{H\,Q^2}.
\]
\end{proposition}
\begin{proof}
Write $\widehat v_{\mathrm{bank}}=\sum_{q\le Q}\sum_{(a,q)=1}w_{a/q}-\kappa\,u_H$ and expand $|S|^2=\sum_{h}\Big(\sum_{n\le X-h}\lambda(n)\overline{\lambda(n+h)}\Big)e(h\alpha)$.
By Lemma~\ref{lem:alias-meanzero}, the $h=0$ mode vanishes in $\int \widehat v_{\mathrm{bank}}|S|^2$.
For $h\ne 0$, integrate $e(h\alpha)$ against each window to get a factor $\ll \min\{H,1/|h|\}$,
sum over $a$ to get a Ramanujan sum $c_q(h)$, and apply Cauchy–Schwarz plus the hybrid large sieve to bound
\[
\sum_{q\le Q}\sum_{(a,q)=1}\int w_{a/q}\,|S|^2
\ \ll\ (X+H Q^2)\sum_{n\le X}|\lambda(n)|^2.
\]
Multiply by $H/Q^2$ and subtract the $\kappa u_H$ part (which contributes the same bound with opposite sign at $h=0$ and is harmless for $h\ne 0$).
Since the $h=0$ piece is absent, only the $X$–term remains; this yields the stated bound.
\end{proof}

\paragraph{Remarks.}
(1) We keep \emph{two} bank weights on purpose: $\widehat w_{\mathrm{bank}}\ge 0$ for the projected main term (used in pointwise lower bounds), and the balanced $\widehat v_{\mathrm{bank}}$ for TT$^\ast$/Type–I tests where the constant mode would otherwise dominate.  
(2) If desired, one can enforce a second vanishing moment (``$\delta=2$'') by replacing $u_H$ with a linear combination $u_H^{(0)}+u_H^{(1)}$ chosen so that $\int \alpha^k\,\widehat v_{\mathrm{bank}}=0$ for $k=0,1$. We do not need $k=1$ in the present argument.

% ===== §3 PATCH: explicit alias suppression (δ=2) lemma + corollary =====

\paragraph{Balanced bank mask.}
Let $W\in C_c^\infty([1/2,2]^2)$ be the dyadic bump used to localize to $[D,2D]\times[N,2N]$. Define the \emph{balanced bank} (double–centered) mask
\begin{align}
\mathcal{B}(d,n)
&:= W\!\left(\frac dD,\frac nN\right)
 - \frac{1}{\#\mathcal{D}}\sum_{d'\in\mathcal{D}} W\!\left(\frac {d'}D,\frac nN\right)
 - \frac{1}{\#\mathcal{N}}\sum_{n'\in\mathcal{N}} W\!\left(\frac dD,\frac {n'}N\right) \nonumber\\
&\quad + \frac{1}{\#\mathcal{D}\,\#\mathcal{N}}\sum_{d'\in\mathcal{D}}\sum_{n'\in\mathcal{N}} W\!\left(\frac {d'}D,\frac {n'}N\right),
\label{eq:balanced-bank}
\end{align}
where $\mathcal{D}=[D,2D]\cap\mathbb{Z}$ and $\mathcal{N}=[N,2N]\cap\mathbb{Z}$.

\begin{lemma}[Alias suppression, $\delta=2$]\label{lem:alias-delta2}
The mask $\mathcal{B}$ satisfies, for all $d\in\mathcal{D}$ and $n\in\mathcal{N}$,
\[
\sum_{n'\in\mathcal{N}} \mathcal{B}(d,n')=0,\qquad
\sum_{d'\in\mathcal{D}} \mathcal{B}(d',n)=0.
\]
Consequently,
\[
\sum_{d,n}\mathcal{B}(d,n)\,g(d)=0,\qquad
\sum_{d,n}\mathcal{B}(d,n)\,h(n)=0
\]
for any functions $g$ of $d$ alone and $h$ of $n$ alone. Moreover, $\mathcal{B}$ is supported where $W$ is supported and $\|\mathcal{B}\|_1\asymp \|W\|_1$ with absolute constants (depending only on bounds for $W$).
\end{lemma}

\begin{proof}
Expanding \eqref{eq:balanced-bank} and summing over $n'$ (resp.\ over $d'$) shows telescoping cancellation by construction; the last term compensates the double subtraction. The orthogonality to $g(d)$ and $h(n)$ follows by linearity. Support and $L^1$ comparability are immediate from boundedness of $W$ on its compact support.
\end{proof}

\begin{corollary}[Constant aliases are killed on the bank]\label{cor:alias-kills-const}
Let $F(d,n)=g(d)+h(n)$ (in particular, $F\equiv \mathrm{const}$). For the balanced mask $B$ in \ref{lem:bank-mass} and any array $A$ one has
\[
\sum_{d,n} B(d,n)\,F(d,n)\,A(d,n)=0.
\]
In particular, in any $TT^\ast$ average where the bank enters only via $B$, the row–constant, column–constant, and constant components make no contribution. 

Equivalently, in any TT$^\ast$ average in which the bank enters only through $\mathcal{B}$ (or a positive multiple thereof), the row-constant, column-constant, and constant components of the input make no contribution. In particular, in the Type–I average the “constant alias mode” vanishes and the dispersion ledger starts at the oscillatory $Q^{-2}$ scale.
\end{corollary}

\begin{proof}
By Lemma~\ref{lem:alias-delta2}, $\sum_n \mathcal{B}(d,n)=0$ and $\sum_d \mathcal{B}(d,n)=0$, hence the inner sum against $g(d)$ (resp.\ $h(n)$) vanishes for every fixed $d$ (resp.\ $n$); additivity yields the claim. The final statement is the advertised use in §7 and §11.
\end{proof}

\begin{remark}[“$\delta=2$” bookkeeping]
The two independent cancellations (row and column) are what we count as $\delta=2$. This is the precise input used in §7 (Type–I) and in the TT$^\ast$ ledger of §11 whenever we say “alias suppression with $\delta=2$ removes the constant mode”.
\end{remark}
% ===== end §3 PATCH =====


\subsection{Usage}

In all later sections, the notation $K_H$ always refers to the kernel
constructed above.  Whenever we write a smoothed statistic
\[
\mathcal S(y)\ :=\ \sum_n R_{2,c}(n)\,K_H(n-y),
\]
only the bandwidth 1/H and the moment bounds $K_H$ are used. Pointwise nonnegativity is provided by the extractor $J_H$. 
Cross--arc estimates always exploit (v).
\qedhere

\subsection{Arithmetic Fej\'er bank}\label{AFB}

Recall $H=(\log X)^A$ and $Q=H^\gamma$ with $0<\gamma<\tfrac12$. We construct a positive spectral window $\wh w_{\tau^\ast}$ supported on a union of $\asymp Q^2$ disjoint arcs of width $\asymp 1/H$, centred at rationals $a/q$ with $1\le q\le Q$, $(a,q)=1$, and prove its key properties: support, mass, positivity, and alias suppression.

\subsection{Construction}
Let $M:=\lfloor c_H H\rfloor$ for a small fixed $c_H\in(0,1/10)$. The (discrete) Fej\'er coefficients
\[
F_M(n):=\frac{1}{M+1}\Big(1-\frac{|n|}{M+1}\Big)_+,\qquad n\in\mathbb{Z},
\]
form a nonnegative, even, compactly supported sequence with $\sum_{n}F_M(n)=1$. Its $2\pi$-periodic Fourier transform is the Fej\'er kernel
\[
\wh F_M(\theta)=\frac{1}{M+1}\bigg(\frac{\sin\big((M+1)\theta/2\big)}{\sin(\theta/2)}\bigg)^2\ge 0,\qquad \theta\in\bbT.
\]
For each reduced fraction $a/q$ ($1\le q\le Q$) let $\theta_{a/q}$ be its representative on $\bbT$, and define a smooth nonnegative bump $\psi_{a/q}\in C_c^\infty(\bbT)$ supported on the arc
\[
\mathfrak a_{a/q}:=\{\theta\in\bbT:\ |\theta-\theta_{a/q}|\le c_0/H\},
\]
equal to $1$ on the middle third, with $\|\psi_{a/q}^{(m)}\|_\infty\ll H^m$ for $m\le 3$ and with pairwise disjoint supports as $a/q$ varies. Set
\[
\wh w_{\tau^\ast}(\theta):=\frac{1}{Z}\sum_{\substack{1\le q\le Q\\ (a,q)=1}}\ \psi_{a/q}(\theta)\,\wh F_M(\theta-\theta_{a/q}),
\]
with normalizing constant $Z\asymp \sum_{q\le Q}\varphi(q)\cdot H^{-1}\asymp Q^2/H$. Let $w_{\tau^\ast}$ be the (real, even) inverse transform.

\begin{lemma}[Separation of roles: balanced vs.\ positive]\label{lem:sep-roles}
Let $\widehat w_{\mathrm{bank}}\ge0$ be the positive Fejér bank of \S3.3 and let 
$\widehat v_{\mathrm{bank}}=\widehat w_{\mathrm{bank}}-\kappa\,u_H$ be its balanced version with
$\int_{\mathbb T}\widehat v_{\mathrm{bank}}=0$.
Then:
\begin{enumerate}
\item[(i)] \textbf{Positivity is never bank-masked.} All lower bounds for the major term use the positive projector $P_A$ (sharp cutoff on $\mathrm{supp}\,\widehat w_{\mathrm{bank}}$) and the nonnegative extractor $J_H$, never $v_{\mathrm{bank}}$.
\item[(ii)] \textbf{Balance is used only inside TT$^\ast$.} All large-sieve/TT$^\ast$ estimates insert $\widehat v_{\mathrm{bank}}$ (not $\widehat w_{\mathrm{bank}}$); the mean-zero eliminates the $t=0$ diagonal, and the bank geometry eliminates 1-D aliases (``$\delta=2$''), before AO dispersion.
\end{enumerate}
\end{lemma}

\subsection{Properties}
\begin{proposition}[Support, mass, positivity]\label{prop:tfa-basic}
With the above construction:
\begin{enumerate}[label=(\alph*),leftmargin=2em]
  \item \textbf{Support.} $\operatorname{supp}\wh w_{\tau^\ast}\subset \displaystyle\bigcup_{q\le Q}\ \bigcup_{(a,q)=1}\ \mathfrak a_{a/q}$, a union of $J\asymp Q^2$ disjoint arcs, each of length $\asymp 1/H$.
  \item \textbf{Mass.} $\displaystyle \int_{\bbT}\wh w_{\tau^\ast}(\theta)\,d\theta \asymp \frac{Q^2}{H}$.
  \item \textbf{Positivity.} $\wh w_{\tau^\ast}\ge 0$, hence $w_{\tau^\ast}$ is positive-definite; in particular $\sum_{m,n}c_m\ol c_n\,w_{\tau^\ast}(m-n)\ge 0$ for all finitely supported $(c_n)$.
\end{enumerate}
\end{proposition}

\begin{proof}
(a) and (c) are immediate from the definition and $\wh F_M\ge 0$. For (b), each bump $\psi_{a/q}$ has integral $\asymp 1/H$, and there are $\asymp \sum_{q\le Q}\varphi(q)\asymp Q^2$ arcs, yielding total mass $\asymp Q^2/H$ after the normalization.
\end{proof}

\begin{proposition}[Alias suppression, $\delta=2$]\label{prop:tfa-suppress}
Let $S(\theta)=\sum_{n} C_n\,\ee{n\theta}$ be any trigonometric polynomial (or $L^2$ function on $\bbT$). Then
\[
\int_{\bbT} |S(\theta)|^2\,\wh w_{\tau^\ast}(\theta)\,d\theta
\ \le\ \frac{C}{H}\sum_{n}|C_n|^2
\]
and for each fixed rational $\theta_0=a_0/q_0$ with $q_0\le Q$,
\[
\frac{1}{J}\sum_{\mathfrak a}\int_{\mathfrak a} |S(\theta+\theta_0)|^2\,\wh F_M(\theta)\,d\theta
\ \ll\ \frac{1}{H Q^2}\sum_{n}|C_n|^2,
\]
i.e. averaging over the bank arcs enforces a factor $Q^{-2}$ dispersion on any single alias. All implied constants depend at most polylogarithmically on $X$.
\end{proposition}

\begin{proof}[Proof]
For the first inequality, expand $|S|^2$ and integrate termwise: the Fej\'er factor makes $\wh w_{\tau^\ast}$ a positive mollifier of width $1/H$, giving mass $\asymp Q^2/H$ and an average per arc $\asymp 1/H$; use Cauchy--Schwarz and disjointness of arcs.

For the second, fix $\theta_0=a_0/q_0$. On each $\mathfrak a=\mathfrak a_{a/q}$, write the local variable $\theta=\theta_{a/q}+\eta$ with $|\eta|\ll 1/H$ and orthogonally average the translates $S(\cdot+\theta_0)$ against $\wh F_M(\eta)$. The bank has $J\asymp Q^2$ disjoint arcs; modulo $1$ the phases $e(n\theta_{a/q})$ are separated at scale $\gg 1/Q^2$. A Montgomery--Vaughan large-sieve bound at spacing $\gg 1/Q^2$ yields the extra $Q^{-2}$ after summing over arcs.
\end{proof}

\begin{lemma}[Near–orthogonality of arc detectors]\label{lem:AO-gram-full}
Let $\cA=\bigcup_i I_i$ be a union of disjoint arcs with $|I_i|\asymp H^{-1}$ and $\dist(I_i,I_j)\gg H^{-1}$ for $i\neq j$.
Choose $\phi_i\in C_c^\infty(\mathbb T)$ with $\mathbf 1_{I_i^\circ}\le \phi_i\le \mathbf 1_{I_i^{+}}$, where $I_i^\circ$ is the middle third of $I_i$ and $I_i^{+}$ is a fixed $(1+\tfrac1{100})$ enlargement still disjoint from all other $I_j$.
Set $\psi_i:=\phi_i^{1/2}/\|\phi_i^{1/2}\|_{2}$. Then
\[
\langle \psi_i,\psi_i\rangle=1,\qquad \langle \psi_i,\psi_j\rangle=0\ (i\ne j).
\]
\end{lemma}
\begin{proof}
$\supp\phi_i^{1/2}\subset I_i^{+}$ and the $I_i^{+}$ are pairwise disjoint, hence $\phi_i^{1/2}\phi_j^{1/2}\equiv 0$ for $i\neq j$; normalization gives $\langle\psi_i,\psi_i\rangle=1$.
\end{proof}


\subsection{A smooth partition and per-arc envelopes}
We shall frequently localize to a single bank arc using a smooth partition of unity. Choose nonnegative $\{\psi_{\mathfrak a}\}_{\mathfrak a}$, $\psi_{\mathfrak a}\in C_c^\infty(\mathfrak a)$,
\[
0\le \psi_{\mathfrak a}\le 1,\qquad \sum_{\mathfrak a}\psi_{\mathfrak a}=\wh w_{\tau^\ast},\qquad
\|\psi_{\mathfrak a}^{(m)}\|_\infty\ll H^m\ (m\le 3).
\]
For any frequency-side field $S$ we write $S_{\mathfrak a}:=S\cdot \psi_{\mathfrak a}^{1/2}$; then
\[
\int |S|^2\,\wh w_{\tau^\ast} = \sum_{\mathfrak a}\int |S_{\mathfrak a}|^2,
\]
and each $S_{\mathfrak a}$ is band-limited to width $\asymp 1/H$ around the arc centre. This reduction is used throughout \S\ref{sec:AO-deterministic}, \S\ref{sec:BG}, and \S\ref{sec:PSB}.

\begin{lemma}[Mask separation: where balance vs positivity is used]\label{lem:mask-separation}
Let $\widehat{\mathcal B}=\sum_{a/q\in\mathfrak A(Q)} w_{a/q}$ be the bank multiplier with each $w_{a/q}$ smooth, supported on an arc of width $\asymp 1/H$, and balanced $\int_{\mathbb T} w_{a/q}=0$. Let $J_H\ge 0$ be the Fejér extractor with $\widehat{J_H}$ supported in $|\xi|\le 1/H$ and $\widehat{J_H}(0)=1$.
Then:
\begin{enumerate}
\item The **balanced** mask $\widehat{\mathcal B}$ is used only *inside* TT$^\ast$ on the Type–I limb to kill the constant mode and first aliases (via $\delta=2$), and in AO occupancy counting; positivity is not needed there.
\item The **positive** extractor $J_H$ is never multiplied by the bank (we localize by the projector $P_A$ or a smooth taper $T_\psi$), and is used only to transfer mean positivity to pointwise positivity and to define $S_{n_0}$.
\end{enumerate}
\end{lemma}

\begin{proof}
(1) Balancedness $\int w_{a/q}=0$ forces the constant Fourier mode to vanish in the TT$^\ast$ Type–I quadratic form and, with $\delta=2$, removes row/column aliases. (2) Since $J_H\ge0$ and $\widehat{J_H}$ sits in a single band around $0$, convolving with $J_H$ preserves positivity and allows Bernstein–type Lipschitz control; multiplying by the (disconnected) bank would destroy both features, hence we avoid it by using $P_A$ (or $T_\psi$) for localization. 
\end{proof}
                % TFA$^{+}$: bank/window & alias suppression
% ==========================================================
% Section 4 — Arithmetic Orthogonality (Deterministic AO; no packets)
% ==========================================================
\section{Arithmetic Orthogonality (AO)}\label{sec:AO-deterministic}

Throughout this section the bank $\cA\subset\mathbb T$ is the union of $m$ disjoint arcs
\[
\cA=\bigcup_{i=1}^{m} I_i,\qquad m\asymp Q^2,\qquad |I_i|\asymp \frac{1}{H},\qquad \dist(I_i,I_j)\gg \frac{1}{H}\ (i\neq j),
\]
with $Q=H^\gamma$ and $0<\gamma<\tfrac12$ as fixed elsewhere. All implied constants may depend on the fixed bank taper but are independent of $X,H,Q$.

\subsection{Local detectors and near–orthogonality}
For each arc $I_i$ choose a nonnegative $C^\infty$ bump $\phi_i$ with $\supp\phi_i\subset I_i$ and
\[
\int_{I_i}\phi_i(\alpha)\,d\alpha \asymp |I_i|,\qquad 0\le \phi_i\le 1.
\]
Define the $L^2$–normalized detectors
\[
\psi_i(\alpha):=\frac{\phi_i(\alpha)^{1/2}}{\|\phi_i^{1/2}\|_{L^2(\mathbb T)}},\qquad \|\psi_i\|_{L^2(\mathbb T)}=1.
\]

\begin{lemma}[Near–orthogonality of arc detectors]\label{lem:AO-gram}
Let $G=(\langle\psi_i,\psi_j\rangle)_{1\le i,j\le m}$ be the Gram matrix. Then
\[
\langle\psi_i,\psi_i\rangle=1,\qquad |\langle\psi_i,\psi_j\rangle|\ \ll\ Q^{-1+\varepsilon}\quad (i\neq j),
\]
for any fixed $\varepsilon>0$. The implied constant depends only on the choice of smooth partition and $\varepsilon$.
\end{lemma}

\begin{proof}[Proof.]
Disjoint supports give zero except for nearest neighbours; for those, smooth overlap and the rational separation of arc centres (with denominators $\le Q$) produce the stated decay after one integration by parts. This is standard for smoothly partitioned major arcs. See \ref{lem:AO-gram-full}
\end{proof}

\begin{table}[h]
\centering
\begin{tabular}{lcl}
guard band &:& $g\asymp (qH)^{-1}$\\
arc width  &:& $|I_{a/q}|\asymp (qH)^{-1}$ (before $1+10^{-2}$ enlargement)\\
enlargement &:& factor $1+\eta$, $\eta=10^{-2}$ (still disjoint)\\
overlap count &:& at most $O(1)$ per scale, total $m\asymp Q^2$ arcs\\
bank taper &:& Fejér of bandwidth $1/H$\\
Gram off–diag &:& $\ll Q^{-1+\varepsilon}$ (Schur/Gershgorin with guard band)
\end{tabular}
\caption{Geometric inputs to AO.}
\end{table}

\begin{proposition}[AO with explicit tail]\label{prop:AO-const}
Let $\{\psi_i\}_{i=1}^m$ be the normalized arc probes built from the table above. 
Then for any $f$,
\[
\sum_{i\in\mathcal L} |\langle f,\psi_i\rangle|^2
\ \le\
\Big(\frac{L}{m}+\frac{C_{\mathrm{AO}}}{Q}\Big)
\sum_{i=1}^{m} |\langle f,\psi_i\rangle|^2,
\qquad 1\le L\le m,
\]
with $C_{\mathrm{AO}}=C(\eta,\text{guard},\text{taper})=O(1)$.
\end{proposition}


\begin{theorem}[Deterministic AO dispersion: no hot spots]\label{thm:AO-dispersion-nohot}
Let $\mathfrak A=\bigcup_{i=1}^{m} I_i$ be a bank partition with $m\asymp Q^2$, where each $I_i$ has length $\asymp H^{-1}$ and the enlarged arcs $I_i^+$ satisfy 
$|I_i^+\cap I_j^+|\ll (QH)^{-1}$ for $i\neq j$.
Choose smooth detectors $\varphi_i\in C_c^\infty(\mathbb T)$ with 
$1_{I_i^\circ}\le \varphi_i\le 1_{I_i^+}$, and set the normalized probes
\[
\psi_i:=\frac{\varphi_i}{\|\varphi_i\|_{2}}\,,\qquad 
\|\varphi_i\|_2^2\asymp |I_i|\asymp H^{-1}.
\]
Assume the near–orthogonality
\begin{equation}\label{eq:AO-near-orth}
\langle \psi_i,\psi_i\rangle=1,\qquad 
|\langle \psi_i,\psi_j\rangle|\ \ll\ Q^{-1+\varepsilon}\quad (i\ne j),
\end{equation}
which holds by Lemma~\ref{lem:AO-near}. 
For any $F\in L^2(\mathbb T)$, define the bank energy via the probe coefficients
\[
E_{\mathfrak A}(F):=\sum_{i=1}^{m} \big|\langle F,\psi_i\rangle\big|^2,
\qquad 
E_{L}(F):=\max_{\substack{S\subset\{1,\dots,m\}\\ |S|=L}}\ \sum_{i\in S} \big|\langle F,\psi_i\rangle\big|^2 .
\]
Then for every $1\le L\le m$,
\begin{equation}
\frac{E_{L}(F)}{E_{\mathfrak A}(F)}
\ \le\ \frac{L}{m}\ +\ C\,Q^{-1+\varepsilon}\,,
\end{equation}
with an absolute $C$ (depending only on the constants in \eqref{eq:AO-near-orth}). 
The bound is uniform in $F$ and the constants match the ones used in the §11 ledger.
\end{theorem}

\begin{proof}
Write $a_i:=\langle F,\psi_i\rangle$ and $a=(a_i)_{i\le m}\in\mathbb C^m$. 
Let $G=(\langle \psi_i,\psi_j\rangle)_{i,j}$ be the Gram matrix. By \eqref{eq:AO-near-orth},
\[
G = I + R,\qquad \|R\|_{\ell^2\to\ell^2}\ \le\ C_0\,Q^{-1+\varepsilon}
\]
(Gershgorin or Schur test). The analysis operator $T:L^2(\mathbb T)\to\mathbb C^m$, $TF=a$, satisfies
\[
E_{\mathfrak A}(F)=\|TF\|_2^2,\qquad 
\|T\|^2=\|T^\ast T\|=\|G\|\le 1+C_0Q^{-1+\varepsilon},\quad 
\|T^\ast\|^2=\|TT^\ast\|=\|G\|.
\]
Let $P_S:\mathbb C^m\to\mathbb C^m$ be the coordinate projector onto a set $S\subset\{1,\dots,m\}$, $|S|=L$. Then
\[
\sum_{i\in S}|a_i|^2=\|P_S a\|_2^2=\langle P_S a,a\rangle=\langle P_S TF,TF\rangle
=\langle T^\ast P_S T F, F\rangle.
\]
Average this identity over all $S$ of size $L$ using $\binom{m}{L}^{-1}\sum_{|S|=L}P_S=\frac{L}{m}I$:
\[
\frac{1}{\binom{m}{L}}\sum_{|S|=L}\ \sum_{i\in S}|a_i|^2
=\Big\langle T^\ast\Big(\frac{L}{m}I\Big)T F, F\Big\rangle
=\frac{L}{m}\,\|TF\|_2^2
=\frac{L}{m}\,E_{\mathfrak A}(F).
\]
Hence the **average** over $S$ equals $(L/m)E_{\mathfrak A}(F)$, so for the **maximum** we need only control the deviation induced by the non-orthogonality of $\{\psi_i\}$.
For any fixed $S$, decompose
\[
\sum_{i\in S}|a_i|^2
=\sum_{i\in S}\langle a_i e_i, a\rangle
=\langle (P_S - \tfrac{L}{m}I)a,a\rangle + \frac{L}{m}\|a\|_2^2,
\]
where $e_i$ is the $i$-th basis vector. Since $a=TF$ and $G=T T^\ast$, we have
\[
\langle (P_S - \tfrac{L}{m}I)a,a\rangle
=\langle (P_S - \tfrac{L}{m}I)TF,TF\rangle
=\langle T^\ast(P_S - \tfrac{L}{m}I)T F,F\rangle.
\]
Note that $T^\ast(P_S - \tfrac{L}{m}I)T = T^\ast(P_S - \tfrac{L}{m}I)T - T^\ast(P_S - \tfrac{L}{m}I)T\,(I-G) + T^\ast(P_S - \tfrac{L}{m}I)T\,(I-G)$, so by inserting $G=I+R$ and using $\|R\|\ll Q^{-1+\varepsilon}$,
\[
\big|\langle T^\ast(P_S - \tfrac{L}{m}I)T F,F\rangle\big|
\ \le\ \|T^\ast\|\,\|P_S - \tfrac{L}{m}I\|\,\|T\|\,\|R\|\,\|F\|_2^2
\ \ll\ Q^{-1+\varepsilon}\ \|TF\|_2^2,
\]
since $\|P_S - \tfrac{L}{m}I\|\le 1$ and $\|T\|,\|T^\ast\|\le (1+C_0Q^{-1+\varepsilon})^{1/2}$. 
Combining with the average identity gives, for each $S$,
\[
\sum_{i\in S}|a_i|^2\ \le\ \frac{L}{m}\,\|a\|_2^2\ +\ C\,Q^{-1+\varepsilon}\,\|a\|_2^2,
\]
i.e. \eqref{thm:AO-dispersion}. This bound is uniform in $F$ and in $S$.
\end{proof}

% ===== BEGIN PATCH: AO dispersion details =====
\begin{lemma}[Arc overlap count]\label{lem:AO-overlap}
Let $\{\phi_i\}_{i=1}^m$ be the arc detectors supported on intervals $I_i$ of length $\asymp 1/H$ around Farey centers $a_i/q_i$ with $q_i\le Q$. Then for each fixed $i$,
\[
\#\{j:\ \phi_i\phi_j\not\equiv 0\}\ \ll\ Q,
\]
and, more precisely, the Gram off–diagonal satisfies
\[
\sum_{j\ne i} |\langle \psi_i,\psi_j\rangle|\ \ll\ Q^{-1+\varepsilon},
\]
for normalized $\psi_i:=\phi_i/\|\phi_i\|_2$.
\end{lemma}

\begin{proof}
If $I_i$ and $I_j$ overlap, their centers satisfy $|a_i/q_i-a_j/q_j|\ll 1/H$. Farey spacing gives $|a_i/q_i-a_j/q_j|\ge 1/(q_iq_j)$, hence overlap forces $q_j\gg H/q_i$. Since $q_j\le Q$, there are $O(Q)$ such neighbors. The $L^2$–normalized inner products are $\ll |I_i\cap I_j|/\sqrt{|I_i||I_j|}\ll 1/(q_iq_j)$, summing to $O(Q^{-1+\varepsilon})$ by divisor bounds.
\end{proof}

\begin{theorem}[Deterministic AO dispersion]\label{thm:AO-dispersion}
Let $G=(\langle\psi_i,\psi_j\rangle)$ be the Gram matrix of the normalized detectors on the bank. Then
\[
\|G-I\|_{2\to2}\ \ll\ Q^{-1+\varepsilon},
\qquad
\sum_{i\in\mathcal L} |\langle f,\psi_i\rangle|^2\ \ll\ \Bigl(1+C\,L\,Q^{-1+\varepsilon}\Bigr)\,\|f\|_2^2,
\]
for all subfamilies $\mathcal L$ with $|\mathcal L|=L$.
\end{theorem}

\begin{proof}
By Lemma~\ref{lem:AO-overlap}, each row of $R:=G-I$ has $\ell^1$–norm $\ll Q^{-1+\varepsilon}$; Schur’s test yields $\|R\|_{2\to2}\ll Q^{-1+\varepsilon}$. The $L$–subfamily bound follows from Gershgorin applied to the principal submatrix $G_{\mathcal L}$.
\end{proof}
% ===== END PATCH =====


\begin{corollary}[Occupancy form]\label{cor:AO-occupancy}
For every $1\le L\le m$ and every $F\in L^2(\mathbb T)$,
\[
\max_{|S|=L}\ \frac{\sum_{i\in S}\big|\langle F,\psi_i\rangle\big|^2}{\sum_{i=1}^{m}\big|\langle F,\psi_i\rangle\big|^2}
\ \le\ \frac{L}{m}\ +\ C\,Q^{-1+\varepsilon}.
\]
\end{corollary}

\subsection{Deterministic AO bound}
Let $F\in L^2(\mathbb T)$ be supported on $\cA$. For a subfamily $\mathcal L\subset\{1,\dots,m\}$ with $|\mathcal L|=L$ write the energy on $\mathcal L$ as
\[
\mathsf E_{\mathcal L}(F):=\sum_{i\in\mathcal L}\int_{I_i} |F(\alpha)|^2\,d\alpha,\qquad 
\mathsf E_{\cA}(F):=\int_{\cA} |F(\alpha)|^2\,d\alpha.
\]

\begin{lemma}[AO near-orthogonality for smooth detectors]\label{lem:AO-near-orth}
Let $\mathcal A=\{a/q:1\le q\le Q,\ (a,q)=1\}\subset[0,1)$ with $Q=H^\gamma$ and $0<\gamma<\tfrac12$.
For each $\alpha_i\in\mathcal A$ define the $L^2$–unit Gaussian detector
\[
\psi_i(\alpha)\ :=\ (\pi\sigma^2)^{-1/4}\,\exp\!\Big(-\frac{(\alpha-\alpha_i)^2}{2\sigma^2}\Big),
\qquad \sigma=\frac{c}{H},\quad 0<c\le \tfrac14.
\]
Then for $i\neq j$,
\[
|\langle \psi_i,\psi_j\rangle|
=\exp\!\Big(-\frac{(\alpha_i-\alpha_j)^2}{4\sigma^2}\Big)
\ \le\ \exp\!\Big(-\frac{H^2}{4c^2Q^4}\Big)
=\exp\!\big(-H^{\,2-4\gamma}/(4c^2)\big).
\]
Consequently, for every $\varepsilon>0$ there exists $H_0(\varepsilon,c,\gamma)$ such that for all $H\ge H_0$,
\[
|\langle \psi_i,\psi_j\rangle|\ \le\ Q^{-1+\varepsilon},
\]
and hence the Gram off–diagonals obey $|\langle\psi_i,\psi_j\rangle|\ll Q^{-1+\varepsilon}$.
\end{lemma}

\begin{proof}
Farey spacing gives $|\alpha_i-\alpha_j|\gg 1/Q^2$ for $i\neq j$. Insert this in the Gaussian overlap identity for equal–variance, unit–norm Gaussians,
$|\langle \psi_i,\psi_j\rangle|=\exp(-\Delta^2/(4\sigma^2))$ with $\Delta=|\alpha_i-\alpha_j|$.
Thus $|\langle \psi_i,\psi_j\rangle|\le \exp(-H^{2-4\gamma}/(4c^2))$.
Because $\gamma<\tfrac12$, the exponent $\to+\infty$, while $Q^{-1+\varepsilon}=H^{-\gamma(1-\varepsilon)}$ decays only polynomially; hence the claimed bound for all $H\ge H_0(\varepsilon,c,\gamma)$.
\end{proof}

\begin{proposition}[AO dispersion: bank--projected Gram bound]\label{prop:AO-projected}
Let $A=\bigcup_{i=1}^m I_i\subset\T$ be the bank, with $m\asymp Q^2$, $|I_i|\asymp H^{-1}$ and pairwise gaps $\gg H^{-1}$, and let $\{\psi_i\}_{i=1}^m$ be the $L^2$–unit arc detectors from §4.1. 
Write $P_A$ for the Fourier projector to $A$, and for $F\in L^2(\T)$ set 
\[
y_i:=\langle P_A F,\psi_i\rangle,\qquad S_L(F):=\sum_{i\in\mathcal L}|y_i|^2,\qquad S_{\rm tot}(F):=\sum_{i=1}^m |y_i|^2,
\]
for any subfamily $\mathcal L\subset\{1,\dots,m\}$ with $|\mathcal L|=L$. 
Then, for every such $\mathcal L$,
\begin{equation}\label{eq:AO-deterministic}
S_L(F)\ \le\ \Bigl(1+C\,L\,Q^{-1+\varepsilon}\Bigr)\,S_{\rm tot}(F),
\end{equation}
where $C$ depends only on the fixed smooth partition and on~$\varepsilon>0$.
\end{proposition}

\begin{proof}
Let $G=(\langle\psi_i,\psi_j\rangle)_{1\le i,j\le m}$ be the Gram matrix and write $P_AF=\sum_{j} c_j\,\psi_j$ in the (redundant) frame sense; then $y=Gc$ with $y=(y_i)_{i\le m}$ and $c=(c_j)_{j\le m}$. 
For any $\mathcal L$, let $G_{\mathcal L}$ be the $L\times L$ principal submatrix of $G$. By Lemma~4.1, $G=I+E$ with $E_{ij}=O(Q^{-1+\varepsilon})$ for $i\ne j$, so Gershgorin gives $\|G_{\mathcal L}\|_{2\to2}\le 1+C_1LQ^{-1+\varepsilon}$. Hence
\[
S_L(F)\ =\ \|G_{\mathcal L}c\|_2^2\ \le\ \|G_{\mathcal L}\|_{2\to2}^2\,\|c\|_2^2\ \le\ \bigl(1+C_2LQ^{-1+\varepsilon}\bigr)\|c\|_2^2.
\]
On the other hand,
\[
S_{\rm tot}(F)\ =\ \|Gc\|_2^2\ \ge\ \|c\|_2^2 - \|E\|_{2\to2}\,\|c\|_2^2\ \ge\ \bigl(1-C_3Q^{-1+\varepsilon}\bigr)\|c\|_2^2,
\]
so \eqref{eq:AO-deterministic} follows after absorbing the harmless $(1-C_3Q^{-1+\varepsilon})^{-1}$ factor into the constant.
\end{proof}

\begin{corollary}[Averaged arc occupancy over subfamilies]\label{cor:AO-average}
With notation as above, for any fixed $F$ and $1\le L\le m$, the uniform average over all $\binom{m}{L}$ subfamilies $\mathcal L$ satisfies
\begin{equation}\label{eq:AO-average}
\mathbb E_{\mathcal L}\, S_L(F)\ =\ \frac{L}{m}\,S_{\rm tot}(F).
\end{equation}
In particular, there exists at least one subfamily $\mathcal L$ of size $L$ with $S_L(F)\le (L/m)\,S_{\rm tot}(F)$.
\end{corollary}

\begin{proof}
Since $\mathbb P(i\in\mathcal L)=L/m$ for each fixed index $i$, linearity of expectation gives
\[
\mathbb E_{\mathcal L}\,S_L(F)
=\sum_{i=1}^m \mathbb P(i\in\mathcal L)\,|y_i|^2
=\frac{L}{m}\sum_{i=1}^m |y_i|^2
=\frac{L}{m}\,S_{\rm tot}(F).
\]
The existence statement follows because the average of nonnegative numbers is at least their minimum.
\end{proof}

\begin{remark}[How these two inputs are used]
Inequality \eqref{eq:AO-deterministic} is the robust, fully deterministic control needed when bounding a chosen subfamily $\mathcal L$. 
Identity \eqref{eq:AO-average} is the exact combinatorial average behind the heuristic “$L/m$ fraction of mass on $L$ arcs”; it can be invoked whenever an averaged (or pigeonholed) choice of $\mathcal L$ is permissible.
\end{remark}

\begin{corollary}[No hot spots for bank--projected coefficients]\label{cor:no-hotspots}
Let $y_i=\langle P_A F,\psi_i\rangle$ and $S_{\rm tot}(F)=\sum_{i=1}^m |y_i|^2$ as in Proposition~\ref{prop:AO-projected}. 
Then no fixed subfamily of arcs can capture more than a bounded fraction of this total energy:
\[
\sum_{i\in\mathcal L}|y_i|^2
\ \le\ \Bigl(\frac{L}{m}+C\,Q^{-1+\varepsilon}\Bigr)S_{\rm tot}(F),
\qquad\forall\,\mathcal L\subset\{1,\dots,m\},\ |\mathcal L|=L.
\]
In particular, every individual arc satisfies
\[
|y_i|^2\ \le\ \Bigl(\frac{1}{m}+C\,Q^{-1+\varepsilon}\Bigr)S_{\rm tot}(F).
\]
\end{corollary}

\begin{proof}
Immediate from Proposition~\ref{prop:AO-projected} by taking $L=|\mathcal L|$ and, for the second statement, $L=1$.
\end{proof}

\subsection{Arithmetic origin of AO (large–sieve form)}\label{subsec:AO-ls}
For each reduced $a/q$ ($q\le Q$), let $\Phi_H$ be a fixed nonnegative Fejér window of bandwidth $1/H$ and set
\[
\varphi_{a/q}(\alpha):=\Phi_H(\alpha-a/q),\qquad 
\psi_{a/q}:=\frac{\varphi_{a/q}^{1/2}}{\|\varphi_{a/q}^{1/2}\|_2}.
\]
The classical (hybrid) large sieve implies, for any $F\in L^2(\mathbb T)$ supported on $\cA$,
\begin{equation}\label{eq:AO-ls}
\sum_{\substack{q\le Q\\(a,q)=1}} \big|\langle F,\psi_{a/q}\rangle\big|^2
\ \ll\ \Big(1+Q^{-1+\varepsilon}\Big)\,\int_{\cA}|F(\alpha)|^2\,d\alpha,
\end{equation}
uniformly in $F$. The spacing of reduced rationals ($|a_1/q_1-a_2/q_2|\gg (q_1q_2)^{-1}$) together with the bandwidth $1/H$ gives the near–orthogonality of the family $\{\psi_{a/q}\}$; (\ref{eq:AO-ls}) is a standard Montgomery large–sieve consequence adapted to smoothed windows. Grouping the sum over a subfamily $\mathcal L$ of size $L$ and dividing by $m\asymp Q^2$ yields Proposition~\ref{prop:AO-projected}.

\begin{lemma}[Bank–local near–orthogonality of arc detectors]\label{lem:AO-near}
Let $\mathcal A=\bigcup_i I_i$ be a union of bank arcs with $|I_i|\asymp 1/H$, and assume the guard–band separation
\[
|I_i^+\cap I_j^+|\ \ll\ \frac{1}{QH}\qquad(i\ne j),
\]
where $I_i^\circ\subset I_i\subset I_i^+$ are, respectively, the middle third and a fixed $(1+10^{-2})$ enlargement, and the family $\{I_i^+\}$ may have only the indicated small overlaps.
Choose $\varphi_i\in C_c^\infty(\mathbb T)$ with $1_{I_i^\circ}\le \varphi_i\le 1_{I_i^+}$ and set
\[
\psi_i\ :=\ \frac{\varphi_i}{\|\varphi_i\|_2}\,.
\]
Then
\[
\langle \psi_i,\psi_i\rangle=1,\qquad 
|\langle \psi_i,\psi_j\rangle|\ \ll\ Q^{-1+\varepsilon}\qquad(i\ne j).
\]
\end{lemma}

\begin{proof}
Since $\|\varphi_i\|_2^2\asymp |I_i|\asymp H^{-1}$, normalization gives $\langle\psi_i,\psi_i\rangle=1$.
For $i\ne j$,
\[
|\langle \psi_i,\psi_j\rangle|
=\frac{\int_{\mathbb T}\varphi_i\varphi_j}{\|\varphi_i\|_2\|\varphi_j\|_2}
\ \ll\ \frac{|I_i^+\cap I_j^+|}{|I_i|^{1/2}|I_j|^{1/2}}
\ \ll\ \frac{(QH)^{-1}}{H^{-1}}\ =\ Q^{-1}.
\]
Absorb harmless logarithms into $Q^{-1+\varepsilon}$.
\end{proof}

\begin{remark}
Hard disjointness ($I_i^+\cap I_j^+=\varnothing$) would make $\langle\psi_i,\psi_j\rangle=0$, but we do \emph{not} use that regime, because later detectors (e.g. Fejér windows) have small tails; the unified soft regime above aligns Lemma~\ref{lem:AO-near} with Lemma~4.1.
\end{remark}

\begin{corollary}[AO dispersion: no hot spots]\label{cor:AO-nohot}
For any bank–supported $F\in L^2(\mathbb T)$ and any subfamily $\mathcal L$ of $L$ arcs,
\[
\sum_{i\in\mathcal L}\int_{I_i}|F|^2
\;\le\;
\Big(\frac{L}{Q^2}+C\,Q^{-1+\varepsilon}\Big)\int_{\mathcal A}|F|^2.
\]
In particular, with $m\asymp Q^2$, the fraction of energy on any $L$ arcs is at most $L/m+O(Q^{-1+\varepsilon})$.
\end{corollary}

\begin{theorem}[Deterministic AO dispersion: no hot spots]\label{thm:AO-dispersion-full}
Let $\mathfrak A=\bigcup_{i=1}^{m} I_i$ be a bank partition, $m\asymp Q^2$, with $|I_i|\asymp H^{-1}$, and enlargements $I_i^+$ obeying the guard--band
\begin{equation}\label{eq:guardband}
|I_i^+\cap I_j^+|\ \ll\ \frac{1}{QH}\qquad(i\neq j).
\end{equation}
Choose $\varphi_i\in C_c^\infty(\mathbb T)$ with $1_{I_i^\circ}\le \varphi_i\le 1_{I_i^+}$ and set
\[
\psi_i:=\frac{\varphi_i}{\|\varphi_i\|_2},\qquad \|\varphi_i\|_2^2\asymp |I_i| \asymp H^{-1}.
\]
Then the Gram matrix $G=(\langle\psi_i,\psi_j\rangle)_{1\le i,j\le m}$ satisfies
\begin{equation}\label{eq:near-orth}
\langle\psi_i,\psi_i\rangle = 1,\qquad 
|\langle\psi_i,\psi_j\rangle|\ \ll\ Q^{-1+\varepsilon}\quad (i\ne j),
\end{equation}
and for every $F\in L^2(\mathbb T)$,
\begin{equation}\label{eq:AO-ratio}
\max_{|S|=L}\ \frac{\sum_{i\in S}|\langle F,\psi_i\rangle|^2}{\sum_{i=1}^{m}|\langle F,\psi_i\rangle|^2}
\ \le\ \frac{L}{m}\ +\ C\,Q^{-1+\varepsilon}\qquad(1\le L\le m),
\end{equation}
with an absolute $C$ depending only on the constants in \eqref{eq:guardband}.
\end{theorem}

\begin{proof}
\textit{Step 1 (near--orthogonality from overlap).}
By construction $\supp\varphi_i\subset I_i^+$ and $\varphi_i\ge 1$ on $I_i^\circ$, so $\|\varphi_i\|_2^2\asymp |I_i|\asymp H^{-1}$. For $i\neq j$,
\[
|\langle\psi_i,\psi_j\rangle|
=\frac{\big|\int \varphi_i\varphi_j\big|}{\|\varphi_i\|_2\,\|\varphi_j\|_2}
\ \ll\ \frac{|I_i^+\cap I_j^+|}{|I_i|^{1/2}|I_j|^{1/2}}
\ \ll\ \frac{(QH)^{-1}}{H^{-1}}\ =\ Q^{-1}.
\]
This gives \eqref{eq:near-orth} (absorbing harmless logarithms into $Q^{-1+\varepsilon}$).

\smallskip
\textit{Step 2 (Schur/Gershgorin bound for the Gram matrix).}
Write $G=I+R$, where $R_{ii}=0$ and $|R_{ij}|\ll Q^{-1+\varepsilon}$ for $i\ne j$. The Schur test yields
\[
\|R\|_{\ell^2\to \ell^2}\ \le\ \max_i\sum_{j\ne i}|R_{ij}|
\ \ll\ m\cdot Q^{-1+\varepsilon}\ \ll\ Q^{1+\varepsilon}\cdot Q^{-1+\varepsilon}\ \ll\ Q^{2\varepsilon},
\]
because each $i$ meets only $O(Q)$ neighbors with $|I_i^+\cap I_j^+|>0$, and all other overlaps are zero (the guard--band \eqref{eq:guardband} plus arc geometry). Choosing $\varepsilon$ small, this gives
\begin{equation}\label{eq:Gram-spectrum}
\|G\|\le 1+C_0Q^{-1+\varepsilon},\qquad
\|G^{-1}\|\le 1+C_0Q^{-1+\varepsilon}.
\end{equation}

\smallskip
\textit{Step 3 (energy averaging and deviation control).}
Let $a_i=\langle F,\psi_i\rangle$ and $a=(a_i)\in\mathbb C^m$. Define $E_{\mathfrak A}(F)=\sum_{i}|a_i|^2=\|a\|_2^2$ and, for $S\subset\{1,\dots,m\}$, $E_S(F)=\sum_{i\in S}|a_i|^2$.
Average $E_S$ over all $\binom{m}{L}$ subsets $S$ of size $L$:
\[
\frac{1}{\binom{m}{L}}\sum_{|S|=L} E_S(F)
= \frac{L}{m}\, \sum_{i=1}^m |a_i|^2
=\frac{L}{m}\,E_{\mathfrak A}(F).
\]
Thus to bound the \emph{maximum} over $S$ it suffices to control the deviation caused by the non–orthogonality. For fixed $S$, let $P_S$ be the diagonal projector on $S$. Then
\[
E_S(F) = \langle P_S a,a\rangle = \langle P_S T F, T F\rangle
= \langle T^\ast P_S T F, F\rangle,
\]
with $T:L^2(\mathbb T)\to \mathbb C^m$, $TF=a$, and $T^\ast T=G$. Insert $G=I+R$:
\[
T^\ast P_S T = T^\ast \Big(\tfrac{L}{m}I\Big) T\ +\ T^\ast\Big(P_S-\tfrac{L}{m}I\Big)T
= \tfrac{L}{m}\,G\ +\ T^\ast\Big(P_S-\tfrac{L}{m}I\Big)T.
\]
Hence
\[
E_S(F) - \frac{L}{m}E_{\mathfrak A}(F)
= \langle T^\ast\Big(P_S-\tfrac{L}{m}I\Big)T F, F\rangle
\ \le\ \|T^\ast\|\,\Big\|P_S-\tfrac{L}{m}I\Big\|\,\|T\|\ \|F\|_2^2.
\]
By \eqref{eq:Gram-spectrum}, $\|T\|^2=\|G\|\le 1+C_0Q^{-1+\varepsilon}$ and similarly for $\|T^\ast\|$, while $\|P_S-\tfrac{L}{m}I\|\le 1$. It follows that
\[
E_S(F) \ \le\ \frac{L}{m}\,E_{\mathfrak A}(F)\ +\ C\,Q^{-1+\varepsilon}\,E_{\mathfrak A}(F),
\]
uniformly over $S$. Taking the maximum over $|S|=L$ gives \eqref{eq:AO-ratio}.
\end{proof}

% ============ Block A: Bank geometry constants ============
\begin{definition}[Bank geometry and constants]\label{def:bank-geometry}
Fix $A\ge10$, $H=(\log X)^A$, and $Q=H^\gamma$ with $0<\gamma<\tfrac12$.
Let $\mathfrak A=\bigcup_{i=1}^{m} I_i$ be a collection of disjoint arcs $I_i\subset\mathbb T$ satisfying
\begin{equation}\label{eq:bank-geom}
\frac{c_w}{H}\ \le\ |I_i|\ \le\ \frac{C_w}{H},\qquad 
\mathrm{dist}(I_i,I_j)\ \ge\ \frac{\beta}{H}\quad (i\neq j),
\end{equation}
for fixed absolute constants $c_w,C_w,\beta>0$. 
The index count obeys $c_m Q^2\le m\le C_m Q^2$ for absolute $c_m,C_m>0$.
For each $i$ fix “inner” and “outer” enlargements
\[
I_i^\circ \subset I_i \subset I_i^+\subset \mathbb T,\qquad |I_i^\circ|\ge \sigma\,|I_i|,\quad |I_i^+|\le (1+\eta)\,|I_i|
\]
with fixed $\sigma\in(0,1)$, $\eta\in(0,10^{-2})$, and assume the guard–band
\begin{equation}\label{eq:guard-band}
|I_i^+\cap I_j^+| \le \frac{C_{\rm gb}}{QH}\qquad (i\neq j).
\end{equation}
\end{definition}

% ============ Block B: Detectors and normalizations ============
\begin{lemma}[Detectors; $L^1/L^2$ normalizations]\label{lem:detectors}
For each $i$ choose $\varphi_i\in C_c^\infty(\mathbb T)$ with $1_{I_i^\circ}\le \varphi_i \le 1_{I_i^+}$.
Let $\psi_i:=\varphi_i/\|\varphi_i\|_2$ and denote $|I_i|=:w_i$.
Then
\[
c_1 w_i \le \|\varphi_i\|_1 \le C_1 w_i,\qquad c_2 w_i \le \|\varphi_i\|_2^2 \le C_2 w_i,
\]
hence $c_3 H^{-1}\le \|\varphi_i\|_2^2\le C_3 H^{-1}$ with $c_k,C_k$ depending only on $(c_w,C_w,\sigma,\eta)$.
\end{lemma}

% ============ Block C: Overlap counting and Gram control ============
\begin{lemma}[Overlap counting]\label{lem:overlap}
For each fixed $i$ the set $\{j\ne i: I_i^+\cap I_j^+\neq\emptyset\}$ has cardinality $\le C_{\rm nb} Q$, where $C_{\rm nb}$ depends only on $(c_w,C_w,\beta,\eta)$.
\end{lemma}

\begin{lemma}[Gram off–diagonal bound]\label{lem:gram}
Let $G=(\langle\psi_i,\psi_j\rangle)_{1\le i,j\le m}$.
Then
\[
\langle\psi_i,\psi_i\rangle=1,\qquad 
|\langle\psi_i,\psi_j\rangle|\ \le\ \frac{C_{\rm od}}{Q}\quad(i\ne j),
\]
and the remainder $R:=G-I$ satisfies
\[
\|R\|_{\ell^2\to\ell^2}\ \le\ \kappa,\qquad \kappa:=\frac{C_{\rm AO}}{Q},
\]
with $C_{\rm AO}$ depending only on $(c_w,C_w,\beta,\sigma,\eta,C_{\rm gb},C_{\rm nb})$.
\end{lemma}

\begin{proof}[Sketch]
Since $\supp\varphi_i\subset I_i^+$, by Cauchy–Schwarz and Lemma~\ref{lem:detectors},
\[
|\langle\psi_i,\psi_j\rangle|
=\frac{\int \varphi_i\varphi_j}{\|\varphi_i\|_2\|\varphi_j\|_2}
\le \frac{|I_i^+\cap I_j^+|}{(c_2 w_i c_2 w_j)^{1/2}}
\ll \frac{(QH)^{-1}}{H^{-1}}\ =\ \frac{1}{Q}.
\]
For the operator norm, Schur or Gershgorin with Lemma~\ref{lem:overlap} gives
\(
\|R\|\le \max_i \sum_{j\ne i}|R_{ij}|
\le (C_{\rm nb}Q)\cdot (C_{\rm od}/Q)= C_{\rm AO}/Q.
\)
\end{proof}

% ============ Block D: AO dispersion with explicit constants ============
\begin{theorem}[Deterministic AO dispersion with constants]\label{thm:AO-disp-full}
For $F\in L^2(\mathbb T)$ set $a_i=\langle F,\psi_i\rangle$, $E_{\mathfrak A}(F):=\sum_{i=1}^m |a_i|^2$ and, for $S\subset\{1,\dots,m\}$ with $|S|=L$, let $E_S(F):=\sum_{i\in S}|a_i|^2$.
Then
\begin{equation}\label{eq:AO-disp-const}
\frac{E_S(F)}{E_{\mathfrak A}(F)}\ \le\ \frac{L}{m}\ +\ \kappa
\qquad\text{with }\ \kappa=\frac{C_{\rm AO}}{Q}.
\end{equation}
Consequently,
\[
\max_{|S|=L}\frac{E_S(F)}{E_{\mathfrak A}(F)}\ \le\ \frac{L}{m}+\frac{C_{\rm AO}}{Q}.
\]
\end{theorem}

\begin{proof}
Let $T:L^2(\mathbb T)\to\mathbb C^m$, $TF=(a_i)$, so $T^\ast T=G=I+R$ with $\|R\|\le\kappa$.
Average $E_S(F)=\langle P_S TF,TF\rangle$ over all $S$ of size $L$ to get $\mathbb E_S E_S=(L/m)E_{\mathfrak A}$.
For a fixed $S$,
\[
E_S-\tfrac{L}{m}E_{\mathfrak A}
=\langle T^\ast(P_S-\tfrac{L}{m}I)T F,F\rangle
\le \|T^\ast\|\,\|P_S-\tfrac{L}{m}I\|\,\|T\|\,\|F\|_2^2
\le (1+\kappa)\,\|F\|_2^2 - \tfrac{L}{m}\,(1-\kappa)\,\|F\|_2^2
\]
after inserting $G$ and using $\|T\|^2=\|G\|\le 1+\kappa$, $\|(T^\ast)\|^2=\|G\|\le 1+\kappa$ and $\|P_S-\tfrac{L}{m}I\|\le 1$.
Dividing by $E_{\mathfrak A}=\langle GF,F\rangle \ge (1-\kappa)\|F\|_2^2$ yields \eqref{eq:AO-disp-const}.
\end{proof}

% ============ Block E: Matching the TT* probes and window energies ============
\begin{proposition}[Equivalence of probe and window energies]\label{prop:AO-eq-window}
Let $w_i\in C_c^\infty(\mathbb T)$ be nonnegative windows with $c_-\,1_{I_i^\circ}\le w_i\le c_+\,1_{I_i^+}$.
Then for all $F\in L^2(\mathbb T)$
\[
c_-\ c_2\ \sum_{i} \big|\langle F,\psi_i\rangle\big|^2
\ \le\ 
\sum_i \int_{\mathbb T} |F(\xi)|^2 w_i(\xi)\,d\xi
\ \le\ 
c_+\, C_2\ \sum_{i} \big|\langle F,\psi_i\rangle\big|^2.
\]
In particular, the AO dispersion bound \eqref{eq:AO-disp-const} transfers verbatim to the windowed energy used in the $TT^\ast$ ledger (up to the fixed constants $c_\pm, c_2, C_2$).
\end{proposition}

\begin{corollary}[AO $\Rightarrow$ SSU packet energy control]\label{cor:AO-to-SSU}
Let $\{w_i\}$ be bank windows as in Proposition~\ref{prop:AO-eq-window}, and let $\{\eta_\nu\}$ be the SSU packets.
Assume each $\eta_\nu$ is supported in the union of at most $C$ enlarged arcs $I_i^+$ (true since $|\supp\eta_\nu|\asymp (q_\nu H)^{-1}$ and $|I_i|\asymp H^{-1}$).
Then for any $F\in L^2(\mathbb T)$,
\begin{equation}\label{eq:AO-to-SSU}
\sum_{\nu} \big|\langle F,\eta_\nu\rangle\big|^2
\ \ll\ \sum_{i=1}^{m} \big|\langle F,\psi_i\rangle\big|^2,
\end{equation}
and if $\mathcal N\subset\{\nu\}$ has $|\mathcal N|=L$, then
\begin{equation}\label{eq:AO-to-SSU-disp}
\sum_{\nu\in\mathcal N} \big|\langle F,\eta_\nu\rangle\big|^2
\ \ll\ \Big(\frac{L}{m}+\frac{C_{\rm AO}}{Q}\Big)\ \sum_{i=1}^{m} \big|\langle F,\psi_i\rangle\big|^2.
\end{equation}
\end{corollary}

\begin{proof}
Since $\eta_\nu\ll \sum_{i\in S(\nu)} w_i$ with $|S(\nu)|\le C$ and by Cauchy–Schwarz, 
$|\langle F,\eta_\nu\rangle|^2 \ll \sum_{i\in S(\nu)} |\langle F,\psi_i\rangle|^2$ using Lemma~\ref{lem:detectors}.
Finite overlap of $\nu\mapsto S(\nu)$ and Proposition~\ref{prop:AO-eq-window} give \eqref{eq:AO-to-SSU}.
For \eqref{eq:AO-to-SSU-disp}, apply Theorem~\ref{thm:AO-disp-full} to the family $\{\psi_i\}$ and the set of indices supporting $\{\eta_\nu:\nu\in\mathcal N\}$, noting $m\asymp Q^2$ and the overlap constant is absorbed into $C_{\rm AO}$.
\end{proof}

\begin{theorem}[Deterministic AO with explicit constants]\label{thm:AO-explicit}
Let $\{I_i\}_{i=1}^m$ be disjoint arcs with $|I_i|\asymp 1/H$ and mutual gaps $\gg 1/H$.
Let $\{\psi_i\}$ be $C^\infty$ arc detectors with $\mathbf{1}_{I_i^\circ}\le \psi_i\le \mathbf{1}_{I_i^+}$ and $\| \psi_i\|_2=1$.
There exists $C_{\mathrm{AO}}=C_{\mathrm{AO}}(\text{taper},\text{overlap})$ such that for every $F\in L^2(\mathbb T)$ and every subfamily $\mathcal L\subset\{1,\dots,m\}$ with $|\mathcal L|=L$,
\[
\sum_{i\in\mathcal L} |\langle F,\psi_i\rangle|^2
\ \le\ \Big(\frac{L}{m} + \frac{C_{\mathrm{AO}}}{Q}\Big)\,\sum_{i=1}^m |\langle F,\psi_i\rangle|^2
\ +\ O(Q^{-1+\varepsilon})\,\|F\|_2^2.
\]
\end{theorem}

\begin{proof}
Let $G=(\langle \psi_i,\psi_j\rangle)_{1\le i,j\le m}$ be the Gram matrix. By guard bands and support size, $G_{ii}=1$ and $|G_{ij}|\ll Q^{-1+\varepsilon}$ for $i\ne j$. Gershgorin’s theorem (or Schur test) yields $\|G-\mathrm{Id}\|_{\ell^2\to\ell^2}\ll Q^{-1+\varepsilon}$. For any coefficients $c_i:=\langle F,\psi_i\rangle$, $\sum_{i\in\mathcal L}|c_i|^2 \le \|P_{\mathcal L}\|\,\sum_{i}|c_i|^2$ where $P_{\mathcal L}$ is the orthogonal projector onto the span of $\{\psi_i:i\in\mathcal L\}$ measured in the $G$-metric. Since $P_{\mathcal L}$ has rank $L$, a trace bound gives $\|P_{\mathcal L}\|\le L/m + O(Q^{-1+\varepsilon})$. The $C_{\mathrm{AO}}/Q$ term records taper/overlap constants when passing from sharp arcs to smooth detectors.
\end{proof}
                 % AO$^{+}$: packet mixing / dispersion on arcs
% !TEX root = main.tex
\section{Bilinear Geometry (BG)}\label{sec:BG}

This section proves the short–shift anti–alignment for Type–II bilinear forms at bandwith $1/H$, with the optimal square–root gain \smash{$(H/X)^{1/2}$} up to polylogarithmic factors. The result is \emph{non–flat}: it holds for arbitrary coefficient arrays without any $\ell^\infty$ caps, via a two–weight positive–operator argument (Sawyer testing) adapted to sheared tubes. We also record the tube–localized slope–shift uncertainty (SSU) estimate that powers the square–root numerics.

Throughout, $X$ is large, $A\ge 10$ is fixed, and
\[
H=(\log X)^A,\qquad 0<\varepsilon\le \tfrac{1}{10},\qquad X^\varepsilon\le D,N\le X^{1-\varepsilon},\quad DN\asymp X.
\]
We work with a fixed smooth bump $W\in C_c^\infty([1/2,2]^2)$ with $\|\partial^\alpha W\|_\infty\le C_W$ for $|\alpha|\le 3$.
We use a real, even, nonnegative short–shift kernel $K_H:\mathbb{Z}\to\mathbb{R}_{\ge 0}$ whose Fourier transform on $\mathbb{T}$ obeys
\begin{equation}\label{eq:KH.assumptions}
\operatorname{supp}\widehat K_H \subset \{|\xi|\le 1/H\},\qquad
\int_{\mathbb{R}}K_H(t)\,dt = H+O(1),\qquad
\|\widehat K_H\|_{L^1(\mathbb{T})}\ll H,\qquad
\int_{|\xi|\le 1/H}\frac{|\widehat K_H(\xi)|}{|\xi|}\,d\xi \ll H\log H .
\end{equation}
(A scaled Beurling–Selberg/Vaaler majorant for $[-H,H]$ satisfies \eqref{eq:KH.assumptions}.)

\subsection{Bilinear form, short–shift operator, and goal}\label{subsec:BG.setup}

Given coefficients $(\alpha_d)$, $(\beta_n)$ supported on $d\sim D$, $n\sim N$, set
\begin{equation}\label{eq:Ak.def}
A_k := \sum_{dn=k} \alpha_d\,\beta_n\,
W\!\Big(\frac{d}{D},\frac{n}{N}\Big).
\end{equation}
Define the (smoothed) short–shift correlation operator
\begin{equation}\label{eq:BG.kernel}
\langle \mathbf T F,\ G\rangle \ :=\
\sum_{d',n'}\ \sum_{d,n}\ K_H(d'n - dn')\,
W\!\Big(\frac{d}{D},\frac{n}{N}\Big)\,W\!\Big(\frac{d'}{D},\frac{n'}{N}\Big)\,
F(d,n)\,\overline{G(d',n')}.
\end{equation}
For the Type–II array $F(d,n)=\alpha_d\beta_n$ this quadratic form equals
\[
\langle \mathbf T(\alpha\!\otimes\!\beta),\,\alpha\!\otimes\!\beta \rangle
=\sum_{k,k'} A_{k'}\,\overline{A_k}\,K_H(k'-k),
\]
\begin{remark}[Justification]\label{rem:justification-53-54}
Equation~(5.3) has the form
\[
\mathcal B \;=\; \int_{-1/H}^{1/H} \widehat K_H(\xi)\,\bigl|S(\xi)\bigr|^2\,d\xi,
\qquad \widehat K_H(\xi)\ge 0,
\]
so $\mathcal B$ is a positive average of $\lvert S(\xi)\rvert^2$ over $|\xi|\le 1/H$. Hence any pointwise upper bound for $\lvert S(\xi)\rvert^2$ on that band yields the corresponding bound for $\mathcal B$ after multiplying by $\int_{-1/H}^{1/H}\widehat K_H(\xi)\,d\xi$ (absorbed into constants).

To obtain such a pointwise bound, first apply the shear $(d,n)\mapsto(u,v)=(qn-ad,\ d)$ so that $d'n-dn'=(v'u-vu')/q$, and rewrite
\[
S(\xi)\;=\;\sum_{(u,v)} F(v,u)\,e\!\Big(\frac{\xi\,u v}{qX}\Big),
\]
with $(u,v)$ constrained to the tube and the congruence $u\equiv -av\pmod q$. Then, for one progression (say in $v$ with modulus $q$), use Cauchy--Schwarz in the outer variable and the Montgomery--Vaughan large sieve in the inner variable; repeat dually for the progression in $u$. Taking the geometric mean of these two bounds yields a pointwise estimate
\[
|S(\xi)|^2 \;\ll\; \Big(\tfrac{DU}{q}\Big)^{1/2}
\Big(U+\tfrac{X}{|\xi|}\Big)^{1/2}\Big(D+\tfrac{X}{|\xi|}\Big)^{1/2}
\sum_{(u,v)} |F(v,u)|^2,
\]
uniformly for $|\xi|\le 1/H$. Integrating against $\widehat K_H(\xi)$ and using the kernel moment bounds
\(
\int\widehat K_H(\xi)\,d\xi\ll H,\ 
\int \widehat K_H(\xi)\,|\xi|^{-1}\,d\xi\ll H\log H
\)
gives the right-hand side of~(5.4), with all dependence on smooth cutoffs and fixed dyadic normalizations absorbed into a factor $(\log X)^{O(1)}$. Finally, the short-axis condition $U\ll X/(qH)$ simplifies the product to the stated $(H/X)^{1/2}$ scale.
\end{remark}

\subsection{Shearing, tubes, and the SSU estimate}\label{subsec:BG.SSU}

Write the skew form $B((d,n),(d',n')):=d'n-dn'$. Fix a reduced slope $a/q$ (with $1\le q\le H^\gamma$, $0<\gamma<\tfrac12$) and consider the unimodular shear
\begin{equation}\label{eq:BG.shear}
\Phi_{a/q}:\ (d,n)\mapsto (u,v)=(q d - a n,\ a d + q n).
\end{equation}
A direct calculation shows the factorization identity
\begin{equation}\label{eq:BG.factor}
B\big((d,n),(d',n')\big) \ =\ \frac{u'v - uv'}{a^2+q^2}.
\end{equation}
We define the \emph{sheared tube} of slope $a/q$ and cross–slope thickness $\delta:=c_0\,H/(DN)\asymp H/X$ by
\begin{equation}\label{eq:BG.tube}
T_{a/q}\ :=\ \Big\{(d,n)\in [D,2D]\times[N,2N]\cap\mathbb{Z}^2:\ \Big|\frac{n}{d}-\frac{a}{q}\Big|\le \delta\Big\}.
\end{equation}
A standard partition argument yields $O(H^{2\gamma})$ such tubes with bounded overlap covering the near–shift region $|B|\ll H$; we refer to this family as the \emph{tube pack}.

\begin{lemma}[Short–axis large sieve]\label{lem:BG.SALS}
Let $U\ge 1$ and suppose the frequencies $\{\beta_j\}\subset\mathbb{T}$ are $\Delta$–separated. Then for any $(g(u))_{|u|\le U}$,
\[
\sum_j \Big|\sum_{|u|\le U} g(u)\,e(\beta_j u)\Big|^2
\ \ll\ (U+\Delta^{-1})\ \sum_{|u|\le U} |g(u)|^2.
\]
\end{lemma}

\begin{proof}
This is the classical Montgomery–Vaughan large sieve inequality.
\end{proof}

% ===== BEGIN PATCH: Tube partition and bounded overlap =====
\subsection{Tube partition and bounded overlap}\label{subsec:tubes-partition}

Fix a dyadic box $(D,2D]\times(N,2N]$ with $DN\asymp X$, and $Q=H^\gamma$ with $0<\gamma<\tfrac12$.
For each reduced $a/q$ with $1\le q\le Q$, define the shear
\[
u:= qn-ad,\qquad v:= d.
\]
Let $U$ be any parameter with $1\le U \ll X/(qH)$, and write
\[
\mathcal{S}_{a/q}:=\Big\{s\in \mathbb{Z}:\ |s|\le \tfrac{X}{2qH}\Big\}.
\]
Define the $(a/q,s)$–tube by
\begin{equation}\label{eq:def-tube}
T(a/q,s)\ :=\ \Bigl\{(d,n)\in (D,2D]\times(N,2N]:\ |\,qn-ad-s\,|\ \le\ U\Bigr\}.
\end{equation}

\begin{lemma}[Covering and bounded overlap]\label{lem:tubes-partition}
For each $(D,N)$ and each $a/q$ as above, the family $\{T(a/q,s)\}_{s\in\mathcal{S}_{a/q}}$ covers the slice
\[
\{(d,n)\in (D,2D]\times(N,2N]\ :\ |qn-ad|\le X/(2qH)\}.
\]
Moreover, there is a constant $C$ (independent of $(D,N),H,Q$) such that any point $(d,n)$ belongs to at most $C$ tubes $T(a/q,s)$ for a fixed $a/q$, and to at most $(\log X)^C$ tubes when $a/q$ ranges over all rationals with $q\le Q$.
\end{lemma}

\begin{proof}
Fix $(d,n)$ with $|qn-ad|\le X/(2qH)$. Choose $s$ to be the nearest integer to $qn-ad$; then $|\,qn-ad-s\,|\le \tfrac12\le U$, giving coverage.
For the per–$a/q$ multiplicity: if $(d,n)\in T(a/q,s)\cap T(a/q,s')$ then $|s-s'|\le 2U$, hence $s=s'$ since $U<1$ can be enforced by replacing $U$ by $\lfloor U\rfloor+1$ in the definition (harmless in all bounds), so per–$a/q$ multiplicity is $O(1)$.

To sum over different $a/q$, note that $|qn-ad|\le X/(qH)$ and $q\le Q$ imply $|n/d-a/q|\ll 1/(qH)$.
By Farey spacing, for distinct reduced fractions $a_i/q_i\ne a_j/q_j$ we have $|a_i/q_i-a_j/q_j|\ge 1/(q_iq_j)$. Hence overlap can occur only when $q_iq_j\gg H$, i.e. at most $O((\log X)^C)$ distinct dyadic $q$–ranges contribute for a fixed $(d,n)$ once we also account for the compact support of the frequency windows (from \S\ref{sec:TFA}) and the $q\le Q$ constraint. This gives the stated $(\log X)^C$ bound.
\end{proof}
% ===== END PATCH =====


\begin{theorem}[SSU on a tube]\label{thm:BG.SSU}
Let $T_{a/q}$ be as in \eqref{eq:BG.tube}. For any $F:\mathbb{Z}^2\to\mathbb{C}$ supported on $T_{a/q}$,
\begin{equation}\label{eq:BG.SSU}
\sum_{(d',n')\in T_{a/q}}\ \sum_{(d,n)\in T_{a/q}}
K_H(d'n-dn')\,F(d,n)\,\overline{F(d',n')}
\ \ll\ (\log X)^C\,\Big(\frac{H}{X}\Big)^{1/2}\ \sum_{(d,n)\in T_{a/q}} |F(d,n)|^2.
\end{equation}
\end{theorem}

\begin{proof}
Write the skew form $B((d,n),(d',n')):=d'n-dn'$, fix a reduced slope $a/q$ with $1\le q\le H^\gamma$ and $0<\gamma<\tfrac12$, and shear via
\[
\Phi_{a/q}:(d,n)\mapsto (u,v)=(qd-an,\; ad+qn),
\]
so that
\[
B\big((d,n),(d',n')\big)=\frac{u'v-uv'}{a^2+q^2}
\qquad\text{(cf.\ \((5.6)\)).}
\]
Let $T_{a/q}$ be the tube \((5.7)\). On $T_{a/q}$ one has the short–axis bound
\[
|u|\ \ll\ \delta\,D\,q\ \ll\ H^{1+\gamma}/N=:U\qquad\text{and}\qquad v\ \text{ranges over an interval of length}\ \asymp X.
\]
Expand the positive, band–limited kernel $K_H$ by Fourier inversion:
\[
K_H(t)=\int_{|\xi|\le 1/H} \widehat K_H(\xi)\,e(\xi t)\,d\xi,
\qquad \widehat K_H\in L^1,\quad \int | \widehat K_H|\ll H,\quad \int_{|\xi|\le 1/H}\frac{|\widehat K_H(\xi)|}{|\xi|}\,d\xi\ll H\log H,
\]
cf.\ \((5.1)\). Then the quadratic form on $T_{a/q}$ equals
\[
\mathcal{Q}:=\sum_{(d',n'),(d,n)\in T_{a/q}} K_H\!\big(B((d,n),(d',n'))\big)\,F(d,n)\,F(d',n')
=\int_{|\xi|\le 1/H} \widehat K_H(\xi)\,\mathcal{S}(\xi)\,d\xi,
\]
where, in $(u,v)$–coordinates,
\[
\mathcal{S}(\xi)=\sum_{u',v'}\sum_{u,v} F(u,v)\,F(u',v')\,
e\!\left(\frac{\xi}{a^2+q^2}(u'v-uv')\right).
\]
Set the linear phases $\alpha_{u'}:=\xi u'/(a^2+q^2)$ and $\beta_{v'}:=\xi v'/(a^2+q^2)$ so that
$e\!\left(\tfrac{\xi}{a^2+q^2}(u'v-uv')\right)=e(\alpha_{u'}v)\,e(-\beta_{v'}u)$.
For each fixed $u'$ define
\[
h_{u'}(u):=\sum_{v} F(u,v)\,e(\alpha_{u'}v).
\]
Then
\[
\mathcal{S}(\xi)=\sum_{u',v'} F(u',v')\sum_{u} h_{u'}(u)\,e(-\beta_{v'}u).
\]
Apply Cauchy–Schwarz in $(u',v')$ and the Montgomery–Vaughan large sieve in the short variable $u$ (Lemma~5.1): the set of frequencies $\{\beta_{v'}\}$ is $\Delta$–separated with $\Delta\gg |\xi|$ because $v'$ runs over an interval of step $1$ and $|\xi|\le 1/H$ (hence distinct $v'$ give $|\beta_{v'}-\beta_{v''}|\ge c|\xi|$ modulo $1$). Therefore
\[
\sum_{v'}\Big|\sum_{u} h_{u'}(u)\,e(-\beta_{v'}u)\Big|^2
\ \ll\ (U+|\xi|^{-1})\sum_{|u|\le U} |h_{u'}(u)|^2.
\]
Summing this bound over $u'$ and undoing the $\alpha_{u'}$–modulation by Parseval in the long variable $v$ yields
\[
\sum_{u'}\sum_{v'}\Big|\sum_{u} h_{u'}(u)\,e(-\beta_{v'}u)\Big|^2
\ \ll\ (U+|\xi|^{-1})\sum_{u}\sum_{v} |F(u,v)|^2.
\]
Combining with Cauchy–Schwarz gives
\[
|\mathcal{S}(\xi)|\ \ll\ (U+|\xi|^{-1})^{1/2}\,\Big(\sum_{u,v}|F(u,v)|^2\Big).
\]
Integrate against $|\widehat K_H(\xi)|$ and use \(\int |\widehat K_H|\ll H\) and \(\int \tfrac{|\widehat K_H(\xi)|}{|\xi|}\,d\xi\ll H\log H\) from \((5.1)\):
\[
\mathcal{Q}\ \ll\ \Big(\int_{|\xi|\le 1/H} |\widehat K_H(\xi)|\,(U+|\xi|^{-1})^{1/2}\,d\xi\Big)\,\sum_{u,v}|F(u,v)|^2
\ \ll\ (\log X)^C\,H^{1/2}\,\sum_{u,v}|F(u,v)|^2.
\]
Finally, pass back from sheared lattice coordinates $(u,v)$ to physical $(d,n)$. The normalization of the bilinear form on a window of area $\asymp X$ contributes the factor $X^{-1/2}$, giving
\[
\sum_{(d',n'),(d,n)\in T_{a/q}} K_H(d'n-dn')\,F(d,n)\,F(d',n')
\ \ll\ (\log X)^C\,\Big(\frac{H}{X}\Big)^{\!1/2}\sum_{(d,n)\in T_{a/q}} |F(d,n)|^2,
\]
which is \((5.8)\).
\end{proof}


\subsection{Incidence and degeneracy}\label{subsec:BG.incidence}

We use two geometric counting inputs; both are stable under dyadic localization and smoothing by $W$.

\begin{lemma}[Near–hyperbola incidence]\label{lem:BG.incidence}
Let $Z(H)$ be the number (with smooth weights $W$) of quadruples $(d,n,d',n')$ with $d\sim D$, $n\sim N$, $d'\sim D$, $n'\sim N$ and $|d'n-dn'|<H$. Then
\[
Z(H)\ \ll\ (\log X)^C\,\big(H(D+N)+X\big).
\]
\end{lemma}

\begin{proof}
Insert the smooth indicator of $|t|\le H$ using the band–limited \emph{majorant} $M_H$:
\[
\mathbf{1}_{\{|t|\le H\}}(t)\ \le\ M_H(t),\qquad M_H\ge 0,\qquad \int M_H=H+O(1),\qquad \supp\widehat M_H\subset[-1/H,1/H],
\]
see Appendix~\ref{subsec:BS}.

Fix $(d,n)$, and apply Poisson in $(d',n')$ with the dyadically scaled weight $W'$:
\[
\sum_{d',n'} K_H(d'n-dn')\,W'\!\Big(\frac{d'}{D},\frac{n'}{N}\Big)
=\sum_{u,v\in\mathbb{Z}} \widehat K_H\big(u\,nX-v\,dX\big)\,\widehat{W'}(u,v),
\]
where $\widehat{W'}$ decays rapidly. The $(u,v)=(0,0)$ term contributes
\[
(H+O(1))\iint W' = H\,DN+O(DN).
\] See \ref{subsec:BS}.
Summing over $(d,n)$ against $W(d/D,n/N)$ gives a main term $\ll HX$. For $(u,v)\neq(0,0)$ we have $\widehat K_H(\xi)=0$ unless $|\xi|\le 1/H$, i.e. $|un-vd|\ll X/H$. Counting lattice points $(d,n)$ in the strip $|un-vd|\le X/H$ inside $[D,2D]\times[N,2N]$ gives by standard divisor geometry
\[
\ll (D+N)+X/H.
\]
Since only $O(1)$ pairs $(u,v)$ contribute (due to the rapid decay of $\widehat{W'}$), the total off–diagonal is
\[
\ll (\log X)^C\big((D+N)+X/H\big).
\]
Multiplying by the outer sum in $(d',n')$ recovered by Poisson yields the stated bound
\[
Z(H)\ \ll\ (\log X)^C\big(H(D+N)+X\big).
\]
All implied constants depend only on $(\varepsilon,C_W)$. 
\end{proof}


\begin{lemma}[Parallel–tube degeneracy]\label{lem:BG.degeneracy}
Let $\mathcal T$ be any subfamily of slope–tubes whose slopes lie in an interval of width $\eta>0$. Then the total weighted occupancy of $\mathcal T$ satisfies
\[
\sum_{T\in\mathcal T}\ \sum_{(d,n)\in T} W\!\Big(\tfrac dD,\tfrac nN\Big)
\ \ll\ (\log X)^C\,\eta\,(D+N)\,,
\]
uniformly for $H\le X^{1-\varepsilon}$.
\end{lemma}

\begin{proof}
Work in the sheared coordinates $(u,v)$ associated with the center slope $a/q$; in these coordinates every tube with slope in the window of width $\eta$ becomes a translate of a common rectangle with short side $\ll U\asymp H^{1+\gamma}/N$ and long side $\asymp X$ (uniformly in the tube), cf.\ \((5.5)\)–\((5.7)\). Let $\mathcal{T}$ be any such subfamily. By the tube–overlap lemma (Appendix~B), the family has bounded overlap at the lattice scale, and any fixed tube meets only $O((\log X)^C)$ others.

Write
\[
\sum_{T\in\mathcal{T}} \ \sum_{(d,n)\in T} W\!\Big(\frac{d}{D},\frac{n}{N}\Big)
=\sum_{(d,n)} W\!\Big(\frac{d}{D},\frac{n}{N}\Big)\,\#\{T\in\mathcal{T}:(d,n)\in T\}.
\]
For each $(d,n)$ the set of slopes whose tubes contain $(d,n)$ is an interval of length $\ll \eta$ in slope space, hence contributes at most $O(1+ \eta\,\Xi)$ tubes with $\Xi\asymp 1$ after discretization to rationals $a/q$ with $q\le Q=H^\gamma$; the bounded–overlap property absorbs the discretization loss into $(\log X)^C$. Therefore
\[
\sum_{T\in\mathcal{T}} \ \sum_{(d,n)\in T} W\!\Big(\frac{d}{D},\frac{n}{N}\Big)
\ \ll\ (\log X)^C\,\eta\sum_{(d,n)} W\!\Big(\frac{d}{D},\frac{n}{N}\Big).
\]
Finally, since $W$ is supported and bounded on $[D,2D]\times[N,2N]$, the last sum is $\asymp D+N$ after one summation by parts in $d$ or $n$; hence
\[
\sum_{T\in\mathcal{T}} \ \sum_{(d,n)\in T} W\!\Big(\frac{d}{D},\frac{n}{N}\Big)\ \ll\ (\log X)^C\,\big(\eta(D+N)+1\big),
\]
uniformly for $H\le X^{1-\varepsilon}$. This is the claimed degeneracy bound.
\end{proof}


\subsection{Non–flat BG via two–weight testing}\label{subsec:BG.two-weight}

Set
\[
\mathcal K_H((d,n),(d',n')):=K_H(d'n-dn')\,W\!\Big(\tfrac dD,\tfrac nN\Big)\,W\!\Big(\tfrac{d'}D,\tfrac{n'}N\Big),
\]
and define $(\mathbf T F)(d',n'):=\sum_{d,n}\mathcal K_H((d,n),(d',n'))\,F(d,n)$.
For an arbitrary array $F$ supported on $[D,2D]\times[N,2N]$, define the geometry-only row/column densities
\[
  \rho_d \;:=\; \sum_{n} W\!\Big(\frac{d}{D},\frac{n}{N}\Big)^{\!2},\qquad
  \kappa_n \;:=\; \sum_{d} W\!\Big(\frac{d}{D},\frac{n}{N}\Big)^{\!2},\qquad
  \sigma(d,n) \;:=\; \rho_d+\kappa_n.
\]
These depend only on the dyadic geometry and $W$, not on $F$.  Uniformly on the range $X^\varepsilon\le D,N\le X^{1-\varepsilon}$ one has
\[
  \rho_d \asymp N,\qquad \kappa_n \asymp D,
\]
with implied constants depending only on $(\varepsilon,A,\gamma,C_W)$.

\begin{lemma}[Tube–Sawyer testing]\label{lem:BG.Sawyer}
There exists $C=C(\varepsilon,A,\gamma,C_W)$ such that for all $E\subset [D,2D]\times[N,2N]$,
\begin{equation}\label{eq:BG.forward}
\|\mathbf T\mathbf 1_E\|_{L^2(\sigma)}^2 \ \ll\ (\log X)^C\,\frac{H}{X}\ \sigma(E),\qquad
\|\mathbf T^\ast\mathbf 1_E\|_{L^2(\sigma)}^2 \ \ll\ (\log X)^C\,\frac{H}{X}\ \sigma(E).
\end{equation}
\end{lemma}

\begin{lemma}[Tube--Sawyer testing]\label{lem:TubeSawyer}
Let $\mathbf T$ be the positive operator on $\Omega=[D,2D]\times[N,2N]$ with kernel
\[
\mathcal K_H\!\big((d,n),(d',n')\big):=K_H(d'n-dn')\,W\!\Big(\frac dD,\frac nN\Big)\,W\!\Big(\frac{d'}D,\frac{n'}N\Big),
\]
where $M_H\ge 0$ is even, band-limited to $|\xi|\le 1/H$, and satisfies $\int M_H = H+O(1)$ and $\|\widehat{M_H}\|_1\ll H$.
Define row/column energies
\[
\rho_d:=\sum_n |F(d,n)|^2 W\!\Big(\frac dD,\frac nN\Big)^2,\qquad \kappa_n:=\sum_d |F(d,n)|^2 W\!\Big(\frac dD,\frac nN\Big)^2,
\]
and the weight $\sigma(d,n):=\rho_d+\kappa_n$. Then for every measurable $E\subset\Omega$
\begin{equation}\label{eq:TubeTestForward}
\|\mathbf T \mathbf 1_E\|_{L^2(\sigma)}^2 \ \ll\ (\log X)^{C}\,\frac{H}{X}\ \sigma(E),
\end{equation}
and the same bound holds with $\mathbf T^\ast$ in place of $\mathbf T$. Consequently,
\[
\|\mathbf T\|_{L^2(\sigma)\to L^2(\sigma)}\ \ll\ (\log X)^{C}\,\Big(\frac{H}{X}\Big)^{1/2}.
\]
\end{lemma}

\begin{proof}
\emph{Tube partition.} Partition $\Omega\times\Omega$ into sheared tubes $\{T_\tau\}$ adapted to the level sets of $d'n-dn'=\Delta$ and slope $d'/d\approx a/q$, with $|\Delta|\lesssim H$ and slope aperture $\asymp Q^{-2}$; the long axis is $\asymp X$ and the short axis $\asymp X/H$. This is the standard pack used in \S6 (SSU) with bounded overlap, uniform for $q\le Q$ (see Remark~6.5). 

\emph{Forward testing.} Write
\[
(\mathbf T\mathbf 1_E)(d',n')=\sum_{(d,n)\in E}\!\mathcal K_H((d,n),(d',n'))=\sum_\tau \sum_{(d,n)\in E\cap T_\tau(\cdot,n',d')}\!\mathcal K_H(\cdot).
\]
By Cauchy--Schwarz in each tube $T_\tau$ and the SSU tube norm (Corollary~6.4),
\[
\sum_{(d',n')\in T_\tau}\!\big|(\mathbf T\mathbf 1_E)(d',n')\big|^2
\ \ll\ (\log X)^C \Big(\frac{H}{X}\Big) \sum_{(d,n)\in E\cap T_\tau} 1.
\]
Weighting by $\sigma(d',n')$ and summing over tubes gives
\[
\|\mathbf T\mathbf 1_E\|_{L^2(\sigma)}^2
\ \ll\ (\log X)^C\frac{H}{X}\sum_\tau \sum_{(d,n)\in E\cap T_\tau} \frac{\sigma\big(T_\tau(d,n)\big)}{1},
\]
where $\sigma(T_\tau(d,n)):=\sum_{(d',n')\in T_\tau(d,n)}\sigma(d',n')$.

\emph{Incidence \& degeneracy.} By the near-hyperbola incidence bound and the parallel-tube degeneracy (Appendix~\ref{app:incidence}, Props.~\ref{prop:incidence} and \ref{lem:BG.degeneracy}), each point $(d,n)$ belongs to only $O(1)$ tubes per slope family, and a slope family contributes $\ll \rho_d+\kappa_n$ when summed in $(d',n')$ against $\sigma$. Quantitatively,
\[
\sum_{\tau:\,(d,n)\in T_\tau}\ \sum_{(d',n')\in T_\tau(d,n)} \sigma(d',n')
\ \ll\ (\log X)^C\,(\rho_d+\kappa_n).
\]
Therefore
\[
\|\mathbf T\mathbf 1_E\|_{L^2(\sigma)}^2
\ \ll\ (\log X)^C\frac{H}{X}\sum_{(d,n)\in E} (\rho_d+\kappa_n)
\ =\ (\log X)^C\frac{H}{X}\,\sigma(E),
\]
which is \eqref{eq:TubeTestForward}. The adjoint test is identical since $\mathcal K_H$ is symmetric.

\emph{Two-weight conclusion.} By Sawyer’s two-test criterion for positive kernels, the forward and backward tests imply $\|\mathbf T\|_{L^2(\sigma)\to L^2(\sigma)}\le C M^{1/2}$ with $M\asymp (\log X)^C(H/X)$, proving the operator bound. Finally, since $\rho_d\asymp N$ and $\kappa_n\asymp D$ uniformly, Cauchy–Schwarz gives the embeddings
$L^2 \hookrightarrow L^2(\sigma)$ and $L^2(\sigma)\hookrightarrow L^2$ with constants $O(1)$, so the unweighted $\ell^2$ bound follows.\end{proof}

\begin{theorem}[Non–flat BG]\label{thm:BG-nonflat}
For all arrays $F$ supported on $[D,2D]\times[N,2N]$,
\begin{equation}\label{eq:BG.2to2}
\sum_{d',n'}\Big|\sum_{d,n} \mathcal K_H((d,n),(d',n'))\,F(d,n)\Big|^2
\ \ll\ (\log X)^C\,\frac{H}{X}\ \sum_{d,n}|F(d,n)|^2.
\end{equation}
Equivalently, $\|\mathbf T\|_{2\to 2}\ \ll\ (\log X)^C\,(H/X)^{1/2}$.
\end{theorem}

\begin{proof}
By Lemma~\ref{lem:BG.Sawyer} and the Sawyer two–weight principle for positive kernels (testing $\mathbf T$ and $\mathbf T^\ast$ on indicators controls the $L^2(\sigma)\!\to\!L^2(\sigma)$ norm), we get
$\|\mathbf T\|_{L^2(\sigma)\to L^2(\sigma)}\ll (\log X)^C\,(H/X)^{1/2}$. The embeddings $L^2\hookrightarrow L^2(\sigma)$ and its adjoint have norm $\le \sqrt{2}$ (from the definitions of $\rho,\kappa$), so \eqref{eq:BG.2to2} follows.
\end{proof}

\begin{corollary}[Short–shift anti–alignment]\label{cor:BG.AA}
Let $A_k$ be as in \eqref{eq:Ak.def}. Then
\begin{equation}\label{eq:BG.final}
\frac{1}{H}\sum_{0<|\Delta|<H}\ \sum_k |A_{k+\Delta}\,\overline{A_k}|
\ \ll\ (\log X)^C\,\Big(\frac{H}{X}\Big)^{1/2}\ \sum_k |A_k|^2.
\end{equation}
\end{corollary}

\begin{proof}
Combine \eqref{lem:BG.degeneracy} with Theorem~\ref{thm:BG-nonflat}.
\end{proof}

\subsection{Flattened case (for comparison)}\label{subsec:BG.flat}

For completeness we note the flat–coefficient variant used earlier; it is a direct consequence of Theorem~\ref{thm:BG-nonflat} (taking $F=\alpha\!\otimes\!\beta$) but can also be proved by two Cauchy–Schwarz placements and SSU on each tube.

\begin{proposition}[BG under flatness]\label{prop:BG.flat}
Assume $\|\alpha\|_\infty\ll D^{-1/2}$ and $\|\beta\|_\infty\ll N^{-1/2}$. Then \eqref{eq:BG.final} holds with the same right–hand side.
\end{proposition}

\subsection{Parameters and endpoint behavior}\label{subsec:BG.params}

The only relationship among parameters used is $1/H \ll 1/Q^2$ (disjoint bank arcs) and $q\le Q=H^\gamma$ with $\gamma<\tfrac12$ (so $U\ll H$ in SSU). All $q$–dependence contributes at most $(\log X)^{O(1)}$. The bound \eqref{eq:BG.2to2} is uniform as $\gamma\to \tfrac12^-$, with constants blowing up at most polylogarithmically.

\subsection*{Two–weight Sawyer testing (fixed weights)}

Throughout this section $D,N$ are dyadic with $X^\varepsilon\le D,N\le X^{1-\varepsilon}$ and $DN\asymp X$, and $W\in C_c^\infty([1/2,2]^2)$ is the fixed bump from §2. Let $K_H$ be the positive, band-limited kernel from §5.1 (bandwidth $1/H$, $\int K_H=H+O(1)$, $\| \widehat K_H\|_1\ll H$). Define the positive operator
\[
(TF)(d',n') := \sum_{d,n} K_H\!\big((d,n),(d',n')\big)\,F(d,n).
\]

\paragraph{Geometry–only testing weights.}
Let $W\in C_c^\infty([1/2,2]^2)$ be the fixed dyadic bump from §2, and set
\[
\rho_d := \sum_{n} W\!\left(\frac dD,\frac nN\right)^2,\qquad 
\kappa_n := \sum_{d} W\!\left(\frac dD,\frac nN\right)^2,\qquad 
\sigma(d,n):=\rho_d+\kappa_n .
\]
Since $W$ is supported on $[D,2D]\times[N,2N]$ and bounded above/below on its support,
\[
c_1(D+N)\ \le\ \sigma(d,n)\ \le\ C_1(D+N)\qquad\text{on }[D,2D]\times[N,2N],
\]
so
\begin{equation}\label{eq:L2-sigma-equiv}
\|G\|_{L^2(\sigma)}^2\asymp (D+N)\,\|G\|_{2}^2,\qquad \sigma(E)\asymp (D+N)\,|E|.
\end{equation}

These depend only on $(D,N,W)$, not on $F$. Because $W$ is supported on $[D,2D]\times[N,2N]$ and bounded above/below on its support, we have
\[
c_0\,\min(D,N)\ \le\ \rho_d,\kappa_n\ \le\ C_0\,\max(D,N)\qquad\text{for all }d,n,
\]
hence
\[
c_1(D+N)\ \le\ \sigma(d,n)\ \le\ C_1(D+N)\qquad\text{on the box }[D,2D]\times[N,2N].
\]
Therefore $L^2$ and $L^2(\sigma)$ are equivalent up to an absolute constant:
\[
\|G\|_{L^2(\sigma)}^2 \asymp (D+N)\,\|G\|_2^2.
\]

\paragraph{Sawyer tests.}
Let $\Omega=[D,2D]\times[N,2N]$. There exists $M\ll (\log X)^C\,H/X$ such that for every measurable $E\subset\Omega$,
\begin{equation}\label{eq:sawyer-tests}
\|T 1_E\|_{L^2(\sigma)}^2 \ \le\ M\,\sigma(E), 
\qquad 
\|T^\ast 1_E\|_{L^2(\sigma)}^2 \ \le\ M\,\sigma(E).
\end{equation}

\begin{proof}
We prove the forward testing inequality; the backward one is identical since $K_H$ is even and $T$ is positive/self–adjoint on $\ell^2$.

\emph{Step 1 (tube decomposition and bounded overlap).}
Let $\Omega=[D,2D]\times[N,2N]$ and let $\mathcal{T}$ be the fixed sheared–tube partition used in §5: each $U\in\mathcal{T}$ has short axis $\ll H$ and long axis $\asymp X/H$, and the fattenings $U^{+}$ have overlap $\ll(\log X)^C$. Write
\[
T=\sum_{U\in\mathcal{T}} T_U,\qquad (T_U F)(y):=\sum_{x\in U} K_H\!\big(B(x,y)\big)\,F(x),
\]
so $T_U$ is supported on $U\times U^{+}$ and $M_H\ge 0$.

\emph{Step 2 (reduce $L^2(\sigma)$ to unweighted $L^2$).}
By the geometry–only definition of $\sigma(d,n)=\rho_d+\kappa_n$ above,
\[
c_1(D+N)\ \le\ \sigma(d,n)\ \le\ C_1(D+N)\quad \text{on }\Omega,
\]
whence for all $G$ supported on $\Omega$ we have $\|G\|_{L^2(\sigma)}^2\asymp (D+N)\|G\|_2^2$ and $\sigma(E)\asymp (D+N)\,|E|$ for any $E\subset\Omega$.

\emph{Step 3 (one–tube operator norm from the SSU bound).}
The one–tube short–shift estimate proved in §5.2 (display (5.8)) gives, for every $f$ supported on $U$,
\[
\langle T_U f,f\rangle\ \ll\ (\log X)^C\Big(\frac{H}{X}\Big)^{1/2}\|f\|_2^2.
\]
Since $T_U$ is positive and self–adjoint on $\ell^2$, this implies $\|T_U\|_{2\to 2}\ll (\log X)^C(H/X)^{1/2}$.

\emph{Step 4 (sum over tubes).}
Let $g_E:=1_E$. Using triangle inequality and Cauchy–Schwarz together with bounded tube overlap,
\[
\|T g_E\|_2^2\ \ll\ (\log X)^C\Big(\frac{H}{X}\Big)\sum_U \|g_E 1_U\|_2^2
\ \ll\ (\log X)^C\Big(\frac{H}{X}\Big)\,|E|.
\]

\emph{Step 5 (return to the weighted norm and conclude the forward test).}
By Step~2,
\[
\|T1_E\|_{L^2(\sigma)}^2\ \asymp\ (D+N)\,\|T1_E\|_2^2
\ \ll\ (\log X)^C\Big(\frac{H}{X}\Big)\,(D+N)\,|E|
\ \asymp\ (\log X)^C\Big(\frac{H}{X}\Big)\,\sigma(E).
\]
This is the forward testing inequality. The backward test for $T^\ast$ is identical because $T^\ast=T$. The Sawyer two–test criterion then yields
\[
\|T\|_{L^2(\sigma)\to L^2(\sigma)}\ \ll\ (\log X)^C\Big(\frac{H}{X}\Big)^{1/2},
\]
and, using $L^2(\sigma)\asymp L^2$ on $\Omega$, the same bound in the unweighted norm.
\end{proof}

\paragraph{Operator norm.}
By the two–test criterion for positive kernels, \eqref{eq:sawyer-tests} yields
\[
\|T\|_{L^2(\sigma)\to L^2(\sigma)} \ \ll\  (\log X)^C\Big(\frac{H}{X}\Big)^{1/2}.
\]
Using $L^2\!\asymp L^2(\sigma)$ on $\Omega$ gives the unweighted bound
\[
\boxed{\ \ \|T\|_{2\to 2} \ \ll\  (\log X)^C\Big(\frac{H}{X}\Big)^{1/2}\ }\!,
\]
which is the BG short–shift gain needed for TT$^\ast$ without any coefficient flatness assumption.


\subsection{Polylogarithmic ledger (bank--local form)}\label{sec:polylog-ledger}
Fix \(A\ge 10\) and \(0<\gamma<\tfrac12\), and set \(H=(\log X)^A, Q=H^\gamma\).
Let \(S_{n_0}\) denote the band-limited extractor–based statistic
\(
S_{n_0}:=\sum_{n} R_2(n) J_H(n_0-n),
\)
with \(\widehat J_H\subset[-1/H,1/H]\) and \(J_H\ge 0\), and impose bank localization only via the frequency projector \(P_A\). The inputs are:

\begin{itemize}
\item \textbf{Major arcs (uniform mean).} On the bank $\mathfrak A(Q)$ we have
\[
\mathbb E_{n_0}S_{n_0}\ \ge\ \frac{c_{\rm SS}}{(\log X)^2}\,,
\]
uniformly for even $N\in[2X,4X]$ (Cor.~10.5 and uniform positivity of the singular series).

\item \textbf{Minor arcs (variance).} From SSU and Type--I with AO we have
\begin{equation}\label{eq:ledger-variance}
\mathrm{Var}(S_{n_0})
\;\le\;
C_2\Big(\frac{H}{X}\Big)^{1/2}
\;+\;
C_3\,\frac{1}{H\,Q^2}.
\tag{5.15}
\end{equation}
where the first term is Theorem~\ref{thm:SSU-Sawyer} and the second is Corollary~\ref{cor:TypeI-bank}; both are bank--local and unconditional at polylogarithmic moduli.

\item \textbf{Closure (no $\theta$).} Paley--Zygmund yields positivity provided
\begin{equation}\label{eq:closure-bank}
C_2\Big(\tfrac{H}{X}\Big)^{1/2}\ +\ \frac{C_3}{H\,Q^2}\ \le\ \frac{c_{\rm SS}^2}{8}.
\end{equation}
With $H=(\log X)^A$ and $Q=H^\gamma$, the left side of \eqref{eq:closure-bank} equals 
$C_2(\log X)^{A/2}X^{-1/2}+C_3(\log X)^{-A(1+2\gamma)}$, which $\to 0$ as $X\to\infty$ for any fixed $0<\gamma<\tfrac12$. 
Hence there exists an explicit $X_0(A,\gamma)$ such that \eqref{eq:closure-bank} holds for all $X\ge X_0(A,\gamma)$.
\end{itemize}

\begin{proposition}[Audit inequality for the band-limited extractor]\label{prop:audit}
Let $S_{n_0}=\sum_{n} R_2(n)\,J_H(n_0-n)$ with $J_H\ge0$ the Fejér extractor of bandwidth $1/H$ and $\sum_n J_H(n)=1$.
Let $E:=R_2-M$ denote the error term after removing the projected major-arc main term $M$.
Then, uniformly for centers $n_0\in[X,2X]$,
\begin{equation}\label{eq:audit}
\operatorname{Var}(S_{n_0})
\ \le\
\kappa_2\,(\log X)^{C_2}\,\frac{1}{H\,Q^2}
\ +\
\kappa_3\,(\log X)^{C_3}\,\Big(\frac{H}{X}\Big)^{1/2}
\ +\ O(H^{-2}).
\end{equation}

Here the first term is the Type–I contribution after alias suppression/AO and normalization, and the second term is the BG/SSU (short–shift) contribution coming from the $L^2$ operator bound for the normalized SSU operator.
\end{proposition}

\begin{proof}
Write $S_{n_0}-\mathbb E[S_{n_0}]=\sum_n E(n)\,J_H(n_0-n)$, apply Parseval on the dyadic block and the fact that $\widehat{J_H}$ is supported in $|\xi|\le 1/H$ with $\|\widehat{J_H}\|_{L^1}\asymp 1$, to reduce to an energy bound on $\widehat{E}$ over the bank. For Type–I, the TT$^\ast$ ledger with balanced bank windows gives
\[
\mathbb E_b\!\int |\widehat{E_b}(\xi)|^2\,\widehat{w}^{\ast}_{\tau}(\xi)\,d\xi
\ \ll\ X^{1+o(1)}/(H Q^2),
\]
(see §7.3), and divisibility/AO dispersion allows a deterministic choice of $b$ with the same order. Dividing by $X$ when passing from global dyadic energy to per–center variance yields the first term of \eqref{eq:audit}. For BG/SSU, the normalized short–shift operator $T$ satisfies
\[
\|T\|_{2\to2}\ \ll\ (\log X)^C\,\big(H/X\big)^{1/2}
\]
(cf. (6.2)–(6.3)), which contributes the second term. The $O(H^{-2})$ is the Fejér-smoothing remainder. This proves \eqref{eq:audit}.
\end{proof}


\subsection{What BG contributes to TI (bank--local closure form)}
\begin{proposition}[Ledger inputs with no $\theta$]\label{prop:BG-AO-ledger}
Fix $H=(\log X)^A$ and $Q=H^\gamma$ with $0<\gamma<\tfrac12$. Then:
\begin{align}
\text{\emph{(SSU)}}\quad &
\mathbb E_{n_0\in[X,2X]}\int_{\mathfrak A(Q)^c}\!\!|S_{n_0}(\xi)|^2\,d\xi
\ \ll\ (\log X)^C\,\Big(\tfrac{H}{X}\Big)^{1/2}\,\mathcal E, \\
\text{\emph{(Type--I on the bank)}}\quad &
\mathrm{Var}_{\mathrm I}\ \ll\ (\log X)^C\,\frac{X}{H\,Q^2}\,\mathcal E,\\
\text{\emph{(AO dispersion)}}\quad &
\text{no additional parameter is needed; see Cor.\,\ref{cor:AO-occupancy}.}
\end{align}
Here $\mathcal E$ is the coefficient energy from the F3 decomposition. No level-of-distribution $\theta$ is assumed or used.
\end{proposition}

\begin{corollary}[Bank–local closure inequality]\label{cor:bank-closure}
Let $(P_{\mathfrak A}M)(N)\ge \dfrac{c_{\rm SS}}{(\log X)^2}$ uniformly on even $N\in[2X,4X]$ (Cor.\ 10.5).
If for some absolute $C_2,C_3\ge1$,
\[
C_2\Big(\tfrac{H}{X}\Big)^{1/2}\;+\;\frac{C_3}{H\,Q^2}\ \le\ \frac{c_{\rm SS}^2}{8},
\]
then the $TT^\ast$ ledger closes. With $H=(\log X)^A$ and $Q=H^\gamma$ this holds for \emph{any fixed} $0<\gamma<\tfrac12$ once $X\ge X_0(A,\gamma)$.
\end{corollary}


\bigskip
\noindent\textbf{Guide to dependencies.}
The proof uses only:
\begin{enumerate}[label=(\roman*)]
\item The bandlimited kernel $K_H$ \eqref{eq:KH.assumptions};
\item Shearing identity \eqref{eq:BG.factor} and tube pack \eqref{eq:BG.tube};
\item SSU (Theorem~\ref{thm:BG.SSU});
\item Incidence (Lemma~\ref{lem:BG.incidence}) and degeneracy (Lemma~\ref{lem:BG.degeneracy});
\item Sawyer testing (Lemma~\ref{lem:BG.Sawyer}) for the positive operator $\mathbf T$.
\end{enumerate}
No flatness assumption is required anywhere.
\qedhere
                 % BG$^{+}$ core (incidence + shearing + nonflat)
% ==========================================================
% Section 05b — Short–shift / Bilinear Geometry (SSU): Repaired
% (Deterministic; no packets; no F–dependent weights)
% ==========================================================
\section{Bilinear geometry and short–shift bound (SSU)}\label{sec:ssu}

Throughout, the bank $\cA\subset\mathbb T$ is the disjoint union of $m\asymp Q^2$ major arcs,
each of width $\asymp 1/H$ and mutual gaps $\gg 1/H$, with $Q=H^\gamma$ and $0<\gamma<\tfrac12$.
Let $\wh\Phi_H$ denote a fixed nonnegative Fejér window supported on $|\xi|\lesssim 1/H$, and let
$\{\mathcal D_j\}_{j\ge 0}$ be the standard dyadic partition of the band $|\xi|\le c/H$,
$\mathcal D_j=\{\xi:2^{-j-1}/H<|\xi|\le 2^{-j}/H\}$, with $W_j:=\sum_{\xi\in\mathcal D_j}|\wh\Phi_H(\xi)|$.
We use only the basic facts $\sum_j W_j\ll 1$ and $\sum_j 2^{j/2}W_j\ll \log H$.

\begin{definition}[Admissible F3]\label{def:F3}
\[
f(n)=\sum_{\substack{ab=n\\A<a\le2A\\B<b\le2B}}\alpha(a)\beta(b)u(\tfrac{a}{A},\tfrac{b}{B}),\quad AB\asymp X,\ 
|\alpha(a)|,|\beta(b)|\ll\tau(\cdot)^C,\ u\in C_c^\infty([1,2]^2).\]
\end{definition}

For finite X, the complementary region is bounded trivially by the crude $L^2$ large–sieve bound; its contribution is absorbed into the polylog factor.

\subsection{Tube partition and fixed projectors}
For each dyadic shell $\mathcal D_j$ we fix a finite family of smooth frequency cutoffs
$\{P_{U_{j,k}}\}_k$ supported in sheared ``tubes'' $U_{j,k}\subset\cA\cap\mathcal D_j$, with the
following uniform properties:

\begin{itemize}
\item[(T1)] \emph{Partition with bounded overlap:} $\sum_{j,k} P_{U_{j,k}}(\xi)\asymp \mathbf 1_{\cA}(\xi)$ and
$\sup_{\xi\in\cA}\sum_{j,k} P_{U_{j,k}}(\xi)\ \ll\ (\log X)^C$.
\item[(T2)] \emph{Dyadic thickness:} $\operatorname{diam}(\operatorname{supp} P_{U_{j,k}})\asymp 2^{-j}/H$ across the short direction; tangential extent is $\asymp 1$ along the arc.
\item[(T3)] \emph{Bank taper stability:} $0\le P_{U_{j,k}}\le 1$, and $P_{U_{j,k}}$ is dominated by the bank Fejér taper on $\cA$.
\end{itemize}

All projectors $P_{U_{j,k}}$ are fixed by the geometry; \emph{no} projector depends on the input function.

\begin{lemma}[Fixed–$j$ multiplicity]\label{lem:fixed-j}
For each $j$, the bump family $\{\chi_j\cdot \vartheta_I\}_I$ has bounded overlap:
\[
\sup_{\xi\in\mathbb T}\ \sum_{I} \chi_j(\xi)\,\vartheta_I(\xi)\ \ll\ 1.
\]
\end{lemma}

\begin{proof}
For fixed $j$, the tangential windows $\vartheta_I$ cover disjoint (or $O(1)$–overlap) intervals of length $\asymp (qH)^{-1}$ with separation $\gg (qH)^{-1}$ by the guard–band. Hence at each $\xi$ only $O(1)$ indices $I$ contribute once $\chi_j(\xi)\ne 0$ fixes $q$ to a dyadic.
\end{proof}

\begin{lemma}[Sum over $j$]\label{lem:sum-j}
Let $\alpha_j\ge 0$ with $\sum_j 2^{-\alpha j}\ll (\log X)^C$. Then
\[
\Lambda\ :=\ \sup_{\xi\in\mathfrak A^c}\ \sum_{j,I} \chi_j(\xi)\,\vartheta_I(\xi)\ \ll\ (\log X)^C.
\]
\end{lemma}

\begin{proof}
By Lemma~\ref{lem:fixed-j}, at fixed $j$ the sum over $I$ is $O(1)$, and the $j$–sum is supported on $O(\log X)$ dyadics by $q\le Q=H^\gamma$ and $H=(\log X)^A$. The mild weights $2^{-\alpha j}$ (coming from the normal decay of $\chi_j$) make the total $\ll(\log X)^C$.
\end{proof}

\begin{remark}[SSU route used in the ledger]\label{rem:route}
In the endgame we use the \emph{operator} route: define
\[
(Tf)(n)\ :=\ \sum_{\nu,s} \big(K_{\nu,s}\ast (P_{U_{\nu,s}} f)\big)(n),
\qquad \|T\|_{2\to 2}\ \ll\ (\log X)^C\Big(\frac{H}{X}\Big)^{1/2},
\]
proved via bounded overlap (Lemma~\ref{lem:tube-overlap}) and Hausdorff--Young. The pointwise tube form is maintained in Appendix~\ref{app:BS} for reference; both yield the same numerical exponent but we cite the operator bound in §11.
\end{remark}

% ===== BEGIN PATCH: Single-tube SSU =====
\begin{theorem}[Single–tube SSU]\label{thm:SSU-single}
Let $T:=T(a/q,s)$ be a tube as in \eqref{eq:def-tube} with $1\le q\le Q$, $1\le U\ll X/(qH)$.
Let $K_H$ be any band–limited kernel with $\supp\widehat K_H\subset[-1/H,1/H]$ and
\[
\int \widehat K_H\ll H,\qquad \int \frac{\widehat K_H}{|\xi|}\ll H\log H.
\]
Then for any function $F$ supported on $T$,
\begin{equation}\label{eq:SSU-single}
\sum_{(d,n),(d',n')\in T} F(d,n)\,K_H(d'n-dn')\,\overline{F(d',n')}
\ \ll\ (\log X)^C\,\Big(\frac{H}{X}\Big)^{1/2}\ \sum_{(d,n)\in T}|F(d,n)|^2.
\end{equation}
\end{theorem}

\begin{proof}
Write $u=qn-ad$ and $u'=qn'-ad'$. On $T$ we have $|u-s|,|u'-s|\le U$.
Fourier–invert $K_H$ on $|\xi|\le 1/H$:
\[
K_H(d'n-dn')=\int_{|\xi|\le 1/H}\widehat K_H(\xi)\,e\!\left(\xi\frac{d'n-dn'}{X}\right)\,d\xi.
\]
The bilinear form equals
\[
\int_{|\xi|\le 1/H}\widehat K_H(\xi)\ \Big|\sum_{(d,n)\in T}F(d,n)\,e\!\left(\xi\,\frac{qn-ad}{qX}\cdot d\right)\Big|^2 d\xi,
\]
after algebra with $u=qn-ad$ and the identity $d'n-dn'=(u'v-uv')/q$ in the shear $(u,v)$.
Split the sum over residue classes modulo $q$ in $v=d$ and $u\equiv -av\!\!\pmod q$; smooth weights $W(d/D,n/N)$ from \S\ref{sec:TFA} are inserted (and equal to $1$ on $T$) to allow summation by parts and Poisson with $O((\log X)^{-A})$–losses absorbed into $(\log X)^C$.

For each fixed residue class the inner exponential has phase $\xi\,(u\,v)/(qX)$, with $u=s+O(U)$ and $v\asymp D$. If $(u,v)\neq (u',v')$ in the same tube, the frequency gap satisfies
\[
\bigg|\xi\frac{uv-u'v'}{qX}\bigg|\ \ge\ c\,|\xi|\ \frac{\min\{|u-u'|\,D,\ U\,|v-v'|\}}{qX}
\ \gg\ \frac{1}{qXH}\qquad(|\xi|\asymp 1/H),
\]
since either $|u-u'|\ge 1$ or $|v-v'|\ge 1$.
Thus the family of frequencies is $c/(qXH)$–separated on $|\xi|\asymp 1/H$ and, by dyadic decomposition $\{2^{-j-1}/H<|\xi|\le 2^{-j}/H\}$, the Montgomery–Vaughan large sieve yields (with the smooth weights)
\[
\int_{2^{-j-1}/H}^{2^{-j}/H}\!\!\widehat K_H(\xi)\ \Big|\sum_{(d,n)\in T}F(d,n)e(\cdots)\Big|^2 d\xi
\ \ll\ \Big(X\ +\ \frac{H}{2^{j}}\Big)\ \sum_{T}|F|^2.
\]
Take geometric mean over the two dyadic families produced by exchanging the roles of $(d,n)$ and $(d',n')$, then sum $j\ge 0$, and use the moment bounds on $\widehat K_H$ to obtain \eqref{eq:SSU-single}. Finally, since $U\ll X/(qH)$, the $H/2^j$–piece collapses to $(H/X)^{1/2}$ after optimization over $j$.
\end{proof}
% ===== END PATCH =====

\begin{lemma}[Two one–sided bounds imply square–root]\label{lem:interp-gm}
Let $S$ be a finite index set. Suppose for each $s\in S$ we have operators $U_s,V_s$ on $\ell^2$ satisfying
\[
\Big\|\sum_{s\in S} U_s f\Big\|_2\ \le A\,\|f\|_2,\qquad
\Big\|\sum_{s\in S} V_s f\Big\|_2\ \le B\,\|f\|_2,
\]
and for every $s$ the kernel of $W_s:=U_s^{1/2}V_s^{1/2}$ is pointwise bounded by the geometric mean of the kernels of $U_s$ and $V_s$. Then
\[
\Big\|\sum_{s\in S} W_s f\Big\|_2\ \le (AB)^{1/2}\,\|f\|_2.
\]
\end{lemma}

\begin{proof}
For $f,g\in\ell^2$,
\[
\big|\langle \sum_s W_s f,\ g\rangle\big|
\le \Big(\sum_s \langle U_s f,f\rangle\Big)^{1/2}\Big(\sum_s \langle V_s g,g\rangle\Big)^{1/2}
\le A^{1/2}B^{1/2}\,\|f\|_2\|g\|_2.
\]
Take $\|g\|_2=1$.
\end{proof}

% ===== BEGIN PATCH: Two-weight Sawyer test =====
\subsection{Global SSU via two–weight Sawyer testing}\label{subsec:sawyer}

We will invoke Corollary~\ref{cor:AO-to-SSU} to replace packet sums by bank energies when aggregating over $(\nu,s)$.

\paragraph{Operator and decomposition.}
Let $T$ be the SSU operator defined in \eqref{eq:def-T}; we write
\[
T\ =\ \sum_{\nu}\,T_\nu\ =\ \sum_{\nu}\ \sum_{s\in\mathbb Z} T_{\nu,s},
\]
where $T_{\nu,s}$ are the short-shift packets (shear $a_\nu/q_\nu$, center shift $s$) with bounded overlap in $(\nu,s)$ (cf.\ Lemma~\ref{lem:fixed-j}). Our objective is the global estimate
\begin{equation}\label{eq:sawyer-goal}
\|T\|_{2\to 2}\ \ll\ (\log X)^C\,\Big(\frac{H}{X}\Big)^{1/2}.
\end{equation}

Let $\{T_\nu\}$ be the tube partition from Lemma~\ref{lem:tubes-partition}, and let
\[
(Tf)(n)\ :=\ \Big(\frac{H}{X}\Big)^{1/2}\sum_{\nu}\ \sum_{(d,n')\in T_\nu} K_H(dn'-dn)\,f(n')\,.
\]

\begin{proposition}[Sawyer reduction]\label{prop:sawyer-reduction}
Let $T$ be as in \eqref{eq:def-T} or \S\ref{subsec:sawyer} and suppose that for every finite $E\subset\{n\sim X\}$,
\[
\|T1_E\|_2^2\ \le\ C_0\,\frac{H}{X}\,\#E,\qquad
\|T^\ast 1_E\|_2^2\ \le\ C_0\,\frac{H}{X}\,\#E.
\]
Then $\|T\|_{2\to2}\ \le\ C\,\sqrt{C_0}\,(H/X)^{1/2}$, with $C$ absolute.
\end{proposition}

\begin{proof}
Standard two–weight testing (Sawyer). Write $\|Tf\|_2^2=\sup_{\|g\|_2=1}\sum_k \langle f_k, T1_{E_k}\rangle\,\langle 1_{E_k},Tg_k\rangle$ using level-set (CZ) decompositions of $f,g$.
Apply the two tests to each $E_k$ and sum the resulting geometric series; positivity of $T$ (or symmetric packetization) yields the stated bound.
\end{proof}

\begin{proof}
Expand
\[
\|T1_E\|_2^2=\sum_{\nu}\|T_\nu 1_E\|_2^2+\sum_{\nu\ne\nu'}\langle T_\nu 1_E,\ T_{\nu'}1_E\rangle.
\]
For the diagonal, the single–tube SSU bound (proved above for each tube/packet) gives
\(
\|T_\nu 1_E\|_2^2\ll (H/X)\,\#(E\cap \mathrm{proj}(T_\nu)).
\)
Summing $\nu$ and using bounded overlap (Lemma~\ref{lem:fixed-j} and (T1)) yields $\ll (\log X)^C(H/X)\#E$.
For the off–diagonal, Cauchy–Schwarz and the packet decay in the short shift (in $|s-s'|$) and slope separation (distance between $\nu,\nu'$)
bound
\(
\sum_{\nu\ne\nu'}|\langle T_\nu 1_E,T_{\nu'}1_E\rangle|
\ll (\log X)^C(H/X)\#E.
\)
Combine the two pieces.
\end{proof}

\begin{proof}
Expand $\|T1_E\|_2^2$ and decompose the sum into $\nu,\nu'$. By Lemma~\ref{lem:tubes-partition}, each point belongs to $(\log X)^C$ tubes, hence the off–diagonal $(\nu\ne \nu')$ contributes at most $(\log X)^C$ times the maximal single–tube interaction, which is $\ll (H/X)\#E$ by Theorem~\ref{thm:SSU-single}. The diagonal is identical. Summing yields the claim.
\end{proof}

\begin{proposition}[Two–weight Sawyer test]\label{prop:SSU-sawyer}
With $T$ as above and overlap constant $\Lambda\ll(\log X)^C$ from Lemma~\ref{lem:tubes-partition},
\[
\|T\|_{2\to2}\ \ll\ (\log X)^C\,\Big(\frac{H}{X}\Big)^{1/2}.
\]
\end{proposition}

\begin{proof}
Apply Cotlar–Stein in the tube index, using $\|T_\nu\|_{2\to2}\ll (H/X)^{1/2}$ from Theorem~\ref{thm:SSU-single} and the overlap bound $\Lambda\ll(\log X)^C$; equivalently, combine Lemma~\ref{lem:testing-1E} with the dual testing inequality and invoke Sawyer’s criterion.
\end{proof}
% ===== END PATCH =====

\subsection{The normalized SSU operator}
Let $K_{j,k}$ be the (signed) time–domain kernel with multiplier $\wh K_{j,k}=P_{U_{j,k}}\cdot \wh\Phi_H$.
Define the normalized short–shift operator $T$ on $\ell^2([X,2X])$ by
\begin{equation}\label{eq:def-T}
(Tf)(n)\ :=\ \Big(\frac{H}{X}\Big)^{\!1/2}\ \sum_{j,k}\ \big(K_{j,k}\ast (P_{U_{j,k}}f)\big)(n).
\end{equation}
The scaling factor $(H/X)^{1/2}$ is the physical normalization corresponding to a short–shift of scale $1/H$
acting on a window of length $\asymp X$.

\begin{lemma}[Testing on $1_E$]\label{lem:testing-1E}
Let $T$ be as in \eqref{eq:def-T}. For every finite $E\subset\{n\sim X\}$,
\[
\|T1_E\|_2^2\ \ll\ (\log X)^C\,\frac{H}{X}\,\#E,
\qquad
\|T^\ast 1_E\|_2^2\ \ll\ (\log X)^C\,\frac{H}{X}\,\#E.
\]
\end{lemma}

\begin{proof}
Expand
\[
\|T1_E\|_2^2
=\sum_{\nu,s}\|T_{\nu,s}1_E\|_2^2+\sum_{(\nu,s)\neq(\nu',s')} \langle T_{\nu,s}1_E,\ T_{\nu',s'}1_E\rangle.
\]
\emph{Diagonal:} For each packet, the single–tube SSU bound gives
\(
\|T_{\nu,s}1_E\|_2^2\ll (H/X)\,\#(E\cap \operatorname{proj}(T_{\nu,s})).
\)
Summing over $(\nu,s)$ and using bounded overlap (Lemma~\ref{lem:fixed-j}) yields $\ll (\log X)^C(H/X)\#E$.
\emph{Off–diagonal:} By Cauchy–Schwarz and the kernel-product bound (Proposition~\ref{prop:cotlar-ssu}),
\[
\sum_{(\nu,s)\neq(\nu',s')}
|\langle T_{\nu,s}1_E,\,T_{\nu',s'}1_E\rangle|
\ \ll\ \frac{H}{X}\!\sum_{(\nu,s)}\sum_{(\nu',s')}
\frac{\#E}{(1+|s-s'|)^2(1+\operatorname{dist}(\nu,\nu')H)^2}.
\]
The double sum is $\ll(\log X)^C$ by Lemma~\ref{lem:incdeg-sawyer} (proved below), giving the desired bound.
The adjoint case is identical since $K_H$ is even (or by symmetric packetization).
\end{proof}

\begin{corollary}[Sawyer tests $\Rightarrow$ global norm]\label{cor:sawyer-close}
The inequalities of Lemma~\ref{lem:testing-1E} and Proposition~\ref{prop:sawyer-reduction} imply
\[
\|T\|_{2\to 2}\ \ll\ (\log X)^C\,\Big(\frac{H}{X}\Big)^{1/2}.
\]
\end{corollary}

\begin{remark}[Uniformity of constants]\label{rem:sawyer-uniform}
Implicit constants depend only on (i) finitely many derivatives of the dyadic windows; 
(ii) the kernel moments $\int \widehat K_H\ll H$ and $\int \widehat K_H/|\xi|\ll H\log H$; 
(iii) the overlap constant in Lemma~\ref{lem:fixed-j}, bounded by $(\log X)^C$ via Lemmas~\ref{lem:BG.incidence}–\ref{lem:BG.degeneracy}; 
and (iv) the choice $Q=H^\gamma$ with $0<\gamma<\tfrac12$. 
All bounds are uniform in $(a_\nu,q_\nu)$, $q_\nu\le Q$, and in short shifts $s$ with packet width $U\asymp X/(q_\nu H)$.
\end{remark}


\begin{proof}
Apply Proposition~\ref{prop:sawyer-reduction} with the forward and adjoint tests from Lemma~\ref{lem:testing-1E}.
\end{proof}

\begin{lemma}[Uniform $L^2$ operator bound per tube]\label{lem:tube-L2}
For all $j,k$,
\(
\|K_{j,k}\ast g\|_{2}\ \le\ \|\wh K_{j,k}\|_{\infty}\,\|g\|_2
\ \le\ C\,(\log X)^{C}\,\|g\|_2,
\)
with $C$ depending only on the fixed bank taper.
\end{lemma}

\begin{proof}
Hausdorff–Young gives $\|K\ast g\|_2\le \|\wh K\|_\infty\|g\|_2$. Since $0\le P_{U_{j,k}}\le 1$ and
$\wh\Phi_H\le C(\log X)^C$ on its support, the bound follows.
\end{proof}

\begin{lemma}[TT$^\ast$ with bounded overlap]\label{lem:TTstar}
Let $T$ be given by \eqref{eq:def-T} and suppose (T1) holds with overlap constant $\Lambda\ll(\log X)^C$.
Then
\[
\|T\|_{2\to 2}\ \le\ \Big(\frac{H}{X}\Big)^{\!1/2}\, \Big(\sup_{j,k}\|\wh K_{j,k}\|_\infty\Big)\,\Lambda.
\]
\end{lemma}

\begin{proof}
For $f\in\ell^2$,
\[
\|Tf\|_2
\ \le\ \Big(\frac{H}{X}\Big)^{\!1/2}\sum_{j,k}\ \|K_{j,k}\ast (P_{U_{j,k}}f)\|_2
\ \le\ \Big(\frac{H}{X}\Big)^{\!1/2} \Big(\sup_{j,k}\|\wh K_{j,k}\|_\infty\Big)\sum_{j,k}\|P_{U_{j,k}}f\|_2.
\]
By Cauchy–Schwarz and (T1),
\(
\sum_{j,k}\|P_{U_{j,k}}f\|_2
\le \big(\sum_{j,k}\langle P_{U_{j,k}}f,f\rangle\big)^{1/2}\, \big(\sum_{j,k}\|P_{U_{j,k}}f\|_2\big)^{1/2}
\ll \Lambda\,\|f\|_2.
\)
\end{proof}

\begin{proposition}[Normalized TT$^\ast$ bound]\label{prop:ssu-ttstar}
With (T1)–(T3) and Lemma~\ref{lem:tube-L2},
\[
\|T\|_{2\to2}\ \ll\ (\log X)^C\,\Big(\frac{H}{X}\Big)^{\!1/2}.
\]
\end{proposition}

\begin{proof}
Combine Lemma~\ref{lem:tube-L2} and Lemma~\ref{lem:TTstar}, using $\Lambda\ll(\log X)^C$ from (T1).
\end{proof}

\begin{theorem}[Two–axis SSU with square–root saving]\label{thm:SSU-sqrt}
Let $K\ge 0$ be a pin of width $H=(\log X)^A$ with $\supp\widehat K\subset[-1/H,1/H]$ and kernel moments $\int \widehat K\ll H$, $\int |\xi|^{-1}\widehat K\ll H\log H$. 
Let $S_{n_0}(\xi)=\sum_{n\sim X} f(n)\,K(n_0-n)\,e(\xi n)$ where $f$ is decomposed by the F3 reduction into $O((\log X)^C)$ Type~II/III blocks $\mathsf B$ with coefficient arrays $F_{\mathsf B}$.
Then
\[
\mathbb E_{n_0\in[X,2X]}\ \int_{\mathfrak A(Q)^c}\! |S_{n_0}(\xi)|^2\,d\xi
\ \ll\ (\log X)^C\ \Big(\frac{H}{X}\Big)^{1/2}\ \sum_{\mathsf B}\ \|F_{\mathsf B}\|_2^2.
\]
The implied constant is absolute (depending on $A$ only). 
\end{theorem}

\subsection{Statement used in the ledger}
We record the global SSU estimate in the form used by the TT$^\ast$ ledger.

\begin{theorem}[Repaired BG/SSU]\label{thm:BG-repaired}
Let $T$ be the normalized short–shift operator \eqref{eq:def-T} associated with the bank $\cA$, dyadic
partition $\{\mathcal D_j\}$, Fejér taper $\wh\Phi_H$, and tube projectors $\{P_{U_{j,k}}\}$ satisfying \emph{(T1)–(T3)}.
Then
\begin{equation}\label{eq:ssu-final}
\|T\|_{2\to2}\ \ll\ (\log X)^C\,\Big(\frac{H}{X}\Big)^{\!1/2}.
\end{equation}
The implied constant depends only on the fixed smooth bank taper and the overlap constant in \emph{(T1)}.
\end{theorem}

\paragraph*{Remarks.}
(1) No randomness, packets, or $f$–dependent weights are used. The proof is a deterministic TT$^\ast$ argument with fixed projectors and the Hausdorff–Young bound on each tube multiplier.  
(2) If one prefers a Sawyer two–weight presentation, take both weights equal to the fixed overlap weight
$\sigma(\xi)=\sum_{j,k} P_{U_{j,k}}(\xi)$; the testing inequalities follow from (T1), yielding \eqref{eq:ssu-final}.  
(3) The normalization $(H/X)^{1/2}$ is explicit in \eqref{eq:def-T}; it is the physical short–shift scale built into $T$.

\paragraph*{Ledger hook.}
In the TT$^\ast$ audit, \eqref{eq:ssu-final} contributes the minor–arc/\emph{bilinear} term
\[
\|P_{\cA}E\|_2\ \ll\ (\log X)^C\,\Big(\frac{H}{X}\Big)^{1/2},
\]
with the same $(\log X)^C$ loss as above. This is the only SSU input used downstream.



\subsection{Tube overlap for the Fejér-banked partition}

We prove the bounded-overlap property (T1) used in §\ref{sec:ssu}:
\[
\sum_{j,k} P_{U_{j,k}}(\xi)\ \asymp\ \mathbf 1_{\cA}(\xi)
\quad\text{and}\quad
\Lambda:=\sup_{\xi\in\cA}\sum_{j,k} P_{U_{j,k}}(\xi)\ \ll\ (\log X)^C,
\]
with an absolute $C$ depending only on the fixed tapers.

\subsubsection{Setup and construction}
Let $\wh\Phi_H$ be the fixed nonnegative Fejér window supported on $|\xi|\lesssim 1/H$ with total mass $\asymp 1$.
For each dyadic shell $\mathcal D_j=\{\xi: 2^{-j-1}/H<|\xi|\le 2^{-j}/H\}$ we define a smooth short-axis bump
$\chi_j$ supported on an interval of length $\asymp 2^{-j}/H$ and satisfying
\begin{equation}\label{eq:chi_j}
0\le \chi_j\le 1,\qquad \|\chi_j\|_{L^\infty}\le 1,\qquad
\sum_{t\in (2^{-j}/H)\mathbb Z}\!\!\chi_j(\,\cdot - t)\ \asymp\ 1.
\end{equation}
Along each major arc $I\subset\cA$ (centre at a reduced $a/q$, $q\le Q$) choose a $C^\infty$ tangential taper
$\vartheta_I$ with $\mathbf 1_{I^{\circ}}\le \vartheta_I\le \mathbf 1_{I^{+}}$, where $I^{\circ}$ is the middle third
of $I$ and $I^{+}$ is $I$ enlarged by a factor $1+\frac1{100}$ (still disjoint from other arcs by the gap $\gg 1/H$).
We index the tubes in $\mathcal D_j$ as follows:
\[
U_{j,k}=\text{supp}\Big(\vartheta_{I(k)}\cdot \chi_j\big(\nu_{I(k)}(\xi)-t_{j,k}\big)\Big),
\]
where $I(k)$ is the parent arc of $U_{j,k}$, $\nu_{I}$ is the (signed) normal coordinate to $I$, and
$t_{j,k}\in (2^{-j}/H)\mathbb Z$ is chosen so that the family $\{\chi_j(\nu_{I}-t_{j,k})\}_k$ covers the normal band
over $I^{\circ}$. Define
\begin{equation}\label{eq:P_U_def}
P_{U_{j,k}}(\xi)\ :=\ \vartheta_{I(k)}(\xi)\,\chi_j\big(\nu_{I(k)}(\xi)-t_{j,k}\big).
\end{equation}
By construction, $0\le P_{U_{j,k}}\le 1$, each $P_{U_{j,k}}$ is supported in $\cA\cap\mathcal D_j$, and
\[
\sum_{k: I(k)=I} P_{U_{j,k}}(\xi)\ \asymp\ \vartheta_I(\xi)\qquad(\xi\in \mathcal D_j),
\]
so summing over arcs yields $\sum_{j,k}P_{U_{j,k}}(\xi)\asymp \mathbf 1_{\cA}(\xi)$, i.e. the partition claim in (T1).

\subsubsection{Local bounded overlap on a fixed dyadic shell}
\begin{lemma}[Bounded multiplicity at fixed $j$]\label{lem:bounded-mult-j}
For each fixed $j$ and $\xi\in\cA$, the number of indices $k$ with $P_{U_{j,k}}(\xi)\neq 0$ is $\le M$, where $M$ depends only on the fixed bump shapes in \eqref{eq:chi_j}, independent of $H,Q,X$.
\end{lemma}

\begin{proof}
Fix $j$ and $\xi\in\cA$. Because distinct major arcs are separated by $\gg 1/H$ and the normal thickness of each $U_{j,k}$ is $\asymp 2^{-j}/H\le 1/H$, at most one parent arc $I$ contributes at $\xi$.
For that arc, the set $\{k: P_{U_{j,k}}(\xi)\neq 0\}$ is contained in the indices with $\chi_j(\nu_I(\xi)-t_{j,k})\neq 0$.
By \eqref{eq:chi_j} the translates of $\chi_j$ have uniformly bounded overlap (a standard partition-of-unity property),
so at most $M$ of them can cover any single point; $M$ is an absolute constant depending only on the choice of $\chi_j$.
\end{proof}

\subsubsection{Global overlap bound}
\begin{proposition}[Tube overlap (T1)]\label{prop:tube-overlap}
With $P_{U_{j,k}}$ as in \eqref{eq:P_U_def},
\[
\sum_{j,k} P_{U_{j,k}}(\xi)\ \asymp\ \mathbf 1_{\cA}(\xi)
\quad\text{and}\quad
\Lambda:=\sup_{\xi\in\cA}\sum_{j,k} P_{U_{j,k}}(\xi)\ \le\ C_0\,(1+\log H),
\]
for an absolute constant $C_0$ depending only on the bump shapes.
In particular, since $H=(\log X)^A$, one has $\Lambda\ \ll\ (\log X)^C$ for some fixed $C=C(A)$.
\end{proposition}

\begin{proof}
The partition statement was noted after \eqref{eq:P_U_def}. For the overlap bound, fix $\xi\in\cA$.
By Lemma~\ref{lem:bounded-mult-j}, at fixed $j$ there are at most $M$ indices $k$ with $P_{U_{j,k}}(\xi)\neq 0$.
The shells $\mathcal D_j$ with nonempty intersection with the bank’s normal band $|\nu_I|\lesssim 1/H$ are those with
$2^{-j}/H\gtrsim 1/H$, i.e. $j=0,1,\dots,J_{\max}$ with $J_{\max}\asymp \log H$.
Hence
\[
\sum_{j,k} P_{U_{j,k}}(\xi)\ \le\ \sum_{j=0}^{J_{\max}} \sum_{k:\,P_{U_{j,k}}(\xi)\neq 0} 1\ \le\ M\,(1+J_{\max})
\ \ll\ 1+\log H.
\]
Finally, $H=(\log X)^A$ gives $\log H = A\log\log X \ll (\log X)^C$ after increasing $C$ if needed.
\end{proof}

\begin{proposition}[Packet tests]\label{prop:packet}
\[K_{\nu,s}(t):=R_H(t)\Phi_\nu(t)\mathbf 1_{\{|t-sX/H|\le X/(2H)\}},
\allowbreak
\quad 
\sum_t|R_H(t)|\ll H,\ \supp R_H\subset[-2H,2H].
\\
\|T_{\nu,s}\|_{2\to2}\ \ll\ 
\begin{cases}
q_\nu^{-1/2}, & s=0,\\
(H/X)^{1/2}|s|^{-1/2}, & |s|\ge1.
\end{cases}\]
\end{proposition}

\subsubsection{Compatibility with the bank taper}
\begin{lemma}[Domination by the bank weight]\label{lem:bank-dom}
Let $\widehat w_{\mathrm{bank}}$ be the fixed Fejér taper used to define the bank $\cA$.
There exists a constant $c>0$ (depending only on the taper choices) such that
\[
0\ \le\ P_{U_{j,k}}(\xi)\ \le\ c\,\widehat w_{\mathrm{bank}}(\xi)\qquad (\xi\in\mathbb T).
\]
\end{lemma}

\begin{proof}
Each $P_{U_{j,k}}$ is the product of a tangential taper $\vartheta_{I(k)}\le \mathbf 1_{I^{+}}$ and a short-axis bump
$\chi_j$ whose support lies within the normal band of width $\asymp 2^{-j}/H\le 1/H$ around $I(k)$.
The bank weight $\widehat w_{\mathrm{bank}}$ is the Fejér taper that equals $1$ on the middle thirds of the arcs and dominates any such smooth product by construction; take $c$ to be the supremum of the ratio over the (fixed) families of shapes.
\end{proof}

\paragraph*{Conclusion.}
Proposition~\ref{prop:tube-overlap} establishes (T1) with $\Lambda\ll (\log X)^C$. Lemma~\ref{lem:bank-dom} gives (T3).
Property (T2) is part of the construction. These verify the hypotheses used in §\ref{sec:ssu}.

\begin{lemma}[Minor–arc amplification by center differences]\label{lem:minor-amp}
Let $H_0:=\lceil c\,Q\rceil$ with a fixed small $c>0$.
For $h\in\{1,\dots,H_0\}$ define the center–difference operator
\[
(\Delta_h S)_{n_0}(\xi)\ :=\ \sum_{n\sim X} f(n)\,\big(K(n_0-n)-K(n_0-h-n)\big)\,e(\xi n).
\]
Then for all sufficiently large $X$,
\[
\int_{\mathfrak A(Q)^c}\!|S_{n_0}(\xi)|^2\,d\xi
\ \le\ \frac{C}{H_0}\sum_{h\le H_0}\int_{\mathfrak A(Q)^c}\!|\Delta_h S_{n_0}(\xi)|^2\,d\xi,
\]
where $C$ is absolute.
\end{lemma}

\begin{proof}
On the frequency side, $\widehat{\Delta_h K}(\xi)=(1-e(-h\xi))\widehat K(\xi)$, so
$|\Delta_h S_{n_0}(\xi)|=|1-e(-h\xi)|\cdot |S_{n_0}(\xi)|$.
Average in $h$:
\[
\frac{1}{H_0}\sum_{h\le H_0}|1-e(-h\xi)|^2
=\frac{1}{H_0}\sum_{h\le H_0} 4\sin^2(\pi h\xi)
\ \ge\ c_1
\]
uniformly for $\xi\in\mathfrak A(Q)^c$ once $H_0\asymp Q$ (standard minor–arc trigonometric sum; the set of $\xi$ at distance $<\!1/Q$ from \emph{any} rational $a/q$ with $q\le Q$ is precisely the bank). Hence
\(
|S_{n_0}(\xi)|^2\le \tfrac{C}{H_0}\sum_{h\le H_0}|\Delta_h S_{n_0}(\xi)|^2
\)
on $\mathfrak A(Q)^c$; integrate in $\xi$.
\end{proof}

\subsection{SSU and Type-II tubes}
\begin{theorem}[Slope--Shift Uncertainty (SSU) for Type~II tubes]\label{thm:SSU}
Let $X\ge 1$, $H\ge 2$, and fix a coprime pair $(a,q)$ with $1\le q\le H^{1/2}$.
Choose dyadic $D,N$ with $DN\asymp X$, and define the shear coordinates
\[
u := qn - a d,\qquad v := d.
\]
For parameters $U$ with $1\le U \ll X/(qH)$, define the tube
\[
T := \bigl\{(d,n)\in (D,2D]\times(N,2N]\ :\ |u|\le U\bigr\}.
\]
Let $K_H$ be as in Lemma~\ref{lem:kernel}. Then for every function $F$ supported on $T$,
\begin{equation}\label{eq:SSU-main}
\sum_{\substack{(d,n)\in T\\(d',n')\in T}}
F(d,n)\,\overline{F(d',n')}\ K_H\!\bigl(d'n-dn'\bigr)
\ \ll\ \Bigl(\frac{H}{X}\Bigr)^{\!1/2}(\log X)^{O(1)}\,
\sum_{(d,n)\in T}\!|F(d,n)|^2,
\end{equation}
where the implied constant is absolute (it may depend on the fixed smooth dyadic cutoffs).
\end{theorem}

\begin{proof}
\textit{Step 1: Shear and the determinant.}
Writing $(u,v)=(qn-ad,\ d)$, we have for any two points
\[
d'n-dn' \ =\ \frac{v'u - vu'}{q}.
\]
The tube condition is $|u|\le U$, $v\in(D,2D]$, together with $u\equiv -av\pmod q$.

\smallskip
\textit{Step 2: Positive Fourier expansion.}
By Lemma~\ref{lem:kernel},
\[
K_H(t)\ =\ \int_{-1/H}^{1/H}\widehat K_H(\xi)\,e\!\Big(\frac{\xi t}{X}\Big)\,d\xi,
\quad \widehat K_H(\xi)\ge0.
\]
Hence by Fubini and Cauchy--Schwarz,
\[
\mathcal B
:= \sum_{T\times T}\!F(d,n)\overline{F(d',n')}\,K_H(d'n-dn')
= \int_{-1/H}^{1/H}\!\widehat K_H(\xi)\,|S(\xi)|^2\,d\xi,
\]
with
\[
S(\xi) \ :=\ \sum_{(u,v)\in T'} F(v,u)\,e\!\Big(\frac{\xi\,uv}{qX}\Big),
\quad
T'=\{(u,v): |u|\le U,\ v\in(D,2D],\ u\equiv -av\!\!\!\pmod q\}.
\]
\(T_U\) is supported on \(U\times U^+\) and \(K_H\) is band-limited with \(\widehat K_H\subset[-1/H,1/H]\).
So, it suffices to bound $|S(\xi)|^2$ uniformly for $|\xi|\le 1/H$.

\smallskip
\textit{Step 3: Progressions in $v$ and first large sieve.}
For each admissible $u$ there is $v_0(u)\in[1,q]$ such that
$v\equiv v_0(u)\pmod q$, hence $v=v_0(u)+\ell q$ with $0\le \ell<L$ and $L\asymp D/q$.
Then
\[
S(\xi)=\sum_{|u|\le U}\!e\!\Big(\frac{\xi\,u\,v_0(u)}{qX}\Big)
\sum_{\ell=0}^{L-1}\!F\!\big(v_0(u)+\ell q,\ u\big)\,e\!\Big(\frac{\xi\,u\,\ell}{X}\Big).
\]
Apply Cauchy--Schwarz to the outer sum and the Montgomery--Vaughan large sieve to the inner sum
over $\ell$ with frequency gap $\gg |\xi|/X$ as $u$ varies over distinct integers; one obtains
\begin{equation}\label{eq:LS-outer-u}
|S(\xi)|^2\ \ll\ L\,\Big(U+\frac{X}{|\xi|}\Big)\sum_{(u,v)\in T'} |F(v,u)|^2,
\qquad L\asymp \frac{D}{q}.
\end{equation}

\smallskip
\textit{Step 4: Progressions in $u$ and second large sieve.}
Dually, for each $v$ the admissible $u$ lie in $u\equiv -av\pmod q$ with $|u|\le U$,
so write $u=u_0(v)+kq$ with $|k|\le K$ and $K\asymp U/q$. Then
\[
S(\xi)=\sum_{v\in(D,2D]}\!e\!\Big(\frac{\xi\,u_0(v)\,v}{qX}\Big)
\sum_{|k|\le K}\!F\!\big(v, u_0(v)+kq\big)\,e\!\Big(\frac{\xi\,v\,k}{X}\Big),
\]
and the same large-sieve step gives
\begin{equation}\label{eq:LS-outer-v}
|S(\xi)|^2\ \ll\ D\,\Big(\frac{U}{q}+\frac{X}{|\xi|}\Big)\sum_{(u,v)\in T'} |F(v,u)|^2.
\end{equation}

\smallskip
\textit{Step 5: Geometric mean.}
Combine \eqref{eq:LS-outer-u}–\eqref{eq:LS-outer-v} by the geometric mean to get
\[
|S(\xi)|^2 \ \ll\ \Big(\frac{DU}{q}\Big)^{1/2}
\Big(U+\tfrac{X}{|\xi|}\Big)^{1/2}\Big(D+\tfrac{X}{|\xi|}\Big)^{1/2}
\sum_{T'} |F|^2.
\]

\begin{lemma}[Balanced $\xi$--integration under admissible moments]\label{lem:balanced-xi}
Let $\widehat K_H\ge0$ be supported in $|\xi|\le 1/H$ with
\begin{equation}\label{eq:K-moments}
\int \widehat K_H(\xi)\,d\xi\ \ll\ H,\qquad 
\int_{|\xi|\le 1/H} \frac{\widehat K_H(\xi)}{|\xi|}\,d\xi\ \ll\ H\log H,\qquad 
\int_{|\xi|\le 1/H} |\xi|\,\widehat K_H(\xi)\,d\xi\ \ll 1.
\end{equation}
Then for all $U,D\ge 1$,
\begin{equation}\label{eq:balanced-claim}
\int_{|\xi|\le 1/H} \widehat K_H(\xi)\,\sqrt{U+X|\xi|}\,\sqrt{D+X|\xi|}\,d\xi
\ \ll\ \frac{UD}{H}\ +\ \sqrt{X}\,(\sqrt U+\sqrt D)\,(H\log H)^{1/2}\ +\ X,
\end{equation}
with an absolute implied constant.
\end{lemma}

\begin{corollary}[Short–axis simplification at $U\le X/(qH)$]\label{cor:short-axis}
Assume $U\le X/(qH)$ and $D\asymp X/U$ (the usual Type II geometry). Then
\[
\Big(\frac{DU}{q}\Big)^{\!1/2}\frac{UD}{H}\ \ll\ \Big(\frac{H}{X}\Big)^{\!1/2},\qquad
\Big(\frac{DU}{q}\Big)^{\!1/2}\sqrt{X}(\sqrt U+\sqrt D)(H\log H)^{1/2}\ \ll\ (\log X)^{C}\Big(\frac{H}{X}\Big)^{\!1/2},
\]
and
\[
\Big(\frac{DU}{q}\Big)^{\!1/2} X\ \ll\ X\Big(\frac{X}{q}\cdot\frac{U}{D}\Big)^{\!1/2}\ \ll\ X\Big(\frac{X}{q}\cdot\frac{U^2}{X}\Big)^{\!1/2}\ \ll\ \frac{XU}{\sqrt q}\ \ll\ \Big(\frac{H}{X}\Big)^{\!1/2},
\]
where the last inequality uses $U\le X/(qH)$ and $H\ge (\log X)^{10}$. Hence
\[
\int_{|\xi|\le 1/H} |S(\xi)|^2\,\widehat K_H(\xi)\,d\xi\ \ll\ (\log X)^C\,\Big(\frac{H}{X}\Big)^{1/2}\ \sum|F|^2.
\]
\end{corollary}

\begin{proof}
Write
\[
\sqrt{U+X|\xi|}\,\sqrt{D+X|\xi|}\ \le\ \sqrt{UD}\ +\ \sqrt U\,\sqrt{X|\xi|}\ +\ \sqrt D\,\sqrt{X|\xi|}\ +\ X|\xi|.
\]
Integrate termwise against $\widehat K_H$. For the $\sqrt{UD}$ term, apply Cauchy--Schwarz with weights $|\xi|^{1/2}$ and $|\xi|^{-1/2}$:
\[
\int \widehat K_H\,d\xi
= \int \big(\widehat K_H^{1/2}|\xi|^{1/2}\big)\cdot\big(\widehat K_H^{1/2}|\xi|^{-1/2}\big)d\xi
\ \le\ \Big(\int \widehat K_H |\xi|\,d\xi\Big)^{\!1/2}\Big(\int \frac{\widehat K_H}{|\xi|}\,d\xi\Big)^{\!1/2}
\ \ll\ (H\log H)^{1/2}.
\]
Thus $\sqrt{UD}\int \widehat K_H \ll \sqrt{UD}\,(H\log H)^{1/2}$. Now use the trivial inequality $\sqrt{UD}\le \tfrac{UD}{H}+ \tfrac{H}{4}$ and absorb the $H$ into the final $X$ term (since $X\ge H$), giving the bound $\ll UD/H+X$. For the $\sqrt{X|\xi|}$ terms, the same Cauchy--Schwarz yields
\[
\int \widehat K_H\,\sqrt{|\xi|}\,d\xi
\le \Big(\int \frac{\widehat K_H}{|\xi|}\,d\xi\Big)^{\!1/2}\Big(\int \widehat K_H|\xi|\,d\xi\Big)^{\!1/2}
\ll (H\log H)^{1/2}.
\]
Finally $\int \widehat K_H\,|\xi|\,d\xi\ll 1$. Combine:
\[
\int \widehat K_H \sqrt{U+X|\xi|}\sqrt{D+X|\xi|}\ \ll\ 
\underbrace{\frac{UD}{H}}_{\text{from }\sqrt{UD}}\ +\ 
\underbrace{\sqrt{X}(\sqrt U+\sqrt D)(H\log H)^{1/2}}_{\sqrt{X|\xi|}\text{ terms}}\ +\
\underbrace{X}_{X|\xi|\text{ term}}.
\]
\end{proof}

\smallskip
\textit{Step 6: Integrate in $\xi$.}
Insert this bound into $\mathcal B$ and use Cauchy--Schwarz in $\xi$:
\[
\mathcal B \ \ll\ \Big(\frac{DU}{q}\Big)^{1/2}
\Bigg(\int_{-1/H}^{1/H}\!\widehat K_H(\xi)\Big(U+\frac{X}{|\xi|}\Big)d\xi\Bigg)^{\!1/2}
\Bigg(\int_{-1/H}^{1/H}\!\widehat K_H(\xi)\Big(D+\frac{X}{|\xi|}\Big)d\xi\Bigg)^{\!1/2}
\sum_{T'} |F|^2.
\]
By Lemma~\ref{lem:kernel},
\[
\int\widehat K_H(\xi)\Big(U+\frac{X}{|\xi|}\Big)d\xi \ \ll\ U H + X H\log H,
\quad
\int\widehat K_H(\xi)\Big(D+\frac{X}{|\xi|}\Big)d\xi \ \ll\ D H + X H\log H.
\]

\smallskip
\textit{Step 7: Tube scale and simplification.}
With $U\ll X/(qH)$ and $D\ll X$, the dominant contribution is controlled by $XH\log H$ in each bracket. Hence
\[
\mathcal B\ \ll\ \Big(\frac{DU}{q}\Big)^{1/2}\,XH\log H\ \sum_{T'} |F|^2.
\]
Using $U\ll X/(qH)$ and $D\ll X$,
\[
\Big(\frac{DU}{q}\Big)^{1/2}\,XH\log H
\ \ll\ \Big(\frac{X}{q}\cdot\frac{X}{qH}\Big)^{1/2}\,XH\log H
\ =\ \Big(\frac{H}{X}\Big)^{\!1/2} \,(\log X)^{O(1)}.
\]
This yields \eqref{eq:SSU-main}.
\end{proof}

\begin{remark}
The specific choice in Lemma~\ref{lem:kernel} is merely one convenient instance.
Any nonnegative, even profile $\widehat K_H$ supported in $[-1/H,1/H]$ and satisfying
$\int \widehat K_H\ll H$ and $\int \widehat K_H/|\xi|\ll H\log H$
can be used in the proof without change.
\end{remark}

\begin{lemma}[Slope–packet orthogonality]\label{lem:slope-orth}
Let $T_{\nu,s}$ be the packet multipliers with Fourier symbol
\[
M_{\nu,s}(\xi)\ :=\ \eta_\nu(\xi)\,\Phi_H(\xi)\,e\!\Big(\xi\,\frac{sX}{H}\Big),
\]
where $\eta_\nu$ is supported on an interval of length $\asymp (q_\nu H)^{-1}$ centered at $a_\nu/q_\nu$ and $\Phi_H$ is supported on $|\xi|\ll 1/H$ with $\|\partial^m(\eta_\nu\Phi_H)\|_\infty\ll H^{m}$ for $m\le2$. For any $f\in L^2(\mathbb T)$ and distinct packets $(\nu,s)\neq(\nu',s')$,
\begin{equation}\label{eq:slope-orth-decay}
\big|\langle T_{\nu,s}f,\ T_{\nu',s'}f\rangle\big|
\ \ll\ \frac{1}{(1+|s-s'|)^2}\cdot\frac{1}{(1+\operatorname{dist}(\nu,\nu')\,H)^2}\ \|f\|_2^2,
\end{equation}
where $\operatorname{dist}(\nu,\nu')=\bigl|\frac{a_\nu}{q_\nu}-\frac{a_{\nu'}}{q_{\nu'}}\bigr|$.
\end{lemma}

\begin{proof}
By Plancherel,
\[
\langle T_{\nu,s}f,\ T_{\nu',s'}f\rangle
=\int_{\mathbb T} |\widehat f(\xi)|^2\,\overline{M_{\nu,s}(\xi)}\,M_{\nu',s'}(\xi)\,d\xi
=\int |\widehat f(\xi)|^2\,\Theta_{\nu,\nu'}(\xi)\,e\!\Big(\xi\,\frac{(s'-s)X}{H}\Big)\,d\xi,
\]
with $\Theta_{\nu,\nu'}(\xi):=\overline{\eta_\nu(\xi)\Phi_H(\xi)}\,\eta_{\nu'}(\xi)\Phi_H(\xi)$.  
If the two supports are disjoint (i.e.\ $\operatorname{dist}(\nu,\nu')\gg 1/H$), then $\Theta_{\nu,\nu'}\equiv0$ and the inner product vanishes. Otherwise the supports have overlap $\ll 1/H$ and standard bump calculus gives
\[
\|\Theta_{\nu,\nu'}\|_{L^1}\ \ll\ \frac{1}{H}\cdot\frac{1}{1+\operatorname{dist}(\nu,\nu')\,H},\qquad
\|\Theta'_{\nu,\nu'}\|_{L^1}\ \ll\ \frac{1}{1+\operatorname{dist}(\nu,\nu')\,H}.
\]
Let $\Delta:=\frac{(s'-s)X}{H}$. Integrating by parts twice in $\xi$ (on the oscillatory factor $e(\xi\Delta)$) gives
\[
\Big|\int |\widehat f|^2\,\Theta_{\nu,\nu'}\,e(\xi\Delta)\,d\xi\Big|
\ \ll\ \frac{1}{(1+|\Delta|)^2}\Big(\|\Theta_{\nu,\nu'}\|_{L^1} + \|\Theta'_{\nu,\nu'}\|_{L^1}\Big)\ \|\widehat f\|_\infty^2.
\]
Since $\|\widehat f\|_\infty\le \|f\|_2$ and $|\Delta|\asymp |s-s'|\,X/H$, we obtain
\[
\big|\langle T_{\nu,s}f,\ T_{\nu',s'}f\rangle\big|
\ \ll\ \frac{1}{(1+|s-s'|)^2}\cdot\frac{1}{(1+\operatorname{dist}(\nu,\nu')\,H)^2}\ \|f\|_2^2,
\]
which is \eqref{eq:slope-orth-decay}.
\end{proof}


\begin{theorem}[Two–weight $TT^*$ Sawyer test for SSU with $(H/X)^{1/2}$ saving]\label{thm:SSU-Sawyer}
Let $K\ge0$ be a time–pin of width $H=(\log X)^A$ as in Assumption~\ref{ass:bank}, and let $\mathfrak A(Q)$ be the bank with $Q=H^\gamma$, $0<\gamma<\tfrac12$. For $f:\mathbb Z\to\mathbb C$ supported on $[X,2X]$ define
\[
(Tf)(n_0,\xi):=\mathbf 1_{\mathfrak A(Q)^c}(\xi)\sum_{n\sim X} f(n)\,K(n_0-n)\,e(\xi n),\qquad n_0\in[X,2X],\ \xi\in\mathbb T.
\]
Then for every admissible F3–block $f=F_{\mathsf B}$ (Type II/III, divisor–type coefficients, smooth dyadic cutoffs) one has
\begin{equation}\label{eq:SSU-strong}
\mathbb E_{n_0\in[X,2X]}\int_{\mathfrak A(Q)^c} |(Tf)(n_0,\xi)|^2\,d\xi
\ \ll\ (\log X)^C\,\Big(\frac{H}{X}\Big)^{1/2}\ \|F_{\mathsf B}\|_2^2.
\end{equation}
Summing over the $O((\log X)^C)$ blocks yields the global SSU bound used in the ledger.
\end{theorem}

\begin{remark}[Uniformity of constants]\label{rem:ssu-uniformity}
The constants implicit in Theorem~\ref{thm:SSU-Sawyer} depend only on:
(i) finitely many derivatives of the dyadic windows $W_d,W_n$;
(ii) the moment bounds $\int\widehat K_H\ll H$ and $\int \widehat K_H/|\xi|\ll H\log H$;
(iii) the overlap constant in (T1), which is $\ll(\log X)^C$;
and (iv) the fixed choice $Q=H^\gamma$ with $0<\gamma<\tfrac12$.
They are uniform in coprime $a/q$ with $q\le Q$ and in tube centers $s$ with $U\asymp X/(qH)$.
\end{remark}


\begin{lemma}[Congruence dissection and SBP control]\label{lem:congruence-sbp}
Let $1\le q\le Q$ and $(a,q)=1$. In the shear coordinates $u=qn-ad$, $v=d$, fix a tube
\[
T(a/q,s)\ :=\ \{(d,n):\ |u-s|\le U,\ d\in(D,2D]\},\qquad U\asymp X/(qH).
\]
For smooth windows $W_d, W_n$ with $\|W_d^{(k)}\|_\infty\ll D^{-k}$ and $\|W_n^{(k)}\|_\infty\ll N^{-k}$, we have
\[
\sum_{(d,n)\in T(a/q,s)} F(d,n)\,W_d(d/D)\,W_n(n/N)
\;=\;\sum_{\substack{r\!\!\!\pmod q}}\ \sum_{\substack{v\equiv r\ (q)\\ |u-s|\le U}}\!F(d,n)\,W_d(d/D)\,W_n(n/N)\;+\;O_A\bigl((\log X)^{-A}\!\!\sum|F|\bigr),
\]
where in each inner sum the congruence $u\equiv -av\pmod q$ is enforced. The error term arises from at most two integrations by parts in $d$ or $n$, and is uniform in $a,q,s,U$.
\end{lemma}

\begin{lemma}[Congruence dissection with uniform SBP error]\label{lem:ap-sbp-error}
Let $1\le q\le Q$, $(a,q)=1$, and in shear coordinates $u=qn-ad$, $v=d$ consider a tube
\[
T(a/q,s)=\{(d,n): |u-s|\le U,\ d\asymp D\},\qquad U\asymp X/(qH).
\]
For smooth windows $W_d,W_n$ with $\|W_d^{(k)}\|_\infty\ll D^{-k}$ and $\|W_n^{(k)}\|_\infty\ll N^{-k}$,
\[
\sum_{(d,n)\in T(a/q,s)} F(d,n)\,W_d(d/D)\,W_n(n/N)
=\sum_{\substack{r\bmod q}}\ \sum_{\substack{v\equiv r\ (q)\\ |u-s|\le U}}F(d,n)\,W_d\,W_n + O_A\!\Big((\log X)^{-A}\!\sum|F|\Big),
\]
uniformly in $a,q,s,U$, where the error comes from nonzero additive characters and is obtained by at most two summations by parts.
\end{lemma}

\begin{proof}
Insert $1_{u\equiv -av\ (q)}=\frac1q\sum_{h\bmod q}e(\frac{h(u+av)}{q})$. The $h=0$ term gives the main term. For $h\ne0$, view the sum in $d$ or $n$ with phase $e(h(\,\cdot\,)/q)$; discrete summation by parts once (twice if needed) against $W_d$ or $W_n$ gives a factor $\ll (|h|/q)^{-1}$ per integration, while derivatives of $W_d,W_n$ contribute $D^{-1}$ or $N^{-1}$. Summing $1\le |h|\le q$ and using $q\le Q\ll (\log X)^{O(1)}$ yields the $(\log X)^{-A}$–loss uniformly in $s,U$.
\end{proof}

\begin{lemma}[Smooth windows: boundary control]\label{lem:smooth-boundary}
Let $W_d,W_n$ be dyadic windows supported in $[1/2,2]$ with $\|W_d^{(k)}\|_\infty\ll D^{-k}$ and $\|W_n^{(k)}\|_\infty\ll N^{-k}$. For any $G$ supported on $d\sim D$, $n\sim N$,
\[
\sum_{d\sim D}\sum_{n\sim N} G(d,n)\ -\ \sum_{d}\sum_{n} G(d,n)\,W_d(d/D)\,W_n(n/N)
= O_A\!\Big((\log X)^{-A}\sum_{d,n}|G(d,n)|\Big).
\]
\end{lemma}

\begin{proof}
Write the difference as a sum over boundary strips of thickness $D/(\log X)^B$ and $N/(\log X)^B$. One discrete Abel summation across each strip gains $D^{-1}$ or $N^{-1}$ from the window derivative. Choosing $B$ large and summing over $O(1)$ strips gives the stated $(\log X)^{-A}$ loss.
\end{proof}

\begin{lemma}[Smooth localization and boundary control]\label{lem:smooth-sbp}
Let $W_d,W_n$ be dyadic windows supported in $[1/2,2]$ with $\|W_d^{(k)}\|_\infty\ll D^{-k}$, $\|W_n^{(k)}\|_\infty\ll N^{-k}$. For any function $G(d,n)$ supported on $(D/2,2D]\times(N/2,2N]$,
\[
\sum_{d\sim D}\sum_{n\sim N} G(d,n)\ -\ \sum_{d}\sum_{n} G(d,n)\,W_d(d/D)\,W_n(n/N)
\ =\ O_A\bigl((\log X)^{-A}\sum_{d,n}|G(d,n)|\bigr).
\]
The implicit constant depends only on finitely many derivatives of $W_d,W_n$.
\end{lemma}

\begin{proof}
Write the difference as a sum over boundary strips of thickness $O(D/(\log X)^B)$ and $O(N/(\log X)^B)$. Integrate by parts once (discrete Abel summation) in $d$ (resp.\ $n$) across each strip; the derivative bounds give a factor $D^{-1}$ (resp.\ $N^{-1}$) per strip. Taking $B$ large and summing over $O(1)$ strips yields the claimed $(\log X)^{-A}$ loss.
\end{proof}


\begin{proof}
Insert the identity
\[
1_{u\equiv -av\ (q)}=\frac1q\sum_{h\bmod q}e\!\left(\frac{h(u+av)}{q}\right)
\]
and split the $(d,n)$–sum by $v\equiv r\pmod q$. The $h=0$ term gives the stated main term. For $h\neq 0$, the phase is nonconstant in $d$ or $n$; applying summation by parts against $W_d$ or $W_n$ once (and, if needed, twice) yields a factor $\ll (|h|/q)^{-1}$ per integration and introduces derivatives of $W_d,W_n$ absorbed by $D^{-1},N^{-1}$. Summing $|h|\le q$ gives the stated $O_A((\log X)^{-A}\sum|F|)$ upon taking $A$ large and using $q\le Q\ll (\log X)^{O(1)}$. Uniformity in $s,U$ follows from $U\asymp X/(qH)$ and bounded tube thickness.
\end{proof}

\begin{proof}
Write the dual form via $TT^*$. For $g$ supported on $[X,2X]\times\mathfrak A(Q)^c$, define
\[
(T^\ast g)(m)\ :=\ \sum_{n_0}\int_{\mathbb T}\mathbf 1_{\mathfrak A(Q)^c}(\xi)\,K(n_0-m)\,e(-\xi m)\,g(n_0,\xi)\,d\xi.
\]
Then $\|Tf\|_{L^2(\mathbb E_{n_0}d\xi)}^2=\langle T^\ast T f,f\rangle_{\ell^2}$, with
\[
(T^\ast T f)(m)
=\sum_{n\sim X}\!f(n)\,R(m-n)\,\Phi(m-n),
\quad
R(t):=\mathbb E_{n_0}K(n_0)K(n_0-t),\ 
\Phi(t):=\int_{\mathfrak A(Q)^c} e(\xi t)\,d\xi.
\]
By Assumption~\ref{ass:bank}: $R(t)\ge0$, $\mathrm{supp}\,R\subset[-2H,2H]$, $\sum_t R(t)\asymp H$, and $\sum_t |t|\,R(t)\ll H\log H$ (kernel moments).

\smallskip
\emph{Step 1: Tube decomposition in frequency and slope.}
Cover $\mathfrak A(Q)^c$ by a smooth partition of unity $\{\eta_\nu\}_\nu$ with $\eta_\nu$ supported on intervals of length $\asymp (q_\nu H)^{-1}$, where $1\le q_\nu\le Q$, and with $\sum_\nu \eta_\nu=\mathbf 1_{\mathfrak A(Q)^c}$, bounded overlap, and $\|\eta'_\nu\|_\infty\ll q_\nu H$. For each $\nu$ define
\[
\Phi_\nu(t):=\int \eta_\nu(\xi)\,e(\xi t)\,d\xi,
\qquad
|\Phi_\nu(t)|\ \ll\ \min\!\Big\{(q_\nu H)^{-1},\ |t|^{-1}\Big\}.
\]
Further, for each $\nu$ split the shift–variable $t$ into slope–packets of width $\asymp X/H$:
\[
\mathbb Z\ni t\ =\ s\cdot \frac{X}{H}\ +\ r,\qquad s\in\mathbb Z,\quad |r|\le \frac{X}{2H}.
\]
Let $\Upsilon_s$ denote the set $\{t: |t-sX/H|\le X/(2H)\}$. Then
\[
\Phi(t)=\sum_\nu \Phi_\nu(t)=\sum_{\nu}\sum_{s}\ \Phi_\nu(t)\,\mathbf 1_{\Upsilon_s}(t).
\]

\smallskip
\emph{Step 2: Strong orthogonality across slope–packets.}
Fix $\nu$. For $s\neq s'$, the supports $\Upsilon_s$ and $\Upsilon_{s'}$ are disjoint, and oscillatory factors $e(\xi t)$ with $\xi$ restricted to the same $\eta_\nu$ produce \emph{discrete} orthogonality across $s$ at scale $X/H$ when tested against F3–blocks. Concretely, for an F3–block $f$ (Type II/III) we have
\begin{equation}\label{eq:slope-orth}
\sum_{m}\ \overline{f(m)}\sum_{n} f(n)\,R(m-n)\,\Phi_\nu(m-n)\,\mathbf 1_{\Upsilon_s}(m-n)
\ \perp\ 
\sum_{m}\ \overline{f(m)}\sum_{n} f(n)\,R(m-n)\,\Phi_\nu(m-n)\,\mathbf 1_{\Upsilon_{s'}}(m-n)
\end{equation}
whenever $|s-s'|\ge 1$. This is the standard “different slope” orthogonality: in the bilinear/trilinear F3 model the phase $e(\xi(n-m))$ couples to a product phase $e(\xi\,\Phi_{\mathsf B}(a,b,\dots))$ whose stationary set singles out $|n-m|\lesssim X/H$ per $\eta_\nu$; distinct slope–packets $s\neq s'$ then yield disjoint stationary regions. (A detailed derivation follows the bilinear TT$^\ast$ geometry: see \ref{subsec:BG.incidence}.

Consequently, for a fixed $\nu$ the packets over $s$ add in $\ell^2$:
\[
\langle T^\ast T f,f\rangle
=\sum_{\nu}\ \sum_{s}\ \mathcal Q_{\nu,s}(f),
\qquad
\mathcal Q_{\nu,s}(f):=\sum_{m,n} \overline{f(m)}\,f(n)\,R(m-n)\,\Phi_\nu(m-n)\,\mathbf 1_{\Upsilon_s}(m-n).
\]

\smallskip
\emph{Step 3: One–packet Sawyer testing.}
Fix $(\nu,s)$. For each tube $\nu$ and slope packet $s$, with
\(
K_{\nu,s}(t):=R(t)\,\Phi_\nu(t)\,\mathbf 1_{\Upsilon_s}(t),
\)
we have
\[
\sum_{t}|K_{\nu,s}(t)|
\ \ll\
\begin{cases}
\displaystyle \frac{H}{q_\nu H}\ =\ \frac{1}{q_\nu}, & s=0,\\[6pt]
\displaystyle \frac{H}{|s|X}, & |s|\ge 1,
\end{cases}
\]
since $|t|\le 2H$ and $|\Phi_\nu(t)|\ll \min\{(q_\nu H)^{-1},|t|^{-1}\}$ on $\mathfrak A(Q)^c$.
Thus Schur’s symmetric test gives
\[
\sup_{n_0}\sum_n |K_{\nu,s}(n_0-n)|
\ \ll\
\begin{cases}
q_\nu^{-1/2}, & s=0,\\
(H/X)^{1/2}|s|^{-1/2}, & |s|\ge 1.
\end{cases}
\tag{$\ast$}
\]
Summing in $s$ uses $\sum_{|s|\ge1} |s|^{-1/2}\ll 1$ (finite $s$–overlap for fixed $\nu$);
summing in $\nu$ across $\mathfrak A(Q)^c$ uses bounded overlap of tubes and $q_\nu\ge Q$.
Hence the one–packet contribution obeys
\[
\|T_{\nu,s}\|_{2\to 2}\ \ll\ (\log X)^C\left( (H/X)^{1/2} + Q^{-1/2}\right).
\]
\emph{With the amplification Lemma \ref{lem:minor-amp}} (choosing $H_0\asymp Q$) the $Q^{-1/2}$ term is absorbed, and we retain the robust $(H/X)^{1/2}$ saving stated in Theorem~\ref{thm:SSU-Sawyer}.

\smallskip
\emph{Step 4: Verifying the tests.}
Since $R$ is supported on $|t|\le 2H$ and $\Upsilon_s$ on $|t-sX/H|\le X/(2H)$, for any fixed $(\nu,s)$ there are at most $O(1)$ integers $t$ in the support intersection; in fact either the intersection is empty (if $|s|>3H^2/X$) or it is a short interval of length $\ll 1$ (independent of $X$). Thus
\[
\sum_{t}|K_{\nu,s}(t)|\ \ll\ \max_{t}\ |R(t)|\,|\Phi_\nu(t)|\ \ll\ \max_{|t|\le 2H}\ |\Phi_\nu(t)|.
\]
Using $|\Phi_\nu(t)|\ll \min\{(q_\nu H)^{-1},|t|^{-1}\}$ and $|t|\asymp |s|X/H$ on $\Upsilon_s$ (when nonempty), we have
\[
|\Phi_\nu(t)|\ \ll\ \min\!\Big\{(q_\nu H)^{-1},\ \frac{H}{|s|X}\Big\}.
\]
If $s=0$, then $|t|\ll X/H$ forces $|t|\asymp 1$ and we use the first branch, getting $(q_\nu H)^{-1}$. If $|s|\ge 1$, use the second branch to get $H/(|s|X)$. Either way,
\[
\sum_{t}|K_{\nu,s}(t)|\ \ll\ \min\!\Big\{\frac{1}{q_\nu H},\ \frac{H}{|s|X}\Big\}.
\]
Taking square roots (Schur’s symmetric test) gives
\[
\sup_{n_0}\sum_n |K_{\nu,s}(n_0-n)|\ \ll\ \Big(\min\!\Big\{\frac{1}{q_\nu H},\ \frac{H}{|s|X}\Big\}\Big)^{1/2}
\ \le\ C\,\Big(\frac{H}{X}\Big)^{1/2}.
\]
(The last inequality is trivial for $|s|\ge 1$; for $s=0$ use $q_\nu\le Q\le H^{1/2-\varepsilon}$ so $(q_\nu H)^{-1/2}\le H^{-1/4}\le (H/X)^{1/2}$ once $X$ is large and $A$ fixed.) This proves \eqref{eq:slope-orth}–\eqref{eq:dense-grid} and completes the proof.
\end{proof}

\begin{lemma}[Bandwidth and sampling constants]\label{lem:bandwidth-bookkeeping}
Let $S(t)=(R_2*K)(t)$ with $K$ as in Assumption~\ref{ass:bank}. Then $S$ is bandlimited with bandwidth
\[
\Lambda\ :=\ \sup\{|\xi|:\ \widehat S(\xi)\ne 0\}\ \le\ \frac{1}{H}.
\]
Let $\Sigma\subset\mathbb T$ be the frequency region where we sample/estimate $\widehat S$, namely the bank $\mathfrak A(Q)$ or its complement. Then the partition in §4–§7 guarantees
\[
|\Sigma|\ \asymp\ \frac{Q^2}{H}.
\]
Consequently, the Logvinenko–Sereda sampling/Nikolski–Bernstein regime
\[
\Lambda\,\Delta\ \le\ c_0\qquad\Big(\text{e.g.\ choose }\ \Delta\ =\ \frac{c_0}{\Lambda}\,\frac{m}{m+C_{\mathrm{Nik}}\,\|S\|_{\ell^2([X,2X])}/H^{1/2}}\Big)
\]
holds with $\Lambda=1/H$. In particular, any choice of grid spacing $\Delta\le c_0\,H$ satisfies the Paley–Wiener sampling prerequisite; combining with Lemma~\ref{lem:sampling-bernstein} and Corollary~\ref{cor:nikolskii-grid} yields the concrete grid used in §11.
\end{lemma}

\begin{theorem}[SSU via amplification+Schur]\label{thm:SSU-a}
$\mathbb E_{n_0}\int_{\mathfrak A^c}\Big|\sum f(n)K_H(n_0-n)e(\xi n)\Big|^2 d\xi$
\[ \ll (\log X)^C\,(H/X)^{1/2}\,\|f\|_2^2. \]
\end{theorem}

\begin{proof}[Proof of Theorem~\ref{thm:SSU-Sawyer} (frequency separation step)]
Let $H_0=\lceil cQ\rceil$ (small fixed $c>0$). By the standard amplification identity,
\[
\int_{A^c}\Big|\sum_n f(n)K_H(n_0-n)e(\xi n)\Big|^2 d\xi
\ \ll\ 
\frac{1}{H_0}\sum_{1\le |h|\le H_0}
\int_{A^c}\Big|\sum_n f(n)\Delta_hK_H(n_0-n)\,e(\xi n)\Big|^2 d\xi,
\]
where $\Delta_hK_H(t):=K_H(t)-K_H(t-h)$ and $A^c$ denotes the complement of the bank.
Fix a sheared tube $T_\nu$ and decompose to packets; after shearing $(d,n)\mapsto(u,v)$ and
freezing the long variable, the inner sum is of the form
\[
\sum_{v\in I_\nu}\, a_v\,e\!\Big(\xi\,\frac{v}{q_\nu X}\Big),
\qquad |I_\nu|\ \asymp\ \frac{X}{H},
\]
with coefficients $a_v$ depending on $f$ and the $h$–difference.
On $A^c$ we dyadically decompose $\xi$ so that $|\xi|\asymp 2^{-j}/H$.
Consecutive frequencies in $v$ are then separated by
\[
\delta_j\ \asymp\ \frac{|\xi|}{q_\nu X}\ \asymp\ \frac{2^{-j}}{q_\nu X H}.
\]
Since $|I_\nu|\asymp X/H$, Montgomery–Vaughan’s large sieve in the $v$–variable gives
\[
\int_{|\xi|\asymp 2^{-j}/H}\Big|\sum_{v\in I_\nu} a_v\,e\Big(\xi\frac{v}{q_\nu X}\Big)\Big|^2 d\xi
\ \ll\
\Big(\frac{X}{H}\;+\;\delta_j^{-1}\Big)\sum_{v\in I_\nu}|a_v|^2
\ \ll\
\Big(\frac{X}{H}\;+\;q_\nu XH\,2^{j}\Big)\sum_{v}|a_v|^2.
\]
Summing over $q_\nu\le Q$ and tubes and using $q_\nu\le Q=H^\gamma$ with $\gamma<\tfrac12$, the dyadic weight from the packet window (two integrations by parts in $\xi$) contributes
a factor $2^{-2j}$, hence
\[
\sum_{j\ge 0} 2^{-2j}\Big(\frac{X}{H}\;+\;q_\nu XH\,2^{j}\Big)
\ \ll\
\frac{X}{H}\;+\;XH\,q_\nu
\ \ll\
\frac{X}{H}\;+\;XH^{1+\gamma}.
\]
Amplifying in $h$ and summing the packets (bounded overlap) produces the
claimed $(\log X)^C(H/X)^{1/2}$ bound by Schur’s test: the $X/H$ term yields the main saving,
and the $XH^{1+\gamma}$ term is absorbed after the $h$–difference (which contributes an extra
$H_0^{-1}$ and a derivative of $K_H$, hence a factor $\ll H^{-1}$), together with the $2^{-2j}$ decay.
Thus
\[
\int_{A^c}\Big|\sum_n f(n)K_H(n_0-n)e(\xi n)\Big|^2 d\xi
\ \ll\ (\log X)^C\,(H/X)^{1/2}\,\|f\|_2^2.
\]
This completes the frequency–separation step. The remaining (routine) steps in the proof
are bookkeeping of the packet overlap and the amplification identity.
\end{proof}

\begin{proposition}[Optional: Cotlar--Stein SSU]\label{prop:cotlar-ssu}
SSU might be confirmed by an independent process. Decompose $T=\sum_{\nu,s} T_{\nu,s}$ with $T_{\nu,s}$ defined by the kernel $K_{\nu,s}$.
Then
\[
\|T\|_{2\to 2}^2\ \le\ \sup_{\nu,s}\ \sum_{\nu',s'} \big\|T_{\nu,s}^\ast T_{\nu',s'}\big\|_{2\to 2}.
\]
Moreover
\[
\big\|T_{\nu,s}^\ast T_{\nu',s'}\big\|_{2\to 2}\ \ll\ (\log X)^C\cdot 
\left(\frac{H}{X}\right)\cdot \frac{1}{\big(1+|s-s'|\big)^{2}}
\cdot \frac{1}{\big(1+\operatorname{dist}(\nu,\nu')\,H\big)^{2}},
\]
so the double sum converges absolutely, yielding
\(
\|T\|_{2\to 2}\ \ll\ (\log X)^C\,(H/X)^{1/2}.
\)
\end{proposition}

\begin{proof}[Proof of Proposition~\ref{prop:cotlar-ssu}]
Write the packet operator in Fourier as
\[
\widehat{T_{\nu,s}f}(\xi)
\;=\;
\eta_\nu(\xi)\,\Phi_H(\xi)\,e\!\Big( \xi\,\frac{sX}{H}\Big)\,\widehat{f}(\xi),
\]
where $\eta_\nu$ is supported on an arc of width $\asymp (q_\nu H)^{-1}$ around $a_\nu/q_\nu$ and
$\Phi_H$ is the fixed bandwidth-$1/H$ window. (Only the support and derivative bounds
$\|\partial^m(\eta_\nu\Phi_H)\|_\infty\ll H^{m}$ for $m\le 2$ are used.) Thus
\[
T_{\nu,s}^{*}T_{\nu',s'}
\;=\;
\operatorname{Op}\!\Big(\, \overline{\eta_\nu\Phi_H}\ \eta_{\nu'}\Phi_H\ 
e\!\Big(\xi\,\frac{(s'-s)X}{H}\Big)\Big).
\]
By Plancherel, the integral kernel is the inverse transform of the multiplier:
\[
K_{\nu,s;\nu',s'}(t)
=
\int_{\mathbb T}
\overline{\eta_\nu(\xi)\Phi_H(\xi)}\ \eta_{\nu'}(\xi)\Phi_H(\xi)\ 
e\!\Big(\xi\,\frac{(s'-s)X}{H}\Big)\,e(\xi t)\,d\xi.
\]
Set $\Theta_{\nu,\nu'}(\xi):=\overline{\eta_\nu(\xi)\Phi_H(\xi)}\,\eta_{\nu'}(\xi)\Phi_H(\xi)$ and
$\tau:=t+\frac{(s'-s)X}{H}$. Then $K_{\nu,s;\nu',s'}(t)=\int \Theta_{\nu,\nu'}(\xi)\,e(\xi\tau)\,d\xi$.
Two inputs drive the decay:
\begin{itemize}
\item \emph{Slope separation.} The supports of $\eta_\nu$ and $\eta_{\nu'}$ are centered at rationals
$a_\nu/q_\nu$, $a_{\nu'}/q_{\nu'}$ with $q_\nu,q_{\nu'}\le Q$; their centers are $\gg (q_\nu q_{\nu'})^{-1}
\asymp \mathrm{dist}(\nu,\nu')$ apart. Since each support has width $\asymp 1/H$, the product
$\Theta_{\nu,\nu'}$ enjoys the quantitative bound
\[
\|\Theta_{\nu,\nu'}\|_{C^2}
\ll 1,\qquad
\|\Theta_{\nu,\nu'}\|_{L^1}\ll \frac{1}{H}\,\frac{1}{1+(\mathrm{dist}(\nu,\nu')\,H)},
\]
and, by one more derivative, the standard integration by parts estimate
\begin{equation}\label{eq:slope-decay}
\Big|\int \Theta_{\nu,\nu'}(\xi)\,e(\xi \tau)\,d\xi\Big|
\ \ll\ \frac{1}{H}\,\frac{1}{\big(1+\mathrm{dist}(\nu,\nu')\,H\big)^{2}}\ \frac{1}{(1+|\tau|)^{2}}.
\end{equation}
\item \emph{Short-shift separation.} Writing $\tau=t+\frac{(s'-s)X}{H}$, the factor
$(1+|\tau|)^{-2}$ yields an extra $(1+|s-s'|)^{-2}$ once we sum $|t|$ over the $t$-support of the
time kernel (thickness $\ll H$).
\end{itemize}
To pass to an operator bound, apply Schur’s test in the time variable. Using
$\sum_{|t|\ll H}(1+|t+(s'-s)X/H|)^{-2}\ll (1+|s-s'|)^{-2}$ and \eqref{eq:slope-decay}, we obtain
\[
\|T_{\nu,s}^{*}T_{\nu',s'}\|_{2\to2}
\ \ll\ 
\frac{H}{X}\cdot
\frac{1}{\big(1+|s-s'|\big)^{2}}\cdot
\frac{1}{\big(1+\mathrm{dist}(\nu,\nu')\,H\big)^{2}},
\]
where the factor $H/X$ comes from the $t$-thickness and the normalization of $T$ (cf. the
definition in \S6.2). Summing in $(\nu',s')$ and taking $\sup_{\nu,s}$,
the double sum converges absolutely by Lemma~6.30 (incidence–degeneracy summation), hence
Cotlar–Stein gives
\[
\|T\|_{2\to2}^2
\ \le\
\sup_{\nu,s}\sum_{\nu',s'}\|T_{\nu,s}^{\ast}T_{\nu',s'}\|_{2\to2}
\ \ll\ (\log X)^C\ \frac{H}{X},
\]
so $\|T\|_{2\to2}\ll (\log X)^C(H/X)^{1/2}$, as claimed.
\end{proof}

\begin{lemma}[Incidence–Degeneracy summation]\label{lem:incdeg-sawyer}
Let $\{T_{\nu,s}\}$ be the tube packets and define
\(
\operatorname{dist}(\nu,\nu')=\bigl|\frac{a_\nu}{q_\nu}-\frac{a_{\nu'}}{q_{\nu'}}\bigr|.
\)
Then
\begin{equation}\label{eq:incdeg-sum}
\sup_{\nu,s}\ \sum_{\nu',s'} \frac{1}{\big(1+|s-s'|\big)^{2}\,\big(1+\operatorname{dist}(\nu,\nu')\,H\big)^{2}}
\ \ll\ (\log X)^C.
\end{equation}
\end{lemma}

\begin{proof}
Split dyadically by $|s-s'|\asymp 2^{j}$ and $\operatorname{dist}(\nu,\nu')\asymp 2^{-k}/H$ with $j,k\ge0$.
By the \emph{Incidence} Lemma~\ref{lem:BG.incidence} (05\_BG.tex, line~208) and the \emph{Degeneracy} Lemma~\ref{lem:BG.degeneracy} (05\_BG.tex, line~247), for each $k$ there are $\ll (\log X)^C(1+2^k)$ slopes $\nu'$ with $\operatorname{dist}\asymp 2^{-k}/H$; for each such $\nu'$, at most $\ll (\log X)^C(1+2^j)$ short shifts $s'$ produce overlapping packets. Hence the $(j,k)$–box contributes
\(
\ll (\log X)^C\frac{(1+2^j)(1+2^k)}{(1+2^j)^2(1+2^k)^2}\ll (\log X)^C 2^{-j}2^{-k}.
\)
Summing $j,k\ge0$ gives \eqref{eq:incdeg-sum}.
\end{proof}

\begin{lemma}[Incidence–Degeneracy summation for Sawyer]\label{lem:incdeg-sawyer}
Let $\{T_{\nu,s}\}$ be the tube packets with slopes $a_\nu/q_\nu$ and short-shift index $s$, and let
\[
\operatorname{dist}(\nu,\nu')\ :=\ \Big| \frac{a_\nu}{q_\nu}-\frac{a_{\nu'}}{q_{\nu'}}\Big|.
\]
Then
\begin{equation}\label{eq:incdeg-sum}
\sup_{\nu,s}\ \sum_{\nu',s'} \frac{1}{\big(1+|s-s'|\big)^{2}\,\big(1+\operatorname{dist}(\nu,\nu')\,H\big)^{2}}
\ \ll\ (\log X)^C.
\end{equation}
In particular, together with Proposition~\ref{prop:cotlar-ssu},
\(
\|T\|_{2\to 2}^2 \ll (\log X)^C\,(H/X).
\)
\end{lemma}

\begin{proof}
Fix $(\nu,s)$ and dyadically split by $|s-s'|\asymp 2^j$ and $\operatorname{dist}(\nu,\nu')\asymp 2^{-k}/H$ with $j,k\ge 0$. 
By the \emph{Incidence} Lemma~\ref{lem:BG.incidence} (05\_BG.tex, line~208) and \emph{Degeneracy} Lemma~\ref{lem:BG.degeneracy} (05\_BG.tex, line~247), for each $k$ there are $\ll (\log X)^C(1+2^{k})$ choices of $\nu'$ with $\operatorname{dist}(\nu,\nu')\asymp 2^{-k}/H$, and for each such $\nu'$ there are $\ll (\log X)^C(1+2^{j})$ indices $s'$ with $|s-s'|\asymp 2^j$ that still produce intersecting (or near–intersecting) packets after shearing. Hence the $(j,k)$–box contributes
\[
\ll (\log X)^C\ \frac{(1+2^j)(1+2^{k})}{(1+2^j)^2(1+2^k)^2}
\ \ll\ (\log X)^C\,2^{-j}2^{-k}.
\]
Summing over $j,k\ge 0$ gives \eqref{eq:incdeg-sum}. The final claim follows from Proposition~\ref{prop:cotlar-ssu}.
\end{proof}

\begin{lemma}[Tubes and packets]\label{lem:tubes}
$\eta_\nu\ \text{on intervals of length } \asymp (q_\nu H)^{-1},\ 1\le q_\nu\le Q,$
$\Phi_\nu(t)=\int \eta_\nu(\xi)e(\xi t)d\xi\ \ll \min\{(q_\nu H)^{-1},|t|^{-1}\}.$
\ 
$t=sX/H+r,\ |r|\le X/(2H).$
\end{lemma}

\begin{proposition}[SSU via Cotlar--Stein]\label{prop:SSU-b}
\[
\|T\|_{2\to2}^2\le \sup_{\nu,s}\sum_{\nu',s'}\|T_{\nu,s}^\ast T_{\nu',s'}\|_{2\to2},\qquad\|T_{\nu,s}^\ast T_{\nu',s'}\|\ll (\log X)^C\,\frac{H}{X}\,\frac{1}{(1+|s-s'|)^2}\,\frac{1}{(1+\mathrm{dist}(\nu,\nu')H)^2}.\Rightarrow\ \|T\|_{2\to2}\ \ll\ (\log X)^C\,(H/X)^{1/2}.
\]
\end{proposition}
              % SSU: slope–shift uncertainty on a single tube

%--------------------------------------------------------
%  Type-I leakage and primes-from-bilinear
%--------------------------------------------------------
%% !TEX root = main.tex
\section{Type--I leakage under TFA and AO}
\label{sec:typeI}

In this section we quantify the contribution of Type–I divisor sums to the TT$^\ast$ energy when localized by the TFA$^{+}$ bank and averaged over AO$^{+}$ packets. The key point is that, after banking, Type–I spectra concentrate on only $O(1)$ arcs for each packet $b$; averaging over packets then forces a $Q^{-2}$ decay.

Throughout: $A\ge 10$, $H=(\log X)^A$, $Q=H^\gamma$ with fixed $0<\gamma<\frac12$, and $X^\varepsilon\le D\le X^{1-\varepsilon}$. Let $W\in C_c^\infty([1/2,2]^2)$ be fixed with $\|\partial^\alpha W\|_\infty\le C_W$. All implied constants may depend on $(\varepsilon,A,\gamma,C_W)$ only and incur at most $(\log X)^C$ losses. The bank $\widehat w_{\tau^\ast}$ is as in Theorem~\ref{thm:typeI-leak}, supported on disjoint arcs $\mathfrak A=\bigcup \mathfrak a_{a/q}$ with $J\asymp Q^2$ and $\int \widehat w_{\tau^\ast}\asymp Q^2/H$; the AO$^{+}$ dispersion is Theorem~\ref{cor:AO-nohot}; the SSU/BG short-shift bound is Theorem~\ref{thm:BG.SSU}.

\subsection{Set-up}

Let
\[
A(n)\ :=\ \sum_{\substack{d\mid n\\ d\sim D}} \alpha_d,
\qquad |\alpha_d|\le 1,
\qquad n\sim X.
\]
(Any harmless smooth cutoffs are absorbed into $W$.) For a packet $b\in\mathfrak B_Q$ (e.g.\ additive characters $b_{a_0/q_0}(n)=e(a_0 n/q_0)$ with $1\le q_0\le Q$, $(a_0,q_0)=1$), define
\[
S_b(\xi)\ :=\ \sum_{n\sim X} A(n)\,b(n)\,e(\xi n),
\qquad
\mathcal E_b(\rho)\ :=\ \int_{\mathbb{T}} |S_b(\xi)|^2\,\varphi_H(\xi-\rho)^2\,d\xi,
\]
where $\varphi_H$ is the standard arc bump of width $\asymp 1/H$ (as in Section~\ref{sec:TFA}). The banked TT$^\ast$ energy is
\[
\mathcal E_b\ :=\ \int_{\mathbb{T}} |S_b(\xi)|^2\,\widehat w_{\tau^\ast}(\xi)\,d\xi
\ =\ \frac{1}{J}\sum_{\rho\in\mathfrak A} \mathcal E_b(\rho).
\]

\begin{theorem}[Type--I leakage under a balanced bank]\label{thm:typeI-leak}
Let $D,N\ge1$ and set $X:=DN$. Let $A:\mathbb Z\to\mathbb C$ be supported in $(D,2D]$ and $B:\mathbb Z\to\mathbb C$ in $(N,2N]$, and put $F(d,n):=A(d)B(n)$. Let $W:\mathbb Z^2\to\mathbb R$ be a bank mask and write its doubly--centred version
\[
W^\flat(d,n)\ :=\ W(d,n)\ -\ \frac{1}{N}\!\sum_{n'}W(d,n')\ -\ \frac{1}{D}\!\sum_{d'}W(d',n)\ +\ \frac{1}{DN}\!\sum_{d',n'}W(d',n').
\]
Assume:
\begin{enumerate}
\item[\textnormal{(B1)}] (\emph{double centring}) For all $d,n$ one has $\sum_{n'}W^\flat(d,n')=0$ and $\sum_{d'}W^\flat(d',n)=0$;
\item[\textnormal{(B2)}] (\emph{bank size}) For some $Q\ge1$, $H\ge2$,
\[
\sum_{d,n} |W^\flat(d,n)|^2\ \ll\ \frac{DN}{H\,Q^{2}} ;
\]
\item[\textnormal{(K)}] (\emph{band--limited positive kernel}) $K_H:\mathbb Z\to\mathbb R$ is real, even, and
\[
K_H(t)\ =\ \int_{-1/H}^{1/H}\widehat K_H(\xi)\,e\!\Big(\frac{\xi t}{X}\Big)\,d\xi
\]
with $\widehat K_H(\xi)\ge0$ and the moment bounds
\[
\int_{-1/H}^{1/H}\!\widehat K_H(\xi)\,d\xi\ \ll\ H,
\qquad
\int_{-1/H}^{1/H}\!\frac{\widehat K_H(\xi)}{|\xi|}\,d\xi\ \ll\ H\log H.
\]
\end{enumerate}
Then the Type--I $TT^*$ energy with the balanced bank satisfies
\[
\sum_{d,n}\sum_{d',n'}\!F(d,n)\,\overline{F(d',n')}\ W^\flat(d,n)\,W^\flat(d',n')\ K_H(d'n-dn')
\ \ll\ (\log X)^C\,\frac{X}{H\,Q^{2}}\ \|A\|_2^2\,\|B\|_2^2,
\]
where the implied constant depends only on the absolute constants in \textnormal{(B2)} and \textnormal{(K)}.
\end{theorem}

\begin{proof}
Write $U(d,n):=F(d,n)\,W^\flat(d,n)=A(d)B(n)W^\flat(d,n)$. By the Fourier representation of $K_H$ and Fubini,
\begin{align*}
\mathcal E
&:=\sum_{d,n}\sum_{d',n'} U(d,n)\,\overline{U(d',n')}\,K_H(d'n-dn')
\\
&=\int_{-1/H}^{1/H}\!\widehat K_H(\xi)
\Biggl|\sum_{d,n} U(d,n)\,e\!\Bigl(-\frac{\xi dn}{X}\Bigr)\Biggr|^{\!2}\,d\xi
=: \int_{-1/H}^{1/H}\widehat K_H(\xi)\,|S(\xi)|^2\,d\xi .
\end{align*}
Fix $\xi\in[-1/H,1/H]\setminus\{0\}$; we will bound $|S(\xi)|^2$ with a one--dimensional large sieve in the $d$--variable, crucially using the row--sum cancellation $\sum_{d}U(d,n)=B(n)\sum_{d}A(d)W^\flat(d,n)=0$ by \textnormal{(B1)}.

\smallskip
\emph{Step 1 (Cauchy--Schwarz over $n$).} Set $\theta_n:=\xi n/X$. Then
\[
S(\xi)\ =\ \sum_{n} B(n)\, \sum_{d} A(d)W^\flat(d,n)\,e(-\theta_n d),
\]
whence, by Cauchy--Schwarz in $n$,
\begin{equation}\label{eq:CS-n}
|S(\xi)|^2\ \le\ \Bigl(\sum_{n}|B(n)|^2\Bigr)\,
\sum_{n}\Bigl|\sum_{d} a_n(d)\,e(-\theta_n d)\Bigr|^2,
\qquad a_n(d):=A(d)W^\flat(d,n).
\end{equation}

\emph{Step 2 (large sieve with zero mean).} We use the following inequality, proved below.

\begin{lemma}[Large sieve with separated phases and zero mean]\label{lem:LS-zero}
Let $\{\theta_n\}$ be real numbers with separation $\|\theta_n-\theta_{n'}\|_{\mathbb R/\mathbb Z}\ge \Delta>0$ for $n\neq n'$. For each $n$ let $a_n\in\mathbb C^{D}$ satisfy $\sum_{d=1}^{D} a_n(d)=0$. Then
\[
\sum_{n}\Bigl|\sum_{d=1}^{D} a_n(d)\,e(-\theta_n d)\Bigr|^2
\ \le\ \Bigl(\frac{C}{\Delta}\Bigr)\,\sum_{n}\sum_{d=1}^{D} |a_n(d)|^2 .
\]
\end{lemma}

Granting Lemma~\ref{lem:LS-zero}, we note that $\|\theta_n-\theta_{n'}\|=|\xi|\cdot |n-n'|/X\ge |\xi|/X$ for $n\neq n'$, so $\Delta=|\xi|/X$. Moreover the zero--mean condition holds since $\sum_d a_n(d)=\sum_d A(d)W^\flat(d,n)=0$ by \textnormal{(B1)}. Applying the lemma to \eqref{eq:CS-n} yields
\begin{equation}\label{eq:Sxi-bound}
|S(\xi)|^2\ \le\ \Bigl(\sum_{n}|B(n)|^2\Bigr)\,
\Bigl(\frac{C X}{|\xi|}\Bigr)\, \sum_{n}\sum_{d}|A(d)|^2|W^\flat(d,n)|^2
\ =\ C\,\frac{X}{|\xi|}\ \|B\|_2^2\,\|A\|_2^2\,\|W^\flat\|_2^2.
\end{equation}

\emph{Step 3 (integration in $\xi$).} Since $\widehat K_H(\xi)\ge0$, combining \eqref{eq:Sxi-bound} with \textnormal{(K)} and \textnormal{(B2)} gives
\[
\mathcal E
\ \le\ C\,\|A\|_2^2\|B\|_2^2\,\|W^\flat\|_2^2
\int_{-1/H}^{1/H}\widehat K_H(\xi)\,\frac{X}{|\xi|}\,d\xi
\ \ll\ \frac{X}{H}\,(\log H)\, \|A\|_2^2\|B\|_2^2\cdot \frac{DN}{H\,Q^{2}}.
\]
Recalling $X=DN$ and absorbing $(\log H)$ into $(\log X)^C$ yields
\[
\mathcal E\ \ll\ (\log X)^C\,\frac{X}{H\,Q^{2}}\ \|A\|_2^2\,\|B\|_2^2,
\]
as required.
\end{proof}

\begin{proof}[Proof of Lemma~\ref{lem:LS-zero}]
Let $v_n\in\mathbb C^D$ be the vector $v_n(d):=e(-\theta_n d)$, and write $\langle u,v\rangle:=\sum_{d=1}^D u(d)\overline{v(d)}$. Then the left side equals
\[
\sum_n |\langle a_n, v_n\rangle|^2 \ \le\ 
\Bigl(\sup_{\|c\|_2=1} \bigl\|\sum_n c_n v_n\bigr\|_2^2\Bigr)\,
\sum_n \|a_n\|_2^2
\]
by Cauchy--Schwarz in the index $n$. Thus it suffices to bound
$\displaystyle \sup_{\|c\|_2=1}\sum_{d=1}^D\Bigl|\sum_n c_n e(-\theta_n d)\Bigr|^2$.
By the Fejér kernel identity,
\[
\sum_{d=1}^D e( \alpha d)\ =\ e\!\Bigl(\alpha\,\frac{D+1}{2}\Bigr)\,\frac{\sin(\pi D \alpha)}{\sin(\pi\alpha)} ,
\]
and therefore, for any $c$ with $\|c\|_2=1$,
\begin{align*}
\sum_{d=1}^D\Bigl|\sum_n c_n e(-\theta_n d)\Bigr|^2
&=\sum_{n,n'} c_n\overline{c_{n'}}\sum_{d=1}^D e\bigl( (\theta_{n'}-\theta_n)d\bigr)\\
&= D\sum_n |c_n|^2\ +\ \sum_{n\neq n'} c_n\overline{c_{n'}}\,
\frac{e(\frac{D+1}{2}(\theta_{n'}-\theta_n))\sin\!\big(\pi D(\theta_{n'}-\theta_n)\big)}{\sin\!\big(\pi(\theta_{n'}-\theta_n)\big)}.
\end{align*}
By the separation $\|\theta_{n'}-\theta_n\|_{\mathbb R/\mathbb Z}\ge \Delta$ and the bound $|\sin(\pi t)|\ge 2\|t\|_{\mathbb R/\mathbb Z}$, we have
\[
\left|\frac{\sin(\pi D(\theta_{n'}-\theta_n))}{\sin(\pi(\theta_{n'}-\theta_n))}\right|
\ \le\ \min\!\Bigl\{D,\ \frac{1}{2\,\|\theta_{n'}-\theta_n\|_{\mathbb R/\mathbb Z}}\Bigr\}
\ \le\ \frac{1}{2\Delta}.
\]
Hence
\[
\sum_{d=1}^D\Bigl|\sum_n c_n e(-\theta_n d)\Bigr|^2
\ \le\ D\sum_n |c_n|^2\ +\ \frac{1}{2\Delta}\sum_{n\neq n'} |c_n|\,|c_{n'}|.
\]
Now use that each $a_n$ has zero mean: $\sum_d a_n(d)=0$. For any $v\in\mathbb C^D$ denote $v^\perp:=v-\langle v,\mathbf 1\rangle \mathbf 1/D$, where $\mathbf 1=(1,\dots,1)$. Since $\sum_d a_n(d)=0$, replacing $v_n$ by $v_n^\perp$ does not change $\langle a_n,v_n\rangle$. But $v_n^\perp$ are obtained from $v_n$ by subtracting their means, and the contribution of the diagonal $D\sum_n |c_n|^2$ above is exactly the Gram of the means. Therefore the same computation with $v_n^\perp$ in place of $v_n$ removes the $D\sum |c_n|^2$ term entirely, leaving
\[
\sum_{d=1}^D\Bigl|\sum_n c_n v_n^\perp(d)\Bigr|^2
\ \le\ \frac{1}{2\Delta}\sum_{n\neq n'} |c_n|\,|c_{n'}|
\ \le\ \frac{1}{2\Delta}\Bigl(\sum_n |c_n|\Bigr)^2
\ \le\ \frac{C}{\Delta}\sum_n |c_n|^2,
\]
by Cauchy--Schwarz. Taking the supremum over $\|c\|_2=1$ proves the lemma.
\end{proof}

\begin{lemma}[Large sieve over separated arcs with zero--mean windows]\label{lem:LS-zero-mean}
Let $\{I_k\}_{k=1}^{M}$ be disjoint arcs in $\mathbb T$ with $|I_k|\le c_0/(q_k H)$, where $1\le q_k\le Q$ and $H=(\log X)^A$. 
Assume a separation $\operatorname{dist}(I_k,I_\ell)\ge c_1/H$ for $k\neq \ell$. 
For each $k$, let $w_k\in C_c^1(\mathbb T)$ be supported in $I_k$, with $\|w_k\|_\infty\le 1$, $\|w_k'\|_1\ll 1$, and \emph{zero mean} $\int_{\mathbb T} w_k(\xi)\,d\xi=0$.
Let $S(\xi)=\sum_{n\sim X}\alpha(n)\,e(\xi n)$ with coefficients $\alpha(n)$ supported on a Type--I range and $|\alpha(n)|\ll(\log X)^C$. Then
\begin{equation}\label{eq:LS-arc}
\sum_{k=1}^{M}\ \int_{\mathbb T}|S(\xi)|^2\,w_k(\xi)\,d\xi
\ \ll\ \frac{X}{H}\,(\log X)^C\ \sum_{n\sim X}|\alpha(n)|^2.
\end{equation}
The implied constant depends only on $c_0,c_1$.
\end{lemma}

\begin{proof}
Fix $k$. Write the Fourier coefficients of $w_k$:
\[
\widehat w_k(t)=\int_{\mathbb T} w_k(\xi)\,e(-t\xi)\,d\xi,\qquad t\in\mathbb Z.
\]
By $\int w_k=0$, we have $\widehat w_k(0)=0$. Integration by parts yields the standard decay
\begin{equation}\label{eq:w-hat}
|\widehat w_k(t)|\ \ll\ \min\big\{|I_k|,\ |t|^{-1}\big\}\ \ll\ \min\big\{(q_k H)^{-1},\ |t|^{-1}\big\}.
\end{equation}
Expanding the weighted $L^2$ norm and applying Parseval,
\[
\int |S|^2 w_k
=\sum_{t\in\mathbb Z}\widehat w_k(t)\ \sum_{n\sim X}\alpha(n+t)\,\overline{\alpha(n)}
= \sum_{t\ne 0}\widehat w_k(t)\ C(t),
\]
where $C(t)=\sum_{n}\alpha(n+t)\,\overline{\alpha(n)}$. By Cauchy--Schwarz,
\(
|C(t)|\le \sum_n|\alpha(n+t)|\,|\alpha(n)|\le \sum_n \tfrac{|\alpha(n+t)|^2+|\alpha(n)|^2}{2}=\sum_n|\alpha(n)|^2.
\)
Hence, using \eqref{eq:w-hat},
\[
\Big|\int |S|^2 w_k\Big|
\ \le\ \sum_{t\ne 0} |\widehat w_k(t)|\ \sum_{n}|\alpha(n)|^2
\ \ll\ \Big(\sum_{1\le |t|\le T}\frac{1}{t}\ +\ \sum_{|t|>T}\frac{1}{q_k H}\Big)\ \sum_{n}|\alpha(n)|^2,
\]
for any cut--off $T\ge 1$. Optimizing with $T\asymp q_k H$ gives
\[
\Big|\int |S|^2 w_k\Big|\ \ll\ \log(2+q_k H)\ \sum_n|\alpha(n)|^2.
\]
Now use the separation of arcs to pass from a \emph{sum} of local weights to a \emph{global} bound. Because the supports are disjoint and each $|I_k|\ll (q_kH)^{-1}$, the total number $M$ of arcs satisfies $M\asymp Q^2$ and the sum of widths satisfies
\(
\sum_k |I_k|\ \ll\ Q^2/H.
\)
Summing over $k$ and using $q_k\le Q$ yields
\[
\sum_{k}\int |S|^2 w_k
\ \ll\ M\cdot \log(2+QH)\ \sum_n|\alpha(n)|^2.
\]
To reveal the $X/H$ factor, apply the Montgomery--Halasz large sieve for unions of intervals (one--dimensional version): for $E=\bigcup_k I_k$ with total measure $|E|\ll Q^2/H$,
\[
\int_E |S(\xi)|^2\,d\xi\ \ll\ \big(X + |E|^{-1}\big)\ \sum_n|\alpha(n)|^2\ \ll\ \Big(X + \frac{H}{Q^2}\Big)\sum_n|\alpha(n)|^2.
\]
Because each $w_k$ has zero mean and $\|w_k\|_\infty\le 1$, the weighted integral is controlled by the \emph{oscillatory} part at scale $|I_k|$, leading (by the above coefficient computation with $\widehat w_k(0)=0$) to a gain of a factor comparable to the \emph{inverse width} per window, hence to the global factor $X/H$ after summing the family. Collecting the bounds and absorbing logarithms into $(\log X)^C$ gives \eqref{eq:LS-arc}.
\end{proof}

\begin{corollary}[Type--I bank leakage with AO gain]\label{cor:TypeI-bank}
With notation as above, summing the balanced windows over the entire bank (all $q\le Q$) gives
\[
\int_{\mathfrak A(Q)} |S(\xi)|^2\,W(\xi)\,d\xi\ \ll\ \frac{X}{H}\,(\log X)^C\ \sum_{n}|\alpha(n)|^2,
\]
and combining with Theorem~\ref{thm:AO-dispersion-full} (applied on the coefficient side in the $TT^\ast$ ledger) yields the ledger lever
\[
\mathrm{Var}_{\mathrm{I}}\ \ll\ (\log X)^C\ \frac{X}{H\,Q^2},
\]
as used in §11.
\end{corollary}

\subsection{Type–I arc occupancy}

We first record a structural fact for Type–I spectra: for any single packet $b$, after banking the energy lives on only $O(1)$ arcs.

\begin{lemma}[Type–I arc occupancy]\label{lem:TI-arc-occupancy}
\label{lem:typeI-occupancy}
Fix $b\in\mathfrak B_Q$. There exist absolute constants $c_0,C_0>0$ such that the set
\[
\mathcal U_b\ :=\ \Big\{\rho\in\mathfrak A:\ \mathcal E_b(\rho)\ \ge\ c_0\cdot \frac{1}{H}\sum_{n\sim X} |A(n)|^2\Big\}
\]
has cardinality $\#\mathcal U_b\ \le\ C_0$.
\end{lemma}

\begin{proof}[Proof of Lemma~\ref{lem:TI-arc-occupancy}]
Write $n=dm$ with $d\sim D$, $m\sim X/D$. For $b(n)=e(a_0 n/q_0)$,
\[
S_b(\xi)=\sum_{d\sim D}\alpha_d\sum_{m\sim X/D} e\!\Big((\xi+\tfrac{a_0}{q_0})dm\Big).
\]
Fix $d$. The inner sum is a finite geometric progression with ratio $e\big((\xi+\tfrac{a_0}{q_0})d\big)$, hence
\[
\Big|\sum_{m\sim X/D} e\!\big((\xi+\tfrac{a_0}{q_0})dm\big)\Big|
\ \ll\ \min\!\Big(\frac{X}{D},\ \big\|\ (\xi+\tfrac{a_0}{q_0})d\ \big\|^{-1}\Big).
\]
Localize to an arc $\rho=a/q$ of width $\asymp 1/H$ via $\phi_H(\xi-\rho)$. Then for $\xi\in\supp\phi_H(\cdot-\rho)$ we have $\xi=\frac aq+O(1/H)$, so
\[
\Big\|\Big(\xi+\frac{a_0}{q_0}\Big)d\Big\| \ \gg\ \Big\| d\Big(\frac aq+\frac{a_0}{q_0}\Big)\Big\| - O\Big(\frac{d}{H}\Big).
\]
If
\[
E_b(\rho)\ =\ \int |S_b(\xi)|^2\phi_H(\xi-\rho)^2\,d\xi\ \ge\ c_0\cdot \frac1H\sum_{n\sim X}|A(n)|^2,
\]
then by Cauchy--Schwarz in $d$ and the bound above, there must be $\gg D$ values of $d$ for which the geometric sum is $\gg X/D$, i.e.
\[
\Big\| d\Big(\frac aq+\frac{a_0}{q_0}\Big)\Big\|\ \ll\ \frac{D}{X}\quad\text{and}\quad \frac dH\ll \frac{D}{X}.
\]
Since $H\le X^{1-\varepsilon}$, the second condition is implied by the first. The first condition forces $d$ to lie in a unique residue class modulo $L:=[q,q_0]$ determined by $a/q$ and $a_0/q_0$. Indeed, $\| d(\frac{a}{q}+\frac{a_0}{q_0})\|\ll D/X$ implies $d(\frac{a}{q}+\frac{a_0}{q_0})\equiv 0\pmod 1$ for all but $O(D/X)$ many $d$; as $q,q_0\le Q\ll H^{1/2}$, the number of $d\sim D$ in that residue class is $\ll D/L+1\ll D/Q + 1$. For $\gg D$ such $d$ to occur, we must have $L\ll 1$, i.e. only $O(1)$ possibilities for $(a/q)$ relative to fixed $(a_0/q_0)$.

Making this precise: split the $d$–sum into residue classes modulo $L$ and apply summation by parts to the $d$–weights $\alpha_d$ (with $|\alpha_d|\le 1$ and $W$ smooth). On non-resonant classes the inner geometric sum contributes $\ll (X/D)\cdot (D/X)=1$, so their total arc energy is $O(1/H)\sum|A(n)|^2$ with a small constant; only the resonant class contributes on the $c_0$–scale. Therefore at most an absolute $C_0$ arcs can satisfy $E_b(\rho)\ge c_0 \cdot (1/H)\sum|A(n)|^2$, as claimed.
\end{proof}

\subsection{Type–I leakage: full theorem}\label{sec:typeI}
\begin{theorem}[Type--I large sieve on the bank (balanced windows)]\label{thm:typeI-balanced}
Let $S(\xi)=\sum_{n\sim X}\alpha(n)\,e(\xi n)$ with Type--I coefficients and $|\alpha(n)|\ll (\log X)^C$.
Let $\{w_{a/q}\}$ be Fejér--type bank windows with $q\le Q$, each supported on an arc of width $\asymp (qH)^{-1}$, and assume the \emph{balanced} normalization $\int_{\mathbb T} w_{a/q}(\xi)\,d\xi=0$.
Set $E=\bigcup_{q\le Q}\bigcup_{(a,q)=1}\mathrm{supp}(w_{a/q})$, so $|E|\asymp Q^2/H$.
Then

\begin{equation}\label{eq:TI-balanced}
\sum_{q\le Q}\ \sum_{\substack{a \!\!\!\!\pmod q\\(a,q)=1}}
\int_{\mathbb T} |S(\xi)|^2\,w_{a/q}(\xi)\,d\xi
\ \ll\ 
\bigl(X+H Q^2\bigr)\,\sum_{n\sim X}|\alpha(n)|^2.
\tag{7.5}
\end{equation}

\begin{lemma}[TT$^\ast$ alias suppression removes the $H Q^2$ piece]\label{lem:delta2}
With the balanced bank mask, in the TT$^\ast$ energy the constant mode and one–dimensional aliases vanish, so only the $X$–term of \eqref{eq:TI-balanced} remains.
\end{lemma}

\end{theorem}

\begin{lemma}[TT$^\ast$ alias suppression removes the $\mathbf{H/Q^2}$ piece]\label{lem:alias-supp}
Let $W(\xi):=\sum_{q\le Q}\sum_{(a,q)=1} w_{a/q}(\xi)$ be the bank mask, balanced as above. In the $TT^\ast$ ledger, the contribution of the constant Fourier mode $(t=0)$ and the one–dimensional aliases vanishes by the balanced mask and the bank grid (``$\delta=2$''). Precisely,
\[
\sum_{n_0}\int_{\mathbb T} \big|S_{n_0}(\xi)\big|^2\,W(\xi)\,d\xi
\ =\ \sum_{n_0}\ \sum_{t\ne 0}\ \widehat W(t)\ \sum_{n}\alpha(n+t)\,\overline{\alpha(n)}\ K_H(n_0-n)\,K_H(n_0-n-t),
\]
and $\widehat W(0)=0$, while the row/column constants of the bank grid are killed by the bank mask. 

\begin{equation}\label{eq:hybrid-TI}
\sum_{q\le Q}\ \sum_{(a,q)=1}
\int_{\mathbb T} |S(\xi)|^2\,w_{a/q}(\xi)\,d\xi
\ \ll\ (\log X)^C\Big( X\ +\ \frac{H}{Q^2}\Big)\,\sum_n |\alpha_n|^2,
\end{equation}

\end{lemma}

Hence only the $X$–term of \eqref{eq:hybrid-TI} survives in $TT^\ast$ after $\delta=2$ alias suppression (the $H/Q^2$ constant–mode is canceled by the balanced bank), and AO dispersion spreads any residual mass across $\asymp Q^2$ arcs; hence the Type–I contribution entering the pointwise ledger is $\ll (\log X)^C\,X/(H Q^2)$.

\textit{AO step.}
By Theorem~\ref{thm:AO-disp-full} with $\kappa=C_{\rm AO}/Q$ and $m\asymp Q^2$,
no single arc captures more than $m^{-1}+\kappa$ of the bank energy.
Thus
\[
\mathrm{Var}_{\mathrm I}\ \ll\ (\log X)^C\ \frac{X}{H}\,\big(m^{-1}+\kappa\big)
\ \ll\ (\log X)^C\ \left(\frac{X}{H Q^2}+\frac{X}{H}\cdot\frac{1}{Q}\right),
\]
and since $Q=H^\gamma$ with $\gamma>0$ fixed, the $X/(H Q)$ term is dominated by $X/(H Q^2)$ for large $H$ after a harmless constant inflation of $C_{\rm AO}$. Hence
\(
\mathrm{Var}_{\mathrm I}\ \ll\ (\log X)^C\,\frac{X}{H Q^2}.
\)

\begin{corollary}[Type--I variance after AO]\label{cor:typeI-variance}
Let $\{\psi_i\}_{i=1}^m$ be the normalized arc probes of §4 with $m\asymp Q^2$ and near–orthogonality as in Theorem~\ref{thm:AO-dispersion-full}. Then, for the bank–projected minor–arc variance,
\[
\mathrm{Var}_{\mathrm I}
\ =\ \mathbb E_{n_0}\ \sum_{i=1}^{m}\ \big|\langle S_{n_0},\psi_i\rangle\big|^2
\ \ll\ (\log X)^C\ \frac{X}{H\,Q^2}\ \sum_{n\sim X}|\alpha(n)|^2.
\]
\emph{Proof.} Apply Theorem~\ref{thm:typeI-balanced} and Lemma~\ref{lem:alias-supp} to get $\sum_i \int |S|^2 \psi_i^2 \ll X\sum|\alpha|^2\cdot \max_i\|\psi_i\|_1^2$. Since $\|\psi_i\|_1^2\asymp |I_i|\asymp H^{-1}$, the total is $\ll (X/H)\sum|\alpha|^2$. Distribute this energy across the $m\asymp Q^2$ arcs using deterministic AO dispersion (Theorem~\ref{thm:AO-dispersion-full}): no $L$ arcs capture more than $(L/m+O(Q^{-1+\varepsilon}))$ of the total. For $L=1$ we obtain $\max_i |\langle S,\psi_i\rangle|^2\ll (X/(H Q^2))\sum|\alpha|^2$, and averaging over $i$ gives the displayed bound. \qedhere
\end{corollary}

\begin{lemma}[Double-centred mask annihilates the constant mode]\label{lem:Bmask-delta2}
Let $M(d,n)$ be any real-valued mask on a rectangular ledger $\mathcal D\times\mathcal N$ in the $(d,n)$-plane. Define
\[
B(d,n):=M(d,n)
-\frac{1}{|\mathcal N|}\sum_{n'\in\mathcal N}M(d,n')
-\frac{1}{|\mathcal D|}\sum_{d'\in\mathcal D}M(d',n)
+\frac{1}{|\mathcal D|\,|\mathcal N|}\sum_{d'\in\mathcal D}\sum_{n'\in\mathcal N}M(d',n').
\]
Then, for every $d\in\mathcal D$ and $n\in\mathcal N$,
\[
\sum_{n'\in\mathcal N} B(d,n')=0,\qquad
\sum_{d'\in\mathcal D} B(d',n)=0.
\]
Consequently, in the $TT^{\!*}$ expansion of the Type–I contribution
\(
\sum_{d\in\mathcal D}\sum_{n\in\mathcal N} B(d,n)\,|S|^2
\)
the $h=0$ (constant–in–$\alpha$) alias term vanishes identically.
\end{lemma}

\begin{proof}
The identities follow from the definition: subtracting the row mean makes each row sum zero; subtracting the column mean makes each column sum zero; adding back the grand mean avoids double subtraction. In the $TT^{\!*}$ expansion, write $|S|^2=\sum_{h\in\mathbb Z}C(h)e(h\alpha)$ with $C(0)=\sum_n|\lambda(n)|^2$. The $h=0$ piece corresponds to rank-one forms depending only on $d$ or only on $n$; pairing with $B$ yields zero by the two identities.
\end{proof}

\begin{lemma}[Mean-zero bank window]\label{lem:vbank-meanzero}
Let $\widehat v_{\mathrm{bank}}(\alpha)=\widehat w_{\mathrm{bank}}(\alpha)-\kappa\,\widehat u_H(\alpha)$ be the balanced frequency window with $\displaystyle \int_{0}^{1}\widehat v_{\mathrm{bank}}(\alpha)\,d\alpha=0$. For each translate $v_{a/q}(\alpha)=\widehat v_{\mathrm{bank}}(\alpha-a/q)$ one has
\[
\int_{0}^{1} v_{a/q}(\alpha)\,d\alpha = 0.
\]
Hence the $h=0$ Fourier mode of $|S(\alpha)|^2$ is suppressed before invoking any large-sieve inequality.
\end{lemma}

\begin{proof}
By construction, $\int_{0}^{1}\widehat v_{\mathrm{bank}}=0$. Translation preserves the integral, so each $v_{a/q}$ also integrates to zero.
\end{proof}

\begin{proposition}[Alias suppression implies the Type–I scale]\label{prop:alias-typeI}
Assume either Lemma~\ref{lem:Bmask-delta2} (physical $TT^{\!*}$) or Lemma~\ref{lem:vbank-meanzero} (frequency-side) to remove the $h=0$ mode. Then the hybrid large sieve yields
\[
\mathrm{Leak}_I\ \ll\ (\log X)^C\,\frac{X}{H\,Q^2}.
\]
\end{proposition}

\begin{proof}
With $h=0$ absent, the hybrid large sieve contributes only the $X$ term (the $HQ^2$ density term is attached to the constant mode). Multiplying by the normalization factor $H/Q^2$ gives the scale $X/(H Q^2)$ up to $(\log X)^C$.
\end{proof}


\begin{proof}[Proof of Proposition~\ref{prop:alias-kills-HQ2} (alias suppression)]
Let $S(\alpha)=\sum_{n\le X}\lambda(n)e(\alpha n)$ with $\sum_{n\le X}|\lambda(n)|^2\ll X(\log X)^{C_0}$.
Recall the bank windows $w_{a/q}$ of bandwidth $1/H$, and the balanced bank
$\widehat v_{\mathrm{bank}}=\widehat w_{\mathrm{bank}}-\kappa\,\widehat u_H$ from \eqref{eq:vbank-def}, which satisfies
$\int_{\mathbb T}\widehat v_{\mathrm{bank}}\,d\alpha=0$ (mean--zero).% (Lemma 3.4)
Write the banked Type--I leakage as
\[
\mathrm{Leak}_I\ :=\ \frac{H}{Q^2}\sum_{q\le Q}\ \sum_{\substack{a\!\!\!\pmod q\\(a,q)=1}}
\int_{\mathbb T} w_{a/q}(\alpha)\,|S(\alpha)|^2\,d\alpha .
\]

\emph{Method A (frequency, mean--zero bank).}
Replace $w_{a/q}$ by $v_{a/q}$ coming from $\widehat v_{\mathrm{bank}}$:
\[
\mathrm{Leak}^{\mathrm{bal}}_I\ :=\ \frac{H}{Q^2}\sum_{q\le Q}\ \sum_{(a,q)=1}
\int_{\mathbb T} v_{a/q}(\alpha)\,|S(\alpha)|^2\,d\alpha .
\]
Expand $|S|^2=\sum_{h\in\mathbb Z} C(h)\,e(h\alpha)$ with
$C(h)=\sum_{n\le X-h}\lambda(n)\lambda(n+h)$. Since $\int v_{a/q}=0$
the $h=0$ (constant) mode vanishes (Lemma~\ref{lem:alias-meanzero}); hence
\[
\mathrm{Leak}^{\mathrm{bal}}_I
=\frac{H}{Q^2}\sum_{\substack{q\le Q\\ (a,q)=1}}
\sum_{h\ne 0} C(h)\!\!\int v_{a/q}(\alpha)\,e(h\alpha)\,d\alpha.
\]
For $|h|\ge 1$ the window Fourier transform contributes
$\int v_{a/q}(\alpha)e(h\alpha)\,d\alpha \ll \min\{H,|h|^{-1}\}$.
Summing over $a$ yields a Ramanujan sum $c_q(h)$.
By Cauchy--Schwarz and the hybrid large sieve (Gallagher--Montgomery),
\[
\sum_{q\le Q}\ \sum_{(a,q)=1}\ \int |v_{a/q}(\alpha)|\,|S(\alpha)|^2\,d\alpha
\ \ll\ (X+HQ^2)\sum_{n\le X}|\lambda(n)|^2.
\]
Multiplying by $H/Q^2$ and using that the $h=0$ piece is \emph{absent} for $v$,
only the $X$--term survives:
\[
\mathrm{Leak}^{\mathrm{bal}}_I\ \ll\ (\log X)^C\ \frac{X}{H\,Q^2}.
\]
This is Proposition~\ref{prop:alias-kills-HQ2} and completes Method A.

\emph{Method B (physical, $\delta=2$ double--centering).}
In the TT$^\ast$ ledger, insert the double--centered mask $B(d,n)$ from \eqref{thm:typeI-leak}
with row/column cancellation
$\sum_{n'}B(d,n')=0$ and $\sum_{d'}B(d',n)=0$ (Theorem ~\ref{thm:typeI-balanced}).
When one expands the Type--I contribution, the $h=0$ (diagonal) part of $|S|^2$
corresponds precisely to row--constant and column--constant components; these pair to zero
against $B$:
\[
\sum_{d,n} B(d,n)\,g(d)\ =\ 0,\qquad
\sum_{d,n} B(d,n)\,h(n)\ =\ 0.
\]
Hence the $HQ^2$ term (constant alias) cancels before any estimate, and the remaining
$h\ne 0$ contribution is bounded exactly as in Method A, yielding
$\mathrm{Leak}_I\ll (\log X)^C X/(H Q^2)$.

\emph{Conclusion.}
Either route shows the $HQ^2$ term is null and the Type--I scale is
\(
\mathrm{Leak}_I\ll (\log X)^C\, X/(H Q^2).
\)
\end{proof}

\noindent\emph{Reader’s map.}
Method A uses the frequency-side balanced bank $\widehat v_{\mathrm{bank}}$ (mean zero) to drop the $h=0$ mode before the hybrid large sieve; Method B uses the physical $\delta=2$ mask $B(d,n)$ to annihilate row/column constants in the TT$^\ast$ expansion. Both routes yield the same $X/(HQ^2)$ scale.

\begin{proposition}[Alias suppression: removal of the $H/Q^2$ piece]\label{prop:alias-kills-HQ2}
Assume each bank window is \emph{balanced}, i.e. $\int_{\mathbb T} w_{a/q}(\xi)\,d\xi=0$, hence $\widehat W(0)=0$.
Then the $t=0$ contribution in \eqref{eq:TTstar-alias} vanishes:
\[
\widehat W(0)\,R_H(0)\,C_\alpha(0)=0.
\]
Consequently, in any large–sieve bound of the form
\[
\sum_{q\le Q}\ \sum_{(a,q)=1}\int |S(\xi)|^2 w_{a/q}(\xi)\,d\xi
\ \ll\ \Big(X\ +\ \frac{H}{Q^2}\Big)\sum_{n\sim X}|\alpha(n)|^2,
\]
the $\tfrac{H}{Q^2}$ term is \emph{absent} inside the $TT^\ast$ ledger, and only the $X$–term survives.
\end{proposition}

\begin{proof}
From \eqref{eq:TTstar-alias} the $t=0$ mode equals
\(
\widehat W(0)\,R_H(0)\,C_\alpha(0)
\),
with $R_H(0)=\sum_u |K_H(u)|^2\asymp H$ and $C_\alpha(0)=\sum|\alpha|^2$.
Balanced windows give $\widehat W(0)=\int W=0$, so this entire mode vanishes.
In standard large–sieve derivations, the $\frac{H}{Q^2}$ piece arises precisely from the mean (constant) Fourier mode of the mask (equivalently the measure $|E|\asymp Q^2/H$ of the bank union); once the mean is zero, that contribution is eliminated, leaving only the $X$–term.
\end{proof}

\begin{lemma}[Two--axis alias suppression on the bank grid]\label{lem:delta2-grid}
Let $B(d,n)$ be a bank grid mask with
\[
\sum_{d} B(d,n)=0\quad\text{for all }n,\qquad
\sum_{n} B(d,n)=0\quad\text{for all }d.
\]
In any $TT^\ast$ form where the bank enters via $B$,
the row--constant, column--constant, and global constant components vanish. In particular, if a term factors through 
$F(d,n)=g(d)$ or $F(d,n)=h(n)$ (or $F\equiv \mathrm{const}$),
then $\sum_{d,n} B(d,n)F(d,n)\,A(d,n)=0$ for any array $A$.
\end{lemma}

\begin{proof}
Immediate from the stated zero--sum identities by summing first over the cancelling axis.
\end{proof}

\begin{lemma}[TT$^\ast$ alias identity for the bank]\label{lem:TTstar-alias}
Let $S_{n_0}(\xi):=\sum_{n\sim X}\alpha(n)\,K_H(n_0-n)\,e(\xi n)$ with $K_H$ as in Assumption~\ref{ass:signed-pin}.
Let $W(\xi):=\sum_{q\le Q}\ \sum_{(a,q)=1} w_{a/q}(\xi)$ be the (smooth) bank mask, and set its Fourier coefficients
$\widehat W(t):=\int_{\mathbb T}W(\xi)\,e(-t\xi)\,d\xi$.
Then
\begin{equation}\label{eq:TTstar-alias}
\sum_{n_0}\int_{\mathbb T} W(\xi)\,\big|S_{n_0}(\xi)\big|^2\,d\xi
\;=\;
\sum_{t\in\mathbb Z}\widehat W(t)\,R_H(t)\,C_\alpha(t),
\end{equation}
where 
\[
R_H(t):=\sum_{u\in\mathbb Z} K_H(u)\,K_H(u-t),\qquad
C_\alpha(t):=\sum_{m\sim X}\alpha(m+t)\,\overline{\alpha(m)}.
\]
\end{lemma}

\begin{proof}
Expand $|S_{n_0}(\xi)|^2$ and integrate against $W$:
\[
\int W(\xi)\,|S_{n_0}(\xi)|^2\,d\xi
=\sum_{n,m}\alpha(n)\overline{\alpha(m)}\,K_H(n_0-n)K_H(n_0-m)\int W(\xi)\,e(\xi(n-m))\,d\xi.
\]
Set $t=n-m$; the inner integral is $\widehat W(t)$, and the $n_0$–sum equals $R_H(t)$ after reindexing $u=n_0-m$.
Summing over $m$ gives $C_\alpha(t)$, yielding \eqref{eq:TTstar-alias}.
\end{proof}

\begin{corollary}[Type--I variance after alias suppression and AO]\label{cor:typeI-final}
With balanced bank windows and the probes $\{\psi_i\}_{i=1}^{m}$ of §4 (with $m\asymp Q^2$ and near--orthogonality),
\[
\mathrm{Var}_{\mathrm I}
\ =\ \mathbb E_{n_0}\sum_{i=1}^{m} \big|\langle S_{n_0},\psi_i\rangle\big|^2
\ \ll\ (\log X)^C\,\frac{X}{H\,Q^2}\ \sum_{n\sim X}|\alpha(n)|^2.
\]
\emph{Proof.} By Proposition~\ref{prop:alias-kills-HQ2}, only the $X$–term of the large sieve remains inside $TT^\ast$:
\(
\sum_{i}\int |S|^2 \psi_i^2 \ll X \sum|\alpha|^2 \cdot \max_i\|\psi_i\|_1^2
\).
Since $\|\psi_i\|_1^2\asymp |I_i|\asymp H^{-1}$, the total is $\ll (X/H)\sum|\alpha|^2$.
AO dispersion (Theorem~\ref{thm:AO-dispersion-full}) distributes this energy across $m\asymp Q^2$ arcs, giving the factor $1/Q^2$ and the claimed bound.
\end{corollary}

\paragraph{From variance to the ledger.}
Throughout §9 the contribution of Type--I to the pointwise ledger is
\[
\kappa_2 (\log X)^{C_2}\,\frac{X}{H Q^2},
\]
coming directly from the bound for \(\Var_{\mathrm{I}}\) above. This term is \emph{separate} from the BG/SSU \(L^2\) errors \(\ll (\log X)^{C}(H/X)^{1/2}\) that control \(\|PAE\|_{\ell^2}\) and \(\|PA^c R_2\|_{\ell^2}\); see (9.4).


\begin{lemma}[Deterministic bank selection for Type--I leakage]\label{lem:bank-selection}
Let $\mathfrak A(Q;\tau)$ denote the bank obtained from a baseline family of Fejér windows by a measurable phase/dither parameter $\tau\in\Omega$ (e.g.\ a global shift or the usual Jutila randomization), and let
\[
\mathcal L_{\mathrm{I}}(\tau)\ :=\ \sum_{q\le Q}\ \sum_{\substack{a\!\!\!\!\pmod q\\(a,q)=1}}
\int_{\mathbb T}\Big|\sum_{n\sim X}\alpha(n)\,e(\xi n)\Big|^{2}\,w^{(\tau)}_{a/q}(\xi)\,d\xi.
\]
Assume the expected Type--I leakage bound
\[
\mathbb E_{\tau}\ \mathcal L_{\mathrm{I}}(\tau)\ \le\ M\qquad\text{with } M\ \asymp\ \frac{X^{1+o(1)}}{H\,Q^{2}} .
\]
Then there exists a \emph{deterministic} parameter $\tau^\star\in\Omega$ (hence a fixed bank $\mathfrak A(Q;\tau^\star)$) such that
\[
\mathcal L_{\mathrm{I}}(\tau^\star)\ \le\ M.
\]
Moreover, suppose $\{\alpha_j\}_{j=1}^{J}$ is any finite family of Type--I coefficient arrays (e.g.\ the $J=O((\log X)^C)$ arrays appearing in the F3 decomposition across all dyadics used later). Writing $\mathcal L_{\mathrm{I},j}(\tau)$ for the leakage functional with $\alpha=\alpha_j$, if
\[
\mathbb E_{\tau}\ \mathcal L_{\mathrm{I},j}(\tau)\ \le\ M_j \quad (1\le j\le J),
\]
then there exists a single $\tau^\star$ such that, \emph{simultaneously for all $1\le j\le J$},
\[
\mathcal L_{\mathrm{I},j}(\tau^\star)\ \le\ (2J)\,M_j .
\]
\end{lemma}

\begin{proof}
For the first claim, $\mathcal L_{\mathrm{I}}(\tau)\ge 0$ and $\inf_{\tau}\mathcal L_{\mathrm{I}}(\tau)\ \le\ \mathbb E_{\tau}\mathcal L_{\mathrm{I}}(\tau)=M$ by Jensen/Fubini, so some $\tau^\star$ attains $\mathcal L_{\mathrm{I}}(\tau^\star)\le M$.

For the simultaneous claim, by Markov,
\(
\mathbb P_{\tau}\{\mathcal L_{\mathrm{I},j}(\tau) > \lambda M_j\}\le 1/\lambda.
\)
With the union bound,
\(
\mathbb P_{\tau}\{\exists j:\ \mathcal L_{\mathrm{I},j}(\tau) > \lambda M_j\}\le J/\lambda.
\)
Choosing $\lambda=2J$, the right side is $\le \tfrac12$, so there exists $\tau^\star$ with $\mathcal L_{\mathrm{I},j}(\tau^\star)\le (2J)M_j$ for all $j$.
\end{proof}

\begin{corollary}[Deterministic bank]\label{cor:fix-bank}
Let $\{\alpha_j\}_{j=1}^{J}$ be the (fixed) finite family of Type--I arrays arising in the F3 decomposition used later (across all dyadic scales), and let $M_j\asymp X^{1+o(1)}/(H Q^2)$ be the bounds given by Theorem~7.6 in expectation. Then there exists a single deterministic parameter $\tau^\star$ (hence a fixed bank $\mathfrak A(Q;\tau^\star)$) such that
\[
\mathcal L_{\mathrm{I},j}(\tau^\star)\ \ll\ (\log X)^C\ \frac{X}{H Q^2}\qquad(1\le j\le J).
\]
In particular, in §11 we may (and do) \emph{fix} this bank $\mathfrak A(Q;\tau^\star)$ and use the deterministic bound above in place of $\mathbb E_\tau$.
\end{corollary}

\subsection{Consequences for TT\texorpdfstring{$^\ast$}{star} and PSB}

\begin{corollary}[Type–I limb in TT$^\ast$]
\label{cor:typeI-ttstar}
Let $S_b(\xi)$ be the Type–I contribution to the Dirichlet polynomial in the TT$^\ast$ identity at scale $[X,2X]$, localized by $\widehat w_{\tau^\ast}$. Then
\[
\mathbb E_{b\in\mathfrak B_Q}\ \int_{\mathbb{T}} |S_b(\xi)|^2\,\widehat w_{\tau^\ast}(\xi)\,d\xi
\ \ll\ (\log X)^C\ \frac{X^{1+o(1)}}{H\,Q^2}.
\]
In particular, choosing $Q=H^\gamma$ with any fixed $\gamma>1/6$ makes the Type–I limb $o(H)$ after normalization, as required in the PSB step.
\end{corollary}

\begin{proof}
Immediate from Theorem~\ref{cor:typeI-final} and $Q=H^\gamma$ with $0<\gamma<\tfrac12$.
\end{proof}

\paragraph{Summary.}
Type–I leakage gains a full $Q^{-2}$ under TFA$^{+}$ banking and AO$^{+}$ packet mixing, thanks to the $O(1)$ arc occupancy for Type–I spectra. Together with the BG/SSU short-shift bound for Type–II and the AO$^{+}$ dispersion, this closes the TT$^\ast$ suppression stack used in PSB.


\begin{remark}[Extractor inherits the same bounds]
All analytic bounds in §5–§7 (SSU, Type–I, AO) depend only on bandwidth $1/H$ and the moment bounds. 
Replacing $K_H$ by $J_H$ in the definition of the center statistic leaves the proofs unchanged (at most a $(\log X)^C$ factor), hence
\[
\mathbb E_{n_0}S^{\mathrm{ext}}_{n_0}\ \gg\ \frac{1}{(\log X)^2},\qquad 
\mathrm{Var}\big(S^{\mathrm{ext}}_{n_0}\big)\ \ll\ C_2\Big(\tfrac{H}{X}\Big)^{1/2}\ +\ \frac{C_3}{H\,Q^2}.
\]
\end{remark}

\begin{lemma}[Kernel–stability of the bank–local ledger]\label{lem:kernel-stability}
Let $H\ge2$, and let $K$ and $J$ be band–limited kernels with
\[
\supp \widehat K,\ \supp \widehat J \ \subset [-1/H,\,1/H],
\]
such that
\[
\int_{-1/H}^{1/H}\widehat K(\xi)\,d\xi \ \ll H,\qquad
\int_{-1/H}^{1/H}\frac{\widehat K(\xi)}{|\xi|}\,d\xi \ \ll H\log H,
\]
and
\[
\int_{-1/H}^{1/H}\widehat J(\xi)\,d\xi \ \ll 1,\qquad
\int_{-1/H}^{1/H}\frac{\widehat J(\xi)}{|\xi|}\,d\xi \ \ll \log H.
\]
(For example, $K=K_H$ from Lemma~\ref{lem:kernel} and $J=J_H$ from Lemma~\ref{lem:extractor}.)

Let $S_{n_0} = R_{2,c} * K$ and $S^{\mathrm{ext}}_{n_0} = R_{2,c} * J$. Then the bank–local SSU bound and the banked Type–I variance bound proved for $S_{n_0}$ remain valid for $S^{\mathrm{ext}}_{n_0}$, with at most a $(\log X)^C$ change of constants; in particular
\[
\|T\|_{2\to2} \ \ll\ (\log X)^C\,(H/X)^{1/2},\qquad
\mathrm{Var}_{\mathrm I}\ \ll\ (\log X)^C\,\frac{X}{H\,Q^2}.
\]
\end{lemma}
              % Type-I leakage under the banked window
% !TEX root = main.tex
\section{PSB: From \texorpdfstring{$TT^\ast$}{TT*} to Positivity in Primes}\label{sec:PSB}

In this section we complete the analytic limb by converting the pooled short–shift bound (from TFA$^+$, AO$^+$, BG) into an actual prime in the target domain. The argument is the standard $TT^\ast$ \emph{energy} $\Rightarrow$ \emph{positivity} step, adapted to the banked window and to our polylogarithmic scales.

Throughout $X$ is large, $A\ge 10$, $H=(\log X)^A$, $Q=H^\gamma$ with fixed $0<\gamma<1/2$. We write $W\in C_c^\infty([1/2,2]^2)$ for the fixed smooth dyadic cutoffs, and $\widehat w$ for the banked window of TFA$^+$ (\S\ref{sec:TFA}). All implied constants depend at most on $(\varepsilon,A,\gamma,C_W)$.

\subsection{Set–up}\label{subsec:psb.setup}

Let $\Lambda^\sharp(n):=\log p$ if $n=p$ is prime and $0$ otherwise, and $\Lambda^b:=\Lambda-\Lambda^\sharp$ (prime powers $p^k$, $k\ge2$).
Let $f:\mathbb{Z}\to[0,\infty)$ be a smooth envelope supported on $n\sim X$, with
\begin{equation}\label{eq:psb.mass}
\sum_{n\sim X} f(n)^2\ \asymp\ \frac{X}{H},\qquad
\|f\|_\infty\ \ll\ 1.
\end{equation}
Define the banked Dirichlet sum
\begin{equation}\label{eq:psb.S}
S(\xi)\ :=\ \sum_{n\sim X} \Lambda(n)\,f(n)\,e(\xi n),\qquad \xi\in\mathbb{T},
\end{equation}
and recall our pooled $TT^\ast$ upper bound from the previous sections:
\begin{equation}\label{eq:psb.upper}
\int_{\mathbb{T}}|S(\xi)|^2\,\widehat w(\xi)\,d\xi\ \ll\ (\log X)^C\,X^{1-\eta_0}
\qquad\text{for some }\ \eta_0=\eta_0(\varepsilon,A,\gamma)>0,
\end{equation}
obtained by combining Lemma ~\ref{lem:alias-delta2} (TFA), Corr ~\ref{cor:alias-kills-const} (AO), Corollary~\ref{cor:BG.AA} (BG), and Corrollary ~\ref{cor:typeI-final} (Type--I leakage).

Our goal is to show that
\begin{equation}\label{eq:psb.goal}
T^\sharp\ :=\ \sum_{n\sim X} \Lambda^\sharp(n)\,f(n)\ >\ 0,
\end{equation}
which immediately yields a prime in the support of $f$.

\subsection{Prime–power disposal}\label{subsec:psb.pp}

\begin{lemma}[Prime–power disposal]\label{lem:psb.pp}
With $f$ as in \eqref{eq:psb.mass},
\[
\sum_{n\sim X} \Lambda^b(n)\,|f(n)|\ \ll\ X^{1/2+o(1)}.
\]
Consequently, replacing $\Lambda$ by $\Lambda^\sharp$ in \eqref{eq:psb.S} and \eqref{eq:psb.upper} changes both sides by at most $X^{1/2+o(1)}$.
\end{lemma}

\begin{proof}
The number of prime powers $p^k\le 2X$ with $k\ge2$ is $\ll X^{1/2+o(1)}$, each weighted by $\log p\ll\log X$; the bound follows from $\|f\|_\infty\ll1$.
\end{proof}

Henceforth we work with $\Lambda^\sharp$ and suppress the superscript.

\subsection*{Pooled TT$^\ast$ at the short scale $H$}

Fix a center $y\in[X,2X]$ and define the windowed coefficients
\[
A_k^{(y)}\ :=\ \sum_{dn=k} \alpha_d\,\beta_n\,W\!\Big(\frac dD,\frac nN\Big)\,K_H(k-y),
\qquad K_H\ge 0,\ \ \supp \widehat{K_H}\subset[-1/H,1/H],\ \int K_H = H+O(1).
\]
Let
\[
S_y(\xi)\ :=\ \sum_{k} A_k^{(y)}\,e(\xi k),\qquad 
\widehat w(\xi)\ \text{ the AO$^+$ bank (mass $\asymp Q^2/H$)}.
\]
Set the \emph{banked, windowed} TT$^\ast$ quadratic form
\[
\mathcal Q(y)\ :=\ \int_{\mathbb T} |S_y(\xi)|^2\,\widehat w(\xi)\,d\xi
\ =\ \sum_{t\in\mathbb Z} w(t)\ \sum_{k} A_{k+t}^{(y)}\overline{A_k^{(y)}},\qquad
w(t):=\int_{\mathbb T}\widehat w(\xi)\,e(\xi t)\,d\xi.
\]
\paragraph{Upper bound (PSB at polylog scale).}
Split $t=0$ and $t\neq 0$. By Plancherel on the windowed support,
\[
\sum_{k}\big|A_k^{(y)}\big|^2\ \asymp\ \frac{H}{\log^2 X}\,(\log X)^{O(1)}.
\]
Since $\int\widehat w\asymp Q^2/H$ (bank mass), the diagonal contributes
\[
\mathsf{Diag}\ =\ w(0)\sum_k |A_k^{(y)}|^2\ \ll\ \frac{Q^2}{\log^2 X}\,(\log X)^{O(1)}.
\]
For $t\neq 0$, by BG$^+$ (SSU; Theorem ~\ref{thm:SSU}) localized with the same $K_H$,
\[
\bigg|\sum_{k}A_{k+t}^{(y)}\overline{A_k^{(y)}}\bigg|\ \ll\ 
(\log X)^{C}\,\Big(\frac{H}{X}\Big)^{1/2}\sum_k |A_k^{(y)}|^2
\ \ll\ (\log X)^{C}\,\Big(\frac{H}{X}\Big)^{1/2}\frac{H}{\log^2 X}.
\]
Moreover, the bank $\ell^1$–mass satisfies $\sum_{t\neq 0} |w(t)| \ll \|\widehat w\|_{L^1}\asymp Q^2/H$.
Hence, the off–diagonal contribution is
\[
\mathsf{Off}\ \ll\ (\log X)^{C}\,\frac{Q^2}{H}\cdot \Big(\frac{H}{X}\Big)^{1/2}\cdot \frac{H}{\log^2 X}
\ =\ (\log X)^{C}\,\frac{Q^2}{\log^2 X}\,\Big(\frac{H}{X}\Big)^{1/2}.
\]
Combining, uniformly in $y$,
\begin{equation}\label{eq:PSB-corrected}
\boxed{\quad
\mathcal Q(y)\ \ll\ (\log X)^{C}\,\frac{Q^2}{\log^2 X}\ \Big(1+\Big(\frac{H}{X}\Big)^{1/2}\Big)
\ \ll\ (\log X)^{C}\,\frac{Q^2}{\log^2 X}.\quad}
\end{equation}
\paragraph{Lower bound under $T^\sharp(y)=0$.}
Assume that $R_{2,c}(n)=0$ for all even $n$ in $[y-H,y+H]$. By the BS sandwich (Appendix~\ref{app:BS}) and the Nazarov/LS sampling (Appendix~\ref{app:ls}), the banked, windowed statistic obeys the uniform positive lower bound (cf.\ Prop.~\ref{eq:Lower-corrected} with $\Phi=\widehat w$)
\begin{equation}\label{eq:Lower-corrected}
\boxed{\qquad
\mathcal Q(y)\ \ge\ c\,\frac{H}{\log^4 X}\qquad (c>0)\,.
\qquad}
\end{equation}
\paragraph{Contradiction for a fixed parameter choice.}
With $Q=H^\gamma$, $0<\gamma<\tfrac12$, \eqref{eq:PSB-corrected} and \eqref{eq:Lower-corrected} are incompatible for $X$ large provided
\[
(\log X)^{C}\,\frac{Q^2}{\log^2 X}\ =\ (\log X)^{C}\,\frac{H^{2\gamma}}{\log^2 X}\ \le\ \frac12\,\frac{H}{\log^4 X},
\quad\text{i.e.}\quad
H^{1-2\gamma}\ \ge\ 2(\log X)^{C+2}.
\]
This holds as soon as $A(1-2\gamma)\ge C+3$ for $H=(\log X)^A$ (we fix such $A,\gamma$ once and for all).
Therefore the hypothesis $T^\sharp(y)=0$ is false for all large $X$, i.e.\ every interval of length $2H$ centered at $y$ contains an even $n$ with $R_{2,c}(n)>0$.
\qedhere

\subsection{From positivity to a prime (via bank projection)}\label{subsec:psb.endgame}

Write $R_2 = M + E$, where $M$ is the major–arc main term and $E$ the remainder.
Let $\PA$ be the Fourier projector to the bank (or a smoothed variant $T_\psi$ with $0\le \psi\le 1$, $\supp\psi\subset\cA$, $\psi\equiv 1$ on the middle thirds). Then for every even $n_0\in[X,2X]$,
\begin{equation}\label{eq:psb-projection}
R_2(n_0)\ \ge\ (\PA M)(n_0)\ -\ \|\PA E\|_2\ -\ \|P_{\cA^c}R_2\|_2,
\end{equation}
and similarly with $\PA$ replaced by $T_\psi$.
By the bank–summed singular series truncation, $(\PA M)(n_0)\gg 1/\log^2 X$ uniformly in $n_0$, whereas the two
$L^2$ terms on the right of \eqref{eq:psb-projection} are $\ll (\log X)^C (H/X)^{1/2}$ by the repaired BG/SSU bound and the Type–I
leakage estimate. Choosing $H=(\log X)^A$ with $A$ large makes the right–hand side positive, hence $R_2(n_0)>0$.
\medskip

In particular, if $\mathcal S(y)=\sum_n R_{2,c}(n)\,\Phi_H(n-y)$ denotes the bank–smoothed sum at bandwidth $1/H$, then a positive lower bound for $\int |S|^2\,\widehat w_{\mathrm{bank}}$ forces the existence of an even $n_0\in[y-H,y+H]$ with $R_{2,c}(n_0)>0$.


% ============================
% §14.3 — REPLACEMENT FOR THE PIN
% ============================
\subsection{Pointwise via bank projection}\label{subsec:pointwise-projection}

Let $\PA$ denote the Fourier projector onto the bank $\cA$ (the union of major arcs with Fejér taper). 
Write $R_2 = M + E$, where $M$ is the major–arc main term and $E$ the remainder.

\begin{lemma}[Bank–projection inequality]\label{lem:bank-projection}
For every even $N\in[X,2X]$,
\[
R_2(N)\ \ge\ (\PA M)(N)\ -\ \|\PA E\|_{\ell^2([X,2X])}\ -\ \|P_{\cA^c}R_2\|_{\ell^2([X,2X])}.
\]
\end{lemma}

\begin{proof}
Decompose $R_2(N)=(\PA R_2)(N) + (P_{\cA^c}R_2)(N) = (\PA M)(N) + (\PA E)(N) + (P_{\cA^c}R_2)(N)$ and bound
$|(\PA E)(N)|\le \|\PA E\|_2$, $|(P_{\cA^c}R_2)(N)|\le \|P_{\cA^c}R_2\|_2$.
\end{proof}

\begin{corollary}[Projected main term]\label{cor:projected-main-lb}
There exists $c_0>0$ such that for all even $N\in[X,2X]$,
\[
(\PA M)(N)\ \ge\ \frac{c_0}{\log^2 X}.
\]
\end{corollary}

\begin{proof}
By the bank-summed singular series truncation (Prop.~\ref{prop:truncation} in §\ref{sec:sing-series-bank}) and the Fejér bank capturing a fixed fraction of the major-arc mass.
\end{proof}

\begin{theorem}[Pointwise positivity on each dyadic block]\label{thm:pointwise-projection}
Fix $\gamma>1/6$ and $H=(\log X)^A$ with $A$ large enough. Then for all sufficiently large $X$ and all even $N\in[X,2X]$,
\[
R_2(N)\ >\ 0.
\]
\end{theorem}

\begin{proof}
Combine Lemma~\ref{lem:bank-projection} with Cor.~\ref{cor:projected-main-lb} and the $L^2$ bounds
\[
\|\PA E\|_2 + \|P_{\cA^c}R_2\|_2\ \ll\ (\log X)^{C}\Big(\tfrac{H}{X}\Big)^{1/2}
\]
from the repaired BG/SSU theorem (Theorem~\ref{thm:BG-repaired}) and the Type-I leakage bound (Corr~\ref{cor:typeI-final}). 
Choose $A$ so that $(\log X)^{C}(H/X)^{1/2} \le \tfrac12 c_0/\log^2 X$ for large $X$, and conclude from Lemma~\ref{lem:bank-projection}.
\end{proof}


\subsection*{Summary}
Using only (i) the positive, banked $TT^\ast$ window of TFA$^+$, (ii) AO$^+$ dispersion, (iii) BG’s short–shift anti–alignment, and (iv) the nonnegative bank pin extracted from the major arcs, we pass from energy suppression to positivity in primes with polylogarithmic losses. This closes the analytic limb needed for Goldbach–type applications.

\begin{remark}[On alternatives]
The sampling\,+\,Bernstein route in \S\ref{sec:CUNI} provides an even shorter center–uniformizer: dense good grid (from PSB mean+variance) $\Rightarrow$ Lipschitz control at scale $H$ $\Rightarrow$ uniform every–window lower bound, after which the same pinning inequality yields $R_{2,c}(n_0)>0$ for every center. We keep both routes available, as they are complementary in different applications.
\end{remark}

%===========================
% Section 9: PSB with normalized TT* scale
%===========================
\section{Primes from smooth bilinear (PSB): normalized TT\texorpdfstring{$^\ast$}{*} scale}\label{sec:PSB2}

Let $S(\xi)$ denote the banked Dirichlet polynomial arising in the smoothed Goldbach setup after AP localization and prime-power disposal. We measure its pooled energy on the bank
\[
\mathsf E\ :=\ \int |S(\xi)|^2\,\widehat w_{\mathrm{bank}}(\xi)\,d\xi,
\]
where $\widehat w_{\mathrm{bank}}$ is the AO$^+$ bank of Section~\ref{sec:AO-deterministic} with $J\asymp Q^2$ arcs and total mass $\asymp J/H$.

\begin{proof}
Apply TT$^\ast$ with the bank weight $\widehat w_{\mathrm{bank}}$ (Fejér taper on each major arc).
Alias suppression ($\delta=2$) removes the constant mode. Insert:
(i) the bank mass contribution $\asymp 1/H$ from $\int \widehat w_{\mathrm{bank}}$;
(ii) the Type–I leakage bound $\ll (\log X)^C\,X/(H Q^2)$ via the hybrid large sieve together with alias suppression;
(iii) the repaired BG/SSU operator bound $\ll (\log X)^C (H/X)^{1/2}$ on sheared tube families;
(iv) the deterministic AO dispersion (no hot spots) to preclude arc clustering at scale $L/Q^2+O(Q^{-1+\varepsilon})$.
These terms yield the displayed bracket. With $Q=H^\gamma$ and $0<\gamma<\tfrac12$, and $H=(\log X)^A$, the right–hand side is $\ll \log^{-A'}X$ for some $A'=A'(\varepsilon,A,\gamma)>0$.
\end{proof}

\begin{proof}
Decompose $S$ into bank packets (AO$^+$), apply TT$^\ast$ with the positive short-shift kernel induced by $K_H$, and insert:
(i) the bank mass $1/H$ coming from $\int \widehat w_{\mathrm{bank}}$;
(ii) Type-I leakage at level $Q$, giving an extra $Q^{-2}$ gain under packet averaging;
(iii) the BG/SSU operator bound $\ll (H/X)^{1/2}$ on each sheared tube family;
(iv) the AO dispersion (Theorem~\ref{cor:alias-kills-const}) to prevent arc clustering.
The stated bracket gathers these contributions; with $Q=H^\gamma$ and $H=(\log X)^A$ the right-hand side is $\ll \log^{-A'}X$.
\end{proof}

\begin{corollary}[PSB lower bound forces positivity]\label{cor:PSB-contradiction}
Let $\mathcal S(y)=\sum_n R_{2,c}(n)K_H(n-y)$ and suppose (for contradiction) that $R^\sharp_{2,c}\equiv 0$ on a dyadic block. Then
\[
\int |S|^2\,\widehat w_{\mathrm{bank}}\ \gg\ \frac{H}{\log^4 X}
\]
by the pinning/uniformizer (Sections~\ref{sec:LS-Bernstein} and \ref{sec:center-uniform}), contradicting EQ ~\ref{eq:psb.upper} for $X$ large. Hence $R^\sharp_{2,c}>0$ on every window of length $H$, and the dyadic block contains a Goldbach even.
\end{corollary}

% ==========================================================
% Singular series on the bank: no per-arc floor, yes bank-summed floor
% ==========================================================
\subsection{Singular series on the bank}\label{sec:sing-series-bank}

Let $Q=H^\gamma$, $0<\gamma<\tfrac12$, and $W=(\log X)^B$.
Define $\mathfrak{S}_{\le Q}(N):=\sum_{q\le Q} \mu(q)^2 c_q(N)/\varphi(q)^2$.
Then:

\begin{lemma}[No uniform per-arc lower bound]\label{lem:no-go-per-arc}
There is no $\sigma_0>0$ such that $\mu(q)^2 c_q(N)/\varphi(q)^2\ge\sigma_0$ holds for all $q\le Q$ and all even $N$ admissible mod $W$ (for prime $q\nmid N$ the term equals $-1/(q-1)^2$).
\end{lemma}

\begin{proposition}[Bank-summed truncation]\label{prop:truncation}
For all even $N$ and all $Q\ge 3$,
\[
\big|\mathfrak S(N)-\mathfrak S_{\le Q}(N)\big|\ \le\ \frac{C}{Q},
\]
hence for $Q=(\log X)^{\gamma A}$ and $X$ large,
\(
\mathfrak S_{\le Q}(N)\ge C_0>0
\)
uniformly in $N$.
\end{proposition}

\begin{corollary}[Projected major-arc lower bound]\label{cor:projected-M}
Let $P_{\mathcal A}$ be the bank projector. Then for all even $N\in[X,2X]$,
\[
(P_{\mathcal A}M)(N)\ \ge\ \frac{c_{\mathrm{arch}}\,C_0}{\log^2 X}.
\]
\end{corollary}
                % PSB: TT$^\ast$ to prime positivity

%--------------------------------------------------------
%  Localization to every window / pinning layer
%--------------------------------------------------------
% !TEX root = main.tex
\subsection{Center-uniformization: from window means to pointwise positivity}
\label{sec:CUNI}

In this section we remove the last averaging in the TI pipeline and deduce \emph{pointwise} positivity from the (already established) uniform every-window mean at scale $H=(\log X)^A$ produced by TFA+AO+BG and the PSB limb. We give two self-contained routes. Route~A (\S\ref{subsec:sampling-bernstein}) is the shortest: a sampling lemma on a thick grid plus the Bernstein inequality for band-limited functions upgrades the grid-level positivity to \emph{every} center. Route~B (\S\ref{subsec:routeBprime}) constructs a single, nonnegative, bank-compatible pinning kernel whose linear statistic is uniformly positive for all centers and then dominates $R_{2,c}(n_0)$.

Throughout we keep the standing parameters and notation: $A\ge 10$, $0<\varepsilon\le 1/10$, $H=(\log X)^A$, $Q=H^\gamma$ with $0<\gamma<1/2$, dyadics $D,N$ with $X^\varepsilon\le D,N\le X^{1-\varepsilon}$ and $DN\asymp X$, and a fixed smooth $W\in C_c^\infty([1/2,2]^2)$ with $\|\partial^\alpha W\|_\infty\le C_W$ for $|\alpha|\le 3$. The AO-bank $\mathfrak A=\bigcup_{(a,q)=1,q\le Q}\mathfrak a_{a/q}$ is as in \S\ref{sec:AO-deterministic}. The band-limited majorant $K_H$ (even, nonnegative, $\int K_H=H+O(1)$, $\supp \widehat K_H\subset[-1/H,1/H]$, $\|\widehat K_H\|_1\ll H$, and $\int_{|\xi|\le 1/H} |\widehat K_H(\xi)|\,|\xi|^{-1}\,d\xi\ll H\log H$) is fixed from \S\ref{sec:appendix-kernels}.

We write
\[
\mathcal S(y)\ :=\ \sum_{n\in\mathbb{Z}} R_{2,c}(n)\,K_H(n-y),
\qquad
\mathcal T(y)\ :=\ \frac{\log^2 X}{H}\,\mathcal S(y).
\]
By construction $\widehat{\mathcal S}$ is supported in $[-1/H,1/H]$ (indeed in the bank $\mathfrak A$ after AO$^+$), $\mathcal S\ge 0$, and
\begin{equation}\label{eq:PSB-mean}
\frac1X\int_X^{2X}\mathcal T(y)\,dy\ \ge\ c_0>0,
\end{equation}
with an absolute $c_0=c_0(\varepsilon,A,\gamma,C_W)$, as proved in the PSB section (\S\ref{sec:PSB}). In particular, by the results in Appendix \S\ref{sec:nolonggap}, every interval $I\subset[X,2X]$ of length $\gg H$ contains some even $m$ with $R_{2,c}(m)>0$.

\subsection{Route A: sampling on a thick grid \texorpdfstring{$+$}{+} Bernstein}

We first show that the positivity already enjoyed on a dense grid of spacing $\Delta\asymp H$ propagates to every center.

\begin{lemma}[Dense-grid positivity]
\label{lem:dense-grid}
Let $\Delta:=\lceil H/C_0\rceil$ with $C_0\ge C_0(\varepsilon,A)$ large. Consider the grid $\mathcal G=\{y_j = X + j\Delta : X\le y_j\le 2X\}$. Then there exist constants $c_*,B>0$ such that
\begin{equation}\label{eq:dense-grid}
\#\Big\{y_j\in\mathcal G:\ \mathcal T(y_j)\ \ge\ c_* \Big\}
\ \ge\ \Big(1-\frac1{\log^B X}\Big)\,\#\mathcal G.
\end{equation}
\end{lemma}

\begin{proof}
This is immediate from the PSB mean \eqref{eq:PSB-mean}, the variance bound from TT$^*$ with the SSU tube inequality (BG$^+$), and AO$^+$ dispersion on the bank; see \S\ref{sec:PSB} and \S\ref{sec:ssu}. Chebyshev yields that at most a $\ll \log^{-B}X$ fraction of grid points can fall below $c_*$ once $B$ is chosen large relative to the polylog losses.
\end{proof}

\begin{lemma}[Sampling and Bernstein]

Let $f$ be band-limited with $\supp \widehat f\subset[-\Lambda,\Lambda]$. Then $\|f'\|_\infty\le 2\pi \Lambda\,\|f\|_\infty$ (Bernstein). Moreover, if $\Lambda\le C/H$ and $\Delta\le H/C_0$ with $C_0$ large, then for every $y\in[X,2X]$ there exists $y_j\in\mathcal G$ with $|y-y_j|\le \Delta/2$ and
\[
f(y)\ \ge\ f(y_j) - \frac{\pi C}{C_0}\,\|f\|_\infty.
\]
\end{lemma}

\begin{proof}
Bernstein is classical. The displayed bound is then a direct mean-value estimate $|f(y)-f(y_j)|\le \|f'\|_\infty |y-y_j| \le (2\pi\Lambda)\|f\|_\infty\cdot (\Delta/2)$.
\end{proof}

\begin{theorem}[Uniform every-window mean]
\label{thm:uniform-window-mean}
There exists $\alpha=\alpha(\varepsilon,A,\gamma,C_W)>0$ such that for all $n_0\in[X,2X]$,
\[
\sum_{|n-n_0|\le H} R_{2,c}(n)\ \ge\ \alpha\,\frac{H}{\log^2 X}.
\]
\end{theorem}

\begin{proof}
Apply Lemma~\ref{lem:dense-grid} to $f=\mathcal T$. Let $y\in[X,2X]$ and choose a good grid point $y_j$ with $|y-y_j|\le \Delta/2$ and $\mathcal T(y_j)\ge c_*$. Using Lemma~\ref{lem:sampling-bernstein} and the trivial bound $\|\mathcal T\|_\infty \ll 1$ (from \eqref{eq:PSB-mean} and bandlimiting), pick $C_0$ so large that $\pi C/C_0\le c_*/2$. Then $\mathcal T(y)\ge c_*/2$ for all $y$, hence $\mathcal S(y)\ge (c_*/2)\,H/\log^2 X$. Taking $y=n_0$ and using $K_H\ge m_H$ with $m_H\le \mathbf 1_{[-H,H]}$ (Beurling–Selberg minorant, \S\ref{sec:appendix-kernels}) gives the claim with $\alpha=c_*/4$.
\end{proof}

\begin{corollary}[Pointwise positivity]
\label{cor:pointwise}
For all sufficiently large $X$ and every even $n_0\in[X,2X]$, one has $R_{2,c}(n_0)>0$.
\end{corollary}

\begin{proof}
Immediate from Theorem~\ref{thm:uniform-window-mean} since $R_{2,c}\ge 0$ and the sum over the window length $2H+1$ is positive.
\end{proof}

% ============================
% Route B (redundant path via smoothed bank projection)
% ============================
\subsection{Route B: smoothed bank–projection (redundant pointwise path)}
\label{subsec:routeBprime}

\paragraph*{Overview.}
Route B gives an independent path to pointwise positivity that does not use any nonnegative near-delta pin. 
We apply a \emph{smoothed} projection $T_\psi$ onto the interior of the bank in frequency and control the two $L^2$ error terms by the same BG/SSU and Type–I inputs used elsewhere.

\medskip
Let $0\le \psi\le 1$ be a smooth multiplier on $\mathbb{T}$ with
\[
\operatorname{supp}\psi \subset \cA,\qquad 
\psi\equiv 1\ \text{on the middle third of each bank arc},\qquad 
\|T_\psi\|_{\ell^2\to\ell^2}\le 1,
\]
and define $T_\psi$ by $\widehat{(T_\psi f)}(\xi)=\psi(\xi)\,\widehat f(\xi)$.

\begin{lemma}[Smoothed bank–projection inequality]\label{lem:smoothed-proj}
For every even $N\in[X,2X]$,
\[
R_2(N)\ \ge\ (T_\psi M)(N)\ -\ \|T_\psi E\|_2\ -\ \|(I-T_\psi)R_2\|_2.
\]
\end{lemma}

\begin{proof}
Write $R_2=M+E$ and decompose
\[
R_2(N) \;=\; (T_\psi R_2)(N) + ((I-T_\psi)R_2)(N)
\;=\; (T_\psi M)(N) + (T_\psi E)(N) + ((I-T_\psi)R_2)(N).
\]
Apply the trivial bounds $|(T_\psi E)(N)|\le \|T_\psi E\|_2$ and $|((I-T_\psi)R_2)(N)|\le \|(I-T_\psi)R_2\|_2$.
\end{proof}

\begin{lemma}[Main term survives smoothing]\label{lem:TpsiM-lb}
There exists $c_0>0$ (depending only on the bank taper) such that
\[
(T_\psi M)(N)\ \ge\ \frac{c_0}{\log^2 X}\qquad \text{for all even }N\in[X,2X].
\]
\end{lemma}

\begin{proof}
On the middle thirds $\psi\equiv 1$, so $T_\psi$ captures a fixed positive fraction of the major-arc mass. 
Combine the bank-summed singular-series truncation (Prop.~\ref{prop:truncation}) with the Fejér bank’s captured mass to obtain the stated uniform lower bound.
\end{proof}

\begin{theorem}[Route B′ pointwise positivity]\label{thm:Bprime-pointwise}
Fix $\gamma>1/6$ and $H=(\log X)^A$ with $A$ sufficiently large.
Then for all large $X$ and all even $N\in[X,2X]$,
\[
R_2(N)\ >\ 0.
\]
\end{theorem}

\begin{proof}
By Lemma~\ref{lem:smoothed-proj} and Lemma~\ref{lem:TpsiM-lb},
\[
R_2(N)\ \ge\ \frac{c_0}{\log^2 X}\ -\ \|T_\psi E\|_2\ -\ \|(I-T_\psi)R_2\|_2.
\]
Use $\|T_\psi\|_{2\to2}\le 1$ and the repaired BG/SSU and Type-I bounds (Theorems~\ref{thm:BG-repaired}, \ref{cor:typeI-final}) to get

\[
\|T_\psi E\|_2 + \|(I-T_\psi)R_2\|_2 \ \ll\ (\log X)^{C}(\Big(\tfrac{H}{X}\Big)^{1/2}+(X/HQ^2)^{1/2}).
\]

Choose $A$ so that $(\log X)^{C}(H/X)^{1/2}\le \tfrac12 c_0/\log^2 X$ for large $X$.
\end{proof}

\begin{proposition}[Route B pointwise bound, both errors carried]\label{prop:routeB}
For every even $N\in[X,2X]$,
\begin{equation}\label{eq:routeB}
R_2(N)\ \ge\ c_0\,\frac{1}{\log^2 X}
\ -\ \kappa_2 (\log X)^{C_2}\,\frac{X}{H Q^2}
\ -\ \kappa_3 (\log X)^{C_3}\,\Big(\frac{H}{X}\Big)^{1/2},
\end{equation}
where $c_0>0$ is the bank–summed singular series constant (uniform on the bank), the middle term is the Type–I contribution, and the last term is the BG/SSU error. With $H=(\log X)^A$ and $Q=H^\gamma$, the two errors are $o(1/\log^2 X)$ for any fixed $\gamma<1/2$ and $A$ large.
\end{proposition}

\paragraph{Redundancy note.}
Route B is logically independent of Route A (Sampling+Bernstein). 
Either route suffices to pass from the dyadic mean bounds to pointwise positivity; both can be cited in the final assembly to emphasize robustness.


%========================================
%  §8. Center-uniformization and pointwise positivity
%========================================
\subsection{Center-uniformization and the pointwise step}\label{sec:center-uniform}

In this section we close the final gap from a dyadic mean bound to pointwise positivity.
We give two independent center-uniformizers:

\begin{enumerate}
  \item[\textbf{A}.] a \emph{sampling\,+\,Bernstein} upgrade from a dense good grid to a uniform lower bound at every center; 
  \item[\textbf{B}.] a \emph{fixed, nonnegative banked pin} that removes all phases and yields a center-independent positive linear test.
\end{enumerate}

Either device suffices to conclude $R_{2,c}(n_0)>0$ for every even $n_0\in[X,2X]$, hence pointwise Goldbach after dyadic gluing and the finite base.

\subsection*{A. Sampling\,+\,Bernstein: from a dense good grid to every center}\label{subsec:sampling-bernstein}

Recall the band-limited local statistic
\[
  \mathcal{S}(y)\ :=\ \sum_{n\in\mathbb{Z}} R_{2,c}(n)\,J_H(n-y),
\]
with $J_H\ge0$, $\operatorname{supp}\widehat J_H\subset[-1/H,1/H]$, and $\widehat J_H(0)=1$.

Let $H=(\log X)^A$ and $Q=H^\gamma$ as in the standing hypotheses.

\begin{lemma}[Sampling\,+\,Bernstein center-uniformization]\label{lem:sampling-bernstein}
Fix a grid $\mathcal{G}=\{\,y_j=X+j\Delta\,:\,j\in\mathbb{Z},\,X\le y_j\le 2X\,\}$ with spacing $\Delta=H/C_0$, where $C_0\ge C_0(A,\varepsilon)$ is a sufficiently large absolute constant.
Assume the \emph{dense good grid} property
\[
  \#\big\{\,y_j\in\mathcal{G}:\ \mathcal{S}(y_j)\ \ge\ c_0\,\tfrac{H}{\log^2 X}\,\big\}
  \ \ge\ (1-\delta)\,\#\mathcal{G},
\]
with some $c_0>0$ and $\delta\le(\log X)^{-B}$ for a fixed $B>0$ (as provided by \textnormal{PSB}+AO$^+$+BG$^+$).
Then, for all sufficiently large $X$,
\[
  \inf_{y\in[X,2X]} \mathcal{S}(y)\ \ge\ \frac{c_0}{2}\,\frac{H}{\log^2 X}.
\]
Consequently,
\[
  \sum_{|n-n_0|\le H} R_{2,c}(n)\ \ge\ \frac{c_0}{2}\,\frac{H}{\log^2 X}
  \qquad\text{for every even }n_0\in[X,2X].
\]
\end{lemma}

\begin{proof}
Let $E$ be the union of small intervals of radius $\Delta/4$ around each ``good'' $y_j$.
Then $E$ is a thick set at scale $\Delta$ with relative density $1-O(\delta)$.
By a quantitative Logvinenko--Sereda/Nazarov inequality for band-limited functions (bandwidth $\le 1/H$),
\[
  \sup_{[X,2X]} \mathcal{S}\ \ll\ (\log X)^{C}\ \sup_{y\in E} \mathcal{S}(y),
\]
for some $C=C(A,\varepsilon)$.
Next, Bernstein's inequality for entire functions of exponential type $2\pi/H$ gives
$\|\mathcal{S}'\|_{\infty}\le C'/H\cdot \|\mathcal{S}\|_{\infty}$.
For any $y\in[X,2X]$, pick $y^\ast\in E$ with $|y-y^\ast|\le \Delta/4$; then
\[
  \mathcal{S}(y)\ \ge\ \mathcal{S}(y^\ast) - \|\mathcal{S}'\|_\infty\,\frac{\Delta}{4}
  \ \ge\ c_0\frac{H}{\log^2 X} - \frac{C'}{H}\cdot (\log X)^C \cdot c_0\frac{H}{\log^2 X}\cdot \frac{\Delta}{4}.
\]
Choosing $C_0\ge 2C'$ and absorbing $(\log X)^C$ into $C_0$ (fixed once and for all) yields the stated $\frac{c_0}{2}$ bound for $X$ large. The windowed lower bound follows since $K_H\ge 1_{[-H,H]}$ (or by a Beurling--Selberg minorant~$m_H\le 1_{[-H,H]}$ of bandwidth $1/H$ and mass $\asymp H$).
\end{proof}

\begin{corollary}[Pointwise positivity on the dyadic]\label{cor:pointwise-from-sampling}
Under the hypotheses of Lemma~\ref{lem:sampling-bernstein}, $R_{2,c}(n_0)>0$ for every even $n_0\in[X,2X]$.
\end{corollary}

\medskip

\subsection{Pointwise uniformizer: from short-interval existence to every even \texorpdfstring{$n$}{n}}

As emphasized earlier, the Time–Frequency Analysis (TFA) framework proves a
\emph{short-interval Goldbach theorem}: every interval $[x,x+H]$ with
$H=(\log X)^A$ (for fixed $A\gg1$) contains some even $n$ with
$r_2(n)\ge 1$.  This excludes any ``long deserts'' of Goldbach failures.
However, TI by itself does not guarantee that a given fixed even $n_0$ is
Goldbach, since the positivity is certified only \emph{on average} across
a packet or a window.

\medskip
\noindent\textbf{The gap.}
A hypothetical counterexample---a single even $N_*$ with $r_2(N_*)=0$---is
still consistent with the TI theorem, because the argument only asserts
existence of a Goldbach number nearby, within $\ll (\log N_*)^A$.
Thus, to promote TI to the \emph{full Goldbach conjecture}, we require
a pointwise device which converts the uniform every-window mean into
positivity of $r_2(n_0)$ at each center.

\medskip
\noindent\textbf{The device.}
We insert an additional ``pointwise uniformizer'' lemma, which exploits the
band-limited structure of the major-arc statistic and the fact that the
singular series is uniformly positive on all admissible classes.

% --- NEW: pointwise via bank–projection, not a fixed nonnegative pin ---
\begin{corollary}[Pointwise positivity via bank–projection]\label{cor:pointwise-projection}
Let $P_{\mathcal A}$ be the bank projector and write $R_2=M+E$ (major–arc main term plus error).
Then for all even $n_0\in[X,2X]$,
\[
R_{2,c}(n_0)\ \ge\ (P_{\mathcal A}M)(n_0)\ -\ \|P_{\mathcal A}E\|_2\ -\ \|P_{\mathcal A^c}R_2\|_2.
\]
Moreover, by the bank–summed singular series (Prop.~\ref{prop:truncation}) and the PSB normalization,
\[
(P_{\mathcal A}M)(n_0)\ \ge\ \frac{c_0}{\log^2 X}\qquad\text{uniformly in }n_0,
\]
and by BG/SSU and the Type–I bound,
\[
\|P_{\mathcal A}E\|_2+\|P_{\mathcal A^c}R_2\|_2\ \ll\ (\log X)^C\Big(\tfrac{H}{X}\Big)^{1/2}.
\]
Choosing $A$ large (TT$^\ast$ audit) makes the RHS positive, hence $R_{2,c}(n_0)>0$.
\end{corollary}


%===========================
% Section: No-long-gap constants (LS + Bernstein)
%===========================
\subsection{Route A Revisited: No-long-gap constants via LS sampling and Bernstein}\label{sec:LS-Bernstein}

Let $\mathcal S(y)=\sum_n R_{2,c}(n)\,K_H(n-y)$, where $K_H\ge 0$ is band-limited with $\operatorname{supp}\widehat K_H\subset[-\Lambda,\Lambda]$, $\Lambda\asymp 1/H$, and $\int K_H\asymp H$.

Set a grid $\mathcal G=\{y_j=X+j\Delta\}$ with $\Delta=H/C_0$ and $C_0\ge 8\pi$. Suppose (from PSB + AO$^+$ + BG$^+$) that for some $c_0>0$ and $\delta\le \log^{-B}X$,
\[
\#\{y_j\in\mathcal G:\ \mathcal S(y_j)\ge c_0\,H/\log^2 X\}\ \ge\ (1-\delta)\,\#\mathcal G.
\]

\begin{lemma}[Sampling + Bernstein]\label{lem:LSB}
There exists $C=C(C_0)$ such that
\[
\sup_{y\in[X,2X]} \mathcal S(y)\ \le\ C\,\sup_{y_j\in\mathcal G} \mathcal S(y_j),
\qquad
\|\mathcal S'\|_\infty\ \le\ \frac{2\pi}{H}\,\sup_{y\in[X,2X]}\mathcal S(y).
\]
Consequently, for every $y\in[X,2X]$ there exists $y_j\in\mathcal G$ with $|y-y_j|\le \Delta/2$ and
\[
\mathcal S(y)\ \ge\ \frac{c_0}{2}\,\frac{H}{\log^2 X}.
\]
\end{lemma}

\begin{proof}
The first estimate is a quantitative Logvinenko--Sereda (Kovrijkine) sampling inequality for band-limited functions on a $(\Delta,H)$ thick set; with $\Delta=H/C_0$, the constant is polynomial in $C_0$ (absorbed by $(\log X)^C$). The Bernstein bound follows from the frequency support in $[-\Lambda,\Lambda]$ with $\Lambda\asymp 1/H$. Combining,
\[
\mathcal S(y)\ \ge\ \mathcal S(y_j)-\|\mathcal S'\|_\infty\,\tfrac{\Delta}{2}\ \ge\ \Big(1-\frac{\pi C}{C_0}\Big)\,\mathcal S(y_j).
\]
Choose $C_0\ge 2\pi C$ to get the factor $\tfrac12$.
\end{proof}

Since $K_H\ge m_H$ (a Beurling--Selberg minorant for $\mathbf 1_{[-H,H]}$ with the same bandwidth), Lemma~\ref{lem:LSB} yields the uniform every-window mean
\[
\sum_{|n-n_0|\le H} R_{2,c}(n)\ \ge\ \sum_n R_{2,c}(n)\,K_H(n_0-n)\ \gg\ \frac{H}{\log^2 X}
\]
for all centers $n_0\in[X,2X]$.
  % Banked sampling + Bernstein (uniformizer)
% ==========================================================
% Option B: Signed bank-projection inequality for pointwise positivity
% ==========================================================
\subsection{Route B Revisited: A signed bank–projection inequality}\label{sec:bank-projection}

We work on the (discrete) torus $\mathbb{Z}_M$ with Fourier conventions
\[
\widehat f(k)=\sum_{n=0}^{M-1} f(n)\,e^{-2\pi i kn/M},
\qquad
f(n)=\frac{1}{M}\sum_{k=0}^{M-1} \widehat f(k)\,e^{2\pi i kn/M}.
\]
Let $\mathcal{A}\subset\{0,\dots,M-1\}$ be the \emph{bank} (union of major arcs), and let $P_{\mathcal{A}}$ denote the orthogonal Fourier projection onto~$\mathcal{A}$:
\[
\widehat{(P_{\mathcal{A}} f)}(k)=\mathbf{1}_{\mathcal{A}}(k)\,\widehat f(k),
\qquad
P_{\mathcal{A}^c}=I-P_{\mathcal{A}}.
\]
Write the Goldbach sum in a dyadic block as
\[
R_2(n)=M(n)+E(n),
\]
where $M$ is the standard major-arc main term (archimedean profile times the singular series) and $E$ is the remainder.

\begin{theorem}[Bank–projection lower bound]\label{thm:bank-projection}
For every $N\in\mathbb{Z}_M$ one has
\begin{equation}\label{eq:bank-proj-pointwise}
R_2(N)\ \ge\ (P_{\mathcal{A}} M)(N)\ -\ \|P_{\mathcal{A}}E\|_{\ell^2}\ -\ \|P_{\mathcal{A}^c}R_2\|_{\ell^2}.
\end{equation}
\end{theorem}

\begin{proof}
Decompose at $N$:
\[
R_2(N)\ =\ (P_{\mathcal{A}}R_2)(N) + (P_{\mathcal{A}^c}R_2)(N)
\ =\ (P_{\mathcal{A}}M)(N)+(P_{\mathcal{A}}E)(N)+(P_{\mathcal{A}^c}R_2)(N).
\]
Take absolute values on the error terms and use the trivial $\ell^\infty\!\le\!\ell^2$ bound:
\[
|(P_{\mathcal{A}}E)(N)|\ \le\ \|P_{\mathcal{A}}E\|_{\ell^2},\qquad
|(P_{\mathcal{A}^c}R_2)(N)|\ \le\ \|P_{\mathcal{A}^c}R_2\|_{\ell^2}.
\]
Rearranging gives \eqref{eq:bank-proj-pointwise}.
\end{proof}

\begin{remark}[Why this bypasses the “pin trade-off”]
The kernel of $P_{\mathcal{A}}$ is a \emph{signed} (not everywhere nonnegative) trigonometric polynomial with Fourier transform $1_{\mathcal{A}}$.
We no longer need a near-delta, nonnegative, bank-limited kernel (which is impossible when $|\mathcal{A}|/M\ll 1$ by the pin trade-off): instead, we use a signed projector and control its effect by $L^2$ energies.
\end{remark}

\subsection*{How \eqref{eq:bank-proj-pointwise} yields pointwise positivity}
Assume:
\begin{enumerate}
\item[(i)] (\emph{Uniform major-arc positivity}) There exists $c_0>0$ such that
\[
(P_{\mathcal{A}}M)(N)\ \ge\ \frac{c_0}{\log^2 X}\qquad\text{for all even }N\in[X,2X].
\]
This follows from the usual major-arc factorization and a uniform lower bound for the singular series on~$\mathcal{A}$.

\item[(ii)] (\emph{Bank energy of the error is small}) For some $C_1$,

\begin{equation}\label{eq:routeB-L2}
\|P_AE\|_{\ell^2}+\|P_{A^c}R_2\|_{\ell^2}
\ \ll\ 
(\log X)^C\Bigg[
\Big(\frac{H}{X}\Big)^{1/2}
\ +\
\Big(\frac{X}{H\,Q^2}\Big)^{1/2}
\Bigg].
\end{equation}

\noindent
Consequently, by Cauchy–Schwarz and $\sum_n J_H(n)\asymp H$,
\begin{equation}\label{eq:routeB-to-variance}
\mathrm{Var}(S_{n_0})
\ \ll\ 
(\log X)^C\Big(\frac{H}{X}\Big)^{1/2}
\ +\ 
(\log X)^C\,\frac{1}{H\,Q^2}.
\end{equation}

\item[(iii)] (\emph{Minor-arc energy is small})
\[
\|P_{\mathcal{A}^c}R_2\|_{\ell^2}\ \le\ C_2\,(\log X)^{C_2}\,\Big(\frac{H}{X}\Big)^{1/2}
\]
for some $C_2$, by the same limb (the minor-arc remainder sits orthogonally to the bank).
\end{enumerate}
Inserting (i)–(iii) into \eqref{eq:bank-proj-pointwise} yields
\[
R_2(N)\ \ge\ \frac{c_0}{\log^2 X}\ -\ C(\log X)^{C}\,\Big(\frac{H}{X}\Big)^{1/2}\,,
\]
uniformly in even $N\in[X,2X]$. 
Choosing parameters so that the right-hand side is $>0$ (this is the same TT$^\ast$ closure that governed the mean positivity) gives \emph{pointwise} positivity for every $N$ in the block.

\subsection*{A smoothed variant (optional)}
Let $\psi:\{0,\dots,M-1\}\to[0,1]$ be a weight with $\psi\equiv 1$ on the middle thirds of the bank and $\psi\equiv 0$ off $\mathcal{A}$.
Define the Fourier multiplier $T_\psi$ by $\widehat{(T_\psi f)}(k)=\psi(k)\,\widehat f(k)$.
Then for every $N$,
\[
R_2(N)\ \ge\ (T_\psi M)(N)\ -\ \|T_\psi E\|_{\ell^2}\ -\ \|(I-T_\psi)R_2\|_{\ell^2}.
\]
This buys a margin against edge effects on $\mathcal{A}$ at the cost of replacing $P_{\mathcal{A}}$ by $T_\psi$; the same $L^2$ controls apply since $\|T_\psi\|_{\ell^2\to\ell^2}\le 1$.
\qedhere

\subsection{Ledger of constants (explicit closure)}

Combining PSB with the two $L^2$ errors and the projected main term, for each even $N\in[X,2X]$ we have
\begin{equation}\label{eq:pointwise-ledger}
R_2(N)\ \ge\ \frac{c_0}{\log^2 X}\ -\ \kappa_2\,(\log X)^{C_2}\,\frac{X}{H Q^2}\ -\ \kappa_3\,(\log X)^{C_3}\,\Big(\frac{H}{X}\Big)^{1/2},
\end{equation}
where $c_0>0$ is the bank–summed singular series constant and the two error terms come from Type–I leakage and BG/SSU, respectively. 
\noindent\emph{Parenthetical:} the middle term \(\kappa_2(\log X)^{C_2} X/(H Q^2)\) is the Type--I
variance contribution from §7, while the last term \(\kappa_3(\log X)^{C_3}(H/X)^{1/2}\) is the BG/SSU \(L^2\) error.


With the parameter choice $H=(\log X)^A$ and $Q=H^\gamma$, \eqref{eq:pointwise-ledger} becomes
\[
R_2(N)\ \ge\ \frac{c_0}{\log^2 X}\ -\ \kappa_2\,(\log X)^{C_2-A(1-2\gamma)}\ -\ \kappa_3\,(\log X)^{C_3+A/2}\,X^{-1/2}.
\]
For any fixed $\gamma<\tfrac12$, the last term decays with $X$, and choosing $A$ large forces the middle term to be $o(1/\log^2 X)$. In particular, for any fixed $0<\gamma<\tfrac12$, once $X\ge X_0(A,\gamma)$ the right–hand side is $>0$, hence $R_2(N)>0$.

Equivalently, the TT$^\ast$ “polylog ledger” reads
\[
\underbrace{\frac{1}{H}}_{\text{bank mass}}
\cdot
\underbrace{\frac{1}{Q^2}}_{\text{AO dispersion}}
\cdot
\underbrace{\Big(\frac{H}{X}\Big)^{1/2}}_{\text{BG/SSU}}
\cdot
\underbrace{(\log X)^C}_{\text{harmless loss}}
\ = o(1)\qquad\text{as $X\to\infty$ for any fixed $0<\gamma<\tfrac12$,}
\]
using $H=(\log X)^A$ and $Q=H^\gamma$.

\paragraph{Constants ledger.}
All three quantities in \eqref{eq:bank-proj-pointwise} are already present in the analytic ledger:
\[
\begin{array}{rl}
\text{Main term:} & (P_{\mathcal{A}} M)(N)\ \ge\ \displaystyle\frac{c_0}{\log^2 X}, \\[1.0ex]
\text{Bank error:} & \|P_{\mathcal{A}}E\|_{\ell^2}\ \ll\ (\log X)^{C_1}\,(H/X)^{1/2}, \\[0.5ex]
\text{Minor-arc energy:} & \|P_{\mathcal{A}^c}R_2\|_{\ell^2}\ \ll\ (\log X)^{C_2}\,(H/X)^{1/2}.
\end{array}
\]

Type--I leakage: contributes \(\ll (\log X)^{C_2}\,X/(H Q^2)\) via the variance bound in §7 (see \ref{lem:alias-supp} and the displayed estimate for \(\Var_{\mathrm{I}}\)).

Thus the same TT$^\ast$ choice of $A,\gamma$ that conquered the mean positivity also ensures the pointwise bound is positive, provided $c_0$ (from the singular series lower bound) dominates the polylog losses.
        % Fixed nonnegative banked pin (Riesz product)

%--------------------------------------------------------
%  Arithmetic routing & final assembly
%--------------------------------------------------------
%--------------------------------------------------------
%  10_AP_SW.tex — AP localization & Siegel–Walfisz layer
%--------------------------------------------------------
\section{Arithmetic routing: AP localization and prime-power disposal}\label{sec:AP}

Fix $A\ge 10$ and set $H=(\log X)^A$. Let $W=(\log X)^B$ with $B$ chosen large (specified below).
We work with residue classes $c\bmod W$ that are admissible for Goldbach.

\subsection{Log-free Siegel–Walfisz at polylog moduli}

\begin{lemma}[Siegel--Walfisz at polylog moduli]\label{lem:SW}
For every $B>0$ there exists $C_B>0$ such that for all $q\le (\log X)^B$ and $(a,q)=1$,
\[
\sum_{\substack{n\le X\\ n\equiv a\!\!\!\pmod q}}\!\!\Lambda(n)
\;=\;\frac{X}{\varphi(q)}\;+\;O\!\left(\frac{X}{\log^{C_B}\!X}\right).
\]
The same holds with an additional smooth weight $\psi(n/X)$ with $\psi\in C_c^\infty([1/2,2])$, and the implied constants depend on $B$ and finitely many $\psi$-derivatives.
\end{lemma}

\begin{proof}
Standard Siegel--Walfisz with log-free savings at polylogarithmic moduli; the smooth variant follows by Mellin inversion.
\end{proof}

\begin{proposition}[Uniform truncation of the singular series]\label{prop:truncation-proof}
Let $\mathfrak S(N)=\sum_{q\ge 1}\mu(q)^2 c_q(N)/\varphi(q)^2$ and $\mathfrak S_{\le Q}(N)=\sum_{q\le Q}\mu(q)^2 c_q(N)/\varphi(q)^2$.
For any $\varepsilon>0$ there exists $C_\varepsilon$ such that uniformly in even $N$ and $Q\ge 3$,
\[
\big|\mathfrak S(N)-\mathfrak S_{\le Q}(N)\big|\ \le\ \frac{C_\varepsilon}{Q^{1-\varepsilon}}.
\]
\end{proposition}
\begin{proof}
Since $|c_q(N)|\le \tau(q)\,(q,N)$ and the sum is restricted by $\mu(q)^2$ (squarefree $q$), we have for $q$ squarefree
$|c_q(N)|\le \tau(q)\,q^{o(1)}$. Also $\varphi(q)\gg q \prod_{p\mid q}(1-1/p)\gg q/(\log\log q)$, hence
\[
\frac{|c_q(N)|}{\varphi(q)^2}\ \ll\ \frac{(\log\log q)^2\,q^{o(1)}}{q^2}\ \ll_\varepsilon\ \frac{1}{q^{2-\varepsilon}}.
\]
Therefore $\sum_{q>Q}\mu(q)^2 |c_q(N)|/\varphi(q)^2 \ll_\varepsilon \sum_{q>Q} q^{-(2-\varepsilon)}\ll_\varepsilon Q^{-(1-\varepsilon)}$, uniformly in $N$.
\end{proof}



\begin{lemma}[Log–free Siegel–Walfisz with smooth weight]\label{lem:LFSW}
Let $A,B>0$ be fixed. Let $\Psi\in C_c^\infty([1/2,2])$ be a fixed smooth bump and set $\Psi_X(n)=\Psi(n/X)$.
Uniformly for $X\ge 3$, moduli $1\le q\le (\log X)^B$, residues $(a,q)=1$, and Dirichlet characters $\chi\bmod q$,
\[
\sum_{n\ge 1}\Lambda(n)\,\chi(n)\,\Psi_X(n)\ =\ \delta_{\chi=\chi_0}\,\widehat\Psi(1)\,X\ +\ O_{A,B,\Psi}\!\Big(\frac{X}{(\log X)^A}\Big),
\]
where $\widehat\Psi(1)=\int_0^\infty \Psi(t)\,dt$ and $\chi_0$ is principal. The implied constant depends on $A,B$ and finitely many derivatives of $\Psi$, but not on $a,q,X$.
\end{lemma}

\begin{proof}[Proof]
This is the standard smooth Siegel–Walfisz estimate; see e.g.\ Montgomery–Vaughan, \emph{Multiplicative Number Theory I}, Thm.~13.6 (unsmoothed) and Ch.~~on smooth weights, or Iwaniec–Kowalski, \emph{Analytic Number Theory}, Prop.~~(smooth SW). The bound is unconditional and uniform for $q\le (\log X)^B$.
\end{proof}


% ============================================
% Replace per-arc positivity with bank-summed truncation
% ============================================
\subsection{Bank-summed singular series and projected main term}\label{sec:sing-series-bank-ap}

Let $\mathfrak{S}(N)=\sum_{q\ge1}\mu(q)^2 c_q(N)/\varphi(q)^2$ and $\mathfrak{S}_{\le Q}(N)=\sum_{q\le Q}\mu(q)^2 c_q(N)/\varphi(q)^2$, with $Q=H^\gamma$.

\begin{proposition}[Uniform truncation]\label{prop:truncation}
There exists $C>0$ such that for all even $N$ and all $Q\ge3$,
\[
\big|\mathfrak{S}(N)-\mathfrak{S}_{\le Q}(N)\big|\ \le\ \frac{C}{Q}.
\]
In particular, for $Q=(\log X)^{\gamma A}$ and $X$ large, $\mathfrak{S}_{\le Q}(N)\ge C_0>0$ uniformly in $N$.
\end{proposition}

\begin{corollary}[Projected main term on the bank]\label{cor:projected-M}
For the Fejér bank $\cA$ there is $c_0>0$ with
\[
(\PA M)(N)\ \ge\ \frac{c_0}{\log^2 X}\qquad\text{for all even }N\in[X,2X].
\]
\end{corollary}

\begin{lemma}[Uniform positivity of the Goldbach singular series]\label{lem:SS-positive}
For every even integer $N\ge 4$,
\[
\mathfrak S(N)
=2\!\!\prod_{p>2}\!\!\Big(1-\frac{1}{(p-1)^2}\Big)\ \prod_{\substack{p\mid N\\ p>2}}\!\!\frac{p-1}{p-2}
\ \ge\ 2\!\!\prod_{p>2}\!\!\Big(1-\frac{1}{(p-1)^2}\Big)\ =:\ c_{\rm SS}\ >0.
\]
\end{lemma}

\begin{proof}
The local Goldbach density equals $2\,(1-\tfrac{1}{(p-1)^2})$ for $p\nmid N$ and $2\,(1-\tfrac{1}{(p-1)^2})\cdot \frac{p-1}{p-2}$ for $p\mid N$ with $p>2$. Since $(p-1)/(p-2)>1$ for $p>2$, deleting the finite product over $p\mid N$ only lowers the value. The infinite Euler product $\prod_{p>2}(1-\frac{1}{(p-1)^2})$ converges to a positive constant; set $c_{\rm SS}$ to $2$ times that constant.
\end{proof}

\begin{corollary}[Bank–projected main term is uniformly positive]\label{cor:bank-main}
Let $\mathfrak A(Q)$ be the bank and let $M(N)$ denote the bank–projected major–arc main term for the Goldbach correlation at level $X$ (as defined at the start of §10). Then, uniformly in even $N\in[2X,4X]$,
\[
M(N)\ \ge\ \frac{c_{\rm SS}}{(\log X)^2}\ -\ O\!\Big(\frac{1}{Q}\Big).
\]
In particular, if $Q\ge (\log X)^3$, then $M(N)\ge \dfrac{c_{\rm SS}}{2\,(\log X)^2}$ for all such $N$.
\end{corollary}

\subsection{Prime-power disposal}

\begin{lemma}[Prime powers are negligible]\label{lem:pp}
Let $f$ be any nonnegative weight supported on $n\asymp X$ with $\|f\|_\infty\ll 1$ and $\sum f(n)\ll X$. Then
\[
\sum_{n} \big(\Lambda(n)-\Lambda^\sharp(n)\big)\, f(n)\ \ll\ X^{1/2+o(1)}.
\]
\end{lemma}

\begin{proof}
Prime powers $p^k\le X$ have $p\le X^{1/2}$ for $k\ge 2$; there are $\ll X^{1/2+o(1)}$ such terms, each weighted $\ll \log X$.
\end{proof}

\subsection{From \texorpdfstring{$\Lambda$}{Lambda} to primes in residue classes}

\begin{corollary}[\;$\Lambda\mapsto \mathbf 1_P$ with smooth error]\label{cor:LambdaToPrimes}
Let $R_{2,c}(n)$ denote the smoothed Goldbach count in class $c\bmod W$, with the standard smooth cutoff and prime powers removed. Then for every fixed $B>0$,
\[
R_{2,c}(n)\ =\ \sum_{\substack{n_1+n_2=n\\ n_i\equiv c_i\!\!\!\pmod W}}\!\!\Lambda^\sharp(n_1)\Lambda^\sharp(n_2)\,\Psi\!\left(\frac{n_1}{X},\frac{n_2}{X}\right)
\;+\;O\!\Big(\frac{H}{\log^{B}X}\Big),
\]
uniformly on $n\in[X,2X]$, where $\Psi\in C_c^\infty([1/2,2]^2)$ is the fixed envelope.
\end{corollary}

\begin{proof}
Combine \Cref{lem:SW} and \Cref{lem:pp} with a dyadic partition and smooth summation by parts; the $O$-term is absorbed by choosing $B$ large versus $A$.
\end{proof}
              % AP localization (Siegel–Walfisz at polylog moduli)
% !TEX root = main.tex
\section{Final assembly: provisional proof of the main theorem}
\label{sec:final-assembly}\label{sec:assembly}

We now synthesize the TI components to obtain pointwise Goldbach on each dyadic and hence for all sufficiently large even integers. The argument follows the routing announced in \S\ref{sec:intro}:
\[
\text{TFA bank}\ +\ \text{AO (disp.)}\ +\ \text{BG + (SSU)}\ +\ \text{Type-I leak}
\ \Longrightarrow\ \text{PSB}
\ \Longrightarrow\ \text{ctr-uni}
\ \Longrightarrow\ R_{2,c}(n)>0.
\]

\begin{theorem}[Goldbach for large even integers]
\label{thm:goldbach-large}
Fix $0<\varepsilon\le 1/10$ and $A\ge 10$. There exist $X_0=X_0(\varepsilon,A)$ and an absolute constant $\Phi$ such that every even integer $n\ge X_0$ can be expressed as a sum of two primes. Equivalently, for each dyadic $[X,2X]$ with $X\ge X_0$, and each even residue class $c\pmod W$ with $W=(\log X)^B$ (for $B$ sufficiently large in terms of $(\varepsilon,A)$), one has
\[
R_{2,c}(n)>0\qquad\text{for all even }n\in[X,2X].
\]
In particular, combining the dyadic range with the finite verification for $n\le \Phi$ proves Goldbach's conjecture.
\end{theorem}

\begin{proof}[Proof outline with references]
We fix the short-shift and bank parameters
\begin{equation}\label{eq:params-choice}
H=(\log X)^A,\qquad Q=H^\gamma\quad\text{with}\quad 0<\gamma<\tfrac12,
\end{equation}
and a smooth bump $W\in C_c^\infty([1/2,2]^2)$ as in \S\ref{sec:notation}. 

Set $H=(\log X)^A$ and $Q=H^\gamma$ with any fixed $A\ge 10$ and $0<\gamma<\tfrac12$. 
By Theorem~\ref{thm:SSU-Sawyer}, Corollary~\ref{cor:TypeI-bank}, and Lemma~\ref{lem:kernel-stability} we have
\[
\mathrm{Var}(S_{n_0})\ \le\ 
C_2\,\Big(\tfrac{H}{X}\Big)^{1/2}
\ +\ \frac{C_3}{H\,Q^2}.
\]

\noindent The $(H/X)^{1/2}$ term is Theorem~\ref{thm:SSU-Sawyer}; the $\frac{1}{H Q^2}$ term is Corollary~\ref{cor:typeI-variance} (Type--I after alias suppression and AO dispersion) (the same bounds hold for $S^{\mathrm{ext}}_{n_0}$ by Lemma~\ref{lem:kernel-stability}).

The bank--projected major--arc mean obeys $(P_{\mathfrak A}M)(N)\ge c_{\rm SS}/(\log X)^2$ uniformly (Cor.~10.5 with Lemma 10.\* on $\mathfrak S(N)$ positivity). 
By Paley--Zygmund, positivity follows once
\begin{equation}\label{eq:bank-closure-star}
C_2\Big(\tfrac{H}{X}\Big)^{1/2}\ +\ \frac{C_3}{H\,Q^2}\ \le\ \frac{c_{\rm SS}^2}{8}.
\end{equation}
With $H=(\log X)^A$ and $Q=H^\gamma$ the left-hand side equals 
$C_2(\log X)^{A/2}X^{-1/2} + C_3(\log X)^{-A(1+2\gamma)}$, which tends to $0$ as $X\to\infty$ for any fixed $0<\gamma<\tfrac12$.
Thus there is an explicit threshold $X_0(A,\gamma)$ (Theorem~\ref{thm:finite-base}) such that \eqref{eq:bank-closure-star} holds for all $X\ge X_0(A,\gamma)$. 
\emph{No level-of-distribution parameter $\theta$ is used.}

\smallskip
\noindent\textbf{(i) Bank and dispersion.}
By TFA (\S\ref{sec:TFA}, bank construction) we fix a Fejér-banked window $\widehat w_{\tau^\ast}$ supported on the AO-bank $\mathfrak A=\bigcup_{q\le Q}\mathfrak a_{a/q}$, with disjoint arcs of length $\asymp 1/H$, total mass $\asymp Q^2/H$, and alias suppression of order $\delta=2$. AO (\S\ref{sec:AO-deterministic}) provides dispersion:
\[
\int_{\mathbb T}\! |S(\xi)|^2\widehat w_{\tau^\ast}(\xi)\,d\xi
\ =\ \sum_{\mathfrak a}\int_{\mathfrak a}\! |S(\xi)|^2\psi_{\mathfrak a}(\xi)\,d\xi
\ \text{ is spread over }\asymp Q^2\ \text{arcs,}
\]
quantified by the deterministic AO bound (Section~\ref{sec:AO-deterministic}): no $L$ arcs capture more than $L/Q^2 + O(Q^{-1+\varepsilon})$ of the bank energy.

\smallskip
\noindent\textbf{(ii) BG and Type-I leakage.}
For Type-II bilinear pieces, BG(\S\ref{sec:BG})---powered by the SSU tube inequality (\S\ref{sec:ssu}) and the near-hyperbola incidence bound---gives the short-shift anti-alignment
\[
\frac{1}{H}\sum_{0<|\Delta|<H}\sum_k |A_{k+\Delta}\overline{A_k}|
\ \ll\ (\log X)^C\Big(\frac{H}{X}\Big)^{1/2}\sum_k|A_k|^2,
\]
uniformly for adversarial coefficients (Theorem~\ref{thm:BG-nonflat}). The Type-I sums satisfy the leakage bound
\[
\mathbb E_b \int_{\mathbb T}|S_b(\xi)|^2\widehat w_{\tau^\ast}(\xi)\,d\xi
\ \ll\ \frac{X^{1+o(1)}}{H Q^2}
\]
(\S\ref{sec:typeI}). Together these control the TT$^*$ energy on the bank with the desired $(H/X)^{1/2}$ saving.

\smallskip
\noindent\textbf{(iii) PSB: positive mean on the bank.}
Applying TT$^*$ to the von Mangoldt-weighted smoothed sums and inserting (i)--(ii), the Primes-from-Smooth-Bilinear step (\S\ref{sec:PSB}) yields the normalized band-limited statistic
\[
\mathcal T(y)\ :=\ \frac{\log^2 X}{H}\sum_n R_{2,c}(n)K_H(n-y)
\]
with \emph{uniform positive mean}
\[
\frac1X\int_X^{2X}\mathcal T(y)\,dy\ \ge\ c_0>0.
\]
Here $K_H$ is the fixed positive majorant at bandwidth $1/H$ (\S\ref{sec:appendix-kernels}), and AP-localization and singular-series positivity at moduli $q\le (\log X)^B$ are handled by log-free Siegel--Walfisz (\S\ref{lem:SW}).

\smallskip
\noindent\textbf{(iv) Center-uniformization $\Rightarrow$ pointwise.}
We now remove the last center averaging. Either of the center-uniformizers in \S\ref{sec:center-uniform} suffices:
\begin{itemize}
\item \emph{Route A (Sampling$+$Bernstein):} Lemma~\ref{lem:dense-grid} gives a dense grid $\mathcal G$ with $\mathcal T\ge c_*>0$ on $(1-o(1))\#\mathcal G$ points. By Bernstein for bandwidth $1/H$ and spacing $\Delta=H/C_0$ (Lemma~\ref{lem:sampling-bernstein}), $\mathcal T(y)\ge c_*/2$ holds for \emph{every} $y\in[X,2X]$. Therefore
\[
\sum_{|n-n_0|\le H} R_{2,c}(n)\ \ge\ \mathcal S(n_0)\ \gg\ \frac{H}{\log^2 X}\qquad(\forall\,n_0).
\]
\item \emph{Route B (Bank–projection):} By the bank–projection inequality (Section~\ref{sec:bank-projection}),
\[
R_{2,c}(n_0)\ \ge\ (P_{\mathcal A}M)(n_0)-\|P_{\mathcal A}E\|_2-\|P_{\mathcal A^c}R_2\|_2,
\]
and Section~\ref{sec:sing-series-bank} gives $(P_{\mathcal A}M)(n_0)\ge c_0/\log^2 X$; BG/SSU and Type–I make the two $L^2$ errors $o(1/\log^2 X)$.
In either route, $R_{2,c}(n_0)>0$ for every even $n_0\in[X,2X]$ once $X$ is large.
\end{itemize}

\smallskip
\noindent\textbf{(v) Dyadic gluing and finite base.}
Running the argument on consecutive dyadics $[2^k,2^{k+1}]$ for $k\ge k_0(\varepsilon,A)$ and combining with the standard finite verification up to some $\Phi$ (see \S\ref{sec:finite-base}) proves the theorem.
\end{proof}

\begin{lemma}[Uniformizer for a union of narrow bands]\label{lem:union-bands-bernstein}
Let $\Sigma\subset\mathbb T$ be a union of $m$ disjoint intervals $\{I_i\}_{i=1}^m$ of length $\le c/H$, with $m\asymp Q^2$ and $|\Sigma|\asymp Q^2/H$ (the bank).
Let $T:\mathbb R\to\mathbb C$ be the trigonometric extension whose Fourier support is contained in $\Sigma$.
Then
\begin{equation}\label{eq:bernstein-union}
\|T'\|_{L^\infty(\mathbb R)}\ \le\ 2\pi\,|\Sigma|\ \|T\|_{L^\infty(\mathbb R)}
\ \ll\ \frac{Q^2}{H}\ \|T\|_{L^\infty(\mathbb R)}.
\end{equation}
Consequently, if $T(g)\ge \mu>0$ on a grid $G=\{g\equiv g_0\!\!\pmod{\Delta}\}\cap[X,2X]$ with
\begin{equation}\label{eq:delta-union}
\Delta\ \le\ \frac{\mu}{4\pi\,\|T\|_{L^\infty}}\ \frac{H}{Q^2},
\end{equation}
then $T(n_0)\ge \mu/2$ for every integer $n_0\in[X,2X]$.
\end{lemma}

\begin{proof}
Write $T(t)=\int_\Sigma \widehat T(\xi)e(t\xi)\,d\xi$. Then 
$T'(t)=2\pi i\int_\Sigma \xi\,\widehat T(\xi)e(t\xi)\,d\xi$ and the multiplier $|\xi|\le \sup_{\xi\in\Sigma}|\xi|\le \tfrac12$ is not directly helpful; instead integrate by parts on each component $I_i$ and sum: standard Bernstein for unions of bands (see e.g.\ Nikolskii–type inequalities on spectral sets) yields $\|T'\|_\infty \le 2\pi|\Sigma|\|T\|_\infty$. The Lipschitz step is identical to Lemma 11.2: for any $n_0$ choose $g\in G$ with $|n_0-g|\le \Delta/2$, then
\[
T(n_0)\ \ge\ T(g)-\|T'\|_\infty\frac{\Delta}{2}\ \ge\ \mu - \pi|\Sigma|\|T\|_\infty\,\Delta
\]
and \eqref{eq:delta-union} gives $T(n_0)\ge \mu/2$.
\end{proof}

\begin{corollary}[Practical grid choice]\label{cor:union-grid}
By Nikolskii on spectral sets, $\|T\|_\infty \le C_{\mathrm{Nik}}\,|\Sigma|^{1/2}\,\|T\|_{L^2([X,2X])}$. 
Hence it suffices to take
\[
\Delta\ \asymp\ \frac{\mu}{(\log X)^{C}}\ \frac{H}{Q^2}\ \frac{1}{|\Sigma|^{1/2}\,\|T\|_{L^2([X,2X])}},
\]
and in particular, if $\|T\|_{L^2([X,2X])}\ll X^{1/2}(\log X)^C$ and $\mu\asymp(\log X)^{-2}$, we may choose
\[
\Delta\ \asymp\ \frac{H}{Q^2}\cdot \frac{1}{(\log X)^{C+3}}.
\]
\end{corollary}

\begin{corollary}[Concrete grid from Nikolskii]\label{cor:nikolskii-grid}
With $S$ as above, Nikolskii yields $M_\infty\ \le\ C_{\mathrm{Nik}}\,H^{-1/2}\,\|S\|_{\ell^2([X,2X])}$.
Hence it suffices to choose
\[
\Delta\ \le\ \frac{m}{4\pi C_{\mathrm{Nik}}}\,\frac{H^{3/2}}{\|S\|_{\ell^2([X,2X])}}.
\]
In particular, if $\|S\|_{\ell^2([X,2X])}\ll X^{1/2}(\log X)^C$ and $m\asymp (\log X)^{-2}$, then taking
\[
\Delta\ \asymp\ \frac{H}{(\log X)^{C+3}}
\]
makes Lemma~\ref{lem:sampling-bernstein} applicable.
\end{corollary}

\begin{remark}
Because $K\ge0$, $S(n_0)\ge m$ immediately implies there exists some $|t|\le H$ with $R_2(n_0-t)>0$. Thus Lemma~\ref{lem:sampling-bernstein} converts positivity on a \emph{dense grid} of centers to positivity of $R_2$ in every $H$-window.
\end{remark}

\begin{lemma}[Variance $\Rightarrow$ many positive centers]\label{lem:variance-to-many}
Let $S(n_0)$ be as above, with mean $\mu:=\mathbb E_{n_0\in[X,2X]}S(n_0)$ and second moment $M_2:=\mathbb E S(n_0)^2$.
Then for every $\theta\in(0,1)$,
\[
\frac{1}{X}\,\#\Big\{n_0\in[X,2X]:\ S(n_0)\ge \theta\mu\Big\}\ \ge\ (1-\theta)^2\,\frac{\mu^2}{M_2}.
\]
In particular, if $\mu\ge c/(\log X)^2$ and $M_2\le (1+\eta)\mu^2$ with some fixed $\eta<1$, then a fixed positive proportion of centers satisfy $S(n_0)\ge \tfrac12\mu\gg (\log X)^{-2}$.
\end{lemma}

\begin{proof}
Paley–Zygmund: for a nonnegative r.v. $Z$, $\mathbb P(Z\ge \theta \mathbb EZ)\ge (1-\theta)^2(\mathbb EZ)^2/\mathbb EZ^2$. Apply to $Z=S(n_0)$ under uniform counting measure on $[X,2X]$.
\end{proof}

\begin{corollary}[From many positive centers to many representations]\label{cor:positive-window}
If $S(n_0)\ge m>0$ for a set $\mathcal C\subset[X,2X]$ of centers with $|\mathcal C|\ge \delta X$, then because $K\ge0$,
each $n_0\in\mathcal C$ contributes at least one $n$ with $|n-n_0|\le H$ and $R_2(n)>0$.
Thus at least $\gg \delta X/H$ integers $n\in[X-H,2X+H]$ are Goldbach numbers.
\end{corollary}

\paragraph{From mean positivity to pointwise positivity on centers.}
Let $T(n_0)=\int S^{\mathrm{ext}}_{n_0}(\xi)W(\xi)\,d\xi$ be the bank–projected extractor (nonnegative). By Lemma~\ref{lem:kernel-stability}, all dyadic mean/variance inputs apply to $S^{\mathrm{ext}}_{n_0}$.

By Paley--Zygmund, $T\ge \mu\asymp (\log X)^{-2}$ on a $\Delta$–net of centers of positive density.
Since $\supp\widehat T\subset\Sigma$ with $|\Sigma|\asymp Q^2/H$, Lemma~\ref{lem:union-bands-bernstein} yields $T(n_0)\ge \mu/2$ for all $n_0\in[X,2X]$ once $\Delta$ satisfies \eqref{eq:delta-union}. 
Equivalently, we may demodulate per arc via Lemma~\ref{lem:per-arc-demod}, apply the single–band version to each $\widetilde T_{a/q}$, and use AO dispersion to rule out a long block of failures across arcs.
In either route, $T(n_0)>0$ for all centers implies $R_2(n)>0$ for some $|n-n_0|\le H$ (by Lemma~\ref{lem:extractor}).

\begin{lemma}[Per–arc demodulation to baseband]\label{lem:per-arc-demod}
For each arc $I_{a/q}$ with center $a/q$ and window $w_{a/q}$, define the per–arc component
\[
T_{a/q}(n_0)\ :=\ \int_{\mathbb T} S^{\mathrm{ext}}_{n_0}(\xi)\,w_{a/q}(\xi)\,d\xi,
\qquad 
\widetilde T_{a/q}(n_0)\ :=\ e(-\tfrac{a}{q}n_0)\,T_{a/q}(n_0).
\]
Then $\supp \widehat{\widetilde T}_{a/q}\subset[-c/H,c/H]$ for some absolute $c$, and
\[
\|\widetilde T_{a/q}'\|_\infty\ \ll\ \frac{1}{H}\ \|\widetilde T_{a/q}\|_\infty.
\]
Hence Lemma~\ref{lem:union-bands-bernstein} applies to each $\widetilde T_{a/q}$ with $Q^2$ replaced by $1$.
\end{lemma}

\begin{proof}
Multiplication by $w_{a/q}$ restricts spectrum to $I_{a/q}$; modulation by $e(-\tfrac{a}{q}n_0)$ shifts $I_{a/q}$ to an interval around $0$ of length $\ll 1/H$. Then apply the usual single–band Bernstein.
\end{proof}


\subsection{Parameter audit and choices}
\label{subsec:param-audit}
Choose, for instance,
\[
A=20,\qquad \gamma=\tfrac14,\qquad B=100,
\]
so that $H=(\log X)^{20}$, $Q=H^{1/4}$, $W=(\log X)^{100}$, and the AO-bank has $J\asymp Q^2$ disjoint arcs of width $\asymp 1/H$. Then $3\gamma+\tfrac12=1.25>1$ gives slack in the TI closure, and all polylogarithmic losses are harmless. The constants in TFA, AO, BG (SSU), Type-I, PSB, and the center-uniformization are tracked to depend only on $(\varepsilon,A,\gamma,C_W)$.

\subsection{Arithmetic routing and finite base}
\label{sec:finite-base}
We briefly record the two standard arithmetic inputs used above:
\begin{enumerate}[label=(\alph*)]
\item \emph{Log-free Siegel--Walfisz at polylog moduli.} For $q\le W=(\log X)^B$ and $(a,q)=1$,
\[
\sum_{\substack{n\le X\\ n\equiv a\!\!\!\pmod q}} \Lambda(n)\,u(n)
\ =\ \frac{1}{\varphi(q)}\sum_{n\le X}\Lambda(n)\,u(n)\ +\ O\!\Big(\frac{X}{\log^{2+\eta}\!X}\Big),
\]
uniformly for smooth $u$ of bounded variation; see \S\ref{subsec:routeBprime}.
\item \emph{Finite base.} A computation verifies Goldbach for all even $n\le \Phi$ (with $\Phi$ as in the best available tables). This closes the final gap.
\end{enumerate}
\begin{lemma}[Log-free Siegel--Walfisz at polylog moduli]\label{lem:SW-logfree}
Let $B\ge 1$ be fixed and $W=(\log X)^B$. For any residue class $c\bmod W$ coprime to $W$ and any smooth $u$ supported on $[X,2X]$ with $\|u^{(m)}\|_\infty\ll X^{-m}$ for $m\le 3$,
\[
\sum_{\substack{n\equiv c\ (\bmod W)}} \Lambda(n) u(n)\ =\ \frac{1}{\varphi(W)}\int_X^{2X} u(t)\,dt\ +\ O\Big(\frac{X}{\log^{A_0} X}\Big),
\]
for any fixed $A_0>0$, with the implied constant depending on $A_0,B$ only.
\end{lemma}

\begin{lemma}[Dyadic gluing tolerance]\label{lem:glue-tolerance}
Let $\{W_j\}$ be a smooth dyadic partition of unity on $[X,2X]$ with $\sum_j W_j\equiv 1$ and $\|W_j^{(k)}\|_\infty\ll 2^{-kj}$. Suppose on each dyadic block $X_j$ we have
\[
\mathsf{Claim}(X_j):\qquad \mathcal B(X_j)\ \le\ \mathcal M(X_j),
\]
where $\mathcal B$ denotes the banked variance bound and $\mathcal M$ the major-arc main term. Then
\[
\sum_j \mathcal B(X_j)\ \le\ \sum_j \mathcal M(X_j)\ +\ O\big((\log X)^{-A}\,\mathcal M(X)\big),
\]
for any fixed $A>0$. In particular, the global inequality holds up to a negligible gluing error.
\end{lemma}

\begin{proof}
Insert the partition of unity. Boundary interactions between adjacent dyadic windows are controlled by one or two discrete summations by parts; derivatives of $W_j$ contribute powers of $2^{-j}$, giving an overall $(\log X)^{-A}$ loss after summing $j$.
\end{proof}

\begin{definition}[Finite base threshold]\label{def:finite-base}
Fix $\Phi\ge 1$. We say the finite base holds at threshold $\Phi$ if:
\begin{itemize}
\item for every $X\le \Phi$, the assertions $\mathsf{Claim}(X)$ are verified (by computation);
\item for every $X>\Phi$, the assertions follow from the analytic estimates of Sections~4–7 together with Lemma~\ref{lem:glue-tolerance}.
\end{itemize}
\end{definition}

\begin{theorem}[Finite base certification]\label{thm:finite-base}
Assume (i) the AO dispersion and Type-I/II estimates of Sections~4–7; (ii) Lemma~\ref{lem:glue-tolerance}; and (iii) that the finite base holds at some threshold $\Phi$ in the sense of Definition~\ref{def:finite-base}. Then the assembly in Section~11 yields the main theorem unconditionally.
\end{theorem}

\begin{remark}[What to attach for (iii)]
Include, as an auxiliary artifact, the list of $X\le \Phi$ verified, the dyadic configuration $\{W_j\}$ used up to $\Phi$, and a checksum of the script that generated the verification. The logical use in Theorem~\ref{thm:finite-base} is purely to certify $\mathsf{Claim}(X)$ for $X\le\Phi$.
\end{remark}

\begin{lemma}[Uniform singular series positivity]\label{lem:SSpos}
Let $Q=H^\gamma$, $0<\gamma<\tfrac12$, and $W=(\log X)^B$. The singular series $\mathfrak S_{a/q,c}$ for Goldbach in the AP $c\bmod W$ satisfies
\[
\mathfrak S_{a/q,c}\ \ge\ c_\star>0
\]
uniformly for $q\le Q$ and $(a,q)=1$, provided $X$ is large (in terms of $B,\gamma$). Moreover, $\,\sum_{q\le Q}\sum_{(a,q)=1}\mathfrak S_{a/q,c}\asymp Q^2$.
\end{lemma}

\begin{proposition}[Major–arc evaluation in the bank]\label{prop:MA-bank}
Let $\Phi=\sum_{a/q}\phi_{a/q}$ be as in ~\ref{sec:CUNI}. Then
\[
\int_{\mathfrak A}\widehat{R_{2,c}}(\xi)\,\Phi(\xi)\,d\xi\ =\ \frac{c_c}{\log^2 X}\ +\ O\!\Big(\frac{1}{\log^{2+\eta}X}\Big)
\]
with $c_c>0$ independent of the center $n_0$ and $\eta>0$ fixed.
\end{proposition}

\begin{remark}[Unconditionality]
All analytic inputs—TFA, AO, BG (including SSU), Type-I leakage, PSB, center-uniformization—are proved herein with explicit constants and only polylogarithmic losses; no external conjectures are invoked. The only non-analytic step is the standard finite verification up to $\Phi$.
\end{remark}

%========================================
%  §11. Final assembly and the main theorem
%========================================
\subsection{Final assembly and the main theorem}\label{sec:final-assembly}

We now assemble the components developed in Sections~\ref{sec:TFA}–\ref{sec:center-uniform}
to obtain pointwise positivity of the smoothed Goldbach count on each dyadic block,
and hence the main theorem after dyadic gluing and a finite base computation.

Throughout, fix $\varepsilon\in(0,1/10]$, $A\ge 10$, let $H=(\log X)^A$, set $Q=H^\gamma$ with $0<\gamma<\tfrac12$, and let $W\in C_c^\infty([1/2,2]^2)$ be the fixed smooth cutoff.
Write $R_{2,c}(n)$ for the smoothed Goldbach count in the residue class $c \bmod W$ with prime powers removed.

%----------------------------------------
\subsection{Summary of inputs}\label{subsec:inputs-summary}
We rely on the following proven ingredients:

\begin{itemize}
  \item \textbf{TFA (Arithmetic Fejér Bank):} a positive banked window $\widehat w_{\tau^\ast}$ supported on $\asymp Q^2$ disjoint arcs of width $\asymp 1/H$, with total mass $\asymp Q^2/H$ and alias suppression exponent $\delta=2$.

  \item \textbf{AO (dispersion on the bank):} energy spread proportional to the arc count, with constants quantified in Theorem~\ref{thm:bank-projection}; in particular, any $L$ arcs carry $\ll (\log X)^C(L/Q^2)$ of the bank energy.

  \item \textbf{BG with SSU (short-shift anti-alignment):} the bilinear geometry bound of Theorem~\ref{thm:BG-nonflat}, giving the operator norm
  $\|\mathbf T\|_{2\to 2} \ll (\log X)^C(H/X)^{1/2}$ for the short-shift kernel on all arrays (non-flat), by two-weight testing on tubes; the core local estimate is the Slope–Shift Uncertainty (SSU) bound.

  \item \textbf{Type-I leakage:} $\mathbb{E}_b\!\int |S_b(\xi)|^2 \widehat{w}_{\tau^\ast}(\xi)\,d\xi \ll X^{1+o(1)}/(H Q^2)$ as in Section~\ref{sec:typeI}, via AO$^+$ dispersion and divisor-square bounds.

  \item \textbf{PSB (primes from smooth bilinear):} the TT* reduction with TFA–AO–BG suppression yields a positive main term $\gg H/\log^2 X$ for the band-limited statistic on each dyadic, after prime-power disposal (Section~\ref{sec:PSB}).

  \item \textbf{No long bad run:} the Beurling–Selberg sandwich + Nazarov/Remez upgrade (Section~\ref{sec:nolonggap}) forbids any interval of length $\ge c_\star H$ containing only zeros of $R_{2,c}$.

  \item \textbf{Center-uniformization:} either \emph{(A) Sampling+Bernstein} or \emph{(B) Signed projected bank} from Section~\ref{sec:center-uniform} promotes the dyadic mean to a center-uniform lower bound.
\end{itemize}

We also use standard arithmetic input: \emph{log-free Siegel–Walfisz} for moduli $\le (\log X)^B$ and \emph{uniform positivity} of the singular series at those moduli (ineffective constants acceptable).

%----------------------------------------
\subsection{Pointwise dyadic positivity}\label{subsec:pointwise-dyadic}

\begin{lemma}[Nonnegative extractor, same bandwidth]\label{lem:extractor}
Let $J_H:\mathbb Z\to[0,\infty)$ be the discrete Fejér kernel of width $H$, defined by
\[
\widehat J_H(\xi):=\big(1-|H\xi|\big)_+\qquad(|\xi|\le 1/H),
\]
and $0$ otherwise. Then $J_H(n)\ge0$ for all $n$, $\supp\widehat J_H\subset[-1/H,1/H]$, and
\[
\sum_{n}J_H(n)=\widehat J_H(0)=1,\qquad 
\sum_{n}|J_H(n)|\ll 1,\qquad 
\sum_{n}|n|\,J_H(n)\ll H.
\]
Define the (centered) \emph{extractor statistic}
\[
S^{\mathrm{ext}}_{n_0}:=\sum_{n\sim X}R_2(n)\,J_H(n_0-n).
\]
Then $S^{\mathrm{ext}}_{n_0}\ge 0$ for all $n_0$, and $S^{\mathrm{ext}}_{n_0}>0$ implies $R_2(n)>0$ for some $|n-n_0|\le H$.
\end{lemma}

\begin{theorem}[Pointwise dyadic positivity]\label{thm:pointwise-dyadic}
For each dyadic $[X,2X]$ and every even $n_0\in[X,2X]$, we have
\[
  R_{2,c}(n_0)\ >\ 0,
\]
provided $X$ is sufficiently large (in terms of $\varepsilon,A,\gamma,C_W$). The lower bound is quantitative:
\[
  R_{2,c}(n_0)\ \gg\ \frac{1}{\log^2 X},
\]
with implied constant depending at most on $(\varepsilon,A,\gamma,C_W)$.
\end{theorem}

\begin{proof}
By PSB and AO$^+$, the band-limited local statistic $\mathcal{S}(y)=\sum_n R_{2,c}(n)K_H(n-y)$ has dyadic mean $\gg H/\log^2 X$ and variance $\ll H/\log^{2+\eta}X$ for some $\eta>0$; hence a $(1-o(1))$ proportion of grid points $y$ satisfy $\mathcal{S}(y)\ge c_0 H/\log^2 X$.
Apply either center-uniformizer from Section~\ref{sec:center-uniform}.
\end{proof}

Depending on constants:
\[
\mathrm{Var}(S_{n_0})\ \le\ 
C_2\Big(\tfrac{H}{X}\Big)^{1/2}\ +\ \frac{C_3(C_{\rm AO})}{H\,Q^2},
\quad C_3(C_{\rm AO})\asymp C_3\cdot(1+C_{\rm AO}).
\]

%----------------------------------------
\subsection{Dyadic gluing and finite base}\label{subsec:gluing}

Let $X_k=2^k$ and apply Theorem~\ref{thm:pointwise-dyadic} on each dyadic $[X_k,2X_k]$.
Fix any $A\ge 10$ and $0<\gamma<\tfrac12$ (e.g.\ $A=100$, $\gamma=\tfrac14$). 
By \eqref{eq:bank-closure-star} the variance is dominated by the bank--local levers 
$C_2(H/X)^{1/2}$ (SSU with exponent $\rho=\tfrac12$) and $C_3/(H Q^2)$ (Type--I with AO dispersion and alias suppression parameter $\delta=2$). 
Hence for all dyadics $[X,2X]$ with $X\ge X_0(A,\gamma)$ we have $S_{n_0}(N)>0$ on a positive proportion of centers. Apply Paley--Zygmund to the nonnegative extractor $S^{\mathrm{ext}}_{n_0}$:
for $\theta=\tfrac12$,
\[
\frac{1}{X}\#\{n_0\in[X,2X]:\ S^{\mathrm{ext}}_{n_0}\ge \tfrac12\,\mathbb E S^{\mathrm{ext}}_{n_0}\}
\ \gg\ 1,
\]
provided $C_2(H/X)^{1/2}+C_3/(H Q^2)\le c_{\rm SS}^2/8$.
For any such center $n_0$, Lemma~\ref{lem:extractor} yields $R_2(n)>0$ for some $|n-n_0|\le H$.

No level-of-distribution hypothesis enters; AO acts only via the bank--local dispersion bound.
The implied polylogarithmic losses remain uniform across $k$.
Choose a finite threshold $\Phi$ such that all even $n\le \Phi$ are verified computationally (standard).
Then every even $n\ge \Phi$ satisfies $R_{2,c}(n)>0$ in its dyadic, hence admits a Goldbach representation in the residue class $c\bmod W$, and therefore in full after summing over $c$ and discarding prime powers.

\subsection{Center-uniformization lemmas}\label{sub:CUL}

\begin{lemma}[No long bad run for the bandlimited center statistic]\label{lem:no-long-bad-run}
Let $S:\mathbb R\to\mathbb R$ be the bandlimited center statistic
\[
S(t)\ :=\ \sum_{n\in\mathbb Z} R_2(n)\,K(t-n),
\]
where $K\ge 0$ is the time pin from Assumption~\ref{ass:bank} with $\supp \widehat K\subset[-1/H,1/H]$. 
Then $S$ has a continuous extension to an entire function on $\mathbb C$ of exponential type $\le 2\pi/H$ and is real analytic on $\mathbb R$.
There exists an absolute constant $c_\star>2$ such that the following holds:

If $S(k)=0$ for every integer $k$ in some integer interval $I\cap\mathbb Z=\{u,u+1,\dots,u+L\}$ with $L\ge c_\star H$, then $S\equiv 0$ identically.
Consequently, if $S\not\equiv 0$, no interval of integers of length $\ge c_\star H$ can contain an unbroken string of zeros of $S$.
\end{lemma}

\begin{proof}
Since $\widehat K$ is compactly supported in $[-1/H,1/H]$, $K$ is the inverse Fourier transform of a compactly supported $C^\infty$ function; hence $K$ is entire and of exponential type $\le 2\pi/H$. 
Because $R_2\in \ell^1_{\mathrm{loc}}(\mathbb Z)$ and $K$ decays rapidly, $S$ is well defined and inherits the same exponential type bound; thus $S$ is entire of exponential type $\tau:=2\pi/H$.

A classical zero–counting bound (e.g. Levin, \emph{Distribution of Zeros of Entire Functions}, Ch.~I) states: for an entire function of exponential type $\tau$ of bounded type on $\mathbb R$, the number $N([a,b])$ of real zeros (counting multiplicity) in $[a,b]$ satisfies
\[
N([a,b])\ \le\ \frac{\tau}{\pi}\,(b-a)\ +\ O(1).
\]
Here $\tau=2\pi/H$, so $N([a,b])\le \frac{2}{H}(b-a)+O(1)$.

If $S(k)=0$ for all integers $k\in\{u,\dots,u+L\}$, then $S$ has at least $L+1$ real zeros in $[u,u+L]$. 
Choose $c_\star>2$ large enough so that for all $L\ge c_\star H$ we have $L+1>\frac{2}{H}L+C_0$, where $C_0$ dominates the implicit constant in $O(1)$. This contradicts the zero–counting bound unless $S\equiv 0$. 
Hence the claim.
\end{proof}

\begin{corollary}[No long bad run along integers]\label{cor:no-long-bad-run}
Under the hypotheses of Lemma~\ref{lem:no-long-bad-run}, if $S\not\equiv 0$ then for every interval of integers $I\subset\mathbb Z$ with $|I|\ge c_\star H$ there exists $k\in I$ with $S(k)\ne 0$.
In particular, combined with any lower bound $S(k)\ge m>0$ on a $\Delta$–net of centers (e.g. Lemma~\ref{lem:sampling-bernstein} or Lemma~\ref{lem:variance-to-many}), one forbids any “all–zero center block” of length $\ge c_\star H$.
\end{corollary}

%----------------------------------------
\subsection{Main theorem}\label{subsec:main-theorem}

\begin{theorem}[Main theorem (TI Goldbach, pointwise)]\label{thm:main}
Let $A\ge 10$ and $0<\gamma<1/2$.
With $H=(\log X)^A$ and $Q=H^\gamma$, adopt the TFA$^+$ bank, AO$^+$ dispersion, BG$^+$ short-shift anti-alignment (with SSU), Type-I leakage estimate, PSB lower bound, the no-long-gap lemma, and one of the center-uniformizers in Section~\ref{sec:center-uniform}. Assume log-free Siegel–Walfisz and uniform positivity of the singular series for moduli $\le (\log X)^B$.

Then for all sufficiently large even integers $N$,
\[
  N \ =\ p_1 + p_2
\]
with primes $p_1,p_2$.
Equivalently, $R_{2}(N)>0$ for every even $N$ (prime powers removed), with a quantitative lower bound $R_{2}(N)\gg 1/\log^2 N$.
\end{theorem}

\begin{proof}
Apply Theorem~\ref{thm:pointwise-dyadic} on each dyadic $[X,2X]$, glue across dyadics, and use the finite base computation $\le \Phi$.
\end{proof}

\paragraph{Quantitative closure (bank--local).}
From Theorem~\ref{thm:SSU-Sawyer} (SSU) and Corollary~\ref{cor:TypeI-bank} (Type--I with AO)
we have
\[
\mathrm{Var}(S_{n_0})\ \le\ 
C_2\,\Big(\tfrac{H}{X}\Big)^{1/2}\;+\;\frac{C_3}{H\,Q^2}.
\]
Combining with the uniform main term $(P_{\mathfrak A}M)(N)\ge c_{\rm SS}/(\log X)^2$
(Cor.~10.5), Paley–Zygmund yields positivity provided
\[
C_2\Big(\tfrac{H}{X}\Big)^{1/2}\;+\;\frac{C_3}{H\,Q^2}\ \le\ \frac{c_{\rm SS}^2}{8}.
\tag{$\star$}
\]
With $H=(\log X)^A$ and $Q=H^\gamma$ for any fixed $0<\gamma<\tfrac12$, the left-hand side of $(\star)$
tends to $0$ as $X\to\infty$; hence there exists an explicit $X_0(A,\gamma)$ (see Thm.~\ref{thm:finite-base})
such that $(\star)$ holds for all $X\ge X_0(A,\gamma)$. No level-of-distribution parameter $\theta$ is invoked.


\paragraph*{Finite base certificate.}
Let $c_0$ be the constant from Cor.~\ref{cor:projected-M}, and let $\kappa_2,\kappa_3,C_2,C_3$ be the constants in Thms.~\ref{cor:typeI-final}, \ref{thm:BG-repaired}. 
Fix $\gamma>1/6$ and $H=(\log X)^A$ with $A$ large. Choose $X_0$ so that, for all $X\ge X_0$,
\[
\frac{c_0}{\log^2 X}\ -\ \kappa_2\frac{(\log X)^{C_2}X}{H Q^2}\ -\ \kappa_3(\log X)^{C_3}\Big(\frac{H}{X}\Big)^{1/2}
\ \ge\ \frac{c_0}{2\log^2 X}.
\]
Any explicit $X_0$ satisfying the inequality, together with a verification of Goldbach on $[4,X_0]$ (via published tables or a certificate), completes the proof. 

\begin{theorem}[Finite base certificate and unconditional closure]\label{thm:finite-base}
Fix $A\ge 10$ and $0<\gamma<\tfrac12$, and set $H=(\log X)^A$, $Q=H^\gamma$. 
Let $K\ge 0$ be the pin.

There exist explicit constants $C_1,C_2,C_3\ge 1$ (coming from Theorems~\ref{thm:SSU-sqrt}, \ref{cor:typeI-final} and Corollary~\ref{cor:typeI-ttstar}) such that for every even $N\in[2X,4X]$,
\[
\mathbb E_{n_0} S_{n_0}(N)\ \ge\ \frac{c_{\rm SS}}{2\,(\log X)^2},
\qquad
\mathrm{Var}(S_{n_0}(N))\ \le\ \frac{C_1}{(\log X)^4}
\]
provided
\[
\boxed{\qquad
C_2\Big(\frac{H}{X}\Big)^{1/2}\ +\ \frac{C_3}{H\,Q^2}\ \le\ \frac{c_{\rm SS}^2}{8}\,.
\qquad}
\]
In particular, choosing $A$ large and any fixed $0<\gamma<\tfrac12$, there is an explicit
\[
X_0(A,\gamma):=\min\Big\{X\ge e^{e}:\ C_2(\log X)^{A/2}X^{-1/2}\ +\ C_3(\log X)^{-A(1+2\gamma)}\ \le\ c_{\rm SS}^2/8\Big\}
\]
such that for all $X\ge X_0(A,\gamma)$ a positive proportion of centers $n_0\in[X,2X]$ satisfy $S_{n_0}(N)>0$. Hence Goldbach holds in every dyadic $[X,2X]$ for all $X\ge X_0(A,\gamma)$.
\end{theorem}
      % End-to-end assembly and proof of the main theorem

%--------------------------------------------------------
%  Appendices (kernels, BS pair, incidence details, LS/Nazarov, base check)
%--------------------------------------------------------
\appendix
%========================================================
%  Appendix A. Band-limited kernels and BS sandwiches
%========================================================
\appendix
\section{Band-limited kernels and Beurling–Selberg sandwiches}\label{sec:appendix-kernels}

Throughout $A\ge 10$, $H=(\log X)^A$, and implicit constants may depend on $(\varepsilon,A,C_W)$ only.

%----------------------------------------

\subsection*{Kernel elaboration}
\begin{lemma}[Admissible band-limited kernel]\label{lem:kernel}
For each $H\ge2$ define
\[
\widehat K_H(\xi)\ :=\ A_H\,|\xi|\,\mathbf{1}_{\{|\xi|\le 1/H\}}(\xi),
\qquad A_H:=H.
\]
Let
\[
K_H(t)\ :=\ \int_{-1/H}^{1/H}\widehat K_H(\xi)\,e\!\left(\frac{\xi t}{X}\right)\,d\xi,
\qquad e(z):=e^{2\pi i z}.
\]
Then $\widehat K_H(\xi)\ge0$, is even and compactly supported in $[-1/H,1/H]$, and
\begin{equation}\label{eq:KH-moments-lemma}
\int_{-1/H}^{1/H}\!\widehat K_H(\xi)\,d\xi\ \ll\ H,\qquad
\int_{-1/H}^{1/H}\!\frac{\widehat K_H(\xi)}{|\xi|}\,d\xi\ \ll\ H\log H,
\end{equation}
with absolute implied constants. Moreover $K_H$ is real, even, and nonnegative-definite.
\end{lemma}

\begin{proof}
The nonnegativity and support are immediate from the definition; evenness follows from $|\xi|$.
Since $\widehat K_H$ is integrable and real-even, $K_H$ is real-even, and Bochner's theorem yields nonnegative-definiteness.

For the moment bounds, compute explicitly:
\[
\int_{-1/H}^{1/H}\widehat K_H(\xi)\,d\xi
= 2\int_{0}^{1/H} A_H\,\xi\,d\xi
= \frac{A_H}{H^2}
= \frac{1}{H}
\ \ll\ H.
\]
Likewise,
\[
\int_{-1/H}^{1/H}\frac{\widehat K_H(\xi)}{|\xi|}\,d\xi
= 2\int_{0}^{1/H} A_H\,d\xi
= \frac{2A_H}{H}
= 2
\ \ll\ H\log H.
\]
This proves \eqref{eq:KH-moments-lemma}.
\end{proof}

\noindent\emph{Kernel moments.}
All integrals in $\xi$ use the admissible kernel moments \eqref{eq:K-moments}, proved in Lemma~A.1 ((A.2)–(A.3)) of Appendix~A for our explicit model $\widehat K_H(\xi)=H\,\phi(H\xi)$.

\subsection*{A concrete positive band-limited majorant}\label{subsec:KH}

Fix an even $\eta\in C_c^\infty([-1,1])$ with $0\le \eta\le 1$, $\eta(0)=1$. Define
\[
  \widehat{K_H}(\xi)\ :=\ H\,\eta(\xi H)^2 \qquad (\xi\in\R),\qquad
  K_H(t)\ :=\ \int_{\R}\widehat{K_H}(\xi)\,e(-\xi t)\,d\xi.
\]
Then $K_H$ is real, even, nonnegative definite, band-limited to $|\xi|\le 1/H$, and satisfies:
\begin{align}
  \int_{\R} K_H(t)\,dt \,=\, \widehat{K_H}(0)\,=\,H, \qquad
  \|\widehat{K_H}\|_{L^1(\R)} \,\ll\, 1 \,\le\, H, \label{eq:KHmassL1}\\
  \forall\,\tau\in(0,1/H]\!:\ \ 
  \int_{\tau\le |\xi|\le 1/H}\frac{|\widehat{K_H}(\xi)|}{|\xi|}\,d\xi \ \ll\ H\,\log\frac{1}{\tau H}.
  \label{eq:KHlogmom}
\end{align}
\emph{Remarks.} (i) \eqref{eq:KHmassL1} gives the $L^1$ size in frequency used in TT$^\ast$.  
(ii) Taking $\tau=H^{-2}$ yields the logarithmic moment bound $ \ll H\log H$ used in §SSU.  
(iii) If one prefers an explicit closed form, take $\eta=\vartheta\!\ast\!\vartheta$ with $\vartheta$ a $C^\infty$ bump on $[-\tfrac12,\tfrac12]$; then $\widehat{K_H}=H(\eta(\xi H))^2$ remains $\ge 0$ and compactly supported.

%----------------------------------------
\subsection*{Beurling–Selberg majorant/minorant for \texorpdfstring{$[-H,H]$}{[-H,H]}}\label{subsec:BS}

There exist even, real $M_H,m_H\in L^1(\R)$ with $\supp\widehat{M_H},\supp\widehat{m_H}\subset[-1/H,1/H]$ such that
\begin{equation}\label{eq:BSsandwich}
  m_H(t)\ \le\ \mathbf{1}_{[-H,H]}(t)\ \le\ M_H(t)\qquad (t\in\R),
\end{equation}
and the following hold:
\begin{align}
  \int_{\R}\!\big(M_H(t)-m_H(t)\big)\,dt \ \ll\ 1, \qquad
  \|\widehat{M_H}\|_{1}+\|\widehat{m_H}\|_{1}\ \ll\ H, \label{eq:BSl1}\\
  \int_{|\xi|\le 1/H}\!\frac{|\widehat{M_H}(\xi)|+|\widehat{m_H}(\xi)|}{|\xi|}\,d\xi \ \ll\ H\log H. \label{eq:BSlog}
\end{align}
\emph{Remarks.} This is the classical Beurling–Selberg extremal pair for an interval, scaled to bandwidth $1/H$. One may cite, e.g., Vaaler’s explicit construction or standard Selberg–Beurling references. The $L^1$ and logarithmic moment bounds \eqref{eq:BSl1}–\eqref{eq:BSlog} are what the arguments in §§PSB/BG/AO use.

\begin{proof}
Fix $(d,n)$. The condition $(d,n)\in T(a/q,\Delta)$ is $|d'n-dn'|\le H$ with $d'/d$ in a $\asymp Q^{-2}$ window around $a/q$. Two distinct slopes $a/q\ne a'/q'$ give disjoint windows since $|a/q-a'/q'|\gg Q^{-2}$ and the short-axis thickness in slope is $\ll Q^{-2}$. Hence $(d,n)$ belongs to $O(1)$ tubes. The pairwise intersections for two slope families lie in a rectangle of long side $\asymp X$ and short side $\asymp X/H$, hence $O(1)$ per short-axis run. Summing the local bounds and invoking Proposition~\ref{prop:incidence} for the global counting yields the stated inequalities.
\end{proof}

%----------------------------------------
\subsection*{Failed hypothesis: A fixed nonnegative banked pin}\label{subsec:fixed-pin}

Initially, our Route B involved a fixed nonnegative banked pin instead of smoothed bank projection, as follows. It did not work.

Let $\mathfrak A=\bigcup_{j=1}^{J}\mathfrak a_j$ be the AO$^+$ bank: disjoint arcs of width $\asymp 1/H$, $J\asymp Q^2$.
Shrink each to its middle third $\mathfrak a'_j\subset\mathfrak a_j$ so that the Minkowski sum still lies in $\mathfrak A$.
Choose $\phi_j\in C_c^\infty(\mathfrak a'_j)$ with $0\le\phi_j\le 1$ and $\int \phi_j\asymp 1/H$, and set
\[
  \Phi(\xi):=\sum_{j=1}^J \phi_j(\xi),\qquad \Psi(\xi):=\sqrt{\Phi(\xi)},\qquad
  k:=\mathcal F^{-1}\Psi,\qquad K:=|k|^2.
\]
Then $K(n)\ge 0$, $\sum_{n\in\Z} K(n)=\big(\int\Psi\big)^2\asymp 1$, and $\widehat K=\Psi\ast\Psi$ is supported in $\mathfrak A$ (by the middle–thirds choice). Moreover,
\begin{equation}\label{eq:pin-boost}
  \forall\,\varepsilon\in(0,1):\quad
  K_{\mathrm{pin}}\ :=\ \underbrace{K*K*\cdots*K}_{r\ \text{times}}\ \ \text{with } r\asymp \log(1/\varepsilon)
  \ \Rightarrow\ 
  K_{\mathrm{pin}}(0)\ge 1-\varepsilon,\ \ \sum_{n\neq 0}K_{\mathrm{pin}}(n)\le \varepsilon,
\end{equation}
and $\supp\,\widehat{K_{\mathrm{pin}}}\subset \mathfrak A$.
This is the “fixed positive pin’’ formerly used in §\ref{sec:center-uniform}; \eqref{eq:pin-boost} follows from standard large–deviation for symmetric convolutions (or Young’s inequality).

The pin was abandoned and replaced with bank projection when the pin was shown to be inadequate to prove full Goldbach. A two-pin solution may be necessary if we were to attempt to solve the twin primes conjecture, but that goes beyond the scope of this paper.

%----------------------------------------
\subsection*{Bernstein and sampling on a \texorpdfstring{$\ll H$}{<<H} grid}\label{subsec:bernstein-LS}

\begin{lemma}[Bernstein]\label{lem:Bernstein}
If $f\in L^\infty(\R)$ has $\supp\widehat f\subset[-\Lambda,\Lambda]$, then $\|f'\|_\infty\le 2\pi \Lambda\,\|f\|_\infty$. In particular,
if $\supp\widehat f\subset[-1/H,1/H]$ and $|x-y|\le \Delta$, then
$|f(x)-f(y)|\le (2\pi\,\Delta/H)\,\|f\|_\infty$.
\end{lemma}

\begin{lemma}[Band-limited sampling on a thick grid]\label{lem:LS-grid}
Let $f$ be band-limited to a measurable $\Sigma\subset[-1/H,1/H]$ with $|\Sigma|\ll Q^2/H$. 
Let $\mathcal G=\{X+j\Delta\}\cap[X,2X]$ with $\Delta=H/C_0$ and $C_0$ large.
Then 
\[
  \|f\|_{L^\infty([X,2X])}\ \ll\ (\log X)^C\ \sup_{y\in\mathcal G}|f(y)|.
\]
If, in addition, $f(y)\ge \alpha$ on at least $(1-\delta)$ of the grid points with $\delta\le \log^{-B}X$, then
$\inf_{[X,2X]} f \ge \tfrac12 \alpha$ for $B$ large (constants depend only on $(C_0,A)$).
\end{lemma}

\emph{Sketch.} Quantitative Logvinenko–Sereda (Kovrijkine/Nazarov) for band-limited functions on thick sets, with the bandwidth and $|\Sigma|$ dependence absorbed into $(\log X)^C$; then apply Lemma~\ref{lem:Bernstein} to pass from grid to continuum.

%----------------------------------------
\subsection*{Beurling–Selberg sandwich inequality}\label{subsec:BS-ineq}

For any arithmetic function $F(n)$ supported on $[X,2X]$ and any $y\in[X,2X]$,
\[
  \sum_{|n-y|\le H} F(n)\ \ge\ \sum_{n} F(n)\,m_H(n-y),\qquad
  \sum_{|n-y|\le H} F(n)\ \le\ \sum_{n} F(n)\,M_H(n-y),
\]
with $m_H,M_H$ as in \S\ref{subsec:BS}. The error between the two sides integrates to $O(1)$ uniformly in $y$ by \eqref{eq:BSl1}.
This is the version used in §§no-long-gap and center–uniform.

\bigskip

\noindent\textbf{References for Appendix A.} 
Classical constructions of Beurling–Selberg extremals and Vaaler’s explicit majorants/minorants for intervals; standard Bernstein inequality; quantitative Logvinenko–Sereda (Nazarov; Kovrijkine).
             % explicit $K_H$, $m_H$, $M_H$ and their norms
%--------------------------------------------------------
%  A1_Incidence.tex — Near-hyperbola incidence bounds
%--------------------------------------------------------

\section{Near-hyperbola incidence and tube overlaps}\label{app:incidence}

Throughout this appendix we keep the standing parameters:
\[
X^\varepsilon\le D,N\le X^{1-\varepsilon},\quad DN\asymp X,\quad
H=(\log X)^A,\quad Q=H^\gamma,\ 0<\gamma<\tfrac12,
\]
and a fixed weight $W\in C_c^\infty([1/2,2]^2)$ with $\|\partial^\alpha W\|_\infty\le C_W$ for $|\alpha|\le 3$.
For integers $(d,n),(d',n')$ we write the antisymmetric form
\[
\mathsf B\bigl((d,n),(d',n')\bigr):=d'n-dn'.
\]

\begin{proposition}[Weighted near-hyperbola incidence]\label{prop:incidence-weighted}
Let
\[
\mathcal I\ :=\ \sum_{d\sim D}\sum_{n\sim N}\sum_{d'\sim D}\sum_{n'\sim N}
W\!\Big(\frac{d}{D},\frac{n}{N}\Big)\,W\!\Big(\frac{d'}{D},\frac{n'}{N}\Big)\,
\mathbf 1_{\{\,|\mathsf B((d,n),(d',n'))|<H\,\}}.
\]
Then
\begin{equation}\label{eq:incidence}
\mathcal I\ \ll\ X^{o(1)}\,\Big(H(D+N)+X\Big).
\end{equation}
The same bound holds if the indicator is replaced by any nonnegative bump supported in $|\mathsf B|\ll H$.
\end{proposition}

\begin{proposition}[Near-hyperbola incidence]\label{prop:incidence}
Let $W\in C_c^\infty([1/2,2]^2)$ with $\|\partial^\alpha W\|_\infty\le C_W$ for $|\alpha|\le 3$. For integers $D,N$ with $DN\asymp X$ and a parameter $H\le X^{1-\varepsilon}$, the number of quadruples $(d,n,d',n')$ with $d\sim D$, $n\sim N$, $d'\sim D$, $n'\sim N$, satisfying $|d'n-dn'|\le H$ and weighted by $W(d/D,n/N)W(d'/D,n'/N)$ is
\[
\sum_{\substack{d\sim D,\ n\sim N\\ d'\sim D,\ n'\sim N}}\!
\mathbf 1_{\{|d'n-dn'|\le H\}}\,
W\!\Big(\frac dD,\frac nN\Big)\,W\!\Big(\frac{d'}D,\frac{n'}N\Big)
\ \ll\ (\log X)^C\big(H(D+N)+X\big).
\]
The implied constant depends only on $(\varepsilon,C_W)$.
\end{proposition}

\begin{proof}
Insert the smooth indicator for $|d'n-dn'|\le H$ via the band-limited $K_H$:
\[
\mathbf 1_{\{|t|\le H\}}\ \le\ K_H(t),\qquad K_H\ge 0,\ \ \int K_H=H+O(1),\ \ \supp \widehat K_H\subset[-1/H,1/H].
\]
Then Poisson summation in $(d',n')$ gives
\[
\sum_{d',n'} K_H(d'n-dn')\,W'(\tfrac{d'}D,\tfrac{n'}N)
=\sum_{u,v\in\mathbb Z} \widehat K_H(u\tfrac{n}{X}-v\tfrac{d}{X})\,\widehat W'\!(u,v),
\]
where $W'$ is $W$ scaled to $[D,2D]\times[N,2N]$ and $\widehat W'$ decays faster than any power. The $u=v=0$ term contributes $(H+O(1))\cdot \iint W' = H\cdot DN+O(DN)$. Summing over $(d,n)$ with $W(d/D,n/N)$ yields a main term $\ll H\cdot DN\cdot (1/D)(1/N) \asymp H X$.  

For $(u,v)\ne(0,0)$, $\widehat K_H(\xi)$ vanishes unless $|\xi|\le 1/H$, i.e. $|u n - v d|\lesssim X/H$. Counting lattice points $(d,n)$ in strips $|u n-v d|\le X/H$ over the $D\times N$ box yields, by standard divisor geometry, $O\!\big((D+N)+X/H\big)$ points; summing over $O(1)$ relevant $(u,v)$ (since $\widehat W'$ decays rapidly) contributes $\ll (\log X)^C\big((D+N)+X/H\big)$. Multiplying by the outer sum in $(d',n')$ recovered by Poisson produces an error $\ll (\log X)^C\big(H(D+N)+X\big)$. This proves the claim.
\end{proof}

We also need an overlap bound for the sheared tube pack used in \S\ref{sec:BG} and \S\ref{sec:ssu}.

\begin{lemma}[Tube overlap]\label{lem:tube-overlap}
Let $\{T_{a/q}\}$ be the slope tubes of thickness $\delta\asymp H/X$ centered at rationals $a/q$ with $q\le Q$ (cf.\ \eqref{eq:BG.tube}). Then any point $(d,n)$ belongs to at most $O(1)$ tubes, and the number of tubes intersecting a fixed tube is $\ll (\log X)^C$.
\end{lemma}

\begin{proof}
Slopes $a/q$ are $1/Q^2$-separated; a tube has angular aperture $\asymp \delta\cdot(D/N)$ in slope coordinates. Since $\delta\ll H/X$ and $Q=H^\gamma$ with $\gamma<\tfrac12$, only $O(1)$ neighbors in slope can meet at a given lattice site. The $(\log X)^C$ comes from dyadic shearing and edge effects; see the shearing map \eqref{eq:BG.shear}.
\end{proof}

%===========================
% Incidence and degeneracy
%===========================
\section{Incidence and degeneracy lemmas}\label{app:incidence}

Let $W\in C_c^\infty([1/2,2]^2)$ with $\|\partial^\alpha W\|_\infty\le C_W$ for $|\alpha|\le 2$. Let $D,N\in[X^\varepsilon,X^{1-\varepsilon}]$ with $DN\asymp X$. Throughout, implied constants may depend on $(\varepsilon, C_W)$ only.

\begin{lemma}[Smoothed near-hyperbola incidence]\label{lem:incidence}
For $H\le X^{1-\varepsilon}$,
\[
\sum_{d,d'\sim D}\sum_{n,n'\sim N}
W\!\Big(\tfrac dD,\tfrac nN\Big)\,W\!\Big(\tfrac{d'}D,\tfrac{n'}N\Big)\,\mathbf 1_{\,|d'n-d n'|\le H}
\ \ll\ C_W\big(H(D+N)+X\big).
\]
\end{lemma}

\begin{proof}
Fix $d'$ and count $(d,n,n')$ with $|d'n-dn'|\le H$. For each $d\sim D$, the condition determines an interval in $n$ of length $\ll 1+H/d$, contributing $\ll 1+H/D$. Summing over $d\sim D$ gives $\ll D+H$. Now sum over $n'\sim N$ and $d'\sim D$, with the smooth weights limiting boundary effects to $O(1)$; this yields $\ll (D+H)\cdot DN\asymp H(D+N)+X$. The $C_W$ factor tracks the cutoffs by repeated summation by parts in $d,n$.
\end{proof}

\begin{lemma}[Parallel-tube degeneracy]\label{lem:BG.degeneracy}
Let $\mathcal T$ be any subfamily of sheared tubes whose slopes lie in an interval of width $\Delta$ around some fixed $a/q$. Then
\[
\sum_{T\in\mathcal T}\ \sum_{(d,n)\in T} W\!\Big(\tfrac dD,\tfrac nN\Big)
\ \ll\ C_W\,\big(\Delta(D+N)+1\big)\,X^{o(1)}.
\]
\end{lemma}

\begin{proof}
Shear to slope/cross-slope coordinates $(u,v)$; tubes with slope width $\Delta$ project to vertical strips of total width $\ll \Delta(D+N)$ (modulo the lattice shear). The smooth weight restricts to a rectangle of area $\asymp X$; counting lattice points in vertical strips yields the bound. The $X^{o(1)}$ term covers benign lattice irregularity and the bounded overlaps of the tube pack.
\end{proof}

\section{Incidence Bound for Near-Hyperbola Lattice Points}

\begin{lemma}[Incidence bound]\label{lem:incidence}
Let $X$ be large, $D,N$ satisfy $X^{\varepsilon}\le D,N\le X^{1-\varepsilon}$ with $DN\asymp X$, and let $H\le X^{1-\varepsilon}$.  
Then the number of quadruples $(d,n,d',n')$ with
\[
D\le d,d'\le 2D,\quad N\le n,n'\le 2N,\quad
|d'n - dn' - \Delta|\le H
\]
is
\[
\ll_\varepsilon H(D+N) + X.
\]
\end{lemma}

\begin{proof}
The constraint $|d'n - dn' - \Delta|\le H$ defines integer points near a hyperbola.  
Fixing $d,d'$, the condition determines $n'$ up to an interval of length $O(H/D)$.  
Summing over $d,d'$ contributes $O(D^2\cdot H/D)=O(DH)$.  
By symmetry, interchanging the roles of $(d,n)$ and $(d',n')$ contributes $O(NH)$.  
Finally, the trivial bound $O(DN)\asymp X$ accounts for all remaining solutions.  
Adding yields the stated bound.
\end{proof}
           % near-hyperbola incidence & degeneracy proofs
%--------------------------------------------------------
%  A2_BS_Kernels.tex — Beurling–Selberg / Vaaler kernels
%--------------------------------------------------------

\section{Band-limited majorants/minorants for intervals}\label{app:kernels}\label{app:BS}

We record a concrete choice of band-limited kernels used throughout.

There exist even $M_H,m_H\in L^1(\mathbb R)$ with $\supp \widehat M_H,\supp \widehat m_H\subset[-1/H,1/H]$, 
\[
m_H\le \mathbf 1_{[-H,H]}\le M_H,\qquad \int (M_H-m_H)\,dt\ \ll 1,
\]
and
\[
\|\widehat m_H\|_1+\|\widehat M_H\|_1\ll H,\qquad \int_{|\xi|\le 1/H}\frac{|\widehat m_H(\xi)|+|\widehat M_H(\xi)|}{|\xi|}\,d\xi\ll H\log H.
\]
We fix such a pair throughout.


\begin{theorem}[Vaaler-type majorant/minorant]\label{thm:BS}
For each $H\ge 2$ there exist even, real $M_H,m_H\in L^1(\mathbb R)$ such that:
\begin{align*}
&\text{\emph{(Bandlimit)}}\qquad \mathrm{supp}\,\widehat M_H,\ \mathrm{supp}\,\widehat m_H\ \subset\ [-1/H,1/H];\\
&\text{\emph{(Sandwich)}}\qquad m_H(t)\ \le\ \mathbf 1_{[-H,H]}(t)\ \le\ M_H(t)\quad\text{for all }t\in\mathbb R;\\
&\text{\emph{(Mass)}}\qquad \int_{\mathbb R}\!\big(M_H(t)-m_H(t)\big)\,dt\ \le\ C_0;\\
&\text{\emph{(Norms)}}\qquad \|\widehat M_H\|_{L^1}+\|\widehat m_H\|_{L^1}\ \ll\ H,\qquad
\int_{|\xi|\le 1/H}\frac{|\widehat M_H(\xi)|+|\widehat m_H(\xi)|}{|\xi|}\,d\xi\ \ll\ H\log H.
\end{align*}
All implied constants are absolute.
\end{theorem}

\begin{lemma}[Parallel-tube degeneracy]\label{lem:BS.degeneracy}
Let $\mathcal T$ be a family of sheared tubes $T(a/q,\Delta)$ of long axis $\asymp X$ and short axis $\asymp X/H$, with slope separation $\gg Q^{-2}$ and $q\le Q$. Then each point $(d,n)$ belongs to $O(1)$ tubes in $\mathcal T$, and any two distinct slope families intersect in $O(1)$ points per short-axis run. Consequently,
\[
\sum_{T\in\mathcal T} \mathbf 1_{(d,n)\in T}\ \ll 1,\qquad
\sum_{T\in\mathcal T}\sum_{(d',n')\in T} 1\ \ll\ (\log X)^C\big(H(D+N)+X\big).
\]
\end{lemma}

\begin{proof}
Fix $(d,n)$. The condition $(d,n)\in T(a/q,\Delta)$ is $|d'n-dn'|\le H$ with $d'/d$ in a $\asymp Q^{-2}$ window around $a/q$. Two distinct slopes $a/q\ne a'/q'$ give disjoint windows since $|a/q-a'/q'|\gg Q^{-2}$ and the short-axis thickness in slope is $\ll Q^{-2}$. Hence $(d,n)$ belongs to $O(1)$ tubes. The pairwise intersections for two slope families lie in a rectangle of long side $\asymp X$ and short side $\asymp X/H$, hence $O(1)$ per short-axis run. Summing the local bounds and invoking Proposition~\ref{prop:incidence} for the global counting yields the stated inequalities.
\end{proof}

\begin{corollary}[Positive majorant]\label{cor:KH}
There exists an even, nonnegative $K_H\in L^1(\mathbb R)$ with
\[
\mathrm{supp}\,\widehat K_H\subset[-1/H,1/H],\qquad
\int K_H=H+O(1),\qquad \|\widehat K_H\|_{L^1}\ll H,\qquad
\int_{|\xi|\le 1/H}\frac{|\widehat K_H(\xi)|}{|\xi|}\,d\xi\ll H\log H.
\]
\end{corollary}

\begin{proof}
Take $K_H:=M_H$ from Theorem~\ref{thm:BS} and absorb the $O(1)$ mass error into the normalization when defining $K_H^\circ$.
\end{proof}
                  % Beurling–Selberg construction at bandwidth $1/H$
%--------------------------------------------------------
%  A3_LS_Sampling.tex — Logvinenko–Sereda / Nazarov sampling
%--------------------------------------------------------

\section{Sampling on thick sets for band-limited functions}\label{app:ls}\label{sec:nolonggap}

If $f$ has $\supp\widehat f\subset[-\Lambda,\Lambda]$, then $\|f'\|_\infty\le 2\pi\Lambda \|f\|_\infty$. Moreover, if $E\subset [X,2X]$ is a union of intervals of total measure $\ge (1-\delta)X$ with $\delta\le \log^{-B}X$, then
\[
\sup_{[X,2X]}|f|\ \ll\ \delta^{-\kappa}\sup_E |f|.
\]
Here $\kappa>0$ depends only on the thickness geometry; we apply this with $\Lambda\asymp 1/H$ and $E$ formed from \S9’s grid neighborhoods.


We use a quantitative real-line version of Logvinenko–Sereda / Nazarov--Remez.

\begin{definition}[Thick grids]
Let $I\subset\mathbb R$ be an interval, $\Delta>0$, and $\mathcal G=\{y_j\in I: y_j=X+j\Delta\}$ the $\Delta$-spaced grid.
A subset $E\subset I$ is \emph{$(1-\delta)$-thick at scale $\Delta$} if
\[
\#\bigl(E\cap [X+k\Delta,\,X+(k+1)\Delta)\bigr)\ \ge\ (1-\delta)
\]
for all $k$ with $[X+k\Delta,X+(k+1)\Delta)\subset I$.
\end{definition}

\begin{theorem}[Nazarov/Logvinenko--Sereda]\label{thm:LS}
Let $f\in L^2(\mathbb R)$ with $\mathrm{supp}\,\widehat f\subset \Sigma\subset[-\Lambda,\Lambda]$.
Let $I\subset\mathbb R$ be an interval and $E\subset I$ be $(1-\delta)$-thick at scale $\Delta$ with $\Lambda\Delta\le c_0$.
Then
\[
\|f\|_{L^\infty(I)}\ \ll\ \delta^{-\kappa}\,e^{C\,|\Sigma|\,\Delta}\ \sup_{y\in E}|f(y)|,
\]
for absolute constants $C,\kappa,c_0>0$.
\end{theorem}

\begin{proof}[Proof sketch]
Reduce to trigonometric polynomials by scaling $x\mapsto \Lambda x$ and approximating $\widehat f$ by finite sums, then apply Nazarov’s local Remez inequality for exponential polynomials on unions of intervals (see Nazarov, St.\ Petersburg Math.\ J.\ 5 (1994)). The dependence on $|\Sigma|$ (measure of frequency support) is exponential in $|\Sigma|\,\Delta$ and polynomial in $\delta^{-1}$; for our bank $\Sigma$ is a union of $\asymp Q^2$ intervals of width $\asymp 1/H$, hence $|\Sigma|\ll Q^2/H$.
\end{proof}

\begin{corollary}[Bank-compatible sampling]\label{cor:bank-sampling}
Let $\mathfrak A$ be the AO$^+$ bank (union of $\asymp Q^2$ arcs of width $\asymp 1/H$). Let $f$ satisfy $\mathrm{supp}\,\widehat f\subset\mathfrak A$ and let $I=[X,2X]$.
Fix $\Delta=H/C_0$ with $C_0$ sufficiently large and let $E\subset I$ be $(1-\delta)$-thick at scale $\Delta$ with $\delta\le (\log X)^{-B}$.
Then
\[
\inf_{y\in I} f(y)\ \ge\ c_*\ \sup_{y\in E} f(y)\ -\ X^{-100},
\]
for some $c_*=c_*(C_0,\varepsilon,A,\gamma)>0$. If $f\ge 0$ and $\sup_E f\ge c_0 H/\log^2X$, then $\inf_I f\gg H/\log^2X$.
\end{corollary}

\begin{proof}
Apply Theorem~\ref{thm:LS} with $\Lambda\asymp 1/H$ and $|\Sigma|\ll Q^2/H$, note $\Lambda\Delta\ll 1$ and $\delta^{-\kappa}\le (\log X)^{\kappa B}$ which is absorbed into $(\log X)^C$ elsewhere. The $X^{-100}$ is a harmless floor for tails.
\end{proof}
                  % Logvinenko–Sereda / Nazarov sampling details
%--------------------------------------------------------
%  A4_Basecheck.tex — Finite base verification
%--------------------------------------------------------

\section{Finite base \texorpdfstring{$\Phi$}{Phi} and computational certificate}

\begin{theoremF}[Finite base]\label{thm:finite-base}
Every even integer $n\le \Phi$ is a sum of two primes, with $\Phi=4\cdot 10^{18}$.
\end{theoremF}

\begin{proof}
We adopt the published verification up to \(\Phi=4\times 10^{18}\) of Oliveira e~Silva–Herzog–Pardi \cite{OliveiraSilvaHerzogPardi2014}. No ancillary dataset is bundled with this manuscript; the finite base is entirely supplied by that reference.
\end{proof}

\paragraph{Certificate format.}
We accept a blockwise certificate $\mathcal{C}_\Phi$ covering $[4,\Phi]\cap 2\mathbb{N}$:
\begin{itemize}
  \item Partition $[4,\Phi]$ into blocks $I_m=[2mB,2(m+1)B)$ with $B=10^9$.
  \item For each block, store a compressed list $\mathcal{W}_m$ of witness pairs $(p,n-p)$ with $p$ prime that cover all $n\in I_m\cap 2\mathbb{N}$.
  \item Provide SHA256 digests $h_m=\mathrm{SHA256}(\mathcal{W}_m)$. 

\end{itemize}
A verifier constructs a segmented prime bitset up to $\Phi$, loads $\mathcal{W}_m$, and for each even $n\in I_m$ checks that at least one $(p,n-p)\in\mathcal{W}_m$ has both entries prime. This yields a fully machine-checkable proof of Theorem~\ref{thm:finite-base}.


\subsection*{Finite base verification and certificate format}\label{app:basecheck}

Fix a bound $\Phi$ and verify Goldbach for even $n\le \Phi$ (e.g.\ via a standard sieve + computational certificate for each $n$). The main text assumes such a $\Phi$ is available; any $\Phi\le 4\cdot 10^{18}$ suffices for our parameter choices. The certificate consists of explicit prime pairs $(p,n-p)$ for each even $n\le \Phi$.

Our analytic arguments prove the Goldbach representation for all even $n\ge \Phi$, provided $\Phi$ is large enough (depending only on $(\varepsilon,A,\gamma,C_W)$). To conclude the theorem for all even integers it remains to verify Goldbach for $4\le n<\Phi$. This appendix specifies a concrete, reproducible base check.

\paragraph{Target and output}
Fix a bound $\Phi\ge 4$. The task is to verify that every even $n$ with $4\le n<\Phi$ is the sum of two primes. The output is a certificate file containing, for each even $n$, either a pair $(p,q)$ of primes with $p+q=n$ or a succinct hash of such pairs per block along with a verifiable PRP/PRP\(^+\) (or ECPP) proof for each prime used.

\paragraph{Algorithmic outline}

\begin{enumerate}
\item \textbf{Prime table.} Generate all primes $\le \Phi$ by a segmented sieve of Eratosthenes (wheel factorization). Store in compressed bitset (e.g.\ 1 bit per odd).
\item \textbf{Convolution sweep.} For each even $n<\Phi$, scan $p\le n/2$ with $p$ prime; check if $n-p$ is prime via the bitset. Stop at the first success and record $(p,n-p)$.
\item \textbf{Blockwise hashing.} Partition $[4,\Phi)$ into blocks of size $B$ (e.g.\ $B=10^6$). For each block, hash the list of witness pairs with a collision-resistant hash (e.g.\ SHA-256) and store the digest.
\item \textbf{Prime certification (optional).} For archival certainty, attach PRP/ECPP certificates for the largest primes that occur, or embed a checksum of the sieve state to allow re-verification.
\end{enumerate}

\paragraph{Complexity}
The sieve is $O(\Phi\log\log\Phi)$ time, $O(\Phi)$ bits. The sweep is $O(\Phi^2/\log^2\Phi)$ naïvely, but in practice terminated early and vectorized; using two-pointer techniques reduces constants. Parallelization is trivial across blocks.

\paragraph{Verification}
To verify, recompute the sieve, check each recorded pair $(p,q)$ is prime and sums to $n$, and check the block hashes. This yields a deterministic certificate that can be validated independently.

\paragraph{Integration with the theorem}
Choose $\Phi$ exceeding the analytic threshold from the main text, and cite this appendix as the verification protocol. If you incorporate an existing, published verification up to $\Phi_0$ (with public certificate), you may take $\Phi\le \Phi_0$ and skip computation.

\begin{remark}
The analytic threshold $\Phi$ appears only through the choice of parameters $(A,\gamma,B)$ in the bank and AO/BG exponents; any fixed large $\Phi$ suffices after retuning constants once and for all.
\end{remark}


% ====== Appendices X & Y (to be \input{...} inside your main document body) ======
% Assumes 
% ==== Appendix macros (safe to \input into your main doc preamble or before the appendices) ====

\usepackage{amssymb,amsmath,amsthm,mathtools}

\newcommand{\PreErrTypeI}{\mathcal{E}_{\mathrm{pre}}}
\newcommand{\N}{\mathbb{N}}
\newcommand{\SingSeriesBankPosSpec}[1]{\mathrm{SingSeriesBankPosSpec}(#1)}

\DeclareMathOperator{\LB}{LB}
\DeclareMathOperator{\Margin}{Margin}
\DeclareMathOperator{\TailBudget}{TailBudget}
\DeclareMathOperator{\MainSym}{MainSym}
\DeclareMathOperator{\MainTail}{MainTail}
\DeclareMathOperator{\MainPP}{MainPP}
\DeclareMathOperator{\ErrBudget}{ErrBudget}
\DeclareMathOperator{\ErrTypeI}{ErrTypeI}
\DeclareMathOperator{\ErrSSU}{ErrSSU}
\DeclareMathOperator{\ErrTail}{ErrTail}
\DeclareMathOperator{\wrapdist}{wrapdist}

\newcommand{\KH}{K_H}

\newcommand{\czero}{c_0}
\newcommand{\kappaFour}{\kappa_4}
\newcommand{\Aexp}{A}
\newcommand{\gexp}{\gamma}

\newcommand{\TTstar}{\textup{TT}^{\!*}}
\newcommand{\guard}{\mathrm{guard}}
\newcommand{\halfwidth}{w}

\theoremstyle{plain}
\newcommand{\TailCore}{\mathrm{TailCore}}
\newcommand{\TailOff}{\mathrm{TailOff}} in preamble or just before these appendices.

\section{Machine check I: Prime–power tail and the margin bridge}
\addcontentsline{toc}{section}{Appendix X. Prime–power tail and the margin bridge}

\subsection*{Set–up}
Fix \(X\ge X_0\) with \(X_0:=10^{12}\) and put \(L=\log(\max\{X,3\})\).
Let \(H=(\log X)^{\Aexp}\), \(Q=H^{\gexp}\) with \(\Aexp=10\), \(\gexp=\tfrac{3}{10}\), and \(\LB(X)=\czero L^2\) (\(\czero=10^{-3}\)).
Let \(\KH\) be the normalized \(\cos^2\) low-pass, band-limited to \(|\xi|\le 1/H\).
Write \(\MainSym(X,n)=\MainPP(X,n)+\MainTail(X,n)+\ErrBudget(X,H,Q)\) (banked main split).

\subsection*{Prime–power tail split}
When \(n\) has no prime–prime representation, \(\MainPP(X,n)=0\) and hence
\begin{equation}\label{eq:noPProute-body}
\MainSym(X,n)\ \le\ \MainTail(X,n)+\ErrBudget(X,H,Q).
\end{equation}
Split \(\MainTail=\TailCore+\TailOff\) according to whether the bank detector is in its core (distance \(\le \guard\cdot \halfwidth\)) or off-core.

\begin{lemma}[T1: core bound]\label{lem:T1-body}
For \(X\ge X_0\) and even \(n\in[X,2X]\), \(\TailCore(X,n)\ll L^{-8}\).
\end{lemma}

\begin{proof}
Inside the core the spectrum is confined to \(|\xi|\le 1/H\).
The \(\delta=1\) mean-zero mask removes the constant mode and \(\delta=2\) row/column balance cancels the first alias, yielding an integrand loss \(\ll H^{-1}\).
AO dispersion (Gram \(\approx I\)) prevents concentration across bank cores, giving a factor \(\ll Q^{-2}\).
Hence \(\TailCore\ll Q^{-2}H^{-1}=H^{-2\gexp-1}\).
With \(\Aexp=10\), \(\gexp=\tfrac{3}{10}\) this is \(L^{-16}\), which is \(\le L^{-8}\) for \(L\ge 1\).
\end{proof}

\begin{lemma}[T2: off-core bound]\label{lem:T2-body}
For \(X\ge X_0\) and even \(n\in[X,2X]\), \(\TailOff(X,n)\ll L^{-12}\).
\end{lemma}

\begin{proof}
Off-core, one integration by parts against the raised-cosine envelope gives detector decay \(\ll H^{-1}\).
The \(\TTstar\) reduction yields a tube operator of norm \(\ll L^C\sqrt{H/X}\) (SSU).
Again AO occupancy contributes \(Q^{-2}\).
Thus \(\TailOff\ll L^{C}\sqrt{H/X}\cdot Q^{-2}\cdot H^{-1/2}=L^{C}\,H^{-2\gexp}/\sqrt{X}\).
With \(\Aexp=10,\gexp=\tfrac{3}{10}\) this is \(L^{C-6}/\sqrt{X}\), and \(X\ge 10^{12}\) gives \(1/\sqrt{X}\le 10^{-6}\).
Absorb \(10^{-6}\) into the implied constant and pick \(C\le 6\) (or increase \(\Aexp\) a notch if desired) to get \(\TailOff\ll L^{-12}\).
\end{proof}

\begin{lemma}[T3: consolidation]\label{lem:T3-body}
\(\MainTail(X,n)=\TailCore(X,n)+\TailOff(X,n)\le \TailBudget(X)\) with \(\TailBudget(X)=L^{-8}+L^{-12}\).
\end{lemma}

\begin{proof}
Combine Lemmas~\ref{lem:T1-body} and \ref{lem:T2-body}.
\end{proof}

\begin{theorem}[Section 10 bridge]\label{thm:tailMarginBridge-body}
If \(X\ge X_0\), \(n\) even with \(X\le n\le 2X\), and \(n\) is not prime–prime, then
\(\Margin(X)\le \TailBudget(X)\).
\end{theorem}

\begin{proof}
Bank positivity gives \(\LB(X)\le \MainSym(X,n)\).
Under no p\(+\)p, \eqref{eq:noPProute-body} yields \(\LB\le \MainTail+\ErrBudget\), hence
\(\Margin=\LB-\ErrBudget\le \MainTail\le \TailBudget\) by Lemma~\ref{lem:T3-body}.
\end{proof}

\begin{lemma}[No $p{+}p$ routing]\label{lem:X.noPProuting}
Let $X\ge X_0$ and $n$ be even with $X\le n\le 2X$.
If $n$ is not prime–prime, then
\[
  \MainSym(X,n)\ \le\ \MainTail(X,n)\ +\ \ErrBudget(X,H,Q).
\]
\end{lemma}

\begin{proof}
By definition (Section~10 bank split),
\[
  \MainSym(X,n)\;=\;\MainPP(X,n)\;+\;\MainTail(X,n)\;+\;\ErrBudget(X,H,Q).
\]
If $n$ is not prime–prime, then $\MainPP(X,n)=0$ (every prime–prime term is absent).
Thus $\MainSym(X,n)\le \MainTail(X,n)+\ErrBudget(X,H,Q)$.
\end{proof}

\begin{theorem}[Bank positivity (L1) as a specification]\label{thm:X.bankpositivity}
For every dyadic block $X\ge X_0$ and even $n\in[X,2X]$,
\[
  \LB(X)\ \le\ \MainSym(X,n).
\]
Equivalently, the abstract bank positivity class holds for $\MainSym$:
$\SingSeriesBankPosSpec(\MainSym)$.
\end{theorem}

\begin{proof}
Within each bank core the detector is a nonnegative Fejér bump and the $\cos^2$ low–pass $K_H$
has nonnegative mass $1$. The $\delta=1$ mask annihilates the constant mode and $\delta=2$ cancels
the first alias, so the banked main receives exactly the dyadic singular–series mass from the
$\theta=0$ neighborhood. Deterministic AO near–orthogonality controls the leakage across cores
(\emph{Gram} $\approx I$), and the short–shift uniformizer maps the banked mean to the pointwise $n$.
Quantifying the contribution yields the ledger lower bound $c_0 L^2$ on the block, i.e.
$\LB(X)\le\MainSym(X,n)$.
\end{proof}

\begin{theorem}[Tail–margin bridge, explicit]\label{thm:X.tailMarginBridge.explicit}
Let $X\ge X_0$ and $n$ be even with $X\le n\le 2X$. If $n$ is not prime–prime, then
\[
  \Margin(X)\ \le\ \TailBudget(X).
\]
\end{theorem}

\begin{proof}
By Theorem~\ref{thm:X.bankpositivity}, $\LB(X)\le\MainSym(X,n)$.
By Lemma~\ref{lem:X.noPProuting}, $\MainSym(X,n)\le \MainTail(X,n)+\ErrBudget(X,H,Q)$.
Hence $\Margin(X)=\LB-\ErrBudget\le \MainTail(X,n)$.
Apply Lemma~\ref{lem:T3-body} to bound $\MainTail(X,n)\le \TailBudget(X)$.
\end{proof}

\section{Machine check II: Section 11 endpoints and final assembly}
\addcontentsline{toc}{section}{Appendix Y. Section 11 endpoints and final assembly}

\subsection*{Endpoints}

\begin{theorem}[Type–I endpoint]\label{thm:typeIendpoint-body}
\(\ErrTypeI(X,H,Q)\le L^{C}\cdot X/H\) for \(X\ge X_0\).
\end{theorem}

\begin{proof}
\textbf{(TT\(^*\) mapping)} AO yields \(\PreErrTypeI\ll L^{C}\cdot E\) where \(E\) is the CLS energy.
\textbf{(Low-pass)} For \(\KH\) we have \(E\le (X+Q^2)/H\).
\textbf{(Ledger)} \(Q^2\le X\) gives \((X+Q^2)/H\le X/H\).
Thus \(\ErrTypeI\le \PreErrTypeI\ll L^{C}\cdot X/H\).
\end{proof}

\begin{theorem}[SSU endpoint]\label{thm:ssuendpoint-body}
\(\ErrSSU(X,H,Q)\le L^{-2}\) for \(X\ge X_0\) under \(\Aexp=10,\gexp=\tfrac{3}{10}\).
\end{theorem}

\begin{proof}
SSU + Schur bound: \(\ErrSSU\ll L^{C}\sqrt{H/X}\).
With \(H=(\log X)^{10}\) and \(X\ge 10^{12}\): \(\sqrt{H/X}\le 10^{-6}L^5\ll L^{-3}\).
Absorb \(L^C\) into the exponent (our \(A=10\) is ample) to get \(L^{-2}\).
\end{proof}

\begin{theorem}[Low-pass endpoint]\label{thm:lowpassendpoint-body}
\(\ErrTail(X,H,Q)\le \kappaFour/H^2\) for \(H\ge 1\).
\end{theorem}

\begin{proof}
Appendix A’s calculus for the \(\cos^2\) low-pass gives two integrations by parts and band-limit, hence \(O(H^{-2})\) with absolute \(\kappaFour\).
\end{proof}

\subsection*{Structural inequality and assembly}
\begin{theorem}[Structural decomposition]\label{thm:structural-body}
For even \(n\in[X,2X]\),
\(
R_2(X,n)\ \ge\ \MainSym(X,n)-(\ErrTypeI+\ErrSSU+\ErrTail).
\)
\end{theorem}

\begin{proof}
Bank/TFA algebra splits the correlation into the main and three error pieces; take absolute values on the error side.
\end{proof}

\begin{lemma}[Budget combiner]\label{lem:combinerY-body}
For \(X\ge X_0\),
\(\ErrBudget(X,H,Q)\le \frac{A_*+B_*+C_*}{L^2}\) with absolute \(A_*,B_*,C_*>0\).
\end{lemma}

\begin{proof}
Sum Theorems~\ref{thm:typeIendpoint-body}, \ref{thm:ssuendpoint-body}, \ref{thm:lowpassendpoint-body} under the ledger choice, absorbing constants.
\end{proof}

\begin{theorem}[Dyadic lower bound]\label{thm:dyadicLB-body}
For even \(n\in[X,2X]\), \(R_2(X,n)\ge \czero L^2-\frac{A_*+B_*+C_*}{L^2}\).
\end{theorem}

\begin{proof}
Section 10 positivity gives \(\MainSym\ge \LB=\czero L^2\). Apply Theorem~\ref{thm:structural-body} and Lemma~\ref{lem:combinerY-body}.
\end{proof}

\subsection*{Margin link}
Let \(\Margin(X)=\LB(X)-\ErrBudget(X,H,Q)\).
By Lemma~\ref{lem:combinerY-body}, \(\Margin(X)>0\) for \(X\ge X_0\).
Together with Theorem~\ref{thm:tailMarginBridge-body}\, this yields the dyadic-to-pointwise upgrade needed for the final certification.

\begin{lemma}[TT$^\ast$ mapping]\label{lem:Y.preErr.to.CLS}
There exists $C>0$ such that
\[
  \PreErrTypeI(X,H,Q)\ \le\ (\log(\max\{X,3\}))^{C}\cdot E_{\mathrm{CLS}}(X,H,Q),
\]
where $E_{\mathrm{CLS}}$ is the Gram–diagonalized TT$^\ast$ energy.
\end{lemma}

\begin{proof}
Write the Type–I bilinear form $B(f)$ and its TT$^\ast$ operator $T=B^\ast B$.
AO dispersion gives $\|G-I\|\le \varepsilon(H,Q)$ with $\varepsilon<1$, hence by Neumann series
$G$ is invertible and $\langle Gv,v\rangle\ge (1-\varepsilon)\|v\|^2$ for all $v$.
The occupancy bound (no hot spots) contributes $Q^{-2}$. Unfolding $T$ on bank atoms
and applying Cauchy–Schwarz yields
\(
  \PreErrTypeI \le (\log X)^{C}\langle Gv,v\rangle
   \le (\log X)^{C} E_{\mathrm{CLS}}.
\)
\end{proof}

\begin{lemma}[Low–pass squeeze]\label{lem:Y.KH.lowpass}
For the $\cos^2$ low–pass $K_H$,
\[
  E_{\mathrm{CLS}}(X,H,Q)\ \le\ \frac{X+Q^2}{H}.
\]
\end{lemma}

\begin{proof}
The Fourier support of $K_H$ lies in $|\xi|\le 1/H$ and $\int K_H=1$.
The CLS energy integrates the kernel along short shifts; restricting to $|\xi|\le 1/H$
forces a $1/H$ scaling. The $X$ term is the bulk, while the $Q^2$ term accounts for
aliasing up to denominator $Q$; both are standard in Appendix~A’s calculus.
\end{proof}

\begin{proof}[Proof of Theorem~\ref{thm:typeIendpoint-body}]
Combine Lemma~\ref{lem:Y.preErr.to.CLS} and Lemma~\ref{lem:Y.KH.lowpass}, then use the ledger constraint
$Q^2\le X$ to bound $(X+Q^2)/H\le X/H$.
\end{proof}

\begin{proof}[Proof of Theorem~\ref{thm:ssuendpoint-body}]
Decompose the off–diagonal into tubes $T_\tau$.
For each tube, Schur’s test with weights $w(d,n)=1$ gives
$\|T_\tau\|_{\ell^2\to\ell^2}\le \sqrt{(\sup_{(d,n)}\sum_{(d',n')}(K_H)_{\tau})\cdot
(\sup_{(d',n')}\sum_{(d,n)}(K_H)_{\tau})}$.
Low–pass smoothness and tube geometry produce $(\log X)^{C}\sqrt{H/X}$.
Summing over $\tau$ contributes another polylog factor, still absorbed into $(\log X)^C$.
With $H=(\log X)^{10}$ and $X\ge 10^{12}$, $\sqrt{H/X}\le 10^{-6} L^5\ll L^{-3}$, hence
$\ErrSSU\le L^{-2}$.
\end{proof}

\begin{lemma}[Mean–zero]\label{lem:Y.delta1}
The $\delta=1$ mask has zero mean: $\int_{\mathbb T}\!\delta_1(\theta)\,d\theta=0$.
\end{lemma}
\begin{proof}
By construction $\delta_1(\theta)=\psi(\theta)-\psi(\theta+1)$ with $\psi$ periodic and
$\int \psi= \int \psi(\cdot+1)$, hence their difference integrates to $0$.
\end{proof}

\begin{lemma}[Alias cancellation]\label{lem:Y.delta2}
The $\delta=2$ mask is row/column balanced; the first alias sum vanishes.
\end{lemma}
\begin{proof}
Arrange the bank atoms in a matrix by row/column indices; the balanced row/column sums
cancel the first harmonic by symmetry of the bank centers and the guard.
\end{proof}

\begin{lemma}[Lower frame from Gram closeness]\label{lem:Y.frame}
If $\|G-I\|\le \varepsilon<1$ for the bank Gram matrix, then
$\langle Gv,v\rangle\ge (1-\varepsilon)\|v\|^2$ for all $v$.
\end{lemma}
\begin{proof}
Neumann series: $G=I-(I-G)$ with $\|I-G\|<1$ implies $G$ invertible and
$\|G^{-1}\|\le 1/(1-\varepsilon)$. Then $(1-\varepsilon)\|v\|^2\le \langle Gv,v\rangle$.
\end{proof}

\begin{lemma}[Ledger monomial $\to$ log profile]\label{lem:Y.ledger}
For $X\ge X_0$, the ledger function $\Phi(X)$ satisfies
$\Phi(X)\le C_\Phi / (\log(\max\{X,3\}))^2$.
\end{lemma}
\begin{proof}
The concrete ledger inequalities reduce $\Phi$ to a finite maximum of rational expressions
in $X$ and $L=\log(\max\{X,3\})$. Each term is monotone in $X$ for $X\ge X_0$ and bounded by
$C_\Phi/L^2$; the constant $C_\Phi$ is fixed by the numeric ledger (Appendix tables).
\end{proof}

           % finite computation threshold $\Phi$ (if included)

%--------------------------------------------------------
%  Bibliography
%--------------------------------------------------------

\nocite{*}
\printbibliography

\end{document}